%%
%% main.tex — MatinBook Documentation
%% 
%% This is the main entry point for the MatinBook user guide.
%% It assembles all chapters from the various parts into a
%% single, coherent book.
%%
%% Compilation:
%%   xelatex -shell-escape main.tex
%%   biber main
%%   xelatex -shell-escape main.tex
%%   xelatex -shell-escape main.tex
%%
%% Or use latexmk:
%%   latexmk -xelatex -shell-escape main.tex
%%
%% IMPORTANT:
%%   - The cover requires at least 2 compilation passes.
%%   - The bibliography requires biber.
%%   - The index requires xindy (called automatically via
%%     -shell-escape).
%%

\documentclass{matinbook}

%========================================================================
% BIBLIOGRAPHY
%========================================================================
\addbibresource{references.bib}

%========================================================================
% BOOK METADATA
%========================================================================
\title{راهنمای جامع MatinBook}
\author{مatin محمودی}
\date{مهر ۱۴۰۵}

\booktitle{راهنمای جامع \lr{MatinBook}}

\renewcommand{\repository}{github.com/matinmahmoudi-cs/matinbook}

%========================================================================
% CUSTOM MACROS FOR THIS DOCUMENTATION
%========================================================================

% Mark a command name: \cmd{\makecover}
\newcommand{\cmd}[1]{\lr{\textbackslash #1}}

% Mark an environment name: \env{theorem}
\newcommand{\env}[1]{\lr{\texttt{#1}}}

% Mark a package name: \pkg{minted}
\newcommand{\pkg}[1]{\lr{\texttt{#1}}}

% Mark a file name: \file{main.tex}
\newcommand{\file}[1]{\lr{\texttt{#1}}}

% Mark a menu path: \menu{File > Open}
\newcommand{\menu}[1]{\lr{\texttt{#1}}}

% Mark a key: \key{Ctrl+S}
\newcommand{\key}[1]{\lr{\texttt{#1}}}

% Mark an option: \opt{draft}
\newcommand{\opt}[1]{\lr{\texttt{#1}}}

% Version of MatinBook documented here.
\newcommand{\mbversion}{1.2}

%========================================================================
% DOCUMENT
%========================================================================
\begin{document}

%------------------------------------------------------------------------
% FRONT COVER
%------------------------------------------------------------------------
\makecover
    {راهنمای جامع MatinBook}
    {از نصب تا انتشار کتاب حرفه‌ای فارسی}
    {مatin محمودی}
    {مهر ۱۴۰۵}

%------------------------------------------------------------------------
% FRONT MATTER (symmetric margins)
%------------------------------------------------------------------------
\frontmatter
\frontmattergeometry

% Title page
\maketitle

% Colophon (copyright, ISBN, edition)
%%
%% frontmatter/colophon.tex
%%
%% Copyright page, edition information, and licensing details.
%%

\thispagestyle{empty}
\vspace*{1.5cm}

\begin{center}
    {\Large\bfseries راهنمای جامع \lr{MatinBook}}

    \vspace{0.4cm}

    {\large از نصب تا انتشار کتاب حرفه‌ای فارسی}

    \vspace{1.5cm}

    {\large\bfseries متین محمودی}

    \vspace{2cm}

    \rule{0.5\textwidth}{0.4pt}

    \vspace{1cm}

    {\bfseries شناسنامه‌ی کتاب}

    \vspace{0.6cm}

    \begin{tabular}{rl}
        عنوان: & راهنمای جامع \lr{MatinBook} \\
        زیرعنوان: & از نصب تا انتشار کتاب حرفه‌ای فارسی \\
        نویسنده: & متین محمودی \\
        ناشر: & انتشارات \lr{MatinBook} \\
        نوبت چاپ: & اول \\
        تیراژ: & \pnum{۱۰۰۰} نسخه \\
        شابک: & \lr{978-600-XXXX-XX-X} \\
        قطع: & وزیری \\
        تعداد صفحه: & حدود \pnum{۲۰۰} \\
        قیمت: & \pnum{۵۰۰٬۰۰۰} ریال \\
    \end{tabular}

    \vspace{1cm}

    \rule{0.5\textwidth}{0.4pt}

    \vspace{1cm}

    {\small
        نسخه‌ی \mbversion \\
        \matinbookdate \\
        مهر ۱۴۰۵
    }

    \vfill

    {\small
        \lr{Copyright \copyright\ 2025--2026 Matin Mahmoudi} \\
        \lr{Published under the MIT License} \\
        \lr{https://github.com/matinmahmoudi-cs/matinbook}
    }

    \vspace{0.5cm}

    {\footnotesize\itshape
        تمامی حقوق این اثر برای نویسنده محفوظ است. \\
        این کتاب تحت مجوز \lr{MIT} منتشر شده و استفاده‌ی
        آموزشی و غیرتجاری از آن آزاد است. \\
        برای استفاده‌ی تجاری، با نویسنده تماس بگیرید.
    }
\end{center}

\clearpage
\clearpage

% Dedication
%%
%% frontmatter/dedication.tex
%%
%% Dedication page.
%%
%% NOTE on RTL:
%%   In an RTL document, `flushright` aligns to the START of
%%   the text (the right edge of the page), which is what we
%%   want for the dedication body. `flushleft` aligns to the
%%   END of the text (the left edge), which is what we want
%%   for the signature at the bottom.
%%

\thispagestyle{empty}
\vspace*{5cm}

% Body of the dedication — right-aligned in RTL
\begin{flushright}
    {\Large\itshape
        تقدیم به\\
        تمام نویسندگانی که\\
        با کلمات، دنیایی تازه می‌سازند\\
        و\\
        به تمام خوانندگانی که\\
        با هر صفحه، دانشی می‌آموزند.
    }
\end{flushright}

\vfill

% Signature — left-aligned in RTL
\begin{flushleft}
    {\large\itshape
        متین محمودی\\
        مهر ۱۴۰۵
    }
\end{flushleft}

\clearpage
\clearpage

% Preface
%%
%% frontmatter/preface.tex
%%
%% Preface: Why MatinBook, who is this book for, what is
%% its philosophy.
%%

\chapter*{پیشگفتار}
\addcontentsline{toc}{chapter}{پیشگفتار}
\label{ch:preface}

\section*{چرا \lr{MatinBook}؟}

حروف‌چینی کتاب‌های فنی فارسی، سال‌هاست با چالش‌های
متعددی روبه‌روست. ابزارهای موجود، یا برای زبان فارسی
بهینه نشده‌اند، یا از نظر بصری کیفیت لازم را ندارند،
یا برای پروژه‌های بزرگ مقیاس‌پذیر نیستند. این وضعیت،
باعث شده که بسیاری از کتاب‌های فنی فارسی، با وجود
محتوای ارزشمند، از نظر ظاهری فاصله‌ی زیادی با
استانداردهای نشر جهانی داشته باشند.

\lr{MatinBook} با هدف پر کردن این شکاف طراحی شده است.
این کلاس \lr{LaTeX}، تمام ابزارهای لازم برای تولید
یک کتاب فنی حرفه‌ای فارسی را در یک بسته‌ی منسجم فراهم
می‌کند: از حروف‌چینی راست‌به‌چپ و پشتیبانی کامل از
فونت‌های فارسی، تا محیط‌های علمی رنگ‌آمیزی‌شده،
نمودارهای پیچیده، و مدیریت حرفه‌ای منابع و نمایه.

\section*{این کتاب برای چه کسانی است؟}

این راهنما برای چهار گروه اصلی نوشته شده است:

\begin{itemize}
    \item \textbf{نویسندگان:} کسانی که می‌خواهند کتاب
    فنی یا علمی فارسی بنویسند و به دنبال ابزاری حرفه‌ای
    برای حروف‌چینی هستند.

    \item \textbf{پژوهشگران:} کسانی که می‌خواهند مقالات،
    پایان‌نامه‌ها یا گزارش‌های فنی خود را با کیفیت
    بالا منتشر کنند.

    \item \textbf{دانشجویان:} کسانی که می‌خواهند با
    ابزارهای حرفه‌ای حروف‌چینی آشنا شوند و مهارتی
    ارزشمند برای آینده‌ی حرفه‌ای خود کسب کنند.

    \item \textbf{مدرسان:} کسانی که می‌خواهند جزوه‌ها،
    کتاب‌های درسی یا مطالب آموزشی خود را با استاندارد
    بالایی آماده کنند.
\end{itemize}

\section*{فلسفه‌ی طراحی}

\lr{MatinBook} بر سه اصل بنیادین استوار است:

\subsection*{۱. وضوح بر تزئین}

هر عنصر بصری در \lr{MatinBook}، نقشی معنایی دارد.
رنگ‌ها، جعبه‌ها، و نمودارها، نه برای زیبایی صرف،
بلکه برای کمک به درک بهتر محتوا طراحی شده‌اند.
این اصل، از فلسفه‌ی نشر \lr{MIT Press} و
\lr{Springer} الهام گرفته است.

\subsection*{۲. سازگاری بر نوآوری}

هر عنصر در \lr{MatinBook}، یک شکل استاندارد دارد
و در تمام کتاب، یکسان ظاهر می‌شود. این سازگاری،
باعث می‌شود خواننده بدون حواس‌پرتی، روی محتوا
تمرکز کند.

\subsection*{۳. خویشتن‌داری بر افزونگی}

فضای خالی، یک عنصر طراحی است. \lr{MatinBook}
از پر کردن هر اینچ صفحه پرهیز می‌کند و به متن
اجازه می‌دهد نفس بکشد. این اصل، خوانایی را
افزایش می‌دهد و خستگی چشم را کاهش می‌دهد.

\section*{ساختار این راهنما}

این راهنما در شش بخش اصلی و پنج پیوست سازماندهی
شده است:

\begin{enumerate}
    \item \textbf{شروع به کار:} معرفی، نصب، اولین
    کتاب، ساختار کلی.

    \item \textbf{اجزای کتاب:} جلد، فصل، متن، محیط‌های
    علمی، ریاضیات، کد، نمودار، جدول، یادداشت حاشیه،
    منابع، نمایه، ارجاع.

    \item \textbf{شخصی‌سازی:} آپشن‌ها، رنگ‌ها، فونت‌ها،
    صفحه‌آرایی، تم‌ها، زبان.

    \item \textbf{مباحث پیشرفته:} ماژول‌ها، تم سفارشی،
    دستورات سفارشی، پروژه‌های بزرگ.

    \item \textbf{مثال‌های کامل:} سه فصل نمونه از
    کتاب‌های ریاضی، برنامه‌نویسی و فیزیک.

    \item \textbf{مرجع دستورات:} مرجع کامل کلاس،
    محیط‌ها و دستورات.
\end{enumerate}

هر فصل با یک مثال ساده شروع می‌شود، به تدریج پیچیده‌تر
می‌شود، و با تمرین‌های عملی به پایان می‌رسد. خواننده
می‌تواند فصل‌ها را به ترتیب بخواند، یا هر فصل را
به‌عنوان یک مرجع مستقل استفاده کند.

\section*{سپاسگزاری}

\lr{MatinBook} بدون پشتیبانی جامعه‌ی \lr{LaTeX} فارسی
و تلاش‌های بی‌دریغ توسعه‌دهندگان بسته‌های متن‌باز،
امکان‌پذیر نبود. از \lr{Vafa Khalighi} برای
\lr{XePersian}، از \lr{Hassan Mesgarha} برای
\lr{XePersian-HM}، از \lr{Thomas F. Sturm} برای
\lr{tcolorbox}، از \lr{Geoffrey Poore} برای
\lr{minted}، و از تمام کسانی که در این مسیر
کمک کرده‌اند، صمیمانه سپاسگزارم.

\vspace{1cm}

\begin{flushleft}
    \textbf{متین محمودی} \\
    مهر ۱۴۰۵
\end{flushleft}

\clearpage
\clearpage

% About this guide
%%
%% frontmatter/about.tex
%%
%% About this guide: conventions, notation, how to read it.
%%

\chapter*{درباره‌ی این راهنما}
\addcontentsline{toc}{chapter}{درباره‌ی این راهنما}
\label{ch:about}

\section*{ساختار این راهنما}

این راهنما به شش بخش اصلی و پنج پیوست تقسیم شده
است. هر بخش، موضوعی مستقل را پوشش می‌دهد و می‌تواند
به‌عنوان یک مرجع جداگانه استفاده شود.

\subsection*{بخش‌های اصلی}

\begin{description}
    \item[بخش ۱ — شروع به کار:] معرفی \lr{MatinBook}،
    نصب و راه‌اندازی، ایجاد اولین کتاب، و ساختار
    کلی یک کتاب.

    \item[بخش ۲ — اجزای کتاب:] تمام اجزای یک کتاب،
    از جلد و فصل تا منابع و نمایه.

    \item[بخش ۳ — شخصی‌سازی:] آپشن‌ها، رنگ‌ها،
    فونت‌ها، صفحه‌آرایی، تم‌ها و زبان.

    \item[بخش ۴ — مباحث پیشرفته:] ماژول‌ها، تم سفارشی،
    دستورات سفارشی و پروژه‌های بزرگ.

    \item[بخش ۵ — مثال‌های کامل:] سه فصل کامل از
    کتاب‌های ریاضی، برنامه‌نویسی و فیزیک.

    \item[بخش ۶ — مرجع دستورات:] مرجع کامل کلاس،
    محیط‌ها و دستورات.
\end{description}

\subsection*{پیوست‌ها}

\begin{description}
    \item[پیوست الف — عیب‌یابی:] خطاهای رایج و
    راه‌حل آن‌ها.

    \item[پیوست ب — ابزارهای موردنیاز:] نصب و
    پیکربندی \lr{XeLaTeX}، \lr{biber}، \lr{xindy} و
    \lr{Pygments}.

    \item[پیوست پ — منابع:] کتاب‌ها، وب‌سایت‌ها و
    انجمن‌های مفید.

    \item[پیوست ت — تاریخچه:] نسخه‌های \lr{MatinBook}
    و تغییرات آن‌ها.

    \item[پیوست ث — مجوز و مشارکت:] جزئیات مجوز
    و راهنمای مشارکت.
\end{description}

\section*{قراردادهای نگارشی}

در این راهنما، از قراردادهای زیر برای نشان دادن
انواع مختلف محتوا استفاده می‌شود:

\subsection*{نام دستورات}

نام دستورات \lr{LaTeX} با فونت \lr{monospace} و
علامت \lr{\textbackslash} نمایش داده می‌شوند:

\begin{itemize}
    \item \cmd{makecover} — دستور جلد
    \item \cmd{chapter} — دستور فصل
    \item \cmd{cref} — دستور ارجاع
\end{itemize}

\subsection*{نام محیط‌ها}

نام محیط‌ها با فونت \lr{monospace} و بدون
علامت \lr{\textbackslash} نمایش داده می‌شوند:

\begin{itemize}
    \item \env{theorem} — محیط قضیه
    \item \env{minted} — محیط کد
    \item \env{algorithm} — محیط الگوریتم
\end{itemize}

\subsection*{نام بسته‌ها}

نام بسته‌ها با فونت \lr{monospace} نمایش داده
می‌شوند:

\begin{itemize}
    \item \pkg{tcolorbox} — بسته‌ی جعبه‌ها
    \item \pkg{minted} — بسته‌ی کد
    \item \pkg{tabularray} — بسته‌ی جدول
\end{itemize}

\subsection*{نام فایل‌ها}

نام فایل‌ها با فونت \lr{monospace} نمایش داده
می‌شوند:

\begin{itemize}
    \item \file{main.tex} — فایل اصلی
    \item \file{references.bib} — کتاب‌نامه
    \item \file{matinbook.cls} — کلاس اصلی
\end{itemize}

\subsection*{گزینه‌های کلاس}

گزینه‌های کلاس با فونت \lr{monospace} نمایش داده
می‌شوند:

\begin{itemize}
    \item \opt{draft} — حالت پیش‌نویس
    \item \opt{final} — حالت نهایی
    \item \opt{vazirmatn} — فونت وزیرمتن
\end{itemize}

\subsection*{کدها}

کدهای \lr{LaTeX} در بلوک‌های زیر نمایش داده
می‌شوند:

\begin{latin}
\begin{minted}{latex}
\documentclass{matinbook}
\begin{document}
\chapter{|\pc{فصل نمونه}|}
\end{document}
\end{minted}
\end{latin}

\subsection*{خروجی}

خروجی کدها، در صورت نیاز، در بلوک‌های جداگانه
نمایش داده می‌شود. مثال:

\begin{quote}
    \itshape
    این یک متن نمونه است که خروجی کد بالا را
    نشان می‌دهد.
\end{quote}

\section*{نحوه‌ی خواندن این راهنما}

این راهنما به سه روش قابل استفاده است:

\subsection*{روش ۱ — خواندن متوالی}

اگر تازه با \lr{MatinBook} آشنا می‌شوید، توصیه
می‌شود فصل‌ها را به ترتیب بخوانید. هر فصل، بر
پایه‌ی فصل‌های قبلی بنا شده است.

\subsection*{روش ۲ — مرجع سریع}

اگر با \lr{MatinBook} آشنا هستید و فقط می‌خواهید
یک قابلیت خاص را یاد بگیرید، می‌توانید مستقیماً
به فصل مربوطه مراجعه کنید. هر فصل، مستقل قابل
خواندن است.

\subsection*{روش ۳ — جستجو}

برای پیدا کردن یک دستور یا مفهوم خاص، از نمایه‌ی
پایان کتاب استفاده کنید. نمایه، به‌صورت الفبایی
مرتب شده و شامل تمام دستورات، محیط‌ها و مفاهیم
کلیدی است.

\section*{تمرین‌ها}

در پایان هر فصل، مجموعه‌ای از تمرین‌ها ارائه
شده است. توصیه می‌شود این تمرین‌ها را انجام
دهید، زیرا یادگیری \lr{LaTeX} بدون تمرین
عملی امکان‌پذیر نیست.

\section*{نمادها}

در این راهنما، از نمادهای زیر برای تأکید بر
نکات مهم استفاده می‌شود:

\begin{itemize}
    \item \textbf{نکته:} اطلاعات تکمیلی که
    ممکن است مفید باشد.
    \item \textbf{هشدار:} نکاتی که باید به
    آن‌ها توجه کنید تا از خطا جلوگیری شود.
    \item \textbf{مثال:} نمونه‌های عملی از
    مفاهیم.
\end{itemize}

\vspace{1cm}

\begin{flushleft}
    \textbf{متین محمودی} \\
    مهر ۱۴۰۵
\end{flushleft}

\clearpage
\clearpage

% Table of Contents
\tableofcontents
\clearpage

% List of Figures
\listoffigures
\clearpage

% List of Tables
\listoftables
\clearpage

% List of Algorithms
\listofalgorithms
\clearpage

%------------------------------------------------------------------------
% MAIN MATTER (wide margin for notes)
%------------------------------------------------------------------------
\mainmatter
\mainmattergeometry

%========================================================================
% PART 1 — GETTING STARTED
%========================================================================
\part{شروع به کار}
\label{part:getting-started}

%%
%% part1-getting-started/ch01-introduction.tex
%%
%% Chapter 1: Introduction to MatinBook — What it is, what it
%% can do, and how it compares to other LaTeX classes.
%%

\chapter{معرفی و قابلیت‌ها}
\label{chap:introduction}

\section{چرا این موضوع مهم است؟}
\label{sec:introduction-why}

اگر تا به حال تلاش کرده‌اید یک کتاب فنی فارسی با
ابزارهای حرفه‌ای حروف‌چینی کنید، احتمالاً با
چالش‌های متعددی روبه‌رو شده‌اید: فونت‌های فارسی
که در اندازه‌های کوچک خوانا نیستند، فرمول‌های
ریاضی که با متن فارسی قاطی می‌شوند، کدهای
برنامه‌نویسی که جهت متن را به‌هم می‌ریزند، و
ده‌ها مشکل کوچک و بزرگ دیگر.

\lr{MatinBook} برای حل همین مشکلات طراحی شده است.
این کلاس \lr{LaTeX}، تمام ابزارهای لازم برای
تولید یک کتاب فنی فارسی حرفه‌ای را در یک بسته‌ی
منسجم فراهم می‌کند. هدف این فصل، آشنایی اولیه با
\lr{MatinBook}، قابلیت‌های کلیدی آن، و تفاوت‌هایش
با سایر ابزارهای موجود است.

\mnote{\lr{MatinBook} یک کلاس \lr{LaTeX} است، نه یک بسته‌ی ساده. این یعنی تمام قابلیت‌های کتاب‌سازی در آن یکپارچه شده است.}

\mnote{نام \lr{MatinBook} از نام کوچک توسعه‌دهنده‌ی آن، «متین»، گرفته شده است.}

\section{\lr{MatinBook} چیست؟}
\label{sec:introduction-what}

\lr{MatinBook} یک کلاس \lr{LaTeX} است که به‌طور
اختصاصی برای نوشتن کتاب‌های فنی و علمی فارسی طراحی
شده است. این کلاس توسط متین محمودی توسعه یافته و
تحت مجوز \lr{MIT} منتشر می‌شود.

\begin{definition}[\lr{MatinBook}]
    \lr{MatinBook} یک کلاس \lr{LaTeX} است که با
    موتور \lr{XeLaTeX} کامپایل می‌شود و تمام
    ابزارهای لازم برای تولید کتاب‌های فنی فارسی
    — از حروف‌چینی راست‌به‌چپ تا نمودارهای پیچیده —
    را در یک ساختار ماژولار فراهم می‌کند.
    \label{def:matinbook}
\end{definition}

\lr{MatinBook} بر پایه‌ی کلاس استاندارد \lr{book}
ساخته شده است، ولی قابلیت‌های زیادی به آن اضافه
کرده است:

\begin{itemize}
    \item پشتیبانی کامل از حروف‌چینی راست‌به‌چپ
    فارسی با \lr{XePersian}.
    \item محیط‌های علمی رنگ‌آمیزی‌شده (قضیه، تعریف،
    مثال، اثبات).
    \item پشتیبانی از کدهای برنامه‌نویسی با
    رنگ‌آمیزی نحوی و کامنت فارسی.
    \item نمودارها و شکل‌های حرفه‌ای با \lr{TikZ} و
    \lr{PGFPlots}.
    \item جداول مدرن با \lr{tabularray}.
    \item مدیریت منابع با \lr{biblatex} و \lr{biber}.
    \item نمایه‌ی حرفه‌ای با \lr{xindy} و مرتب‌سازی
    صحیح فارسی.
    \item جلد تمام‌رنگی با طرح حرفه‌ای.
    \item یادداشت‌های حاشیه‌ای برای توضیحات مکمل.
\end{itemize}

\mnote{کلمه‌ی کلیدی اینجا «یکپارچه» است. در ابزارهای دیگر، باید هر قابلیت را دستی پیکربندی کنید.}

\section{قابلیت‌های کلیدی}
\label{sec:introduction-features}

در این بخش، نگاهی گذرا به مهم‌ترین قابلیت‌های
\lr{MatinBook} می‌اندازیم. هر یک از این قابلیت‌ها
در فصل‌های بعدی به‌طور مفصل بررسی خواهد شد.

\subsection{محیط‌های علمی رنگ‌آمیزی‌شده}
\label{sec:introduction-environments}

\lr{MatinBook} ده نوع محیط علمی را با دو خانواده‌ی
بصری متمایز ارائه می‌دهد. محیط‌های نظری (قضیه، لم،
نتیجه، گزاره) از یک نوار رنگی توپر استفاده
می‌کنند، در حالی که محیط‌های توضیحی (مثال، اثبات،
راه‌حل) از یک خط عمودی سمت راست بهره می‌برند.

\mnote{تفکیک دو خانواده، بر اساس نوع محتوا است، نه سلیقه. محیط‌های نظری «رسمی» هستند، محیط‌های توضیحی «غیررسمی».}

\begin{theorem}[قضیه‌ی فیثاغورث]
    در هر مثلث قائم‌الزاویه، مربع وتر برابر است با
    مجموع مربعات دو ضلع دیگر:
    \[
        a^{2} + b^{2} = c^{2}
    \]
    \label{thm:pythagoras-intro}
\end{theorem}

\begin{proof}
    اثبات این قضیه را در فصل \cref{chap:environments}
    به‌طور کامل بررسی خواهیم کرد.
\end{proof}

\begin{example}
    برای یک مثلث با ضلع‌های $a = 3$ و $b = 4$،
    وتر برابر است با:
    \[
        c = \sqrt{3^{2} + 4^{2}} = \sqrt{25} = 5
    \]
    \label{ex:pythagoras-intro}
\end{example}

\begin{remark}
    محیط‌های علمی در \lr{MatinBook} از یک شمارنده‌ی
    مشترک استفاده می‌کنند. یعنی اگر یک تعریف با
    شماره‌ی ۱.۱ ثبت شود، قضیه‌ی بعدی شماره‌ی ۱.۲
    خواهد داشت.
    \label{rem:shared-counter-intro}
\end{remark}

این محیط‌ها در فصل \cref{chap:environments} به‌طور
مفصل بررسی می‌شوند.

\subsection{پشتیبانی از ریاضیات}
\label{sec:introduction-math}

\lr{MatinBook} از تمام قابلیت‌های \lr{amsmath}،
\lr{mathtools} و \lr{unicode-math} پشتیبانی می‌کند
و مجموعه‌ای از عملگرها و جداکننده‌های سفارشی را
نیز اضافه کرده است:

\begin{itemize}
    \item عملگرهای حسابان برداری: \cmd{grad}،
    \cmd{curl}، \cmd{diver}.
    \item عملگرهای جبر خطی: \cmd{rank}، \cmd{trace}،
    \cmd{diag}.
    \item مجموعه‌های اعداد: $\N$، $\Z$، $\Q$، $\R$،
    $\C$.
    \item جداکننده‌های هوشمند: \cmd{abs}، \cmd{norm}،
    \cmd{ceil}، \cmd{floor}.
\end{itemize}

مثال:

\[
    \grad f = \nabla f, \quad
    \diver \vec{F} = \nabla \cdot \vec{F}, \quad
    \norm{\vec{v}} = \sqrt{\inner{\vec{v}}{\vec{v}}}
\]

\mnote{در متن فارسی، فرمول‌های ریاضی باید داخل \lr{\textbackslash lr\{...\}} قرار گیرند تا جهت‌نمایی درست حفظ شود.}

ریاضیات در \lr{MatinBook} در فصل \cref{chap:math}
به‌طور کامل بررسی می‌شود.

\subsection{کد و الگوریتم}
\label{sec:introduction-code}

\lr{MatinBook} از بسته‌ی \lr{minted} برای نمایش کد
با رنگ‌آمیزی نحوی استفاده می‌کند. این بسته از بیش
از ۳۰۰ زبان برنامه‌نویسی پشتیبانی می‌کند و امکان
نوشتن کامنت فارسی داخل کد را نیز فراهم می‌سازد.

برای کامنت فارسی در پایتون، از دستور \cmd{pcc}
استفاده می‌کنیم که به‌طور خودکار علامت \lr{\#} را
قبل از متن فارسی قرار می‌دهد:

\begin{latin}
\begin{minted}{python}
def factorial(n: int) -> int:
    """Calculate n! recursively."""
    if n <= 1:
        return 1
    return n * factorial(n - 1)

\pcc{محاسبه‌ی فاکتوریل ۵}
print(factorial(5))  # 120
\end{minted}
\end{latin}
\codecaption{تابع فاکتوریل بازگشتی در پایتون}
\label{code:factorial-intro}

\mnote{برای کامنت فارسی در پایتون، \lr{\textbackslash pcc} بهترین گزینه است. برای زبان‌های دیگر (مثل \lr{LaTeX} یا \lr{bash}) می‌توانید از \lr{\textbackslash pc} داخل \lr{|...|} استفاده کنید.}

\mnote{اگر \lr{\textbackslash pcc} را دستی باز کنید، به \lr{|\textbackslash pc\{\textbackslash\#~متن\}|} بسط می‌یابد. این یعنی \lr{\#} داخل \lr{\textbackslash pc} قرار می‌گیرد تا \lr{Pygments} آن را به‌عنوان شروع کامنت نبیند.}

الگوریتم‌ها نیز با بسته‌ی استاندارد \lr{algorithm}
نمایش داده می‌شوند. نکته‌ی مهم این است که
\lr{xepersian} عنوان فارسی «الگوریتم» را به‌طور
خودکار اضافه می‌کند، ولی متن داخل الگوریتم باید
لاتین باشد (به‌جز کامنت‌هایی که با \cmd{pc} نوشته
می‌شوند):

\begin{latin}
\begin{algorithm}
\caption{جستجوی دودویی}
\label{alg:binary-intro}
\begin{algorithmic}[1]
    \Require Sorted array $A$ of $n$ elements
    \Require Target value $x$
    \Ensure Index of $x$ in $A$, or $-1$
    \State $left \gets 0$, $right \gets n - 1$
    \While{$left \leq right$}
        \State $mid \gets (left + right) / 2$
        \If{$A[mid] = x$}
            \State \Return $mid$  \Comment{|\pc{یافتن مقدار}|}
        \ElsIf{$A[mid] < x$}
            \State $left \gets mid + 1$  \Comment{|\pc{جستجو در نیمه‌ی راست}|}
        \Else
            \State $right \gets mid - 1$  \Comment{|\pc{جستجو در نیمه‌ی چپ}|}
        \EndIf
    \EndWhile
    \State \Return $-1$  \Comment{|\pc{یافت نشد}|}
\end{algorithmic}
\end{algorithm}
\end{latin}

\mnote{در الگوریتم‌ها، متن داخل \lr{\textbackslash Comment\{...\}} می‌تواند فارسی باشد، ولی باید با \lr{\textbackslash pc} پیچیده شود.}

کد و الگوریتم در \cref{chap:code} به‌طور مفصل
بررسی می‌شوند.

\subsection{نمودارها و شکل‌ها}
\label{sec:introduction-graphics}

\lr{MatinBook} از \lr{TikZ}، \lr{PGFPlots} و
\lr{pgf-pie} برای ترسیم نمودارها و شکل‌های حرفه‌ای
استفاده می‌کند. از نمودارهای ساده‌ی توابع تا
نمودارهای سه‌بعدی، فلوچارت‌ها، درخت‌ها و گراف‌ها،
همه در \lr{MatinBook} پشتیبانی می‌شوند.

\begin{figure}[h]
\centering
\begin{latin}
\begin{tikzpicture}
  \begin{axis}[
    width=0.8\textwidth, height=5cm,
    xlabel={$x$},
    ylabel={$f(x)$},
    grid=major,
    legend pos=north west,
    domain=-2:2,
    samples=100
  ]
    \addplot[matinblue, thick] {x^2};
    \addlegendentry{$x^2$}
    \addplot[matinred, thick] {x^3};
    \addlegendentry{$x^3$}
  \end{axis}
\end{tikzpicture}
\end{latin}
\caption{نمودار دو تابع $x^{2}$ و $x^{3}$}
\label{fig:functions-intro}
\end{figure}

\mnote{برای متن فارسی داخل \lr{tikzpicture}، باید از \lr{\textbackslash rl\{...\}} استفاده کنید.}

نمودار \cref{fig:functions-intro} نمونه‌ای ساده از
قابلیت‌های گرافیکی \lr{MatinBook} است. این موضوع
در \cref{chap:graphics} به‌طور کامل بررسی می‌شود.

\subsection{جداول مدرن}
\label{sec:introduction-tables}

\lr{MatinBook} از بسته‌ی \lr{tabularray} برای
جداول استفاده می‌کند. این بسته، برخلاف \lr{tabularx}،
با \lr{xepersian} کاملاً سازگار است و از
حروف‌چینی راست‌به‌چپ پشتیبانی می‌کند.

\begin{matintable}{مقایسه‌ی سه الگوریتم مرتب‌سازی}{tab:sorting-intro}
\begin{tblr}{
    width = \textwidth,
    colspec = {l X[c] X[c]},
    hlines, vlines,
    colsep = 8pt,
    rowsep = 4pt,
    row{1} = {font=\bfseries},
}
الگوریتم & بهترین حالت & بدترین حالت \\
مرتب‌سازی سریع & \lr{$O(n \log n)$} & \lr{$O(n^2)$} \\
مرتب‌سازی ادغامی & \lr{$O(n \log n)$} & \lr{$O(n \log n)$} \\
مرتب‌سازی حبابی & \lr{$O(n)$} & \lr{$O(n^2)$} \\
\end{tblr}
\end{matintable}

\mnote{بسته‌ی \lr{tabularx} با \lr{xepersian} ناسازگار است و خطا می‌دهد. همیشه از \lr{tabularray} استفاده کنید.}

جدول \cref{tab:sorting-intro} نمونه‌ای از جدول‌های
\lr{MatinBook} است. فصل \cref{chap:tables} به‌طور
کامل به این موضوع می‌پردازد.

\subsection{جلد حرفه‌ای}
\label{sec:introduction-cover}

\lr{MatinBook} یک جلد تمام‌رنگی با طرح حرفه‌ای ارائه
می‌دهد. جلد جلو شامل پس‌زمینه‌ی سرمه‌ای، گوه‌ی
قرمز، بیضی طلایی و نمونه‌کد \lr{C++} است. جلد
عقب شامل دو بلوک ایستاتیک — «درباره این کتاب» و
«درباره این نویسنده» — با فاصله‌ی دقیق
\pnum{۰.۷} سانتی‌متر است.

\begin{figure}[h]
\centering
\begin{latin}
\begin{tikzpicture}[scale=0.6]
  % Simplified front cover mockup
  \fill[matinnavy] (0,0) rectangle (6,8);
  \fill[matincrimson] (0,8) -- (3.5,8) -- (2,6.5) -- cycle;
  \draw[white, opacity=0.5, line width=0.5pt] (0.2,0.2) rectangle (5.8,7.8);
  \node[white, font=\bfseries] at (3,5) {\rl{عنوان کتاب}};
  \node[white, font=\small] at (3,4.3) {\rl{زیرعنوان}};
  \node[matingold, font=\bfseries] at (3,1) {MatinBook};
\end{tikzpicture}
\end{latin}
\caption{نمای ساده‌شده‌ی جلد \lr{MatinBook}}
\label{fig:cover-intro}
\end{figure}

\mnote{در \lr{tikzpicture}، حتی برای متن کوتاه فارسی هم باید از \lr{\textbackslash rl\{...\}} استفاده کنید.}

شکل \cref{fig:cover-intro} نمای ساده‌شده‌ای از
جلد را نشان می‌دهد. فصل \cref{chap:cover} به‌طور
کامل به جلد می‌پردازد.

\section{مقایسه با سایر ابزارها}
\label{sec:introduction-comparison}

چندین کلاس و بسته‌ی \lr{LaTeX} برای نوشتن
کتاب‌های فارسی وجود دارد. در این بخش،
\lr{MatinBook} را با مهم‌ترین آن‌ها مقایسه
می‌کنیم.

\subsection{مقایسه با کلاس \lr{book} استاندارد}
\label{sec:introduction-vs-book}

کلاس \lr{book} استاندارد \lr{LaTeX}، پایه‌ی
\lr{MatinBook} است. ولی استفاده‌ی مستقیم از
آن برای کتاب فارسی، چالش‌های زیادی دارد:

\begin{itemize}
    \item \textbf{حروف‌چینی راست‌به‌چپ:} کلاس
    \lr{book} از RTL پشتیبانی نمی‌کند. باید
    \lr{xepersian} را دستی اضافه کنید.

    \item \textbf{فونت فارسی:} باید فونت فارسی
    را دستی تنظیم کنید.

    \item \textbf{محیط‌های علمی:} هیچ محیط علمی
    رنگ‌آمیزی‌شده‌ای ندارد.

    \item \textbf{کد و الگوریتم:} باید \lr{minted}
    و \lr{algorithm} را دستی پیکربندی کنید.

    \item \textbf{جلد:} هیچ طرح جلدی ندارد.
\end{itemize}

\lr{MatinBook} تمام این مشکلات را حل کرده است.

\subsection{مقایسه با \lr{XePersian} خالص}
\label{sec:introduction-vs-xepersian}

\lr{XePersian} یک بسته‌ی عالی برای حروف‌چینی
فارسی است، ولی یک بسته است، نه یک کلاس. اگر
از \lr{XePersian} خالص استفاده کنید، باید
همه‌ی موارد زیر را دستی پیکربندی کنید:

\begin{itemize}
    \item هندسه‌ی صفحه.
    \item سرصفحه و پاصفحه.
    \item محیط‌های علمی.
    \item سبک عناوین.
    \item کد و الگوریتم.
    \item جلد.
\end{itemize}

\mnote{تفاوت اصلی این است: \lr{XePersian} یک «موتور» است، \lr{MatinBook} یک «کتاب کامل» است.}

\lr{MatinBook} همه‌ی این‌ها را از پیش پیکربندی
شده ارائه می‌دهد.

\subsection{جدول مقایسه}
\label{sec:introduction-comparison-table}

جدول \cref{tab:comparison} مقایسه‌ی خلاصه‌ای
بین \lr{MatinBook}، کلاس \lr{book} استاندارد،
و \lr{XePersian} خالص را نشان می‌دهد.

\begin{matintable}{مقایسه‌ی \lr{MatinBook} با سایر ابزارها}{tab:comparison}
\begin{tblr}{
    width = \textwidth,
    colspec = {X[l] X[c] X[c] X[c]},
    hlines, vlines,
    colsep = 6pt,
    rowsep = 4pt,
    row{1} = {font=\bfseries},
}
قابلیت & \lr{MatinBook} & \lr{book} & \lr{XePersian} \\
حروف‌چینی RTL & \textcolor{matingreen}{بله} & \textcolor{matinred}{خیر} & \textcolor{matingreen}{بله} \\
فونت فارسی & \textcolor{matingreen}{خودکار} & \textcolor{matinred}{خیر} & \textcolor{matinorange}{دستی} \\
محیط‌های علمی & \textcolor{matingreen}{۱۰ نوع} & \textcolor{matinred}{خیر} & \textcolor{matinred}{خیر} \\
کد و الگوریتم & \textcolor{matingreen}{کامل} & \textcolor{matinred}{خیر} & \textcolor{matinred}{خیر} \\
نمودار & \textcolor{matingreen}{TikZ} & \textcolor{matinred}{خیر} & \textcolor{matinred}{خیر} \\
جدول مدرن & \textcolor{matingreen}{tblr} & \textcolor{matinorange}{tabular} & \textcolor{matinorange}{tabular} \\
نمایه‌ی فارسی & \textcolor{matingreen}{xindy} & \textcolor{matinred}{خیر} & \textcolor{matinorange}{دستی} \\
جلد & \textcolor{matingreen}{تمام‌رنگی} & \textcolor{matinred}{خیر} & \textcolor{matinred}{خیر} \\
ماژولار & \textcolor{matingreen}{۱۹ ماژول} & \textcolor{matinred}{خیر} & \textcolor{matinred}{خیر} \\
\end{tblr}
\end{matintable}

\section{معماری \lr{MatinBook}}
\label{sec:introduction-architecture}

\lr{MatinBook} یک کلاس ماژولار است. یعنی به‌جای
یک فایل بزرگ، از ۱۹ ماژول مستقل تشکیل شده است
که هر یک مسئولیت مشخصی دارند. این معماری،
نگهداری، توسعه و شخصی‌سازی کلاس را بسیار آسان‌تر
می‌کند.

\mnote{معماری ماژولار یعنی اگر فقط می‌خواهید یک ماژول را تغییر دهید، نیازی به دست زدن به بقیه‌ی کلاس ندارید.}

\subsection{ماژول‌های اصلی}
\label{sec:introduction-modules}

ماژول‌های \lr{MatinBook} به سه دسته تقسیم می‌شوند:

\textbf{۱. هسته (۲ ماژول):}
\begin{itemize}
    \item \file{mb-core.sty} — بارگذارنده‌ی
    بسته‌های پایه.
    \item \file{mb-utils.sty} — دستورات کمکی.
\end{itemize}

\textbf{۲. ماژول‌های کاربردی (۱۰ ماژول):}
\begin{itemize}
    \item \file{mb-math.sty} — ریاضیات.
    \item \file{mb-boxes.sty} — استایل‌های
    جعبه.
    \item \file{mb-theorem.sty} — محیط‌های
    قضیه.
    \item \file{mb-code.sty} — کد.
    \item \file{mb-layout.sty} — چیدمان.
    \item \file{mb-headings.sty} — عناوین.
    \item \file{mb-references.sty} — ارجاعات.
    \item \file{mb-graphics.sty} — گرافیک.
    \item \file{mb-index.sty} — نمایه.
    \item \file{mb-table.sty} — جداول.
\end{itemize}

\textbf{۳. تم و لوکال (۵ ماژول):}
\begin{itemize}
    \item \file{mb-theme-colors.sty} — پالت
    رنگی.
    \item \file{mb-theme-cover.sty} — جلد.
    \item \file{mb-theme-default.sty} —
    بارگذارنده‌ی تم.
    \item \file{fa-IR.sty} — لوکال فارسی.
    \item \file{en-US.sty} — لوکال انگلیسی.
\end{itemize}

\subsection{ترتیب بارگذاری}
\label{sec:introduction-loading-order}

ماژول‌ها در سه فاز بارگذاری می‌شوند:

\textbf{فاز ۱ (قبل از \lr{xepersian}):}
هسته، تم، و ماژول‌های کاربردی.

\textbf{فاز ۲:}
\lr{xepersian} — باید آخرین بسته باشد.

\textbf{فاز ۳ (بعد از \lr{xepersian}):}
فونت‌ها، تایپوگرافی، و لوکال.

\begin{remark}
    تغییر این ترتیب، باعث خطاهای عجیب و غریب
    می‌شود. اگر می‌خواهید ماژول جدیدی اضافه
    کنید، حتماً آن را در فاز مناسب قرار دهید.
    \label{rem:loading-order}
\end{remark}

معماری \lr{MatinBook} در \cref{chap:modules}
به‌طور کامل بررسی می‌شود.

\section{خطاهای رایج}
\label{sec:introduction-mistakes}

در این بخش، چند خطای رایج که کاربران جدید
\lr{MatinBook} ممکن است مرتکب شوند را
معرفی می‌کنیم:

\subsection{استفاده از \lr{pdfLaTeX} به‌جای \lr{XeLaTeX}}
\label{sec:introduction-mistake-engine}

\lr{MatinBook} فقط با موتور \lr{XeLaTeX} کار
می‌کند. اگر از \lr{pdfLaTeX} استفاده کنید،
خطای زیر را دریافت می‌کنید:

\begin{latin}
\begin{minted}{text}
! Package matinbook Error:
  MatinBook requires XeLaTeX.
\end{minted}
\end{latin}

\mnote{\lr{XeLaTeX} برای پشتیبانی از فونت‌های \lr{OpenType} و حروف‌چینی راست‌به‌چپ ضروری است.}

\textbf{راه‌حل:} همیشه از \lr{XeLaTeX} استفاده
کنید.

\subsection{فراموش کردن \lr{-shell-escape}}
\label{sec:introduction-mistake-shellescape}

\lr{MatinBook} از \lr{minted} و \lr{xindy}
استفاده می‌کند که نیاز به \lr{-shell-escape}
دارند. اگر آن را فراموش کنید، خطاهای زیر
را دریافت می‌کنید:

\begin{itemize}
    \item \lr{minted}: عدم توانایی در اجرای
    \lr{Pygments}.
    \item \lr{xindy}: عدم توانایی در ساخت نمایه.
\end{itemize}

\textbf{راه‌حل:} همیشه از
\lr{xelatex -shell-escape} استفاده کنید.

\subsection{فراموش کردن \lr{biber}}
\label{sec:introduction-mistake-biber}

اگر از \lr{biblatex} و \lr{biber} استفاده
می‌کنید، باید \lr{biber} را دستی اجرا کنید:

\begin{latin}
\begin{minted}{bash}
xelatex -shell-escape main.tex
biber main
xelatex -shell-escape main.tex
xelatex -shell-escape main.tex
\end{minted}
\end{latin}

\mnote{این چهار مرحله را می‌توانید با \lr{latexmk} خودکار کنید: \lr{latexmk -xelatex -shell-escape main.tex}.}

\textbf{راه‌حل:} این چرخه‌ی چهار مرحله‌ای را
همیشه برای کتاب‌های حاوی منابع اجرا کنید.

\subsection{وارد کردن متن فارسی در بلوک کد}
\label{sec:introduction-mistake-persian-code}

بلوک‌های \env{minted} از فونت لاتین استفاده
می‌کنند، پس **نمی‌توانند متن فارسی را مستقیماً
نمایش دهند**. اگر متن فارسی در بلوک کد بنویسید،
نتیجه به‌صورت نامفهوم نمایش داده می‌شود.

\textbf{راه‌حل:} بسته به زبان، از یکی از دو
دستور زیر استفاده کنید:

\begin{itemize}
    \item \textbf{پایتون:} از \cmd{pcc} استفاده
    کنید.
    \item \textbf{سایر زبان‌ها:} از
    \cmd{pc} داخل \lr{|...|} استفاده کنید.
\end{itemize}

\begin{latin}
\begin{minted}{python}
x = 5  |\pcc{مقدار متغیر}|
y = 10  |\pcc{مقدار دیگر}|
\end{minted}
\end{latin}

\mnote{دستور \lr{\textbackslash pcc} توسط \lr{mb-code.sty} تعریف شده و علامت \lr{\#} را به‌طور خودکار اضافه می‌کند.}

خطاهای رایج در \cref{app:troubleshooting}
به‌طور کامل بررسی می‌شوند.

\section{تمرین‌ها}
\label{sec:introduction-exercises}

\begin{exercise}
    سه قابلیت \lr{MatinBook} که برای کار شما
    مهم‌تر هستند را فهرست کنید و توضیح دهید
    چرا.
    \label{exc:introduction-features}
\end{exercise}

\begin{exercise}
    \lr{MatinBook} را با یکی از ابزارهای
    حروف‌چینی که قبلاً استفاده کرده‌اید
    مقایسه کنید. سه مزیت و یک عیب
    \lr{MatinBook} را بنویسید.
    \label{exc:introduction-compare}
\end{exercise}

\begin{exercise}
    در فصل‌های بعدی، چه موضوعاتی را می‌خواهید
    اول یاد بگیرید؟ یک فهرست اولویت‌بندی‌شده
    تهیه کنید.
    \label{exc:introduction-priority}
\end{exercise}

\section{خلاصه}
\label{sec:introduction-summary}

در این فصل، با \lr{MatinBook} آشنا شدیم:

\begin{itemize}
    \item \lr{MatinBook} یک کلاس \lr{LaTeX}
    است که برای کتاب‌های فنی فارسی طراحی
    شده است.

    \item این کلاس از موتور \lr{XeLaTeX}
    استفاده می‌کند و از حروف‌چینی راست‌به‌چپ
    کامل پشتیبانی می‌کند.

    \item قابلیت‌های کلیدی شامل محیط‌های علمی
    رنگ‌آمیزی‌شده، ریاضیات، کد، نمودار، جدول،
    جلد حرفه‌ای و نمایه‌ی فارسی است.

    \item \lr{MatinBook} از ۱۹ ماژول مستقل
    تشکیل شده که در سه فاز بارگذاری
    می‌شوند.

    \item خطاهای رایج شامل استفاده از
    \lr{pdfLaTeX}، فراموش کردن
    \lr{-shell-escape}، فراموش کردن
    \lr{biber} و وارد کردن متن فارسی در
    بلوک کد است.
\end{itemize}

در فصل بعدی (\cref{chap:installation})،
به‌طور گام‌به‌گام یاد می‌گیرید که
\lr{MatinBook} را نصب کنید و اولین کتاب
خود را بسازید.
%%
%% part1-getting-started/ch02-installation.tex
%%
%% Chapter 2: Installation and Setup — How to install XeLaTeX,
%% fonts, Python dependencies, and MatinBook itself.
%%

\chapter{نصب و راه‌اندازی}
\label{chap:installation}

\section{چرا این موضوع مهم است؟}
\label{sec:installation-why}

قبل از اینکه بتوانید اولین کتاب خود را با
\lr{MatinBook} بسازید، باید چند ابزار را روی
سیستم خود نصب کنید. این ابزارها شامل موتور
حروف‌چینی \lr{XeLaTeX}، فونت‌های فارسی، بسته‌های
\lr{LaTeX}، و چند ابزار جانبی مانند \lr{biber} و
\lr{xindy} هستند.

نصب نادرست یا ناقص، شایع‌ترین دلیل خطاهای اولیه
در \lr{MatinBook} است. بنابراین این فصل را
با دقت بخوانید و هر گام را به‌ترتیب انجام دهید.

\mnote{اگر با \lr{LaTeX} و ابزارهای آن آشنایی دارید، می‌توانید مستقیماً به \cref{sec:installation-matinbook} بروید.}

\section{پیش‌نیازها}
\label{sec:installation-prerequisites}

برای کار با \lr{MatinBook}، به ابزارهای زیر
نیاز دارید:

\begin{enumerate}
    \item \textbf{\lr{TeX Live}:} توزیع اصلی
    \lr{LaTeX} که شامل \lr{XeLaTeX} و اکثر بسته‌ها
    است. نسخه‌ی ۲۰۲۳ یا جدیدتر توصیه می‌شود.

    \item \textbf{\lr{XeLaTeX}:} موتور حروف‌چینی
    که از فونت‌های \lr{OpenType} و حروف‌چینی
    راست‌به‌چپ پشتیبانی می‌کند.

    \item \textbf{\lr{biber}:} ابزار پردازش
    کتاب‌نامه برای \lr{biblatex}.

    \item \textbf{\lr{xindy} و \lr{texindy}:}
    ابزار ساخت نمایه با پشتیبانی از مرتب‌سازی
    فارسی.

    \item \textbf{\lr{Python 3} و \lr{Pygments}:}
    برای رنگ‌آمیزی کد در \lr{minted}.

    \item \textbf{فونت‌های فارسی و لاتین:} شامل
    \lr{XB Niloofar}، \lr{XITS Math}، و
    \lr{DejaVu Sans Mono}.
\end{enumerate}

\mnote{اگر از ویندوز استفاده می‌کنید، نصب \lr{TeX Live} ممکن است چند ساعت طول بکشد. صبور باشید!}

\section{نصب \lr{TeX Live}}
\label{sec:installation-texlive}

\lr{TeX Live} توزیع استاندارد \lr{LaTeX} است
که برای \lr{Linux}، \lr{macOS} و \lr{Windows}
در دسترس است.

\subsection{نصب روی \lr{Linux}}
\label{sec:installation-texlive-linux}

روی \lr{Debian} و \lr{Ubuntu}:

\begin{latin}
\begin{minted}{bash}
sudo apt update
sudo apt install texlive-full
\end{minted}
\end{latin}

\mnote{بسته‌ی \lr{texlive-full} تمام بسته‌های \lr{LaTeX} را نصب می‌کند و حدود ۵ گیگابایت فضا می‌گیرد. اگر فضا محدود دارید، از \lr{texlive-xetex} و بسته‌های موردنیاز دیگر استفاده کنید.}

\subsection{نصب روی \lr{macOS}}
\label{sec:installation-texlive-macos}

روش توصیه‌شده، استفاده از \lr{MacTeX} است:

\begin{latin}
\begin{minted}{bash}
# با Homebrew
brew install --cask mactex
\end{minted}
\end{latin}

\subsection{نصب روی \lr{Windows}}
\label{sec:installation-texlive-windows}

از سایت رسمی \lr{TeX Live}، فایل \lr{install-tl-windows.exe}
را دانلود و اجرا کنید.

\section{نصب فونت‌ها}
\label{sec:installation-fonts}

\lr{MatinBook} به چند فونت نیاز دارد که برخی
از آن‌ها ممکن است روی سیستم شما نصب نباشند.

\subsection{فونت \lr{XB Niloofar}}
\label{sec:installation-font-niloofar}

\lr{XB Niloofar} فونت پیش‌فرض متن فارسی در
\lr{MatinBook} است. اگر در مخزن پروژه هستید:

\begin{latin}
\begin{minted}{bash}
mkdir -p ~/.local/share/fonts
cp fonts/niloofar/*.ttf ~/.local/share/fonts/
fc-cache -fv
\end{minted}
\end{latin}

\mnote{در \lr{macOS}، مسیر \lr{\textasciitilde/Library/Fonts} است. در \lr{Windows}، روی فایل \lr{.ttf} راست‌کلیک کرده و \lr{Install} بزنید.}

برای بررسی نصب:

\begin{latin}
\begin{minted}{bash}
fc-list | grep -i niloofar
\end{minted}
\end{latin}

\subsection{فونت \lr{XITS Math}}
\label{sec:installation-font-xits}

فونت ریاضی \lr{XITS Math} برای نمایش فرمول‌های
ریاضی استفاده می‌شود:

\begin{latin}
\begin{minted}{bash}
# Debian/Ubuntu
sudo apt install fonts-xits

# macOS
brew install --cask font-xits
\end{minted}
\end{latin}

\subsection{فونت \lr{DejaVu Sans Mono}}
\label{sec:installation-font-dejavu}

فونت \lr{DejaVu Sans Mono} برای نمایش کد و
متن فارسی داخل کد استفاده می‌شود:

\begin{latin}
\begin{minted}{bash}
# Debian/Ubuntu
sudo apt install fonts-dejavu
\end{minted}
\end{latin}

\mnote{فونت \lr{DejaVu Sans Mono} روی اکثر توزیع‌های \lr{Linux} به‌طور پیش‌فرض نصب است.}

\subsection{فونت‌های اختیاری}
\label{sec:installation-font-optional}

\lr{MatinBook} از چهار فونت فارسی دیگر نیز
پشتیبانی می‌کند که نصب آن‌ها اختیاری است:

\begin{itemize}
    \item \lr{Vazirmatn}
    \item \lr{Sahel}
    \item \lr{IR Lotus}
    \item \lr{B Nazanin} (تجاری)
\end{itemize}

\section{نصب \lr{Python} و \lr{Pygments}}
\label{sec:installation-python}

برای رنگ‌آمیزی کد با \lr{minted}، به \lr{Python 3}
و \lr{Pygments} نیاز دارید:

\begin{latin}
\begin{minted}{bash}
# Debian/Ubuntu
sudo apt install python3 python3-pip
pip3 install pygments

# macOS
brew install python3
pip3 install pygments
\end{minted}
\end{latin}

برای بررسی نصب:

\begin{latin}
\begin{minted}{bash}
pygmentize -V
\end{minted}
\end{latin}

باید نسخه‌ی \lr{Pygments} را نمایش دهد (مثلاً
\lr{2.17.2}).

\mnote{اگر از \lr{TeX Live 2024} یا جدیدتر استفاده می‌کنید، ممکن است به‌جای \lr{pygments} به \lr{latexminted} هم نیاز داشته باشید. با \lr{pip install latexminted} نصب کنید.}

\section{نصب \lr{MatinBook}}
\label{sec:installation-matinbook}

پس از نصب پیش‌نیازها، نوبت به نصب خود
\lr{MatinBook} می‌رسد. سه روش برای این کار
وجود دارد:

\subsection{روش ۱ — نصب در \lr{texmf} (توصیه‌شده)}
\label{sec:installation-texmf}

این روش، استانداردترین و پایدارترین روش است:

\begin{latin}
\begin{minted}{bash}
mkdir -p ~/texmf/tex/latex/matinbook
cp matinbook.cls ~/texmf/tex/latex/matinbook/
find tex -name "*.sty" -exec cp {} ~/texmf/tex/latex/matinbook/ \;
texhash ~/texmf
\end{minted}
\end{latin}

\mnote{مسیر \lr{\textasciitilde/texmf} در \lr{TeX Live} به‌عنوان \lr{TEXMFHOME} شناخته می‌شود و به‌طور خودکار جستجو می‌شود.}

برای بررسی نصب:

\begin{latin}
\begin{minted}{bash}
kpsewhich matinbook.cls
\end{minted}
\end{latin}

باید مسیر فایل را نمایش دهد.

\subsection{روش ۲ — استفاده از \lr{TEXINPUTS}}
\label{sec:installation-texinputs}

اگر نمی‌خواهید فایل‌ها را در \lr{texmf} کپی کنید،
می‌توانید مسیر پروژه را به \lr{TEXINPUTS} اضافه کنید:

\begin{latin}
\begin{minted}{bash}
export TEXINPUTS="/path/to/matinbook//:"
\end{minted}
\end{latin}

\mnote{دو اسلش \lr{//} یعنی «همه‌ی زیرپوشه‌ها به‌صورت بازگشتی». این روش برای توسعه‌ی پروژه مفید است ولی برای استفاده‌ی نهایی توصیه نمی‌شود.}

\subsection{روش ۳ — نصب در پروژه‌ی محلی}
\label{sec:installation-local}

اگر فقط می‌خواهید در یک پروژه‌ی خاص از
\lr{MatinBook} استفاده کنید، می‌توانید فایل‌ها
را کنار \file{main.tex} قرار دهید. اما این
روش، پایداری کمتری دارد.

\section{تست نصب}
\label{sec:installation-test}

برای اطمینان از نصب درست، یک فایل ساده بسازید:

\begin{latin}
\begin{minted}{latex}
\documentclass{matinbook}

\title{|\pc{تست نصب}|}
\author{|\pc{من}|}
\date{|\pc{مهر ۱۴۰۵}|}

\begin{document}

\maketitle

\chapter{|\pc{فصل اول}|}

|\pc{این یک متن فارسی ساده است.}|

\begin{theorem}
    |\pc{این یک قضیه‌ی آزمایشی است.}|
\end{theorem}

\end{document}
\end{minted}
\end{latin}

\mnote{نام فایل را \lr{test.tex} بگذارید و با \lr{xelatex -shell-escape test.tex} کامپایل کنید.}

سپس با دستور زیر کامپایل کنید:

\begin{latin}
\begin{minted}{bash}
xelatex -shell-escape test.tex
\end{minted}
\end{latin}

اگر یک فایل \lr{test.pdf} ساخته شد و شامل
متن فارسی و قضیه‌ی رنگی بود، نصب شما موفق
بوده است.

\section{عیب‌یابی نصب}
\label{sec:installation-troubleshooting}

در این بخش، چند خطای رایج نصب و راه‌حل آن‌ها
را بررسی می‌کنیم.

\subsection{خطای \lr{File matinbook.cls not found}}
\label{sec:installation-error-cls}

\textbf{علت:} فایل \file{matinbook.cls} در مسیر
جستجوی \lr{LaTeX} نیست.

\textbf{راه‌حل:}
\begin{itemize}
    \item اگر از روش ۱ استفاده می‌کنید،
    \lr{texhash} را اجرا کنید.
    \item اگر از روش ۲ استفاده می‌کنید،
    \lr{TEXINPUTS} را بررسی کنید.
\end{itemize}

\subsection{خطای \lr{Font XB Niloofar not found}}
\label{sec:installation-error-font}

\textbf{علت:} فونت نصب نشده یا \lr{fc-cache}
اجرا نشده است.

\textbf{راه‌حل:}
\begin{latin}
\begin{minted}{bash}
fc-cache -fv
fc-list | grep -i niloofar
\end{minted}
\end{latin}

\subsection{خطای \lr{Package minted Error}}
\label{sec:installation-error-minted}

\textbf{علت:} فراموش کردن \lr{-shell-escape}
یا نصب نبودن \lr{Pygments}.

\textbf{راه‌حل:}
\begin{itemize}
    \item از \lr{xelatex -shell-escape} استفاده کنید.
    \item \lr{Pygments} را نصب کنید:
    \lr{pip3 install pygments}.
\end{itemize}

\subsection{خطای \lr{xindy: command not found}}
\label{sec:installation-error-xindy}

\textbf{علت:} \lr{xindy} نصب نیست.

\textbf{راه‌حل:}
\begin{latin}
\begin{minted}{bash}
sudo apt install xindy
\end{minted}
\end{latin}

\mnote{خطاهای بیشتر در \cref{app:troubleshooting} بررسی می‌شوند.}

\section{تمرین‌ها}
\label{sec:installation-exercises}

\begin{exercise}
    \lr{TeX Live} را روی سیستم خود نصب کنید
    و با دستور \lr{xelatex -v} نسخه‌ی آن را
    بررسی کنید.
    \label{exc:installation-texlive}
\end{exercise}

\begin{exercise}
    سه فونت اصلی \lr{MatinBook} (\lr{XB Niloofar}،
    \lr{XITS Math}، \lr{DejaVu Sans Mono}) را نصب
    کنید و با \lr{fc-list} نصب آن‌ها را تأیید کنید.
    \label{exc:installation-fonts}
\end{exercise}

\begin{exercise}
    یک فایل \lr{test.tex} ساده بسازید و آن را
    با \lr{xelatex -shell-escape} کامپایل کنید.
    اگر خطایی گرفتید، آن را در \cref{sec:installation-troubleshooting}
    جستجو کنید.
    \label{exc:installation-test}
\end{exercise}

\section{خلاصه}
\label{sec:installation-summary}

در این فصل، یاد گرفتیم که:

\begin{itemize}
    \item پیش‌نیازهای \lr{MatinBook} شامل
    \lr{TeX Live}، \lr{biber}، \lr{xindy}،
    \lr{Python} و فونت‌های فارسی و لاتین است.

    \item \lr{TeX Live} روی سه سیستم‌عامل
    \lr{Linux}، \lr{macOS} و \lr{Windows} نصب
    می‌شود.

    \item فونت‌های موردنیاز شامل
    \lr{XB Niloofar}، \lr{XITS Math} و
    \lr{DejaVu Sans Mono} هستند.

    \item \lr{MatinBook} به سه روش نصب می‌شود:
    \lr{texmf} (توصیه‌شده)، \lr{TEXINPUTS} و
    پروژه‌ی محلی.

    \item برای تست نصب، یک فایل ساده بسازید
    و با \lr{xelatex -shell-escape} کامپایل کنید.

    \item خطاهای رایج نصب شامل «فایل پیدا نشد»،
    «فونت پیدا نشد»، «خطای \lr{minted}» و
    «\lr{xindy} پیدا نشد» است.
\end{itemize}

در فصل بعدی (\cref{chap:first-book})، اولین
کتاب خود را با \lr{MatinBook} می‌سازید و
با ساختار پایه‌ی کتاب آشنا می‌شوید.
%%
%% part1-getting-started/ch03-first-book.tex
%%
%% Chapter 3: Your First Book — Build a minimal but complete
%% book with MatinBook, from cover to back cover.
%%

\chapter{اولین کتاب}
\label{chap:first-book}

\section{چرا این موضوع مهم است؟}
\label{sec:first-book-why}

پس از نصب موفق \lr{MatinBook}، وقت آن رسیده که
اولین کتاب خود را بسازید. هدف این فصل، نه آموزش
تمام قابلیت‌های \lr{MatinBook}، بلکه آشنایی با
\textbf{ساختار پایه‌ی یک کتاب} است. پس از خواندن
این فصل، شما یک کتاب کامل با جلد، پیشگفتار، فصل،
منابع و نمایه خواهید داشت.

این کتاب کوچک، پایه‌ای خواهد بود که در فصل‌های
بعدی، با افزودن قابلیت‌های پیشرفته‌تر، آن را
گسترش خواهید داد.

\mnote{برای شروع، نیازی به دانش عمیق \lr{LaTeX} نیست. کافی است دستورات این فصل را گام‌به‌گام دنبال کنید.}

\section{ساختار یک کتاب کامل}
\label{sec:first-book-structure}

یک کتاب حرفه‌ای در \lr{MatinBook} از پنج بخش
اصلی تشکیل شده است:

\begin{enumerate}
    \item \textbf{جلد جلو:} با دستور \cmd{makecover}.
    \item \textbf{پیش‌متن:} شامل شناسنامه، تقدیم،
    پیشگفتار، فهرست مطالب.
    \item \textbf{متن اصلی:} شامل فصل‌ها.
    \item \textbf{پس‌متن:} شامل منابع و نمایه.
    \item \textbf{جلد عقب:} با دستور \cmd{makebackcover}.
\end{enumerate}

\mnote{ترتیب این پنج بخش مهم است. اگر جابه‌جا شوند، کتاب به‌درستی کامپایل نمی‌شود.}

\section{ساخت فایل اصلی}
\label{sec:first-book-main}

یک فایل جدید به نام \file{my-first-book.tex} بسازید
و محتوای زیر را در آن قرار دهید:

\begin{latin}
\begin{minted}{latex}
\documentclass{matinbook}

\title{|\pc{اولین کتاب من}|}
\author{|\pc{نام من}|}
\date{|\pc{مهر ۱۴۰۵}|}

\booktitle{|\pc{اولین کتاب من}|}

\begin{document}

% |\pc{جلد جلو}|
\makecover
    {|\pc{اولین کتاب من}|}
    {|\pc{یک کتاب آزمایشی}|}
    {|\pc{نام من}|}
    {|\pc{مهر ۱۴۰۵}|}

% |\pc{پیش‌متن}|
\frontmatter
\frontmattergeometry

\maketitle
\tableofcontents

% |\pc{متن اصلی}|
\mainmatter
\mainmattergeometry

\chapter{|\pc{فصل اول}|}
\label{chap:first}

|\pc{این یک متن ساده است.}|

% |\pc{پس‌متن}|
\backmatter
\frontmattergeometry

% |\pc{جلد عقب}|
\makebackcover{|\pc{توضیحات پشت جلد...}|}

\end{document}
\end{minted}
\end{latin}

\mnote{نام فایل را می‌توانید هر چه می‌خواهید بگذارید. در این کتاب از \lr{my-first-book.tex} استفاده می‌کنیم.}

\section{کامپایل اولین کتاب}
\label{sec:first-book-compile}

برای ساخت فایل \lr{PDF}، از دستور زیر استفاده
کنید:

\begin{latin}
\begin{minted}{bash}
xelatex -shell-escape my-first-book.tex
xelatex -shell-escape my-first-book.tex
\end{minted}
\end{latin}

\mnote{دو بار کامپایل لازم است چون جلد با \lr{TikZ} ساخته می‌شود و به دو مرحله نیاز دارد.}

پس از اجرای این دستورات، یک فایل
\file{my-first-book.pdf} ساخته می‌شود. آن را باز
کنید و نتیجه را ببینید.

\mnote{اگر خطایی گرفتید، به \cref{app:troubleshooting} مراجعه کنید.}

\section{آشنایی با ساختار}
\label{sec:first-book-tour}

حالا که اولین کتاب خود را ساختید، بگذارید
ساختار آن را بخش‌به‌بخش بررسی کنیم.

\subsection{فرمان \cmd{documentclass}}
\label{sec:first-book-documentclass}

اولین خط هر سند \lr{LaTeX}، دستور
\cmd{documentclass} است:

\begin{latin}
\begin{minted}{latex}
\documentclass{matinbook}
\end{minted}
\end{latin}

این دستور به \lr{LaTeX} می‌گوید که از کلاس
\lr{MatinBook} استفاده کند. می‌توانید آپشن‌هایی
مانند \opt{draft} یا \opt{vazirmatn} را هم اضافه
کنید:

\begin{latin}
\begin{minted}{latex}
\documentclass[draft, vazirmatn]{matinbook}
\end{minted}
\end{latin}

\mnote{آپشن‌ها با کاما از هم جدا می‌شوند. ترتیب آن‌ها مهم نیست.}

\subsection{فراداده کتاب}
\label{sec:first-book-metadata}

بعد از \cmd{documentclass}، فراداده‌ی کتاب را
تعریف می‌کنیم:

\begin{latin}
\begin{minted}{latex}
\title{|\pc{اولین کتاب من}|}
\author{|\pc{نام من}|}
\date{|\pc{مهر ۱۴۰۵}|}

\booktitle{|\pc{اولین کتاب من}|}
\end{minted}
\end{latin}

\begin{itemize}
    \item \cmd{title}: عنوان کتاب (برای صفحه‌ی
    عنوان).
    \item \cmd{author}: نام نویسنده.
    \item \cmd{date}: تاریخ.
    \item \cmd{booktitle}: عنوان کتاب برای
    سرصفحه‌ی صفحات زوج.
\end{itemize}

\mnote{فراداده‌ی \cmd{title} و آرگومان اول \cmd{makecover} باید یکسان باشند. برای جلوگیری از ناسازگاری، آن‌ها را هماهنگ نگه دارید.}

\subsection{جلد جلو}
\label{sec:first-book-frontcover}

جلد جلو با دستور \cmd{makecover} ساخته می‌شود:

\begin{latin}
\begin{minted}{latex}
\makecover
    {|\pc{اولین کتاب من}|}
    {|\pc{یک کتاب آزمایشی}|}
    {|\pc{نام من}|}
    {|\pc{مهر ۱۴۰۵}|}
\end{minted}
\end{latin}

این دستور چهار آرگومان می‌گیرد:

\begin{enumerate}
    \item عنوان کتاب.
    \item زیرعنوان.
    \item نام نویسنده.
    \item تاریخ یا نسخه.
\end{enumerate}

\subsection{پیش‌متن}
\label{sec:first-book-frontmatter}

پیش‌متن با دستور \cmd{frontmatter} شروع می‌شود:

\begin{latin}
\begin{minted}{latex}
\frontmatter
\frontmattergeometry
\end{minted}
\end{latin}

دستور \cmd{frontmatter} شماره‌گذاری صفحات را به
حروف ابجد فارسی (الف، ب، ج، ...) تغییر می‌دهد.
دستور \cmd{frontmattergeometry} نیز هندسه‌ی
صفحه را به حالت متقارن تغییر می‌دهد.

\mnote{در پیش‌متن، حاشیه‌ی بیرونی پهن وجود ندارد. این حالت برای فهرست مطالب و پیشگفتار مناسب‌تر است.}

\subsection{صفحه‌ی عنوان}
\label{sec:first-book-titlepage}

صفحه‌ی عنوان با دستور \cmd{maketitle} ساخته می‌شود:

\begin{latin}
\begin{minted}{latex}
\maketitle
\end{minted}
\end{latin}

این دستور از فراداده‌ی \cmd{title}، \cmd{author} و
\cmd{date} استفاده می‌کند.

\subsection{فهرست مطالب}
\label{sec:first-book-toc}

فهرست مطالب با دستور \cmd{tableofcontents} ساخته
می‌شود:

\begin{latin}
\begin{minted}{latex}
\tableofcontents
\end{minted}
\end{latin}

\mnote{برای فهرست تصاویر و جداول، از \lr{\textbackslash listoffigures} و \lr{\textbackslash listoftables} استفاده کنید.}

\subsection{متن اصلی}
\label{sec:first-book-mainmatter}

متن اصلی با دستور \cmd{mainmatter} شروع می‌شود:

\begin{latin}
\begin{minted}{latex}
\mainmatter
\mainmattergeometry
\end{minted}
\end{latin}

دستور \cmd{mainmatter} شماره‌گذاری صفحات را به
اعداد عربی (۱، ۲، ۳، ...) تغییر می‌دهد.
دستور \cmd{mainmattergeometry} نیز هندسه‌ی
صفحه را به حالت حاشیه‌ی پهن برمی‌گرداند.

\subsection{فصل}
\label{sec:first-book-chapter}

فصل با دستور \cmd{chapter} شروع می‌شود:

\begin{latin}
\begin{minted}{latex}
\chapter{|\pc{فصل اول}|}
\label{chap:first}
\end{minted}
\end{latin}

\mnote{دستور \lr{\textbackslash label} برای ارجاع به فصل استفاده می‌شود. همیشه یک برچسب معنادار بدهید.}

\subsection{پس‌متن}
\label{sec:first-book-backmatter}

پس‌متن با دستور \cmd{backmatter} شروع می‌شود:

\begin{latin}
\begin{minted}{latex}
\backmatter
\frontmattergeometry
\end{minted}
\end{latin}

در پس‌متن، معمولاً منابع و نمایه قرار می‌گیرند.

\subsection{جلد عقب}
\label{sec:first-book-backcover}

جلد عقب با دستور \cmd{makebackcover} ساخته
می‌شود:

\begin{latin}
\begin{minted}{latex}
\makebackcover{|\pc{توضیحات پشت جلد...}|}
\end{minted}
\end{latin}

\mnote{دستور \lr{\textbackslash makebackcover} باید حتماً در انتهای سند باشد، پیش از \lr{\textbackslash end\{document\}}.}

\section{افزودن منابع و نمایه}
\label{sec:first-book-bibliography-index}

کتاب شما اکنون کامل است، ولی هنوز منابع و نمایه
ندارد. در این بخش، آن‌ها را اضافه می‌کنیم.

\subsection{افزودن کتاب‌نامه}
\label{sec:first-book-bibliography}

برای افزودن کتاب‌نامه، در بخش پیش از
\cmd{begin\{document\}}، این خط را اضافه کنید:

\begin{latin}
\begin{minted}{latex}
\addbibresource{references.bib}
\end{minted}
\end{latin}

و در پس‌متن، این دستور را اضافه کنید:

\begin{latin}
\begin{minted}{latex}
\printbibliography[title={|\pc{منابع و مراجع}|}]
\end{minted}
\end{latin}

\mnote{کتاب‌نامه را داخل \lr{latin} قرار ندهید. \lr{xepersian} جهت متن را به‌درستی مدیریت می‌کند.}

\subsection{افزودن نمایه}
\label{sec:first-book-index}

برای افزودن نمایه، کافی است یک کلمه را با دستور
\cmd{index} علامت‌گذاری کنید:

\begin{latin}
\begin{minted}{latex}
|\pc{برنامه‌نویسی}|\index{|\pc{برنامه‌نویسی}|}
\end{minted}
\end{latin}

\mnote{دستور \lr{\textbackslash index} متن فارسی را مستقیماً در متن عادی می‌پذیرد و نیازی به \lr{\textbackslash pc} در متن عادی ندارد.}

و در پس‌متن، این دستور را اضافه کنید:

\begin{latin}
\begin{minted}{latex}
\printindex
\end{minted}
\end{latin}

\mnote{نمایه با \lr{xindy} ساخته می‌شود و مرتب‌سازی فارسی را به‌درستی انجام می‌دهد.}

\section{ساختار فایل‌های یک پروژه}
\label{sec:first-book-project-structure}

برای کتاب‌های بزرگ‌تر، بهتر است فایل‌ها را به
چند بخش تقسیم کنید:

\begin{latin}
\begin{minted}{text}
my-book/
├── main.tex
├── references.bib
├── chapters/
│   ├── chapter-01.tex
│   ├── chapter-02.tex
│   └── chapter-03.tex
├── figures/
│   ├── figure-01.pdf
│   └── figure-02.pdf
└── bibliography/
    └── references.bib
\end{minted}
\end{latin}

\mnote{تقسیم فایل‌ها به بخش‌های کوچک، مدیریت پروژه را بسیار آسان‌تر می‌کند. هر فصل در یک فایل جداگانه.}

سپس در \file{main.tex}:

\begin{latin}
\begin{minted}{latex}
\input{chapters/chapter-01}
\input{chapters/chapter-02}
\input{chapters/chapter-03}
\end{minted}
\end{latin}

\mnote{دستور \lr{\textbackslash input} فایل را در همان‌جا وارد می‌کند. برای فایل‌های بزرگ، از \lr{\textbackslash include} استفاده کنید که صفحه‌ی جدید ایجاد می‌کند.}

\section{خطاهای رایج}
\label{sec:first-book-mistakes}

\subsection{فراموش کردن \cmd{frontmatter} یا \cmd{mainmatter}}
\label{sec:first-book-mistake-frontmain}

اگر \cmd{frontmatter} و \cmd{mainmatter} را
فراموش کنید، شماره‌گذاری صفحات به‌درستی انجام
نمی‌شود و ممکن است شماره‌ها به‌صورت عربی از
ابتدای سند شروع شوند.

\textbf{راه‌حل:} ترتیب صحیح را رعایت کنید:
\cmd{frontmatter}، \cmd{mainmatter}، \cmd{backmatter}.

\subsection{قرار دادن \cmd{makebackcover} در ابتدای سند}
\label{sec:first-book-mistake-backcover}

اگر \cmd{makebackcover} را در ابتدای سند قرار
دهید، جلد عقب روی صفحه‌ی دوم ظاهر می‌شود و
بقیه‌ی کتاب به‌هم می‌ریزد.

\textbf{راه‌حل:} همیشه \cmd{makebackcover} را در
انتهای سند، پیش از \cmd{end\{document\}} قرار دهید.

\subsection{فراموش کردن \cmd{addbibresource}}
\label{sec:first-book-mistake-bibresource}

اگر \cmd{addbibresource} را فراموش کنید،
کتاب‌نامه خالی خواهد بود و استنادها به‌صورت
«؟؟» نمایش داده می‌شوند.

\textbf{راه‌حل:} همیشه \cmd{addbibresource} را
در پیش از \cmd{begin\{document\}} قرار دهید.

\section{تمرین‌ها}
\label{sec:first-book-exercises}

\begin{exercise}
    فایل \lr{my-first-book.tex} را بسازید و
    آن را کامپایل کنید. آیا فایل \lr{PDF} بدون
    خطا ساخته می‌شود؟
    \label{exc:first-book-compile}
\end{exercise}

\begin{exercise}
    یک فصل دوم اضافه کنید و در فهرست مطالب
    آن را ببینید. آیا شماره‌گذاری صفحات به‌درستی
    انجام می‌شود؟
    \label{exc:first-book-second-chapter}
\end{exercise}

\begin{exercise}
    یک کتاب‌نامه‌ی کوچک با دو منبع بسازید
    و با \cmd{cite} به آن‌ها استناد کنید.
    \label{exc:first-book-bibliography}
\end{exercise}

\begin{exercise}
    برای کتاب خود یک نمایه با سه مدخل بسازید.
    آیا مرتب‌سازی فارسی به‌درستی انجام می‌شود؟
    \label{exc:first-book-index}
\end{exercise}

\section{خلاصه}
\label{sec:first-book-summary}

در این فصل، اولین کتاب خود را ساختیم:

\begin{itemize}
    \item یک کتاب کامل در \lr{MatinBook} از پنج
    بخش تشکیل شده: جلد جلو، پیش‌متن، متن اصلی،
    پس‌متن، جلد عقب.

    \item ترتیب این بخش‌ها مهم است و باید رعایت
    شود.

    \item کتاب‌نامه با \cmd{addbibresource} و
    \cmd{printbibliography} اضافه می‌شود.

    \item نمایه با \cmd{index} و \cmd{printindex}
    اضافه می‌شود.

    \item برای کتاب‌های بزرگ، بهتر است فایل‌ها
    را به چند بخش تقسیم کنید.

    \item خطاهای رایج شامل فراموش کردن
    \cmd{frontmatter}، \cmd{mainmatter}،
    \cmd{backmatter}، \cmd{makebackcover} در
    انتهای سند، و \cmd{addbibresource} است.
\end{itemize}

در فصل بعدی (\cref{chap:book-structure})،
ساختار کامل یک کتاب را به‌طور مفصل بررسی
می‌کنیم و یاد می‌گیرید که هر بخش را چطور
شخصی‌سازی کنید.
%%
%% part1-getting-started/ch04-book-structure.tex
%%
%% Chapter 4: Book Structure — A detailed look at the
%% frontmatter, mainmatter, and backmatter split, and how
%% each part of a MatinBook project is organized.
%%

\chapter{ساختار یک کتاب}
\label{chap:book-structure}

\section{چرا این موضوع مهم است؟}
\label{sec:book-structure-why}

در فصل قبل، یک کتاب ساده ساختیم و با پنج بخش
اصلی آن آشنا شدیم. در این فصل، نگاهی عمیق‌تر به
ساختار یک کتاب \lr{MatinBook} می‌اندازیم.

درک درست از ساختار کتاب، به شما کمک می‌کند که:

\begin{itemize}
    \item بدانید هر عنصر کتاب کجا قرار می‌گیرد.
    \item بتوانید کتاب‌های بزرگ را به‌راحتی مدیریت کنید.
    \item از خطاهای رایج جلوگیری کنید.
    \item کتاب خود را حرفه‌ای‌تر سازمان‌دهی کنید.
\end{itemize}

\mnote{اگر با ساختار \lr{LaTeX} آشنا هستید، می‌دانید که \lr{\textbackslash frontmatter}، \lr{\textbackslash mainmatter} و \lr{\textbackslash backmatter} از کلاس \lr{book} استاندارد می‌آیند.}

\section{سه بخش اصلی کتاب}
\label{sec:book-structure-three-parts}

هر کتاب \lr{MatinBook} از سه بخش اصلی تشکیل شده
است که با سه دستور مشخص می‌شوند:

\begin{enumerate}
    \item \textbf{پیش‌متن} (\lr{\textbackslash frontmatter}):
    بخش‌هایی که پیش از متن اصلی می‌آیند.
    \item \textbf{متن اصلی} (\lr{\textbackslash mainmatter}):
    فصل‌های اصلی کتاب.
    \item \textbf{پس‌متن} (\lr{\textbackslash backmatter}):
    بخش‌هایی که پس از متن اصلی می‌آیند.
\end{enumerate}

\mnote{هر یک از این سه دستور، شماره‌گذاری صفحات را تغییر می‌دهد: پیش‌متن با حروف ابجد، متن اصلی با اعداد عربی، پس‌متن بدون شماره.}

\section{پیش‌متن}
\label{sec:book-structure-frontmatter}

پیش‌متن، بخش‌هایی است که پیش از فصل‌های اصلی
کتاب قرار می‌گیرند. این بخش‌ها شامل موارد زیر
هستند:

\begin{itemize}
    \item \textbf{صفحه‌ی عنوان:} با \cmd{maketitle}.
    \item \textbf{شناسنامه:} صفحه‌ی حقوق و
    اطلاعات نشر.
    \item \textbf{تقدیم‌نامه:} صفحه‌ی تقدیم.
    \item \textbf{پیشگفتار:} متن مقدمه.
    \item \textbf{فهرست مطالب:} با
    \cmd{tableofcontents}.
    \item \textbf{فهرست تصاویر:} با
    \cmd{listoffigures}.
    \item \textbf{فهرست جداول:} با
    \cmd{listoftables}.
    \item \textbf{فهرست الگوریتم‌ها:} با
    \cmd{listofalgorithms}.
\end{itemize}

\subsection{شماره‌گذاری صفحات در پیش‌متن}
\label{sec:book-structure-frontmatter-numbering}

در پیش‌متن، شماره‌گذاری صفحات با \textbf{حروف
ابجد فارسی} انجام می‌شود: الف، ب، پ، ت، ث، ج،
چ، ح، خ، د، ذ، ر، ز، ژ، س، ش، ص، ض، ط، ظ، ع،
غ، ف، ق، ک، گ، ل، م، ن، و، ه، ی.

این کار، مطابق استاندارد نشر ایران است و باعث
می‌شود که پیش‌متن، مستقل از متن اصلی شماره‌گذاری
شود.

\mnote{اگر پیش‌متن طولانی باشد، ممکن است حروف ابجد تمام شوند. \lr{xepersian} به‌طور خودکار از ترکیب‌های جدید استفاده می‌کند.}

\subsection{هندسه‌ی صفحه در پیش‌متن}
\label{sec:book-structure-frontmatter-geometry}

در پیش‌متن، حاشیه‌ی بیرونی پهن وجود ندارد و
صفحه، حالت \textbf{متقارن} دارد:

\begin{itemize}
    \item حاشیه‌ی داخلی: \pnum{۱.۵} سانتی‌متر.
    \item حاشیه‌ی بیرونی: \pnum{۱.۵} سانتی‌متر.
    \item عرض متن: حدود \pnum{۱۴.۳} سانتی‌متر.
\end{itemize}

برای فعال‌سازی این حالت، از دستور
\cmd{frontmattergeometry} استفاده کنید.

\mnote{دلیل حذف حاشیه‌ی پهن در پیش‌متن این است که این بخش‌ها به یادداشت حاشیه‌ای نیاز ندارند و می‌توانند از عرض کامل متن استفاده کنند.}

\subsection{ساختار یک پیش‌متن نمونه}
\label{sec:book-structure-frontmatter-example}

یک پیش‌متن نمونه در \lr{MatinBook} به این صورت است:

\begin{latin}
\begin{minted}{latex}
\frontmatter
\frontmattergeometry

% |\pc{صفحه‌ی عنوان}|
\maketitle

% |\pc{شناسنامه}|
\thispagestyle{empty}
\vfill
\begin{center}
    |\pc{تمامی حقوق محفوظ است.}|
\end{center}
\clearpage

% |\pc{تقدیم‌نامه}|
\thispagestyle{empty}
\vfill
\begin{flushright}
    |\pc{تقدیم به ...}|
\end{flushright}
\clearpage

% |\pc{پیشگفتار}|
\chapter*{|\pc{پیشگفتار}|}
\addcontentsline{toc}{chapter}{|\pc{پیشگفتار}|}
|\pc{متن پیشگفتار...}|

% |\pc{فهرست‌ها}|
\tableofcontents
\clearpage
\listoffigures
\clearpage
\listoftables
\clearpage
\end{minted}
\end{latin}

\mnote{دستور \lr{\textbackslash chapter*} یک فصل بدون شماره می‌سازد. برای پیشگفتار و پیوست‌های بدون شماره مناسب است.}

\subsection{دستور \cmd{chapter*} در پیش‌متن}
\label{sec:book-structure-chapter-star}

برای بخش‌هایی مانند پیشگفتار که نباید شماره‌ی
فصل داشته باشند، از دستور \cmd{chapter*} استفاده
می‌کنیم. این دستور، یک عنوان بدون شماره می‌سازد.

برای اینکه این بخش در فهرست مطالب هم بیاید،
باید دستور \cmd{addcontentsline} را نیز اضافه
کنیم:

\begin{latin}
\begin{minted}{latex}
\chapter*{|\pc{پیشگفتار}|}
\addcontentsline{toc}{chapter}{|\pc{پیشگفتار}|}
\end{minted}
\end{latin}

\mnote{دستور \lr{\textbackslash addcontentsline} سه آرگومان می‌گیرد: فایل فهرست (\lr{toc})، سطح (\lr{chapter})، و متن.}

\section{متن اصلی}
\label{sec:book-structure-mainmatter}

متن اصلی، بخش اصلی کتاب است و شامل فصل‌های
اصلی می‌شود. این بخش با دستور \cmd{mainmatter}
شروع می‌شود.

\subsection{شماره‌گذاری صفحات در متن اصلی}
\label{sec:book-structure-mainmatter-numbering}

در متن اصلی، شماره‌گذاری صفحات با \textbf{اعداد
عربی} انجام می‌شود: ۱، ۲، ۳، ...

\subsection{هندسه‌ی صفحه در متن اصلی}
\label{sec:book-structure-mainmatter-geometry}

در متن اصلی، حاشیه‌ی بیرونی پهن (به‌طور پیش‌فرض
\pnum{۵.۹} سانتی‌متر) فعال می‌شود:

\begin{itemize}
    \item حاشیه‌ی داخلی: \pnum{۱.۵} سانتی‌متر.
    \item حاشیه‌ی بیرونی: \pnum{۵.۹} سانتی‌متر.
    \item عرض متن: حدود \pnum{۱۰.۳} سانتی‌متر.
\end{itemize}

برای فعال‌سازی این حالت، از دستور
\cmd{mainmattergeometry} استفاده کنید.

\mnote{حاشیه‌ی پهن برای یادداشت‌های حاشیه‌ای، شکل‌های کوچک و فرمول‌های مکمل طراحی شده است.}

\subsection{ساختار یک فصل نمونه}
\label{sec:book-structure-chapter-example}

یک فصل نمونه در \lr{MatinBook} به این صورت است:

\begin{latin}
\begin{minted}{latex}
\chapter{|\pc{فصل اول}|}
\label{chap:first}

\section{|\pc{مقدمه}|}
\label{sec:first-intro}

|\pc{متن مقدمه...}|

\section{|\pc{مفهوم اصلی}|}
\label{sec:first-main}

|\pc{متن بخش اصلی...}|

\begin{theorem}[|\pc{قضیه‌ی نمونه}|]
    |\pc{متن قضیه...}|
    \label{thm:first}
\end{theorem}

\begin{proof}
    |\pc{اثبات قضیه...}|
\end{proof}

\section{|\pc{نتیجه‌گیری}|}
\label{sec:first-conclusion}

|\pc{متن نتیجه‌گیری...}|
\end{minted}
\end{latin}

\mnote{هر فصل باید حداقل یک \lr{\textbackslash label} داشته باشد تا بتوان به آن ارجاع داد.}

\section{پس‌متن}
\label{sec:book-structure-backmatter}

پس‌متن، بخش‌هایی است که پس از فصل‌های اصلی
کتاب قرار می‌گیرند. این بخش‌ها شامل موارد زیر
هستند:

\begin{itemize}
    \item \textbf{منابع و مراجع:} با
    \cmd{printbibliography}.
    \item \textbf{نمایه:} با \cmd{printindex}.
    \item \textbf{پیوست‌ها:} با \cmd{appendix}.
\end{itemize}

\subsection{شماره‌گذاری صفحات در پس‌متن}
\label{sec:book-structure-backmatter-numbering}

در پس‌متن، شماره‌گذاری صفحات همانند متن اصلی
با اعداد عربی ادامه می‌یابد.

\subsection{هندسه‌ی صفحه در پس‌متن}
\label{sec:book-structure-backmatter-geometry}

در پس‌متن، هندسه‌ی صفحه به حالت \textbf{متقارن}
برمی‌گردد (همانند پیش‌متن). برای این کار، از
دستور \cmd{frontmattergeometry} استفاده می‌کنیم.

\mnote{استفاده از \lr{\textbackslash frontmattergeometry} در پس‌متن، کمی گیج‌کننده است. این دستور فقط «هندسه‌ی متقارن» را فعال می‌کند، نه «حالت پیش‌متن» را.}

\subsection{پیوست‌ها}
\label{sec:book-structure-appendices}

پیوست‌ها بخش‌هایی هستند که در انتهای کتاب
می‌آیند و معمولاً شامل مطالب تکمیلی، جدول‌ها،
یا کدهای طولانی هستند. برای شروع پیوست‌ها،
از دستور \cmd{appendix} استفاده می‌کنیم:

\begin{latin}
\begin{minted}{latex}
\appendix

\chapter{|\pc{پیوست الف}|}
\label{app:first}

|\pc{متن پیوست...}|
\end{minted}
\end{latin}

پس از \cmd{appendix}، شماره‌گذاری فصل‌ها به
حروف ابجد لاتین (A, B, C, ...) تغییر می‌کند.

\mnote{پیوست‌ها معمولاً در \lr{\textbackslash mainmatter} قرار می‌گیرند، نه در \lr{\textbackslash backmatter}.}

\subsection{ساختار یک پس‌متن نمونه}
\label{sec:book-structure-backmatter-example}

یک پس‌متن نمونه در \lr{MatinBook} به این صورت
است:

\begin{latin}
\begin{minted}{latex}
% |\pc{پیوست‌ها}| (|\pc{در متن اصلی}|)
\appendix

\chapter{|\pc{پیوست الف}|}
\label{app:a}

|\pc{متن پیوست...}|

% |\pc{پس‌متن}| (|\pc{بدون شماره}|)
\backmatter
\frontmattergeometry

% |\pc{منابع}| (|\pc{بدون}|\lr{latin})
\printbibliography[title={|\pc{منابع و مراجع}|}]

% |\pc{نمایه}|
\printindex
\end{minted}
\end{latin}

\mnote{دستور \lr{\textbackslash printbibliography} را داخل \lr{latin} قرار ندهید. \lr{xepersian} جهت متن را به‌درستی مدیریت می‌کند.}

\section{جدول مقایسه سه بخش}
\label{sec:book-structure-comparison}

جدول \cref{tab:book-structure} مقایسه‌ی خلاصه‌ای
بین سه بخش اصلی کتاب را نشان می‌دهد.

\begin{matintable}{مقایسه‌ی سه بخش اصلی کتاب}{tab:book-structure}
\begin{tblr}{
    width = \textwidth,
    colspec = {X[l] X[c] X[c] X[c]},
    hlines, vlines,
    colsep = 6pt,
    rowsep = 4pt,
    row{1} = {font=\bfseries},
}
ویژگی & پیش‌متن & متن اصلی & پس‌متن \\
دستور شروع & \lr{\textbackslash frontmatter} & \lr{\textbackslash mainmatter} & \lr{\textbackslash backmatter} \\
شماره‌گذاری & حروف ابجد & اعداد عربی & اعداد عربی \\
هندسه & متقارن & حاشیه‌ی پهن & متقارن \\
محتوا & پیشگفتار، فهرست & فصل‌ها & منابع، نمایه \\
شماره‌ی فصل & ندارد & دارد & ندارد \\
\end{tblr}
\end{matintable}

\section{ترتیب دقیق عناصر}
\label{sec:book-structure-order}

ترتیب دقیق عناصر کتاب در \lr{MatinBook} به
این صورت است:

\begin{enumerate}
    \item \textbf{جلد جلو:} \cmd{makecover}.
    \item \textbf{پیش‌متن:} \cmd{frontmatter} +
    \cmd{frontmattergeometry}.
    \begin{enumerate}
        \item صفحه‌ی عنوان: \cmd{maketitle}.
        \item شناسنامه.
        \item تقدیم‌نامه.
        \item پیشگفتار: \cmd{chapter*} +
        \cmd{addcontentsline}.
        \item فهرست مطالب: \cmd{tableofcontents}.
        \item فهرست تصاویر: \cmd{listoffigures}.
        \item فهرست جداول: \cmd{listoftables}.
        \item فهرست الگوریتم‌ها: \cmd{listofalgorithms}.
    \end{enumerate}
    \item \textbf{متن اصلی:} \cmd{mainmatter} +
    \cmd{mainmattergeometry}.
    \begin{enumerate}
        \item بخش‌ها: \cmd{part}.
        \item فصل‌ها: \cmd{chapter}.
        \item پیوست‌ها: \cmd{appendix} +
        \cmd{chapter}.
    \end{enumerate}
    \item \textbf{پس‌متن:} \cmd{backmatter} +
    \cmd{frontmattergeometry}.
    \begin{enumerate}
        \item منابع: \cmd{printbibliography}.
        \item نمایه: \cmd{printindex}.
    \end{enumerate}
    \item \textbf{جلد عقب:} \cmd{makebackcover}.
\end{enumerate}

\mnote{ترتیب دقیق این عناصر، یکی از مهم‌ترین عوامل در تولید کتاب حرفه‌ای است.}

\section{خطاهای رایج}
\label{sec:book-structure-mistakes}

\subsection{فراموش کردن \cmd{frontmattergeometry}}
\label{sec:book-structure-mistake-frontgeo}

اگر \cmd{frontmattergeometry} را فراموش کنید،
پیش‌متن با حاشیه‌ی پهن نمایش داده می‌شود که
برای فهرست مطالب و پیشگفتار مناسب نیست.

\textbf{راه‌حل:} همیشه \cmd{frontmattergeometry} را
پس از \cmd{frontmatter} قرار دهید.

\subsection{فراموش کردن \cmd{mainmattergeometry}}
\label{sec:book-structure-mistake-maingeo}

اگر \cmd{mainmattergeometry} را فراموش کنید،
متن اصلی با حاشیه‌ی متقارن (بدون حاشیه‌ی پهن)
نمایش داده می‌شود که برای یادداشت‌های حاشیه‌ای
مشکل ایجاد می‌کند.

\textbf{راه‌حل:} همیشه \cmd{mainmattergeometry} را
پس از \cmd{mainmatter} قرار دهید.

\subsection{قرار دادن پیوست‌ها در پس‌متن}
\label{sec:book-structure-mistake-appendix}

پیوست‌ها باید در متن اصلی قرار گیرند، نه در
پس‌متن. اگر پیوست‌ها را در پس‌متن قرار دهید،
شماره‌گذاری آن‌ها به‌درستی انجام نمی‌شود.

\textbf{راه‌حل:} پیوست‌ها را پس از فصل‌های اصلی
و پیش از \cmd{backmatter} قرار دهید.

\subsection{فراموش کردن \cmd{addcontentsline} برای \cmd{chapter*}}
\label{sec:book-structure-mistake-addcontentsline}

اگر از \cmd{chapter*} برای پیشگفتار استفاده
کنید و \cmd{addcontentsline} را فراموش کنید،
پیشگفتار در فهرست مطالب ظاهر نمی‌شود.

\textbf{راه‌حل:} همیشه \cmd{addcontentsline} را
پس از \cmd{chapter*} قرار دهید.

\section{تمرین‌ها}
\label{sec:book-structure-exercises}

\begin{exercise}
    یک کتاب با پیش‌متن کامل (شناسنامه، تقدیم،
    پیشگفتار، فهرست مطالب) بسازید. آیا
    شماره‌گذاری صفحات با حروف ابجد انجام
    می‌شود؟
    \label{exc:book-structure-frontmatter}
\end{exercise}

\begin{exercise}
    یک کتاب با دو بخش و چهار فصل بسازید.
    از دستور \cmd{part} برای ایجاد بخش‌ها
    استفاده کنید.
    \label{exc:book-structure-parts}
\end{exercise}

\begin{exercise}
    یک پیوست به کتاب خود اضافه کنید و
    شماره‌گذاری آن را بررسی کنید.
    \label{exc:book-structure-appendix}
\end{exercise}

\begin{exercise}
    جدول \cref{tab:book-structure} را با
    کتاب خود مقایسه کنید. آیا همه‌ی
    ویژگی‌ها را رعایت کرده‌اید؟
    \label{exc:book-structure-compare}
\end{exercise}

\section{خلاصه}
\label{sec:book-structure-summary}

در این فصل، ساختار یک کتاب \lr{MatinBook} را
به‌طور مفصل بررسی کردیم:

\begin{itemize}
    \item یک کتاب از سه بخش اصلی تشکیل شده:
    پیش‌متن، متن اصلی، پس‌متن.

    \item پیش‌متن با حروف ابجد فارسی و هندسه‌ی
    متقارن شماره‌گذاری می‌شود.

    \item متن اصلی با اعداد عربی و حاشیه‌ی پهن
    شماره‌گذاری می‌شود.

    \item پس‌متن با اعداد عربی و هندسه‌ی متقارن
    شماره‌گذاری می‌شود.

    \item پیوست‌ها در متن اصلی قرار می‌گیرند،
    نه در پس‌متن.

    \item برای بخش‌های بدون شماره، از
    \cmd{chapter*} همراه با \cmd{addcontentsline}
    استفاده می‌شود.

    \item خطاهای رایج شامل فراموش کردن
    \cmd{frontmattergeometry}، \cmd{mainmattergeometry}،
    قرار دادن پیوست‌ها در پس‌متن، و فراموش
    کردن \cmd{addcontentsline} است.
\end{itemize}

این فصل، آخرین فصل از بخش اول کتاب (شروع به
کار) بود. در بخش دوم (اجزای کتاب)، هر یک از
اجزای کتاب را به‌طور مفصل بررسی خواهیم کرد.

%========================================================================
% PART 2 — BOOK COMPONENTS
%========================================================================
\part{اجزای کتاب}
\label{part:book-components}

%%
%% part2-book-components/ch05-cover.tex
%%
%% Chapter 5: Cover and Back Cover — Full-color front cover
%% and static two-block back cover with author bio.
%%

\chapter{جلد و پشت جلد}
\label{chap:cover}

\section{چرا این موضوع مهم است؟}
\label{sec:cover-why}

جلد کتاب، اولین چیزی است که خواننده می‌بیند و
نقش مهمی در جذب مخاطب دارد. در کتاب‌های فنی و
علمی، جلد حرفه‌ای نشان‌دهنده‌ی کیفیت محتواست و
اعتماد خواننده را جلب می‌کند.

\lr{MatinBook} یک سیستم جلد کامل ارائه می‌دهد که
شامل جلد جلو تمام‌رنگی و جلد عقب با دو بلوک
ایستاتیک است. هدف این فصل، آموزش نحوه‌ی استفاده
و شخصی‌سازی این سیستم است.

\mnote{جلد \lr{MatinBook} با \lr{TikZ} ساخته می‌شود و به حداقل دو بار کامپایل نیاز دارد.}

\section{جلد جلو}
\label{sec:cover-front}

جلد جلو با دستور \cmd{makecover} ساخته می‌شود.
این دستور چهار آرگومان می‌گیرد:

\begin{latin}
\begin{minted}{latex}
\makecover
    {|\pc{عنوان کتاب}|}
    {|\pc{زیرعنوان کتاب}|}
    {|\pc{نام نویسنده}|}
    {|\pc{مهر ۱۴۰۵}|}
\end{minted}
\end{latin}

\begin{enumerate}
    \item \textbf{عنوان کتاب:} عنوان اصلی.
    \item \textbf{زیرعنوان:} توضیح کوتاه یا شعار.
    \item \textbf{نام نویسنده:} نویسنده یا نویسندگان.
    \item \textbf{تاریخ یا نسخه:} تاریخ یا شماره‌ی نسخه.
\end{enumerate}

\mnote{اگر زیرعنوان خالی باشد، از \lr{\{\}} استفاده کنید: \lr{\textbackslash makecover\{عنوان\}\{\}\{نویسنده\}\{تاریخ\}}.}

\subsection{عناصر بصری جلد جلو}
\label{sec:cover-front-elements}

جلد جلو \lr{MatinBook} از عناصر زیر تشکیل شده
است:

\begin{itemize}
    \item \textbf{پس‌زمینه‌ی سرمه‌ای:} رنگ اصلی جلد
    که با \lr{matinnavy} تعریف شده است.
    \item \textbf{گوه‌ی قرمز:} یک مثلث قرمز در
    گوشه‌ی بالا-چپ.
    \item \textbf{بیضی طلایی:} یک بیضی طلایی
    در گوشه‌ی پایین-راست.
    \item \textbf{شبکه‌ی مختصات:} یک شبکه‌ی
    کم‌کنتراست برای بافت.
    \item \textbf{نمونه‌کد C++:} یک نمونه‌ی کد
    با فریم طلایی.
    \item \textbf{معادلات ریاضی:} چهار معادله‌ی
    معروف.
    \item \textbf{فلش‌ها و نقاط:} عناصر تزئینی.
\end{itemize}

\mnote{همه‌ی این عناصر در فایل \lr{mb-theme-cover.sty} تعریف شده‌اند. برای شخصی‌سازی، این فایل را ویرایش کنید.}

\subsection{نمای ساده‌شده‌ی جلد}
\label{sec:cover-front-mockup}

شکل \cref{fig:cover-mockup} نمای ساده‌شده‌ای
از جلد جلو را نشان می‌دهد:

\begin{figure}[h]
\centering
\begin{latin}
\begin{tikzpicture}[scale=0.5]
  % Background
  \fill[matinnavy] (0,0) rectangle (10,14);
  % Crimson wedge
  \fill[matincrimson] (0,14) -- (5.5,14) -- (3,11.5) -- cycle;
  % White frame
  \draw[white, opacity=0.6, line width=0.6pt]
      (0.4,0.4) rectangle (9.6,13.6);
  % Gold ellipse
  \draw[covergold, opacity=0.4, line width=1pt, rotate=-18]
      (7,3) ellipse (3.5 and 1.8);
  % Title area
  \node[white, font=\Large\bfseries] at (7,9) {\rl{عنوان کتاب}};
  \draw[covergold, line width=1pt] (6,8) -- (9,8);
  \node[white, font=\small] at (7.5,7.5) {\rl{زیرعنوان}};
  % Author
  \node[white, font=\bfseries] at (7,4.5) {\rl{نام نویسنده}};
  % MatinBook label
  \node[white, font=\footnotesize, anchor=west] at (1,0.8) {MatinBook};
  \node[white, font=\footnotesize, anchor=east] at (9,0.8) {github.com/...};
\end{tikzpicture}
\end{latin}
\caption{نمای ساده‌شده‌ی جلد جلو}
\label{fig:cover-mockup}
\end{figure}

\mnote{در \lr{tikzpicture}، متن فارسی باید داخل \lr{\textbackslash rl\{...\}} قرار گیرد.}

\subsection{شخصی‌سازی جلد جلو}
\label{sec:cover-front-custom}

برای شخصی‌سازی جلد جلو، سه راه دارید:

\textbf{راه ۱ — تغییر آرگومان‌ها:} فقط محتوای
\cmd{makecover} را تغییر دهید.

\textbf{راه ۲ — تغییر رنگ‌ها:} در
\file{mb-theme-cover.sty}، رنگ‌های
\lr{coverprimary}، \lr{coveraccent}، \lr{covergold}
و \lr{coverpaper} را تغییر دهید.

\textbf{راه ۳ — تغییر طرح:} فایل
\file{mb-theme-cover.sty} را ویرایش کنید و
عناصر جدید اضافه یا حذف کنید.

\mnote{برای تغییرات ساده، راه ۱ و ۲ کافی است. برای طراحی کاملاً جدید، راه ۳ را انتخاب کنید.}

\section{جلد عقب}
\label{sec:cover-back}

جلد عقب با دستور \cmd{makebackcover} ساخته
می‌شود. این دستور یک آرگومان می‌گیرد:

\begin{latin}
\begin{minted}{latex}
\makebackcover{|\pc{توضیحات پشت جلد...}|}
\end{minted}
\end{latin}

\mnote{دستور \lr{\textbackslash makebackcover} باید در انتهای سند، پیش از \lr{\textbackslash end\{document\}} قرار گیرد.}

\subsection{ساختار دو بلوکی}
\label{sec:cover-back-two-blocks}

جلد عقب \lr{MatinBook} از دو بلوک ایستاتیک
تشکیل شده است:

\begin{enumerate}
    \item \textbf{بلوک اول — درباره این کتاب:}
    با عنوان \lr{\textbackslash coverabout} و متن
    توضیحات.
    \item \textbf{بلوک دوم — درباره این نویسنده:}
    با عنوان \lr{\textbackslash coverauthorbio} و
    متن زندگی‌نامه.
\end{enumerate}

فاصله‌ی بین دو بلوک، دقیقاً \pnum{۰.۷}
سانتی‌متر است.

\mnote{هر دو بلوک در یک \lr{minipage} با ارتفاع ثابت قرار می‌گیرند تا متن طولانی باعث سرریز به صفحه‌ی دوم نشود.}

\subsection{افزودن زندگی‌نامه نویسنده}
\label{sec:cover-back-authorbio}

برای افزودن زندگی‌نامه‌ی نویسنده، از دستور
\cmd{authorbio} استفاده کنید:

\begin{latin}
\begin{minted}{latex}
\authorbio{%
    |\pc{نام نویسنده، پژوهشگر حوزه‌ی الگوریتم و ریاضیات است.}|%
    |\pc{او سال‌ها در زمینه‌ی آموزش برنامه‌نویسی فعالیت کرده است.}|%
}

\makebackcover{%
    |\pc{توضیحات پشت جلد...}|%
}
\end{minted}
\end{latin}

\mnote{اگر \lr{\textbackslash authorbio} را فراخوانی نکنید، بلوک دوم به‌کلی حذف می‌شود و فقط بلوک کتاب نمایش داده می‌شود.}

\subsection{شخصی‌سازی جلد عقب}
\label{sec:cover-back-custom}

برای شخصی‌سازی جلد عقب:

\begin{itemize}
    \item \textbf{تغییر عنوان‌ها:} با
    \cmd{renewcommand} می‌توانید
    \lr{\textbackslash coverabout} و
    \lr{\textbackslash coverauthorbio} را تغییر دهید.
    \item \textbf{تغییر رنگ‌ها:} در
    \file{mb-theme-cover.sty}، رنگ‌های
    \lr{coveraccent} و \lr{coverprimary} را
    تغییر دهید.
    \item \textbf{تغییر متن:} فقط آرگومان‌های
    \cmd{makebackcover} و \cmd{authorbio} را
    تغییر دهید.
\end{itemize}

\mnote{عنوان‌های \lr{\textbackslash coverabout} و \lr{\textbackslash coverauthorbio} در فایل \lr{fa-IR.sty} تعریف شده‌اند.}

\section{اطلاعات مخزن}
\label{sec:cover-repository}

جلد \lr{MatinBook} آدرس مخزن پروژه را نمایش
می‌دهد. برای تغییر آن:

\begin{latin}
\begin{minted}{latex}
\renewcommand{\repository}{github.com/username/repo}
\end{minted}
\end{latin}

\mnote{مقدار پیش‌فرض \lr{github.com/matinmahmoudi-cs/matinbook} است.}

\section{اطلاعات نسخه}
\label{sec:cover-version}

جلد \lr{MatinBook} نسخه‌ی کلاس را نمایش می‌دهد.
این مقدار در \file{matinbook.cls} تعریف شده است:

\begin{latin}
\begin{minted}{latex}
\def\matinbookversion{1.2}
\def\matinbookdate{2026/10/08}
\end{minted}
\end{latin}

\mnote{نسخه‌ی \lr{MatinBook} در صفحه‌ی شناسنامه و روی جلد نمایش داده می‌شود.}

\section{رنگ‌های جلد}
\label{sec:cover-colors}

جلد \lr{MatinBook} از رنگ‌های زیر استفاده می‌کند
که در \file{mb-theme-cover.sty} تعریف شده‌اند:

\begin{matintable}{رنگ‌های جلد \lr{MatinBook}}{tab:cover-colors}
\begin{tblr}{
    width = \textwidth,
    colspec = {l X[l] X[c]},
    hlines, vlines,
    colsep = 8pt,
    rowsep = 4pt,
    row{1} = {font=\bfseries},
}
نام & نقش & مقدار \lr{RGB} \\
\lr{coverprimary} & پس‌زمینه‌ی جلد جلو & \lr{(30, 42, 58)} \\
\lr{coveraccent} & گوه و عناصر تزئینی & \lr{(166, 25, 46)} \\
\lr{covergold} & بیضی و متن طلایی & \lr{(201, 162, 39)} \\
\lr{coverpaper} & پس‌زمینه‌ی جلد عقب & \lr{(242, 238, 227)} \\
\lr{covergray} & متن‌های ثانویه & \lr{(145, 151, 157)} \\
\end{tblr}
\end{matintable}

\mnote{این رنگ‌ها با پالت \lr{CMYK} کتاب هماهنگ هستند تا جلد با محتوای داخلی یکپارچه باشد.}

\section{کامپایل جلد}
\label{sec:cover-compile}

جلد \lr{MatinBook} با \lr{TikZ} ساخته می‌شود و
از \lr{remember picture, overlay} استفاده می‌کند.
این یعنی حداقل \textbf{دو بار کامپایل} لازم است:

\begin{latin}
\begin{minted}{bash}
# |\pc{مرحله‌ی اول: نوشتن موقعیت}| \lr{TikZ}
xelatex -shell-escape document.tex

# |\pc{مرحله‌ی دوم: استفاده از موقعیت}|
xelatex -shell-escape document.tex
\end{minted}
\end{latin}

\mnote{اگر جلد ناقص یا جابه‌جا ظاهر شد، یک بار دیگر کامپایل کنید.}

\section{خطاهای رایج}
\label{sec:cover-mistakes}

\subsection{قرار دادن \cmd{makebackcover} در ابتدای سند}
\label{sec:cover-mistake-position}

اگر \cmd{makebackcover} را در ابتدای سند قرار
دهید، جلد عقب روی صفحه‌ی دوم ظاهر می‌شود و
بقیه‌ی کتاب به‌هم می‌ریزد.

\textbf{راه‌حل:} همیشه \cmd{makebackcover} را در
انتهای سند قرار دهید.

\subsection{فراموش کردن \cmd{authorbio}}
\label{sec:cover-mistake-authorbio}

اگر \cmd{authorbio} را فراموش کنید، بلوک «درباره
این نویسنده» حذف می‌شود و جلد عقب فقط یک بلوک
خواهد داشت.

\textbf{راه‌حل:} اگر می‌خواهید زندگی‌نامه‌ی نویسنده
نمایش داده شود، \cmd{authorbio} را پیش از
\cmd{makebackcover} قرار دهید.

\subsection{سرریز متن جلد عقب به صفحه‌ی دوم}
\label{sec:cover-mistake-overflow}

اگر متن جلد عقب بسیار طولانی باشد، ممکن است
به صفحه‌ی دوم سرریز کند. \lr{MatinBook} این
مشکل را با \lr{minipage} با ارتفاع ثابت حل
کرده است، ولی اگر متن بسیار طولانی باشد،
ممکن است همچنان سرریز رخ دهد.

\textbf{راه‌حل:} متن را کوتاه‌تر کنید یا از
\cmd{small} برای کاهش اندازه‌ی فونت استفاده کنید.

\section{تمرین‌ها}
\label{sec:cover-exercises}

\begin{exercise}
    یک جلد جلو با عنوان، زیرعنوان، نام نویسنده
    و تاریخ بسازید. آیا همه‌ی عناصر به‌درستی
    نمایش داده می‌شوند؟
    \label{exc:cover-front}
\end{exercise}

\begin{exercise}
    یک جلد عقب با دو بلوک بسازید و از
    \cmd{authorbio} استفاده کنید.
    \label{exc:cover-back}
\end{exercise}

\begin{exercise}
    رنگ‌های جلد را تغییر دهید و کتاب را
    دوباره کامپایل کنید. آیا رنگ‌ها اعمال
    می‌شوند؟
    \label{exc:cover-colors}
\end{exercise}

\begin{exercise}
    یک جلد عقب با متن بسیار طولانی بسازید.
    آیا به صفحه‌ی دوم سرریز می‌کند؟
    \label{exc:cover-overflow}
\end{exercise}

\section{خلاصه}
\label{sec:cover-summary}

در این فصل، سیستم جلد \lr{MatinBook} را بررسی
کردیم:

\begin{itemize}
    \item جلد جلو با \cmd{makecover} و چهار
    آرگومان ساخته می‌شود.

    \item جلد جلو شامل پس‌زمینه‌ی سرمه‌ای، گوه‌ی
    قرمز، بیضی طلایی و عناصر تزئینی است.

    \item جلد عقب با \cmd{makebackcover} و یک
    آرگومان ساخته می‌شود.

    \item جلد عقب دو بلوک ایستاتیک دارد:
    «درباره این کتاب» و «درباره این نویسنده».

    \item زندگی‌نامه‌ی نویسنده با \cmd{authorbio}
    اضافه می‌شود.

    \item جلد با \lr{TikZ} ساخته می‌شود و به
    حداقل دو بار کامپایل نیاز دارد.

    \item خطاهای رایج شامل قرار دادن
    \cmd{makebackcover} در ابتدای سند، فراموش
    کردن \cmd{authorbio} و سرریز متن است.
\end{itemize}

در فصل بعدی (\cref{chap:chapters})، با فصل‌ها
و بخش‌ها آشنا می‌شوید و یاد می‌گیرید که چگونه
ساختار کتاب خود را سازمان‌دهی کنید.
%%
%% part2-book-components/ch06-chapters.tex
%%
%% Chapter 6: Chapters and Sections — How to organize a book
%% into parts, chapters, sections, subsections, and
%% subsubsections, and how to cross-reference them.
%%

\chapter{فصل و بخش}
\label{chap:chapters}

\section{چرا این موضوع مهم است؟}
\label{sec:chapters-why}

یک کتاب خوب، فقط مجموعه‌ای از متن نیست؛ بلکه
ساختاری منظم دارد که خواننده را در مسیر یادگیری
هدایت می‌کند. فصل‌ها، بخش‌ها و زیربخش‌ها،
ستون فقرات این ساختار هستند.

\lr{MatinBook} از یک سیستم عنوان‌بندی پنج‌سطحی
استفاده می‌کند که بر اساس مقیاس \lr{Boyer/Stewart}
طراحی شده است. در این فصل، با این سیستم آشنا
می‌شوید و یاد می‌گیرید که چگونه ساختار کتاب
خود را سازمان‌دهی کنید.

\mnote{مقیاس \lr{Boyer/Stewart} استاندارد کتاب‌های درسی دانشگاهی است و در انتشارات معتبری مانند \lr{Wiley} و \lr{Pearson} استفاده می‌شود.}

\section{سطوح عنوان‌بندی}
\label{sec:chapters-levels}

\lr{MatinBook} از پنج سطح عنوان‌بندی پشتیبانی
می‌کند:

\begin{enumerate}
    \item \textbf{بخش} (\cmd{part}): بزرگ‌ترین
    واحد ساختاری.
    \item \textbf{فصل} (\cmd{chapter}): واحد
    اصلی کتاب.
    \item \textbf{بخش} (\cmd{section}): زیرمجموعه‌ی
    فصل.
    \item \textbf{زیربخش} (\cmd{subsection}):
    زیرمجموعه‌ی بخش.
    \item \textbf{زیرزیربخش} (\cmd{subsubsection}):
    عمیق‌ترین سطح.
\end{enumerate}

\mnote{توجه کنید که «بخش» در فارسی برای هر دو \lr{part} و \lr{section} استفاده می‌شود. در این کتاب، برای \lr{part} از «بخش اصلی» و برای \lr{section} از «بخش» استفاده می‌کنیم.}

\section{بخش اصلی (\lr{part})}
\label{sec:chapters-part}

بخش اصلی، بزرگ‌ترین واحد ساختاری در یک کتاب
است. یک کتاب می‌تواند چند بخش اصلی داشته باشد
که هر یک شامل چند فصل است.

برای ایجاد بخش اصلی:

\begin{latin}
\begin{minted}{latex}
\part{|\pc{مبانی}|}
\label{part:fundamentals}
\end{minted}
\end{latin}

\mnote{بخش اصلی در فهرست مطالب ظاهر می‌شود ولی شماره‌ی فصل را تغییر نمی‌دهد.}

\subsection{شماره‌گذاری بخش اصلی}
\label{sec:chapters-part-numbering}

بخش‌های اصلی با اعداد رومی شماره‌گذاری می‌شوند:
بخش اول، بخش دوم، بخش سوم و...

\subsection{نمایش بخش اصلی}
\label{sec:chapters-part-display}

هر بخش اصلی، یک صفحه‌ی جداگانه دارد که فقط
شامل عنوان بخش است. این صفحه، بدون شماره
است و به عنوان یک جداکننده بین بخش‌ها عمل
می‌کند.

\mnote{اگر کتاب شما کوتاه است (کمتر از ۸ فصل)، نیازی به بخش اصلی ندارید.}

\section{فصل (\lr{chapter})}
\label{sec:chapters-chapter}

فصل، واحد اصلی کتاب است. هر فصل، یک موضوع
مشخص را پوشش می‌دهد و معمولاً ۱۵ تا ۲۵ صفحه
است.

برای ایجاد فصل:

\begin{latin}
\begin{minted}{latex}
\chapter{|\pc{مقدمه}|}
\label{chap:introduction}
\end{minted}
\end{latin}

\subsection{شماره‌گذاری فصل}
\label{sec:chapters-chapter-numbering}

فصل‌ها با اعداد عربی شماره‌گذاری می‌شوند: ۱، ۲،
۳ و... این شماره‌گذاری در \cmd{mainmatter}
شروع می‌شود و در \cmd{backmatter} ادامه می‌یابد.

\subsection{سبک نمایش فصل}
\label{sec:chapters-chapter-style}

فصل‌ها در \lr{MatinBook} با یک سبک خاص نمایش
داده می‌شوند:

\begin{itemize}
    \item شماره‌ی فصل با فونت \pnum{۲۴pt} و
    رنگ کم‌رنگ.
    \item عنوان فصل با فونت \pnum{۲۰pt}.
    \item یک خط افقی رنگی زیر عنوان.
\end{itemize}

\mnote{سبک فصل در فایل \lr{mb-headings.sty} تعریف شده است. برای شخصی‌سازی، این فایل را ویرایش کنید.}

\subsection{فصل‌های بدون شماره}
\label{sec:chapters-chapter-star}

برای فصل‌های بدون شماره، از \cmd{chapter*}
استفاده می‌کنیم. این دستور برای پیشگفتار،
پیوست‌های بدون شماره و بخش‌های ویژه مناسب
است:

\begin{latin}
\begin{minted}{latex}
\chapter*{|\pc{پیشگفتار}|}
\addcontentsline{toc}{chapter}{|\pc{پیشگفتار}|}
\end{minted}
\end{latin}

\mnote{دستور \lr{\textbackslash addcontentsline} باعث می‌شود که این فصل در فهرست مطالب ظاهر شود.}

\section{بخش (\lr{section})}
\label{sec:chapters-section}

بخش، زیرمجموعه‌ی فصل است. هر فصل معمولاً
شامل ۳ تا ۵ بخش است. برای ایجاد بخش:

\begin{latin}
\begin{minted}{latex}
\section{|\pc{مقدمه}|}
\label{sec:introduction}
\end{minted}
\end{latin}

\subsection{سبک نمایش بخش}
\label{sec:chapters-section-style}

بخش‌ها در \lr{MatinBook} با یک سبک خاص نمایش
داده می‌شوند:

\begin{itemize}
    \item شماره‌ی بخش در یک \lr{colorbox}
    با پس‌زمینه‌ی آبی روشن.
    \item عنوان بخش با فونت \pnum{۱۵pt}.
\end{itemize}

\mnote{شماره‌ی بخش به صورت «فصل.بخش» نمایش داده می‌شود، مثلاً \lr{3.2} یعنی بخش دوم از فصل سوم.}

\section{زیربخش (\lr{subsection})}
\label{sec:chapters-subsection}

زیربخش، زیرمجموعه‌ی بخش است. برای ایجاد
زیربخش:

\begin{latin}
\begin{minted}{latex}
\subsection{|\pc{پیشینه}|}
\label{sec:background}
\end{minted}
\end{latin}

\subsection{سبک نمایش زیربخش}
\label{sec:chapters-subsection-style}

زیربخش‌ها با فونت \pnum{۱۳pt} و شماره‌ی آبی
کم‌رنگ نمایش داده می‌شوند.

\mnote{شماره‌ی زیربخش به صورت «فصل.بخش.زیربخش» نمایش داده می‌شود، مثلاً \lr{3.2.1}.}

\section{زیرزیربخش (\lr{subsubsection})}
\label{sec:chapters-subsubsection}

زیرزیربخش، عمیق‌ترین سطح عنوان‌بندی در
\lr{MatinBook} است. برای ایجاد زیرزیربخش:

\begin{latin}
\begin{minted}{latex}
\subsubsection{|\pc{مثال اول}|}
\label{sec:example-first}
\end{minted}
\end{latin}

\mnote{استفاده‌ی بیش از حد از زیرزیربخش، ساختار کتاب را پیچیده می‌کند. سعی کنید از سه سطح اول (فصل، بخش، زیربخش) استفاده کنید.}

\section{جدول مقایسه‌ی سطوح}
\label{sec:chapters-comparison}

جدول \cref{tab:heading-levels} مقایسه‌ی خلاصه‌ای
بین پنج سطح عنوان‌بندی را نشان می‌دهد.

\begin{matintable}{مقایسه‌ی سطوح عنوان‌بندی}{tab:heading-levels}
\begin{tblr}{
    width = \textwidth,
    colspec = {l X[c] X[c] X[c]},
    hlines, vlines,
    colsep = 6pt,
    rowsep = 4pt,
    row{1} = {font=\bfseries},
}
دستور & اندازه & شماره‌گذاری & در فهرست \\
\lr{\textbackslash part} & صفحه‌ی جداگانه & رومی & بله \\
\lr{\textbackslash chapter} & \pnum{۲۴/۲۰}pt & عربی & بله \\
\lr{\textbackslash section} & \pnum{۱۵}pt & فصل.بخش & بله \\
\lr{\textbackslash subsection} & \pnum{۱۳}pt & فصل.بخش.زیر & بله \\
\lr{\textbackslash subsubsection} & \pnum{۱۱}pt & فصل.بخش.زیر.زیر & بله \\
\end{tblr}
\end{matintable}

\section{ارجاع به فصل و بخش}
\label{sec:chapters-references}

برای ارجاع به فصل و بخش، از دستور \cmd{cref}
استفاده می‌کنیم. این دستور به‌طور خودکار نوع
مرجع را تشخیص می‌دهد و پیشوند مناسب (فصل،
بخش، زیربخش) را اضافه می‌کند:

\begin{latin}
\begin{minted}{latex}
% |\pc{ارجاع به فصل}|
\cref{chap:introduction}

% |\pc{ارجاع به بخش}|
\cref{sec:introduction}

% |\pc{ارجاع به چند مرجع}|
\cref{chap:introduction,chap:chapters}
\end{minted}
\end{latin}

\mnote{خروجی \lr{\textbackslash cref\{chap:introduction\}} می‌شود «فصل ۱» و خروجی \lr{\textbackslash cref\{sec:introduction\}} می‌شود «بخش ۱.۱».}

\subsection{نام‌گذاری برچسب‌ها}
\label{sec:chapters-label-naming}

برای اینکه برچسب‌ها معنادار باشند، از
پیشوندهای زیر استفاده کنید:

\begin{itemize}
    \item \textbf{فصل:} \lr{chap:xxx} —
    مثال: \lr{chap:introduction}.
    \item \textbf{بخش:} \lr{sec:xxx} —
    مثال: \lr{sec:motivation}.
    \item \textbf{زیربخش:} \lr{subsec:xxx} —
    مثال: \lr{subsec:history}.
    \item \textbf{زیرزیربخش:} \lr{subsubsec:xxx}.
\end{itemize}

\mnote{برچسب‌های معنادار، خواندن فایل‌های \lr{tex} را آسان‌تر می‌کنند و از خطاهای ارجاع جلوگیری می‌کنند.}

\section{مثال کامل: یک فصل با همه‌ی سطوح}
\label{sec:chapters-full-example}

در ادامه، یک فصل کامل با همه‌ی سطوح عنوان‌بندی
نمایش داده می‌شود:

\begin{latin}
\begin{minted}{latex}
\chapter{|\pc{مقدمه}|}
\label{chap:introduction}

\section{|\pc{انگیزه}|}
\label{sec:motivation}

|\pc{متن انگیزه...}|

\subsection{|\pc{پیشینه}|}
\label{subsec:background}

|\pc{متن پیشینه...}|

\subsubsection{|\pc{کارهای مرتبط}|}
\label{subsubsec:related}

|\pc{متن کارهای مرتبط...}|

\section{|\pc{تعریف مسئله}|}
\label{sec:problem}

|\pc{متن تعریف مسئله...}|

\section{|\pc{نتیجه‌گیری}|}
\label{sec:conclusion}

|\pc{متن نتیجه‌گیری...}|
\end{minted}
\end{latin}

\mnote{در این مثال، از پنج سطح عنوان‌بندی استفاده شده است. سعی کنید از چهار سطح اول استفاده کنید و از زیرزیربخش فقط در موارد ضروری.}

\section{خطاهای رایج}
\label{sec:chapters-mistakes}

\subsection{عنوان‌بندی بدون متن}
\label{sec:chapters-mistake-empty}

هرگز دو عنوان را پشت سر هم بدون متن قرار
ندهید. این کار، خواننده را گیج می‌کند:

\begin{latin}
\begin{minted}{latex}
% |\pc{اشتباه}|:
\section{|\pc{مقدمه}|}
\subsection{|\pc{پیشینه}|}

% |\pc{درست}|:
\section{|\pc{مقدمه}|}
|\pc{این فصل به بررسی... می‌پردازد.}|

\subsection{|\pc{پیشینه}|}
|\pc{پیشینه پژوهش...}|
\end{minted}
\end{latin}

\subsection{استفاده‌ی بیش از حد از زیرزیربخش}
\label{sec:chapters-mistake-deep}

اگر از زیرزیربخش بیش از حد استفاده کنید،
ساختار کتاب پیچیده می‌شود و خواننده مسیر
خود را گم می‌کند.

\textbf{راه‌حل:} از سه سطح اول (فصل، بخش،
زیربخش) استفاده کنید و فقط در موارد ضروری
به زیرزیربخش بروید.

\subsection{فراموش کردن \cmd{label}}
\label{sec:chapters-mistake-label}

اگر \cmd{label} را فراموش کنید، نمی‌توانید
به آن بخش ارجاع دهید و ممکن است در آینده
مشکل ایجاد شود.

\textbf{راه‌حل:} برای هر فصل و بخش مهم، یک
\cmd{label} معنادار تعریف کنید.

\section{تمرین‌ها}
\label{sec:chapters-exercises}

\begin{exercise}
    یک فصل با سه بخش و هر بخش با دو زیربخش
    بسازید. آیا ساختار در فهرست مطالب به‌درستی
    نمایش داده می‌شود؟
    \label{exc:chapters-structure}
\end{exercise}

\begin{exercise}
    به فصل و بخش‌های خود ارجاع دهید. آیا
    پیشوندها به‌درستی نمایش داده می‌شوند؟
    \label{exc:chapters-refs}
\end{exercise}

\begin{exercise}
    یک کتاب با دو بخش اصلی و چهار فصل بسازید.
    از دستور \cmd{part} برای ایجاد بخش‌های
    اصلی استفاده کنید.
    \label{exc:chapters-parts}
\end{exercise}

\begin{exercise}
    یک \cmd{chapter*} با \cmd{addcontentsline}
    بسازید. آیا در فهرست مطالب ظاهر می‌شود؟
    \label{exc:chapters-star}
\end{exercise}

\section{خلاصه}
\label{sec:chapters-summary}

در این فصل، با سیستم عنوان‌بندی \lr{MatinBook}
آشنا شدیم:

\begin{itemize}
    \item \lr{MatinBook} از پنج سطح عنوان‌بندی
    پشتیبانی می‌کند: \cmd{part}، \cmd{chapter}،
    \cmd{section}، \cmd{subsection}،
    \cmd{subsubsection}.

    \item مقیاس اندازه‌ها بر اساس استاندارد
    \lr{Boyer/Stewart} است: ۲۴/۲۰، ۱۵، ۱۳، ۱۱pt.

    \item برای فصل‌های بدون شماره، از
    \cmd{chapter*} با \cmd{addcontentsline}
    استفاده می‌شود.

    \item ارجاع به فصل و بخش با \cmd{cref}
    انجام می‌شود که پیشوند مناسب را به‌طور
    خودکار اضافه می‌کند.

    \item برچسب‌ها باید معنادار باشند و از
    پیشوندهای استاندارد (\lr{chap:}،
    \lr{sec:}، \lr{subsec:}) استفاده کنند.

    \item خطاهای رایج شامل عنوان‌بندی بدون
    متن، استفاده‌ی بیش از حد از زیرزیربخش،
    و فراموش کردن \cmd{label} است.
\end{itemize}

در فصل بعدی (\cref{chap:text})، با متن و
پاراگراف آشنا می‌شوید و یاد می‌گیرید که
چگونه متن فارسی و لاتین را با هم ترکیب کنید.
%%
%% part2-book-components/ch07-text.tex
%%
%% Chapter 7: Text and Paragraphs — How to write Persian
%% text, mix it with Latin text, and control paragraph
%% formatting.
%%

\chapter{متن و پاراگراف}
\label{chap:text}

\section{چرا این موضوع مهم است؟}
\label{sec:text-why}

متن، بخش اصلی هر کتاب است. در کتاب‌های فنی
فارسی، متن ترکیبی از فارسی و لاتین است و
مدیریت درست این ترکیب، کیفیت نهایی کتاب را
تعیین می‌کند.

\lr{MatinBook} با استفاده از \lr{xepersian} و
تنظیمات دقیق تایپوگرافی، این ترکیب را به‌درستی
مدیریت می‌کند. هدف این فصل، آموزش نحوه‌ی نوشتن
متن فارسی، ترکیب آن با لاتین، و کنترل پاراگراف‌ها
است.

\mnote{تایپوگرافی خوب، خوانایی را افزایش می‌دهد و خستگی چشم خواننده را کاهش می‌دهد.}

\section{متن فارسی}
\label{sec:text-persian}

نوشتن متن فارسی در \lr{MatinBook}، ساده‌ترین
کار ممکن است: فقط متن را به فارسی بنویسید و
\lr{LaTeX} بقیه‌ی کارها را انجام می‌دهد.

\begin{latin}
\begin{minted}{latex}
|\pc{این یک متن فارسی ساده است. لاتک به‌طور خودکار}|%
|\pc{جهت متن را راست‌به‌چپ تشخیص می‌دهد.}|%
\end{minted}
\end{latin}

\mnote{نیازی به هیچ دستور خاصی برای متن فارسی نیست. فقط متن را بنویسید.}

\subsection{فاصله‌گذاری در فارسی}
\label{sec:text-persian-spacing}

در فارسی، فاصله‌ی میان کلمات و نشانه‌های نگارشی،
قواعد خاصی دارد:

\begin{itemize}
    \item \textbf{فاصله‌ی معمولی:} بین کلمات.
    \item \textbf{نیم‌فاصله (\lr{ZWNJ}):} در کلماتی
    مانند «می‌رود»، «کتاب‌ها»، «برنامه‌نویسی».
    \item \textbf{فاصله‌ی نشکن (\lr{\textasciitilde}):}
    برای جلوگیری از شکستن خط.
\end{itemize}

\mnote{نیم‌فاصله را می‌توانید با \lr{Ctrl+Shift+2} در ویندوز یا \lr{Shift+Space} در لینوکس تایپ کنید.}

\subsection{نشانه‌های نگارشی فارسی}
\label{sec:text-persian-punctuation}

\lr{MatinBook} از تمام نشانه‌های نگارشی فارسی
پشتیبانی می‌کند:

\begin{itemize}
    \item \textbf{گیومه:} «نقل قول مستقیم».
    \item \textbf{گیومه‌ی تکی:} ‹نقل قول درونی›.
    \item \textbf{پرانتز:} (متن داخل پرانتز).
    \item \textbf{کروشه:} [متن داخل کروشه].
    \item \textbf{نقطه‌ویرگول:} نشانه‌ی مکث؛ ادامه.
    \item \textbf{ویرگول:} جداسازی، ادامه.
    \item \textbf{علامت سؤال:} چرا؟
    \item \textbf{علامت تعجب:} شگفت‌انگیز!
\end{itemize}

\mnote{در فارسی، نقطه‌ویرگول «؛» و ویرگول «،» با نسخه‌ی لاتین متفاوت هستند. \lr{xepersian} این تفاوت را به‌درستی مدیریت می‌کند.}

\section{ترکیب فارسی و لاتین}
\label{sec:text-latin}

در کتاب‌های فنی، ترکیب متن فارسی با واژگان
لاتین امری اجتناب‌ناپذیر است. \lr{MatinBook}
سه روش برای این ترکیب ارائه می‌دهد:

\begin{enumerate}
    \item \cmd{lr\{...\}}: برای واژگان و
    عبارات کوتاه لاتین در متن فارسی.
    \item \cmd{rl\{...\}}: برای متن فارسی در
    محیط لاتین (مخصوص \lr{tikzpicture}).
    \item \cmd{begin\{latin\}}: برای پاراگراف‌های
    کامل لاتین.
\end{enumerate}

\subsection{واژگان لاتین در متن فارسی}
\label{sec:text-latin-inline}

برای واژگان و عبارات کوتاه لاتین، از دستور
\cmd{lr\{...\}} استفاده کنید:

\begin{latin}
\begin{minted}{latex}
|\pc{این کتاب درباره}|\lr{machine learning}|\pc{و}|\lr{deep learning}|\pc{است.}|%
\end{minted}
\end{latin}

خروجی: این کتاب درباره \lr{machine learning}
و \lr{deep learning} است.

\mnote{فراموش کردن \lr{\textbackslash lr} باعث می‌شود واژه‌ی لاتین با متن فارسی قاطی شود و به‌صورت نامفهوم نمایش داده شود.}

\subsection{متن فارسی در محیط لاتین}
\label{sec:text-persian-in-latin}

برای متن فارسی در محیط لاتین (مخصوص
\lr{tikzpicture})، از دستور \cmd{rl\{...\}}
استفاده کنید:

\begin{latin}
\begin{minted}{latex}
\begin{tikzpicture}
    \node {|\rl{متن فارسی}|};
\end{tikzpicture}
\end{minted}
\end{latin}

\mnote{در \lr{tikzpicture}، جهت متن به‌طور خودکار لاتین است. برای متن فارسی، از \lr{\textbackslash rl} استفاده کنید.}

\subsection{پاراگراف‌های کامل لاتین}
\label{sec:text-latin-block}

برای پاراگراف‌های کامل لاتین (مثلاً یک متن
انگلیسی چند خطی)، از محیط \env{latin} استفاده
کنید:

\begin{latin}
\begin{minted}{latex}
\begin{latin}
This is a complete English paragraph.
It can span multiple lines and will be
typeset left-to-right.
\end{latin}
\end{minted}
\end{latin}

\mnote{محیط \lr{latin} برای کد، الگوریتم، نمودار و هر متن لاتین طولانی مناسب است.}

\subsection{جدول تصمیم‌گیری}
\label{sec:text-decision-table}

جدول \cref{tab:latin-decision} به شما کمک می‌کند
که در هر موقعیت، از کدام دستور استفاده کنید.

\begin{matintable}{انتخاب دستور مناسب برای ترکیب فارسی و لاتین}{tab:latin-decision}
\begin{tblr}{
    width = \textwidth,
    colspec = {X[l] X[c]},
    hlines, vlines,
    colsep = 8pt,
    rowsep = 4pt,
    row{1} = {font=\bfseries},
}
موقعیت & دستور \\
واژه‌ی لاتین در متن فارسی & \lr{\textbackslash lr\{...\}} \\
عبارت کوتاه لاتین & \lr{\textbackslash lr\{...\}} \\
پاراگراف کامل لاتین & \lr{\textbackslash begin\{latin\}} \\
متن فارسی در \lr{tikzpicture} & \lr{\textbackslash rl\{...\}} \\
کد برنامه‌نویسی & \lr{\textbackslash begin\{latin\}} + \lr{minted} \\
الگوریتم & \lr{\textbackslash begin\{latin\}} + \lr{algorithm} \\
فرمول ریاضی & \lr{\textbackslash lr\{...\}} در متن فارسی \\
\lr{URL} & \lr{\textbackslash lr\{...\}} \\
\end{tblr}
\end{matintable}

\mnote{اگر شک دارید، از \lr{\textbackslash lr} استفاده کنید. این دستور در ۹۰٪ موارد کار می‌کند.}

\section{پاراگراف‌ها}
\label{sec:text-paragraphs}

پاراگراف، واحد اصلی متن است. \lr{MatinBook}
تنظیمات زیر را برای پاراگراف‌ها اعمال می‌کند:

\begin{itemize}
    \item \textbf{تورفتگی:} \pnum{۱em} برای
    پاراگراف‌های دوم به بعد.
    \item \textbf{فاصله‌ی بین پاراگراف‌ها:}
    \pnum{۰pt} (بدون فاصله‌ی اضافی).
    \item \textbf{فاصله‌ی خطوط:} \pnum{۱.۱۵}.
\end{itemize}

\mnote{پاراگراف اول هر بخش، تورفتگی ندارد. این یک استاندارد در حروف‌چینی فارسی است.}

\subsection{شکستن پاراگراف}
\label{sec:text-paragraph-break}

برای شکستن پاراگراف، یک خط خالی بین دو متن
قرار دهید:

\begin{latin}
\begin{minted}{latex}
|\pc{این پاراگراف اول است.}|%

|\pc{این پاراگراف دوم است.}|%
\end{minted}
\end{latin}

\mnote{هرگز از \lr{\textbackslash\textbackslash} برای شکستن پاراگراف استفاده نکنید. این دستور برای شکستن خط در یک پاراگراف است، نه برای ایجاد پاراگراف جدید.}

\subsection{طول پاراگراف}
\label{sec:text-paragraph-length}

طبق اصول نگارش حرفه‌ای:

\begin{itemize}
    \item \textbf{طول ایده‌آل:} ۳ تا ۷ جمله.
    \item \textbf{حداکثر:} ۱۰ جمله.
    \item \textbf{حداقل:} ۲ جمله.
\end{itemize}

\mnote{اگر پاراگراف بیش از ۱۰ جمله است، آن را به دو پاراگراف تقسیم کنید. اگر فقط یک جمله است، آن را با پاراگراف بعدی ادغام کنید.}

\section{متن تأکیدی}
\label{sec:text-emphasis}

\lr{MatinBook} از دستورات استاندارد \lr{LaTeX}
برای تأکید پشتیبانی می‌کند. جدول
\cref{tab:emphasis-commands} این دستورات را
خلاصه می‌کند:

\begin{matintable}{دستورات تأکید در متن}{tab:emphasis-commands}
\begin{tblr}{
    width = 0.7\textwidth,
    colspec = {l X[c]},
    hlines, vlines,
    colsep = 8pt,
    rowsep = 4pt,
    row{1} = {font=\bfseries},
}
نوع تأکید & دستور \\
پررنگ & \lr{\textbackslash textbf\{...\}} \\
مورب & \lr{\textbackslash textit\{...\}} \\
زیرخط‌دار & \lr{\textbackslash underline\{...\}} \\
مونواسپیس & \lr{\textbackslash texttt\{...\}} \\
\end{tblr}
\end{matintable}

در ادامه، هر یک را با یک مثال نشان می‌دهیم.

\subsection{متن پررنگ}

برای پررنگ کردن متن، از دستور \cmd{textbf}
استفاده کنید. مثال:

\begin{quote}
    این متن \textbf{پررنگ} است.
\end{quote}

\mnote{در متن فارسی، \lr{\textbackslash textbf} به‌درستی کار می‌کند و متن را پررنگ نمایش می‌دهد.}

\subsection{متن مورب}

برای مورب کردن متن، از دستور \cmd{textit}
استفاده کنید. مثال:

\begin{quote}
    این متن \textit{مورب} است.
\end{quote}

\mnote{متن مورب در فارسی کمتر استفاده می‌شود، چون خوانایی را کاهش می‌دهد. برای تأکید، از پررنگ استفاده کنید.}

\subsection{متن زیرخط‌دار}

برای زیرخط‌دار کردن متن، از دستور \cmd{underline}
استفاده کنید. مثال:

\begin{quote}
    این متن \underline{زیرخط‌دار} است.
\end{quote}

\mnote{استفاده از زیرخط در متن رسمی توصیه نمی‌شود، چون با لینک‌های \lr{hyperref} اشتباه می‌شود.}

\subsection{متن مونواسپیس}

برای نمایش متن با فونت مونواسپیس، از دستور
\cmd{texttt} استفاده کنید. مثال:

\begin{quote}
    این متن \texttt{monospace} است.
\end{quote}

\mnote{برای متن مونواسپیس فارسی، از \lr{\textbackslash texttt} استفاده نکنید، چون فونت مونواسپیس ممکن است فارسی را پشتیبانی نکند.}

\section{فهرست‌ها}
\label{sec:text-lists}

\lr{MatinBook} از سه نوع فهرست پشتیبانی می‌کند.
هر یک را با یک مثال نشان می‌دهیم.

\subsection{فهرست نشانه‌دار}

فهرست نشانه‌دار، پرکاربردترین نوع فهرست است.
مثال:

\begin{quote}
    \begin{itemize}
        \item مورد اول.
        \item مورد دوم.
        \item مورد سوم.
    \end{itemize}
\end{quote}

\mnote{از فهرست نشانه‌دار برای موارد بدون ترتیب مشخص استفاده کنید.}

\subsection{فهرست شماره‌دار}

فهرست شماره‌دار، برای مواردی که ترتیب مهم
است، استفاده می‌شود. مثال:

\begin{quote}
    \begin{enumerate}
        \item مورد اول.
        \item مورد دوم.
        \item مورد سوم.
    \end{enumerate}
\end{quote}

\mnote{از فهرست شماره‌دار برای مراحل، دستورالعمل‌ها یا اولویت‌ها استفاده کنید.}

\subsection{فهرست توصیفی}

فهرست توصیفی، برای تعریف اصطلاحات استفاده
می‌شود. مثال:

\begin{quote}
    \begin{description}
        \item[اصطلاح اول:] توضیح اول.
        \item[اصطلاح دوم:] توضیح دوم.
    \end{description}
\end{quote}

\mnote{از فهرست توصیفی برای واژه‌نامه یا تعریف اصطلاحات استفاده کنید.}

\mnote{در همه‌ی فهرست‌ها، آیتم‌ها باید ساختار گرامری یکسان داشته باشند. این اصل «ساختار موازی» نام دارد.}

\section{خطاهای رایج}
\label{sec:text-mistakes}

\subsection{فراموش کردن \cmd{lr}}
\label{sec:text-mistake-lr}

اگر \cmd{lr} را برای واژگان لاتین فراموش کنید،
نتیجه به‌صورت نامفهوم نمایش داده می‌شود:

\begin{latin}
\begin{minted}{latex}
% |\pc{اشتباه}|:
|\pc{این کتاب درباره}|machine learning|\pc{است.}|%

% |\pc{درست}|:
|\pc{این کتاب درباره}|\lr{machine learning}|\pc{است.}|%
\end{minted}
\end{latin}

\subsection{استفاده از \cmd{\textbackslash\textbackslash} برای پاراگراف}
\label{sec:text-mistake-parbreak}

دستور \cmd{\textbackslash\textbackslash} برای شکستن
خط در یک پاراگراف است، نه برای ایجاد پاراگراف
جدید. برای پاراگراف جدید، از یک خط خالی
استفاده کنید.

\subsection{پاراگراف بسیار طولانی}
\label{sec:text-mistake-longpara}

پاراگراف‌های بیش از ۱۰ جمله، خواننده را خسته
می‌کنند. آن‌ها را به پاراگراف‌های کوچک‌تر
تقسیم کنید.

\section{تمرین‌ها}
\label{sec:text-exercises}

\begin{exercise}
    یک پاراگراف فارسی با پنج جمله بنویسید
    که شامل سه واژه‌ی لاتین باشد. از
    \cmd{lr} برای واژگان لاتین استفاده کنید.
    \label{exc:text-paragraph}
\end{exercise}

\begin{exercise}
    یک فهرست نشانه‌دار با پنج مورد بنویسید
    که همه‌ی موارد ساختار گرامری یکسان داشته
    باشند.
    \label{exc:text-list}
\end{exercise}

\begin{exercise}
    یک پاراگراف لاتین با محیط \env{latin}
    بنویسید و آن را با متن فارسی ترکیب کنید.
    \label{exc:text-latin}
\end{exercise}

\begin{exercise}
    یک متن فارسی با نشانه‌های نگارشی مختلف
    (گیومه، پرانتز، کروشه، نقطه‌ویرگول) بنویسید.
    \label{exc:text-punctuation}
\end{exercise}

\section{خلاصه}
\label{sec:text-summary}

در این فصل، با متن و پاراگراف در \lr{MatinBook}
آشنا شدیم:

\begin{itemize}
    \item متن فارسی به‌طور ساده نوشته می‌شود و
    \lr{LaTeX} جهت را خودکار تشخیص می‌دهد.

    \item برای واژگان لاتین در متن فارسی، از
    \cmd{lr\{...\}} استفاده کنید.

    \item برای متن فارسی در محیط لاتین، از
    \cmd{rl\{...\}} استفاده کنید.

    \item برای پاراگراف‌های کامل لاتین، از
    محیط \env{latin} استفاده کنید.

    \item پاراگراف‌ها با یک خط خالی از هم جدا
    می‌شوند.

    \item طول ایده‌آل پاراگراف، ۳ تا ۷ جمله
    است.

    \item خطاهای رایج شامل فراموش کردن
    \cmd{lr}، استفاده از \cmd{\textbackslash\textbackslash}
    برای پاراگراف، و پاراگراف‌های بسیار طولانی
    است.
\end{itemize}

در فصل بعدی (\cref{chap:environments})، با
محیط‌های علمی (قضیه، لم، تعریف، مثال، اثبات)
آشنا می‌شوید و یاد می‌گیرید که چگونه از آن‌ها
استفاده کنید.
%%
%% part2-book-components/ch08-environments.tex
%%
%% Chapter 8: Scientific Environments — Theorem, lemma,
%% corollary, proposition, definition, example, remark,
%% exercise, solution, and proof.
%%

\chapter{محیط‌های علمی}
\label{chap:environments}

\section{چرا این موضوع مهم است؟}
\label{sec:environments-why}

در کتاب‌های ریاضی و فنی، محتوا به انواع مختلفی
تقسیم می‌شود: قضیه، لم، تعریف، مثال، اثبات و...
هر نوع، نقش معنایی مشخصی دارد و باید در قالب
بصری متمایز نمایش داده شود تا خواننده بتواند
به‌سرعت نوع محتوا را تشخیص دهد.

\lr{MatinBook} ده نوع محیط علمی را با دو خانواده‌ی
بصری متمایز ارائه می‌دهد. در این فصل، با این
محیط‌ها آشنا می‌شوید و یاد می‌گیرید که چگونه
از آن‌ها استفاده کنید.

\mnote{استفاده‌ی درست از محیط‌های علمی، یکی از عوامل مهم در کیفیت کتاب‌های ریاضی و فنی است.}

\section{دو خانواده‌ی جعبه}
\label{sec:environments-families}

محیط‌های علمی \lr{MatinBook} به دو خانواده
تقسیم می‌شوند:

\begin{enumerate}
    \item \textbf{خانواده‌ی \lr{matinbox}:}
    با یک نوار رنگی توپر در بالای جعبه. این
    خانواده برای محیط‌های نظری و رسمی استفاده
    می‌شود.
    \item \textbf{خانواده‌ی \lr{matin-outline}:}
    با یک خط عمودی در سمت راست و دو خط افقی
    زیر عنوان و زیر محتوا. این خانواده برای
    محیط‌های توضیحی و غیررسمی استفاده می‌شود.
\end{enumerate}

\mnote{تفکیک دو خانواده، بر اساس نوع محتوا است، نه سلیقه. محیط‌های نظری «رسمی» هستند، محیط‌های توضیحی «غیررسمی».}

\subsection{خانواده‌ی \lr{matinbox}}
\label{sec:environments-matinbox}

خانواده‌ی \lr{matinbox} شامل هفت محیط است:

\begin{itemize}
    \item \env{theorem} — قضیه.
    \item \env{lemma} — لم.
    \item \env{corollary} — نتیجه.
    \item \env{proposition} — گزاره.
    \item \env{definition} — تعریف.
    \item \env{remark} — نکته.
    \item \env{exercise} — تمرین.
\end{itemize}

\subsection{خانواده‌ی \lr{matin-outline}}
\label{sec:environments-outline}

خانواده‌ی \lr{matin-outline} شامل سه محیط است:

\begin{itemize}
    \item \env{example} — مثال.
    \item \env{proof} — اثبات.
    \item \env{solution} — راه‌حل.
\end{itemize}

\mnote{خانواده‌ی \lr{matin-outline} در نسخه‌ی \lr{1.2} بازطراحی شده و خط عمودی آن همیشه در سمت راست است.}

\section{محیط‌های خانواده‌ی \lr{matinbox}}
\label{sec:environments-matinbox-detail}

\subsection{قضیه (\lr{theorem})}
\label{sec:environments-theorem}

قضیه، یک حکم ریاضی است که اثبات شده است.
برای ایجاد قضیه:

\begin{latin}
\begin{minted}{latex}
\begin{theorem}[|\pc{قضیه‌ی فیثاغورث}|]
    |\pc{در هر مثلث قائم‌الزاویه، مربع وتر برابر است با}|%
    |\pc{مجموع مربعات دو ضلع دیگر:}|%
    \[
        a^{2} + b^{2} = c^{2}
    \]
    \label{thm:pythagoras}
\end{theorem}
\end{minted}
\end{latin}

خروجی:

\begin{theorem}[قضیه‌ی فیثاغورث]
    در هر مثلث قائم‌الزاویه، مربع وتر برابر است با
    مجموع مربعات دو ضلع دیگر:
    \[
        a^{2} + b^{2} = c^{2}
    \]
    \label{thm:pythagoras-ch8}
\end{theorem}

\mnote{عنوان داخل \lr{[]} اختیاری است. اگر آن را ندهید، فقط «قضیه ۱.۱» نمایش داده می‌شود.}

\subsection{لم (\lr{lemma})}
\label{sec:environments-lemma}

لم، یک قضیه‌ی کمکی است که برای اثبات قضیه‌های
بزرگ‌تر استفاده می‌شود:

\begin{latin}
\begin{minted}{latex}
\begin{lemma}[|\pc{لم اقلیدس}|]
    |\pc{اگر}|\lr{$p$}|\pc{عددی اول باشد و}|\lr{$p \mid ab$}|\pc{، آنگاه}|%
    |\lr{$p \mid a$}|\pc{یا}|\lr{$p \mid b$}|\pc{.}|%
    \label{lem:euclid}
\end{lemma}
\end{minted}
\end{latin}

خروجی:

\begin{lemma}[لم اقلیدس]
    اگر \lr{$p$} عددی اول باشد و \lr{$p \mid ab$}،
    آنگاه \lr{$p \mid a$} یا \lr{$p \mid b$}.
    \label{lem:euclid-ch8}
\end{lemma}

\subsection{نتیجه (\lr{corollary})}
\label{sec:environments-corollary}

نتیجه، یک حکم است که بلافاصله از یک قضیه یا
لم نتیجه می‌شود:

\begin{latin}
\begin{minted}{latex}
\begin{corollary}
    |\pc{هر دنباله‌ی یکنوا و کراندار همگراست.}|%
    \label{cor:monotone}
\end{corollary}
\end{minted}
\end{latin}

خروجی:

\begin{corollary}
    هر دنباله‌ی یکنوا و کراندار همگراست.
    \label{cor:monotone-ch8}
\end{corollary}

\subsection{گزاره (\lr{proposition})}
\label{sec:environments-proposition}

گزاره، یک حکم است که ارزش اثبات دارد ولی
لزوماً به اندازه‌ی قضیه مهم نیست:

\begin{latin}
\begin{minted}{latex}
\begin{proposition}
    |\pc{مجموعه‌ی اعداد گویا}|\lr{$\Q$}|\pc{در}|\lr{$\R$}|\pc{چگال است.}|%
    \label{prop:q-dense}
\end{proposition}
\end{minted}
\end{latin}

خروجی:

\begin{proposition}
    مجموعه‌ی اعداد گویا \lr{$\Q$} در \lr{$\R$}
    چگال است.
    \label{prop:q-dense-ch8}
\end{proposition}

\subsection{تعریف (\lr{definition})}
\label{sec:environments-definition}

تعریف، یک مفهوم جدید را معرفی می‌کند:

\begin{latin}
\begin{minted}{latex}
\begin{definition}[|\pc{گراف}|]
    |\pc{گراف}|\lr{$G = (V, E)$}|\pc{مجموعه‌ای از رئوس}|\lr{$V$}|%
    |\pc{و یال‌های}|\lr{$E$}|\pc{است.}|%
    \label{def:graph}
\end{definition}
\end{minted}
\end{latin}

خروجی:

\begin{definition}[گراف]
    گراف \lr{$G = (V, E)$} مجموعه‌ای از رئوس
    \lr{$V$} و یال‌های \lr{$E$} است.
    \label{def:graph-ch8}
\end{definition}

\mnote{تعریف‌ها معمولاً در فصل‌های اول کتاب می‌آیند و پایه‌ی مفاهیم بعدی هستند.}

\subsection{نکته (\lr{remark})}
\label{sec:environments-remark}

نکته، یک توضیح مکمل یا هشدار است:

\begin{latin}
\begin{minted}{latex}
\begin{remark}
    |\pc{تقسیم بر صفر تعریف نشده است! هنگام ساده‌سازی}|%
    |\pc{عبارات کسری، همواره فرض کنید مخرج غیرصفر است.}|%
    \label{rem:division-zero}
\end{remark}
\end{minted}
\end{latin}

خروجی:

\begin{remark}
    تقسیم بر صفر تعریف نشده است! هنگام ساده‌سازی
    عبارات کسری، همواره فرض کنید مخرج غیرصفر است.
    \label{rem:division-zero-ch8}
\end{remark}

\mnote{استفاده‌ی سنجیده از نکته ضروری است. اگر بیش از حد استفاده شود، اثر هشداردهندگی آن از بین می‌رود.}

\subsection{تمرین (\lr{exercise})}
\label{sec:environments-exercise}

تمرین، یک مسئله برای حل کردن است:

\begin{latin}
\begin{minted}{latex}
\begin{exercise}
    |\pc{ثابت کنید مجموع دو عدد زوج، زوج است.}|%
    \label{exc:sum-even}
\end{exercise}
\end{minted}
\end{latin}

خروجی:

\begin{exercise}
    ثابت کنید مجموع دو عدد زوج، زوج است.
    \label{exc:sum-even-ch8}
\end{exercise}

\section{محیط‌های خانواده‌ی \lr{matin-outline}}
\label{sec:environments-outline-detail}

\subsection{مثال (\lr{example})}
\label{sec:environments-example}

مثال، یک نمونه‌ی عملی از یک مفهوم است:

\begin{latin}
\begin{minted}{latex}
\begin{example}[|\pc{محاسبه‌ی انتگرال}|]
    |\pc{مساحت زیر نمودار}|\lr{$f(x) = x^{2}$}|\pc{از}|\lr{$x = 0$}|%
    |\pc{تا}|\lr{$x = 2$}|\pc{:}|%
    \[
        \int_{0}^{2} x^{2} \, dx = \left[ \frac{x^{3}}{3} \right]_{0}^{2}
        = \frac{8}{3} \approx 2.67
    \]
    \label{ex:integral-x2}
\end{example}
\end{minted}
\end{latin}

خروجی:

\begin{example}[محاسبه‌ی انتگرال]
    مساحت زیر نمودار \lr{$f(x) = x^{2}$} از
    \lr{$x = 0$} تا \lr{$x = 2$}:
    \[
        \int_{0}^{2} x^{2} \, dx = \left[ \frac{x^{3}}{3} \right]_{0}^{2}
        = \frac{8}{3} \approx 2.67
    \]
    \label{ex:integral-x2-ch8}
\end{example}

\subsection{اثبات (\lr{proof})}
\label{sec:environments-proof}

اثبات، استدلال ریاضی برای یک حکم است:

\begin{latin}
\begin{minted}{latex}
\begin{proof}
    |\pc{از برهان خلف استفاده می‌کنیم. فرض کنید}|\lr{$\sqrt{2}$}|%
    |\pc{گویا باشد:}|\lr{$\sqrt{2} = \frac{a}{b}$}|\pc{که}|\lr{$\gcd(a,b) = 1$}|\pc{.}|%

    \[
    2 = \frac{a^{2}}{b^{2}} \implies a^{2} = 2b^{2}
    \]

    |\pc{بنابراین}|\lr{$a$}|\pc{زوج است. پس}|\lr{$b$}|\pc{نیز زوج است که}|%
    |\pc{با فرض}|\lr{$\gcd(a,b) = 1$}|\pc{تناقض دارد.}|%
\end{proof}
\end{minted}
\end{latin}

خروجی:

\begin{proof}
    از برهان خلف استفاده می‌کنیم. فرض کنید \lr{$\sqrt{2}$}
    گویا باشد: \lr{$\sqrt{2} = \frac{a}{b}$} که
    \lr{$\gcd(a,b) = 1$}.

    \[
    2 = \frac{a^{2}}{b^{2}} \implies a^{2} = 2b^{2}
    \]

    بنابراین \lr{$a$} زوج است. پس \lr{$b$} نیز
    زوج است که با فرض \lr{$\gcd(a,b) = 1$}
    تناقض دارد.
\end{proof}

\mnote{محیط اثبات به‌طور خودکار علامت پایان (□) را در انتهای جعبه اضافه می‌کند.}

\subsection{راه‌حل (\lr{solution})}
\label{sec:environments-solution}

راه‌حل، پاسخ تشریحی یک تمرین است:

\begin{latin}
\begin{minted}{latex}
\begin{solution}
    |\pc{فرض کنید}|\lr{$a = 2m$}|\pc{و}|\lr{$b = 2n$}|\pc{.}|%
    |\pc{آنگاه}|\lr{$a + b = 2(m + n)$}|\pc{که زوج است.}|%
\end{solution}
\end{minted}
\end{latin}

خروجی:

\begin{solution}
    فرض کنید \lr{$a = 2m$} و \lr{$b = 2n$}.
    آنگاه \lr{$a + b = 2(m + n)$} که زوج است.
\end{solution}

\mnote{محیط راه‌حل و اثبات شماره‌گذاری نمی‌شوند، چون به‌طور طبیعی به تمرین یا قضیه‌ی قبلی متصل هستند.}

\section{شمارنده‌ی مشترک}
\label{sec:environments-counter}

تمام محیط‌های شماره‌دار (قضیه، لم، نتیجه، گزاره،
تعریف، مثال، نکته، تمرین) از یک \textbf{شمارنده‌ی
مشترک} استفاده می‌کنند. این یعنی شماره‌گذاری
پیوسته است:

\begin{latin}
\begin{minted}{latex}
\begin{definition}[|\pc{گراف}|]   % |\pc{شماره ۱.۱}|%
    ...
\end{definition}

\begin{theorem}[|\pc{اویلر}|]   % |\pc{شماره ۱.۲}|%
    ...
\end{theorem}

\begin{lemma}[|\pc{دست‌دادن}|]   % |\pc{شماره ۱.۳}|%
    ...
\end{lemma}

\begin{example}                % |\pc{شماره ۱.۴}|%
    ...
\end{example}
\end{minted}
\end{latin}

\mnote{این شمارنده‌گذاری مشترک، یک تصمیم آگاهانه است. جامعه‌ی \lr{LaTeX} فارسی ترجیح می‌دهد شماره‌گذاری پیوسته باشد تا خواننده راحت‌تر مفاهیم را پیدا کند.}

\section{ارجاع به محیط‌ها}
\label{sec:environments-references}

برای ارجاع به محیط‌های علمی، از دستور \cmd{cref}
استفاده می‌کنیم:

\begin{latin}
\begin{minted}{latex}
% |\pc{ارجاع به قضیه}|%
\cref{thm:pythagoras}

% |\pc{ارجاع به تعریف}|%
\cref{def:graph}

% |\pc{ارجاع به چند مرجع}|%
\cref{thm:pythagoras,def:graph}

% |\pc{ارجاع بازه‌ای}|%
\crefrange{thm:pythagoras}{def:graph}
\end{minted}
\end{latin}

\mnote{دستور \lr{\textbackslash cref} به‌طور خودکار نوع مرجع را تشخیص می‌دهد و پیشوند مناسب (قضیه، تعریف، لم) را اضافه می‌کند.}

\subsection{نام‌گذاری برچسب‌ها}
\label{sec:environments-labels}

برای برچسب‌ها، از پیشوندهای زیر استفاده کنید:

\begin{itemize}
    \item \textbf{قضیه:} \lr{thm:xxx}.
    \item \textbf{لم:} \lr{lem:xxx}.
    \item \textbf{نتیجه:} \lr{cor:xxx}.
    \item \textbf{گزاره:} \lr{prop:xxx}.
    \item \textbf{تعریف:} \lr{def:xxx}.
    \item \textbf{مثال:} \lr{ex:xxx}.
    \item \textbf{نکته:} \lr{rem:xxx}.
    \item \textbf{تمرین:} \lr{exc:xxx}.
\end{itemize}

\mnote{برچسب‌های معنادار، خواندن فایل‌های \lr{tex} را آسان‌تر می‌کنند و از خطاهای ارجاع جلوگیری می‌کنند.}

\section{جدول خلاصه}
\label{sec:environments-summary-table}

جدول \cref{tab:environments-summary} خلاصه‌ای از
ده محیط علمی \lr{MatinBook} را نشان می‌دهد.

\begin{matintable}{خلاصه‌ی محیط‌های علمی \lr{MatinBook}}{tab:environments-summary}
\begin{tblr}{
    width = \textwidth,
    colspec = {l l l X[c]},
    hlines, vlines,
    colsep = 6pt,
    rowsep = 4pt,
    row{1} = {font=\bfseries},
}
محیط & عنوان & خانواده & شماره‌دار \\
\lr{theorem} & قضیه & \lr{matinbox} & بله \\
\lr{lemma} & لم & \lr{matinbox} & بله \\
\lr{corollary} & نتیجه & \lr{matinbox} & بله \\
\lr{proposition} & گزاره & \lr{matinbox} & بله \\
\lr{definition} & تعریف & \lr{matinbox} & بله \\
\lr{remark} & نکته & \lr{matinbox} & بله \\
\lr{exercise} & تمرین & \lr{matinbox} & بله \\
\lr{example} & مثال & \lr{matin-outline} & بله \\
\lr{proof} & اثبات & \lr{matin-outline} & خیر \\
\lr{solution} & راه‌حل & \lr{matin-outline} & خیر \\
\end{tblr}
\end{matintable}

\section{خطاهای رایج}
\label{sec:environments-mistakes}

\subsection{قرار دادن \cmd{label} بیرون محیط}
\label{sec:environments-mistake-label}

اگر \cmd{label} را بیرون محیط قرار دهید، ارجاع
به آن کار نمی‌کند:

\begin{latin}
\begin{minted}{latex}
% |\pc{اشتباه}|:
\begin{theorem}
    |\pc{متن قضیه...}|%
\end{theorem}
\label{thm:wrong}

% |\pc{درست}|:
\begin{theorem}
    |\pc{متن قضیه...}|%
    \label{thm:correct}
\end{theorem}
\end{minted}
\end{latin}

\subsection{فراموش کردن \cmd{label}}
\label{sec:environments-mistake-no-label}

اگر \cmd{label} را فراموش کنید، نمی‌توانید به
آن محیط ارجاع دهید.

\subsection{استفاده از محیط نامناسب}
\label{sec:environments-mistake-wrong-env}

هر محیط، نقش معنایی مشخصی دارد. استفاده از
\env{theorem} برای یک تعریف یا \env{definition}
برای یک قضیه، ساختار کتاب را به‌هم می‌ریزد.

\section{تمرین‌ها}
\label{sec:environments-exercises}

\begin{exercise}
    یک قضیه، یک تعریف و یک مثال بنویسید و
    به هر یک ارجاع دهید.
    \label{exc:environments-basic}
\end{exercise}

\begin{exercise}
    یک لم بنویسید و یک نتیجه از آن استخراج
    کنید.
    \label{exc:environments-lemma}
\end{exercise}

\begin{exercise}
    یک اثبات برای قضیه‌ی فیثاغورث بنویسید
    و از محیط \env{proof} استفاده کنید.
    \label{exc:environments-proof}
\end{exercise}

\begin{exercise}
    یک تمرین بنویسید و راه‌حل آن را با محیط
    \env{solution} ارائه دهید.
    \label{exc:environments-solution}
\end{exercise}

\section{خلاصه}
\label{sec:environments-summary}

در این فصل، با محیط‌های علمی \lr{MatinBook}
آشنا شدیم:

\begin{itemize}
    \item \lr{MatinBook} ده محیط علمی را در دو
    خانواده‌ی بصری ارائه می‌دهد.

    \item خانواده‌ی \lr{matinbox} شامل قضیه، لم،
    نتیجه، گزاره، تعریف، نکته و تمرین است.

    \item خانواده‌ی \lr{matin-outline} شامل مثال،
    اثبات و راه‌حل است.

    \item تمام محیط‌های شماره‌دار، از یک
    شمارنده‌ی مشترک استفاده می‌کنند.

    \item اثبات و راه‌حل، شماره‌گذاری نمی‌شوند.

    \item ارجاع به محیط‌ها با \cmd{cref} انجام
    می‌شود که پیشوند مناسب را خودکار اضافه
    می‌کند.

    \item خطاهای رایج شامل قرار دادن \cmd{label}
    بیرون محیط، فراموش کردن \cmd{label}، و
    استفاده از محیط نامناسب است.
\end{itemize}

در فصل بعدی (\cref{chap:math})، با ریاضیات
در \lr{MatinBook} آشنا می‌شوید و یاد می‌گیرید
که چگونه معادلات، ماتریس‌ها و عملگرها را
بنویسید.
%%
%% part2-book-components/ch09-math.tex
%%
%% Chapter 9: Mathematics — Equations, matrices, operators,
%% delimiters, and number sets.
%%

\chapter{ریاضیات}
\label{chap:math}

\section{چرا این موضوع مهم است؟}
\label{sec:math-why}

ریاضیات، ستون فقرات کتاب‌های فنی و علمی است.
از معادلات ساده تا ماتریس‌های پیچیده، از
عملگرهای حسابان تا مجموعه‌های عددی، همه در
\lr{MatinBook} پشتیبانی می‌شوند.

\lr{MatinBook} از بسته‌های \lr{amsmath}،
\lr{mathtools} و \lr{unicode-math} استفاده
می‌کند و مجموعه‌ای از عملگرها و جداکننده‌های
سفارشی را نیز اضافه کرده است. در این فصل،
با تمام این قابلیت‌ها آشنا می‌شوید.

\mnote{فونت ریاضی \lr{MatinBook}، \lr{XITS Math} است که از \lr{OpenType} پشتیبانی می‌کند و نمادهای ریاضی را با کیفیت بالا نمایش می‌دهد.}

\section{معادلات درون‌خطی}
\label{sec:math-inline}

معادلات درون‌خطی، برای اشاره به نمادها و
عبارات کوتاه استفاده می‌شوند. در \lr{MatinBook}،
این معادلات باید داخل \cmd{lr} قرار گیرند تا
جهت‌نمایی درست حفظ شود:

\begin{latin}
\begin{minted}{latex}
|\pc{رابطه‌ی معروف اینشتین:}|\lr{$E = mc^{2}$}|\pc{است.}|%
|\pc{قضیه‌ی فیثاغورث:}|\lr{$a^{2} + b^{2} = c^{2}$}|\pc{است.}|%
|\pc{اتحاد اویلر:}|\lr{$e^{i\pi} + 1 = 0$}|\pc{است.}|%
\end{minted}
\end{latin}

خروجی:

رابطه‌ی معروف اینشتین: \lr{$E = mc^{2}$} است.
قضیه‌ی فیثاغورث: \lr{$a^{2} + b^{2} = c^{2}$} است.
اتحاد اویلر: \lr{$e^{i\pi} + 1 = 0$} است.

\mnote{همیشه معادلات درون‌خطی را داخل \lr{\textbackslash lr\{...\}} قرار دهید. این کار از به‌هم‌ریختن جهت متن جلوگیری می‌کند.}

\section{معادلات نمایشی}
\label{sec:math-display}

معادلات نمایشی، در خط جداگانه و وسط‌چین نمایش
داده می‌شوند. \lr{MatinBook} از سه محیط اصلی
پشتیبانی می‌کند:

\subsection{معادله‌ی ساده (\lr{equation})}
\label{sec:math-equation}

برای یک معادله‌ی شماره‌دار، از محیط
\env{equation} استفاده کنید:

\begin{latin}
\begin{minted}{latex}
\begin{equation}
    \int_{0}^{\infty} \frac{x^{3}}{e^{x} - 1} \, dx
    = \frac{\pi^{4}}{15}
    \label{eq:stefan-boltzmann}
\end{equation}
\end{minted}
\end{latin}

خروجی:

\begin{equation}
    \int_{0}^{\infty} \frac{x^{3}}{e^{x} - 1} \, dx
    = \frac{\pi^{4}}{15}
    \label{eq:stefan-boltzmann-ch9}
\end{equation}

\mnote{شماره‌گذاری معادلات در \lr{MatinBook} به‌صورت «شماره-فصل» است: \lr{(۱-۲)} یعنی فرمول اول از فصل دوم.}

\subsection{معادلات هم‌تراز (\lr{align})}
\label{sec:math-align}

برای چند معادله‌ی هم‌تراز، از محیط \env{align}
استفاده کنید:

\begin{latin}
\begin{minted}{latex}
\begin{align}
    (a + b)^{2} &= (a + b)(a + b) \nonumber \\
    &= a^{2} + ab + ba + b^{2} \nonumber \\
    &= a^{2} + 2ab + b^{2} \label{eq:binomial}
\end{align}
\end{minted}
\end{latin}

خروجی:

\begin{align}
    (a + b)^{2} &= (a + b)(a + b) \nonumber \\
    &= a^{2} + ab + ba + b^{2} \nonumber \\
    &= a^{2} + 2ab + b^{2} \label{eq:binomial-ch9}
\end{align}

\mnote{دستور \lr{\textbackslash nonumber} معادله را بدون شماره می‌کند. برای معادلات میانی که به آن‌ها ارجاع نمی‌دهید، از این دستور استفاده کنید.}

\subsection{معادلات غیرهم‌تراز (\lr{gather})}
\label{sec:math-gather}

برای معادلات غیرهم‌تراز، از محیط \env{gather}
استفاده کنید:

\begin{latin}
\begin{minted}{latex}
\begin{gather}
    x = a + b + c + d + e + f \label{eq:long1} \\
    y = \alpha + \beta + \gamma + \delta + \epsilon \label{eq:long2}
\end{gather}
\end{minted}
\end{latin}

خروجی:

\begin{gather}
    x = a + b + c + d + e + f \label{eq:long1-ch9} \\
    y = \alpha + \beta + \gamma + \delta + \epsilon \label{eq:long2-ch9}
\end{gather}

\mnote{تفاوت \lr{align} و \lr{gather}: در \lr{align}، معادلات با \lr{\&} هم‌تراز می‌شوند، ولی در \lr{gather}، هر معادله در خط جداگانه و بدون هم‌ترازی نمایش داده می‌شود.}

\section{ماتریس‌ها}
\label{sec:math-matrices}

\lr{MatinBook} سه دستور برای نمایش ماتریس
ارائه می‌دهد:

\begin{itemize}
    \item \cmd{mat\{...\}}: ماتریس با کروشه.
    \item \cmd{pmat\{...\}}: ماتریس با پرانتز.
    \item \cmd{vmat\{...\}}: ماتریس با خط عمودی
    (دترمینان).
\end{itemize}

\subsection{ماتریس با کروشه (\cmd{mat})}
\label{sec:math-mat-bracket}

\begin{latin}
\begin{minted}{latex}
\[
    A = \mat{1 & 2 & 3 \\ 4 & 5 & 6 \\ 7 & 8 & 9}
\]
\end{minted}
\end{latin}

خروجی:

\[
    A = \mat{1 & 2 & 3 \\ 4 & 5 & 6 \\ 7 & 8 & 9}
\]

\subsection{ماتریس با پرانتز (\cmd{pmat})}
\label{sec:math-mat-paren}

\begin{latin}
\begin{minted}{latex}
\[
    B = \pmat{a_{11} & a_{12} \\ a_{21} & a_{22}}
\]
\end{minted}
\end{latin}

خروجی:

\[
    B = \pmat{a_{11} & a_{12} \\ a_{21} & a_{22}}
\]

\subsection{دترمینان (\cmd{vmat})}
\label{sec:math-mat-vmatrix}

\begin{latin}
\begin{minted}{latex}
\[
    \det(B) = \vmat{a_{11} & a_{12} \\ a_{21} & a_{22}}
    = a_{11}a_{22} - a_{12}a_{21}
\]
\end{minted}
\end{latin}

خروجی:

\[
    \det(B) = \vmat{a_{11} & a_{12} \\ a_{21} & a_{22}}
    = a_{11}a_{22} - a_{12}a_{21}
\]

\mnote{برای ماتریس‌های بزرگ، می‌توانید از \lr{\textbackslash cdots}، \lr{\textbackslash vdots} و \lr{\textbackslash ddots} برای نشان دادن ادامه‌ی الگو استفاده کنید.}

\section{عملگرهای ریاضی}
\label{sec:math-operators}

\lr{MatinBook} مجموعه‌ای از عملگرهای سفارشی
ارائه می‌دهد که در حوزه‌های مختلف ریاضیات
و علوم کامپیوتر کاربرد دارند. جدول
\cref{tab:math-operators} این عملگرها را
خلاصه می‌کند.

\begin{matintable}{عملگرهای ریاضی \lr{MatinBook}}{tab:math-operators}
\begin{tblr}{
    width = \textwidth,
    colspec = {l l X[l]},
    hlines, vlines,
    colsep = 6pt,
    rowsep = 4pt,
    row{1} = {font=\bfseries},
}
دسته & دستور & خروجی \\
حسابان برداری & \lr{\textbackslash grad} & \lr{$\grad f$} \\
حسابان برداری & \lr{\textbackslash curl} & \lr{$\curl \vec{F}$} \\
حسابان برداری & \lr{\textbackslash diver} & \lr{$\diver \vec{F}$} \\
جبر خطی & \lr{\textbackslash rank} & \lr{$\rank(A)$} \\
جبر خطی & \lr{\textbackslash trace} & \lr{$\trace(A)$} \\
جبر خطی & \lr{\textbackslash diag} & \lr{$\diag(a, b, c)$} \\
بهینه‌سازی & \lr{\textbackslash argmax} & \lr{$\argmax_x f(x)$} \\
بهینه‌سازی & \lr{\textbackslash argmin} & \lr{$\argmin_x f(x)$} \\
نظریه اعداد & \lr{\textbackslash gcdop} & \lr{$\gcdop(a, b)$} \\
نظریه اعداد & \lr{\textbackslash lcm} & \lr{$\lcm(a, b)$} \\
توابع خاص & \lr{\textbackslash sinc} & \lr{$\sinc(x)$} \\
توابع خاص & \lr{\textbackslash sgn} & \lr{$\sgn(x)$} \\
توابع خاص & \lr{\textbackslash erf} & \lr{$\erf(x)$} \\
\end{tblr}
\end{matintable}

\mnote{نام \lr{\textbackslash gcdop} از \lr{gcd operator} گرفته شده، چون \lr{\textbackslash gcd} در برخی بسته‌ها تعریف شده است.}

\section{جداکننده‌های هوشمند}
\label{sec:math-delimiters}

\lr{MatinBook} شش جداکننده‌ی هوشمند ارائه می‌دهد
که به‌طور خودکار با محتوای خود مقیاس می‌گیرند.
جدول \cref{tab:math-delimiters} این جداکننده‌ها
را خلاصه می‌کند.

\begin{matintable}{جداکننده‌های هوشمند}{tab:math-delimiters}
\begin{tblr}{
    width = \textwidth,
    colspec = {l l X[l]},
    hlines, vlines,
    colsep = 6pt,
    rowsep = 4pt,
    row{1} = {font=\bfseries},
}
دستور & خروجی & کاربرد \\
\lr{\textbackslash abs\{x\}} & \lr{$\abs{x}$} & قدر مطلق \\
\lr{\textbackslash norm\{x\}} & \lr{$\norm{x}$} & نرم \\
\lr{\textbackslash ceil\{x\}} & \lr{$\ceil{x}$} & سقف \\
\lr{\textbackslash floor\{x\}} & \lr{$\floor{x}$} & کف \\
\lr{\textbackslash inner\{u\}\{v\}} & \lr{$\inner{u}{v}$} & ضرب داخلی \\
\lr{\textbackslash paren\{x\}} & \lr{$\paren{x}$} & پرانتز هوشمند \\
\end{tblr}
\end{matintable}

\subsection{مقیاس خودکار}
\label{sec:math-delimiters-scaling}

این جداکننده‌ها به‌طور خودکار با محتوای خود
مقیاس می‌گیرند:

\begin{latin}
\begin{minted}{latex}
\[
    \abs{\frac{a}{b}}, \quad
    \norm{\sum_{i=1}^{n} x_i}, \quad
    \paren{\frac{a + b}{c + d}}
\]
\end{minted}
\end{latin}

خروجی:

\[
    \abs{\frac{a}{b}}, \quad
    \norm{\sum_{i=1}^{n} x_i}, \quad
    \paren{\frac{a + b}{c + d}}
\]

\mnote{برای مقیاس دستی، می‌توانید از \lr{\textbackslash abs[\textbackslash Big]\{...\}}، \lr{\textbackslash abs[\textbackslash big]\{...\}} و... استفاده کنید.}

\section{مجموعه‌های عددی}
\label{sec:math-numbers}

\lr{MatinBook} پنج مجموعه‌ی عددی استاندارد
ارائه می‌دهد:

\begin{itemize}
    \item \cmd{N}: مجموعه‌ی اعداد طبیعی
    (\lr{$\N$}).
    \item \cmd{Z}: مجموعه‌ی اعداد صحیح
    (\lr{$\Z$}).
    \item \cmd{Q}: مجموعه‌ی اعداد گویا
    (\lr{$\Q$}).
    \item \cmd{R}: مجموعه‌ی اعداد حقیقی
    (\lr{$\R$}).
    \item \cmd{C}: مجموعه‌ی اعداد مختلط
    (\lr{$\C$}).
\end{itemize}

رابطه‌ی شمول میان این مجموعه‌ها:

\begin{latin}
\begin{minted}{latex}
\[
    \N \subset \Z \subset \Q \subset \R \subset \C
\]
\end{minted}
\end{latin}

خروجی:

\[
    \N \subset \Z \subset \Q \subset \R \subset \C
\]

\mnote{این مجموعه‌ها با \lr{\textbackslash providecommand} تعریف شده‌اند، پس می‌توانید آن‌ها را با \lr{\textbackslash renewcommand} بازنویسی کنید.}

\section{نمادگذاری مجموعه}
\label{sec:math-sets}

\lr{MatinBook} سه دستور برای نمادگذاری مجموعه
ارائه می‌دهد:

\begin{itemize}
    \item \cmd{set\{...\}}: مجموعه‌ی ساده.
    \item \cmd{setbuilder\{...\}\{...\}}: مجموعه‌ی
    سازنده.
    \item \cmd{given}: جداکننده‌ی شرط.
\end{itemize}

\begin{latin}
\begin{minted}{latex}
\[
    A = \set{1, 2, 3, 4, 5}, \quad
    B = \setbuilder{x \in \R}{x > 0}
\]
\end{minted}
\end{latin}

خروجی:

\[
    A = \set{1, 2, 3, 4, 5}, \quad
    B = \setbuilder{x \in \R}{x > 0}
\]

\mnote{دستور \lr{\textbackslash setbuilder} دو آرگومان می‌گیرد: شرط و محدوده.}

\section{مشتق و انتگرال}
\label{sec:math-derivatives}

\lr{MatinBook} چهار دستور برای مشتق و انتگرال
ارائه می‌دهد:

\begin{itemize}
    \item \cmd{dd\{x\}}: دیفرانسیل (\lr{$dx$}).
    \item \cmd{D}: دیفرانسیل کوتاه (\lr{$\D$}).
    \item \cmd{deriv\{f\}\{x\}}: مشتق معمولی
    (\lr{$\deriv{f}{x}$}).
    \item \cmd{pderiv\{f\}\{x\}}: مشتق جزئی
    (\lr{$\pderiv{f}{x}$}).
\end{itemize}

\begin{latin}
\begin{minted}{latex}
\[
    \deriv{f}{x} = \lim_{h \to 0} \frac{f(x + h) - f(x)}{h}, \quad
    \pderiv{f}{x} = \frac{\partial f}{\partial x}
\]
\end{minted}
\end{latin}

خروجی:

\[
    \deriv{f}{x} = \lim_{h \to 0} \frac{f(x + h) - f(x)}{h}, \quad
    \pderiv{f}{x} = \frac{\partial f}{\partial x}
\]

\mnote{استفاده از \lr{\textbackslash dd\{x\}} به‌جای \lr{dx} در انتگرال‌ها، فاصله‌ی بهتری ایجاد می‌کند.}

\section{شماره‌گذاری معادلات}
\label{sec:math-equation-numbering}

شماره‌گذاری معادلات در \lr{MatinBook} به‌صورت
\textbf{«شماره-فصل»} است. یعنی معادله‌ی شماره‌ی
\lr{(۵-۱)} به معنای فرمول پنجم از فصل اول است.

\begin{matintable}{نمونه‌ی شماره‌گذاری معادلات}{tab:equation-numbering}
\begin{tblr}{
    width = 0.7\textwidth,
    colspec = {c c c},
    hlines, vlines,
    colsep = 8pt,
    rowsep = 4pt,
    row{1} = {font=\bfseries},
}
فصل & معادله & نمایش \\
۱ & اول & \lr{(۱-۱)} \\
۱ & دوم & \lr{(۲-۱)} \\
۲ & اول & \lr{(۱-۲)} \\
۲ & دوم & \lr{(۲-۲)} \\
\end{tblr}
\end{matintable}

\mnote{این شماره‌گذاری، مطابق استاندارد نشر دانشگاهی ایران (دانشگاه تهران و دانشگاه علوم پزشکی) است.}

\section{شکستن معادلات طولانی}
\label{sec:math-breaking}

\lr{MatinBook} امکان شکستن معادلات طولانی بین
صفحات را فراهم می‌کند:

\begin{latin}
\begin{minted}{latex}
\begin{align}
    f(x) &= a_{0} + a_{1}x + a_{2}x^{2} + \cdots \nonumber \\
         &\quad + a_{n}x^{n} \\
    &= \sum_{i=0}^{n} a_{i}x^{i}
\end{align}
\end{minted}
\end{latin}

خروجی:

\begin{align}
    f(x) &= a_{0} + a_{1}x + a_{2}x^{2} + \cdots \nonumber \\
         &\quad + a_{n}x^{n} \\
    &= \sum_{i=0}^{n} a_{i}x^{i}
\end{align}

\mnote{دستور \lr{\textbackslash allowdisplaybreaks} در \lr{MatinBook} فعال است، پس معادلات طولانی به‌طور خودکار بین صفحات می‌شکنند.}

\section{خطاهای رایج}
\label{sec:math-mistakes}

\subsection{فراموش کردن \cmd{lr} برای معادلات درون‌خطی}
\label{sec:math-mistake-lr}

اگر معادله‌ی درون‌خطی را داخل \cmd{lr} قرار
ندهید، ممکن است جهت آن به‌هم بریزد:

\begin{latin}
\begin{minted}{latex}
% |\pc{اشتباه}|:
|\pc{این معادله}|\lr{$E = mc^{2}$}|\pc{است.}|%

% |\pc{درست}|:
|\pc{این معادله}|\lr{$E = mc^{2}$}|\pc{است.}|%
\end{minted}
\end{latin}

\subsection{استفاده از \cmd{\$...\$} در عنوان بخش}
\label{sec:math-mistake-title}

هرگز از حالت ریاضی در عنوان بخش استفاده نکنید:

\begin{latin}
\begin{minted}{latex}
% |\pc{اشتباه}|:
\section{|\pc{تحلیل}|\lr{$O(n)$}}

% |\pc{درست}|:
\section{|\pc{تحلیل}|\lr{O(n)}}
\end{minted}
\end{latin}

\subsection{فراموش کردن شماره‌گذاری معادلات مهم}
\label{sec:math-mistake-label}

معادلاتی که به آن‌ها ارجاع می‌دهید، باید
شماره‌دار باشند:

\begin{latin}
\begin{minted}{latex}
% |\pc{اشتباه}|:
\[
    E = mc^{2}
\]
|\pc{طبق معادله‌ی بالا...}|%

% |\pc{درست}|:
\begin{equation}
    E = mc^{2}
    \label{eq:einstein}
\end{equation}
|\pc{طبق}|\cref{eq:einstein}|\pc{...}|%
\end{minted}
\end{latin}

\section{تمرین‌ها}
\label{sec:math-exercises}

\begin{exercise}
    یک معادله‌ی درون‌خطی و یک معادله‌ی نمایشی
    بنویسید. به معادله‌ی نمایشی ارجاع دهید.
    \label{exc:math-equation}
\end{exercise}

\begin{exercise}
    یک ماتریس ۳×۳ با کروشه و یک ماتریس ۲×۲
    با پرانتز بنویسید.
    \label{exc:math-matrix}
\end{exercise}

\begin{exercise}
    از سه عملگر سفارشی \lr{MatinBook} در یک
    معادله استفاده کنید.
    \label{exc:math-operators}
\end{exercise}

\begin{exercise}
    یک مجموعه‌ی سازنده بنویسید که اعداد
    حقیقی مثبت را توصیف کند.
    \label{exc:math-set}
\end{exercise}

\section{خلاصه}
\label{sec:math-summary}

در این فصل، با ریاضیات در \lr{MatinBook}
آشنا شدیم:

\begin{itemize}
    \item معادلات درون‌خطی باید داخل \cmd{lr}
    قرار گیرند.

    \item معادلات نمایشی با سه محیط
    \env{equation}، \env{align} و \env{gather}
    نمایش داده می‌شوند.

    \item سه دستور برای ماتریس:
    \cmd{mat}، \cmd{pmat}، \cmd{vmat}.

    \item ۱۳ عملگر سفارشی برای حسابان برداری،
    جبر خطی، بهینه‌سازی، نظریه اعداد و توابع
    خاص.

    \item شش جداکننده‌ی هوشمند که با محتوا
    مقیاس می‌گیرند.

    \item پنج مجموعه‌ی عددی (\lr{$\N$}،
    \lr{$\Z$}، \lr{$\Q$}، \lr{$\R$}، \lr{$\C$}).

    \item شماره‌گذاری معادلات به‌صورت
    «شماره-فصل» است.

    \item خطاهای رایج شامل فراموش کردن
    \cmd{lr}، استفاده از حالت ریاضی در
    عنوان بخش، و فراموش کردن شماره‌گذاری
    معادلات مهم است.
\end{itemize}

در فصل بعدی (\cref{chap:code})، با کد و
الگوریتم در \lr{MatinBook} آشنا می‌شوید و
یاد می‌گیرید که چگونه کدهای برنامه‌نویسی
و شبه‌کدها را نمایش دهید.
%%
%% part2-book-components/ch10-code.tex
%%
%% Chapter 10: Code and Algorithms — Syntax-highlighted
%% code with minted, and pseudocode with algorithm.
%%

\chapter{کد و الگوریتم}
\label{chap:code}

\section{چرا این موضوع مهم است؟}
\label{sec:code-why}

در کتاب‌های برنامه‌نویسی، نمایش حرفه‌ای کد
اهمیت ویژه‌ای دارد. کد باید با رنگ‌آمیزی نحوی
مناسب، فونت خوانا، شماره‌گذاری منظم و کپشن‌های
گویا نمایش داده شود.

\lr{MatinBook} از بسته‌ی \lr{minted} برای نمایش
کد و از بسته‌ی استاندارد \lr{algorithm} برای
شبه‌کد استفاده می‌کند. در این فصل، با هر دو
آشنا می‌شوید و یاد می‌گیرید که چگونه کدهای
فارسی و لاتین را با هم ترکیب کنید.

\mnote{کدها با \lr{minted} و \lr{Pygments} رنگ‌آمیزی می‌شوند که از بیش از ۳۰۰ زبان برنامه‌نویسی پشتیبانی می‌کند.}

\section{نمایش کد با \lr{minted}}
\label{sec:code-minted}

\subsection{بلوک کد پایه}
\label{sec:code-basic}

برای نمایش یک بلوک کد، از محیط \env{minted}
داخل محیط \env{latin} استفاده کنید:

\begin{latin}
\begin{minted}{latex}
\begin{latin}
\begin{minted}{python}
def factorial(n: int) -> int:
    """Calculate n! recursively."""
    if n <= 1:
        return 1
    return n * factorial(n - 1)
\end{minted}
\end{latin}
\codecaption{|\pc{تابع فاکتوریل بازگشتی}|}
\label{code:factorial}
\end{minted}
\end{latin}

خروجی:

\begin{latin}
\begin{minted}{python}
def factorial(n: int) -> int:
    """Calculate n! recursively."""
    if n <= 1:
        return 1
    return n * factorial(n - 1)
\end{minted}
\end{latin}
\codecaption{تابع فاکتوریل بازگشتی}
\label{code:factorial-ch10}

\mnote{محیط \lr{minted} باید همیشه داخل \lr{\textbackslash begin\{latin\}} قرار گیرد. این کار جهت کد را لاتین نگه می‌دارد.}

\subsection{کپشن کد}
\label{sec:code-caption}

برای افزودن کپشن به کد، از دستور \cmd{codecaption}
استفاده کنید:

\begin{latin}
\begin{minted}{latex}
\codecaption{|\pc{کپشن کد}|}
\label{code:my-code}
\end{minted}
\end{latin}

\mnote{دستور \lr{\textbackslash codecaption} یک شمارنده‌ی مستقل برای کدها ایجاد می‌کند که به‌صورت «فصل.شماره» نمایش داده می‌شود.}

\subsection{ارجاع به کد}
\label{sec:code-reference}

برای ارجاع به کد، از دستور \cmd{cref} استفاده
کنید:

\begin{latin}
\begin{minted}{latex}
|\pc{همان‌طور که در}|\cref{code:factorial}|\pc{می‌بینیم...}|%
\end{minted}
\end{latin}

خروجی: همان‌طور که در \cref{code:factorial-ch10}
می‌بینیم...

\mnote{دستور \lr{\textbackslash cref} به‌طور خودکار نوع مرجع را تشخیص می‌دهد و پیشوند «کد» را اضافه می‌کند.}

\section{کامنت فارسی در کد}
\label{sec:code-persian-comments}

یکی از ویژگی‌های مهم \lr{MatinBook}، پشتیبانی
از کامنت فارسی در کد است. برای این کار، از دو
دستور استفاده می‌کنیم:

\begin{itemize}
    \item \cmd{pcc\{...\}}: برای کامنت فارسی
    در پایتون (و زبان‌های مشابه که \lr{\#}
    دارند).
    \item \cmd{pc\{...\}}: برای متن فارسی غیرکامنت
    (مثل \lr{LaTeX}، \lr{bash}).
\end{itemize}

\mnote{دستور \lr{\textbackslash pcc} علامت \lr{\#} را به‌طور خودکار اضافه می‌کند. دستور \lr{\textbackslash pc} فقط متن فارسی را با فونت مناسب نمایش می‌دهد.}

\subsection{کامنت فارسی در پایتون}
\label{sec:code-pcc}

برای کامنت فارسی در پایتون، از \cmd{pcc} داخل
\lr{|...|} استفاده کنید:

\begin{latin}
\begin{minted}{python}
x = 5  |\pcc{مقدار متغیر}|
y = 10  |\pcc{مقدار دیگر}|
result = x + y  |\pcc{جمع دو متغیر}|
\end{minted}
\end{latin}


خروجی:

\begin{latin}
\begin{minted}{python}
x = 5  |\pcc{مقدار متغیر}|
y = 10  |\pcc{مقدار دیگر}|
result = x + y  |\pcc{جمع دو متغیر}|
\end{minted}
\end{latin}

\mnote{اگر به‌جای \lr{\textbackslash pcc} از \lr{\# |\textbackslash pc\{...\}|} استفاده کنید، \lr{Pygments} علامت \lr{\#} را به‌عنوان شروع کامنت می‌بیند و \lr{|} را نادیده می‌گیرد. نتیجه، مربع‌های خالی است.}

\subsection{متن فارسی در \lr{LaTeX}}
\label{sec:code-pc}

برای متن فارسی غیرکامنت (مثلاً در \lr{LaTeX})،
از \cmd{pc} استفاده کنید:

\begin{latin}
\begin{minted}{latex}
\begin{latin}
\begin{minted}{latex}
\documentclass{matinbook}
\begin{document}
\chapter{|\pc{فصل نمونه}|}
\end{document}
\end{minted}
\end{latin}
\end{minted}
\end{latin}

خروجی:

\begin{latin}
\begin{minted}{latex}
\documentclass{matinbook}
\begin{document}
\chapter{|\pc{فصل نمونه}|}
\end{document}
\end{minted}
\end{latin}

\mnote{دستور \lr{\textbackslash pc} باید همیشه داخل \lr{|...|} قرار گیرد. اگر بیرون از \lr{|...|} باشد، به‌عنوان متن خام نمایش داده می‌شود.}

\subsection{جدول تصمیم‌گیری}
\label{sec:code-decision-table}

جدول \cref{tab:code-persian} به شما کمک می‌کند
که در هر موقعیت، از کدام دستور استفاده کنید.

\begin{matintable}{انتخاب دستور مناسب برای فارسی در کد}{tab:code-persian}
\begin{tblr}{
    width = \textwidth,
    colspec = {X[l] X[c] X[c]},
    hlines, vlines,
    colsep = 6pt,
    rowsep = 4pt,
    row{1} = {font=\bfseries},
}
موقعیت & زبان & دستور \\
کامنت فارسی & پایتون & \lr{|\textbackslash pcc\{...\}|} \\
کامنت فارسی & \lr{C++}، \lr{Java} & \lr{|\textbackslash pcc\{...\}|} \\
کامنت فارسی & \lr{LaTeX}، \lr{bash} & \lr{|\textbackslash pc\{\textbackslash\# ...\}|} \\
متن فارسی غیرکامنت & همه & \lr{|\textbackslash pc\{...\}|} \\
کد کاملاً لاتین & همه & بدون \lr{\textbackslash pc} \\
\end{tblr}
\end{matintable}

\section{کد درون‌خطی}
\label{sec:code-inline}

برای ارجاع به یک تابع یا متغیر در متن، از
دستور \cmd{inlcode} استفاده کنید:

\begin{latin}
\begin{minted}{latex}
|\pc{تابع}|\inlcode{factorial()}|\pc{یک تابع بازگشتی است.}|%
\end{minted}
\end{latin}

خروجی: تابع \inlcode{factorial()} یک تابع
بازگشتی است.

\mnote{دستور \lr{\textbackslash inlcode} فونت مونواسپیس را در متن اعمال می‌کند و برای کدهای کوتاه (کمتر از ۴۰ کاراکتر) مناسب است.}

\section{الگوریتم‌ها}
\label{sec:code-algorithms}

\subsection{محیط \lr{algorithm}}
\label{sec:code-algorithm-basic}

برای نمایش شبه‌کد، از بسته‌ی استاندارد
\lr{algorithm} استفاده می‌کنیم. این محیط
باید داخل \env{latin} قرار گیرد:

\begin{latin}
\begin{minted}{latex}
\begin{latin}
\begin{algorithm}
\caption{|\pc{جستجوی دودویی}|}
\label{alg:binary}
\begin{algorithmic}[1]
    \Require Sorted array $A$ of $n$ elements
    \Require Target value $x$
    \Ensure Index of $x$ in $A$, or $-1$
    \State $left \gets 0$, $right \gets n - 1$
    \While{$left \leq right$}
        \State $mid \gets (left + right) / 2$
        \If{$A[mid] = x$}
            \State \Return $mid$
        \ElsIf{$A[mid] < x$}
            \State $left \gets mid + 1$
        \Else
            \State $right \gets mid - 1$
        \EndIf
    \EndWhile
    \State \Return $-1$
\end{algorithmic}
\end{algorithm}
\end{latin}
\end{minted}
\end{latin}

خروجی:

\begin{latin}
\begin{algorithm}
\caption{جستجوی دودویی}
\label{alg:binary-ch10}
\begin{algorithmic}[1]
    \Require Sorted array $A$ of $n$ elements
    \Require Target value $x$
    \Ensure Index of $x$ in $A$, or $-1$
    \State $left \gets 0$, $right \gets n - 1$
    \While{$left \leq right$}
        \State $mid \gets (left + right) / 2$
        \If{$A[mid] = x$}
            \State \Return $mid$
        \ElsIf{$A[mid] < x$}
            \State $left \gets mid + 1$
        \Else
            \State $right \gets mid - 1$
        \EndIf
    \EndWhile
    \State \Return $-1$
\end{algorithmic}
\end{algorithm}
\end{latin}

\mnote{عنوان الگوریتم (\lr{\textbackslash caption}) می‌تواند فارسی باشد، ولی متن داخل \lr{algorithmic} باید لاتین باشد.}

\subsection{کامنت فارسی در الگوریتم}
\label{sec:code-algorithm-comment}

برای کامنت فارسی در الگوریتم، از \cmd{pc} داخل
\lr{|...|} استفاده کنید:

\begin{latin}
\begin{minted}{latex}
\begin{algorithmic}[1]
    \Require Sorted array $A$
    \While{$left \leq right$}
        \State $mid \gets (left + right) / 2$
        \If{$A[mid] = x$}
            \State \Return $mid$  \Comment{|\pc{یافتن مقدار}|}
        \EndIf
    \EndWhile
\end{algorithmic}
\end{minted}
\end{latin}

خروجی:

\begin{latin}
\begin{algorithm}
\caption{جستجو با کامنت فارسی}
\label{alg:binary-comment}
\begin{algorithmic}[1]
    \Require Sorted array $A$
    \State $left \gets 0$, $right \gets n - 1$
    \While{$left \leq right$}
        \State $mid \gets (left + right) / 2$
        \If{$A[mid] = x$}
            \State \Return $mid$  \Comment{|\pc{یافتن مقدار}|}
        \ElsIf{$A[mid] < x$}
            \State $left \gets mid + 1$  \Comment{|\pc{نیمه‌ی راست}|}
        \EndIf
    \EndWhile
\end{algorithmic}
\end{algorithm}
\end{latin}

\mnote{کامنت‌های الگوریتم با \lr{\textbackslash Comment\{|\textbackslash pc\{...\}|\}} نوشته می‌شوند.}

\subsection{دستورات \lr{algorithmic}}
\label{sec:code-algorithm-commands}

جدول \cref{tab:algorithm-commands} دستورات
اصلی \env{algorithmic} را خلاصه می‌کند.

\begin{matintable}{دستورات اصلی \lr{algorithmic}}{tab:algorithm-commands}
\begin{tblr}{
    width = \textwidth,
    colspec = {l X[l]},
    hlines, vlines,
    colsep = 8pt,
    rowsep = 4pt,
    row{1} = {font=\bfseries},
}
دستور & کاربرد \\
\lr{\textbackslash Require} & پیش‌نیاز الگوریتم \\
\lr{\textbackslash Ensure} & خروجی الگوریتم \\
\lr{\textbackslash State} & یک دستور ساده \\
\lr{\textbackslash Statex} & دستور بدون شماره \\
\lr{\textbackslash If} / \lr{\textbackslash ElsIf} / \lr{\textbackslash Else} & شرط \\
\lr{\textbackslash EndIf} & پایان شرط \\
\lr{\textbackslash For} / \lr{\textbackslash EndFor} & حلقه‌ی \lr{for} \\
\lr{\textbackslash While} / \lr{\textbackslash EndWhile} & حلقه‌ی \lr{while} \\
\lr{\textbackslash Function} / \lr{\textbackslash EndFunction} & تعریف تابع \\
\lr{\textbackslash Call\{name\}\{args\}} & فراخوانی تابع \\
\lr{\textbackslash Return} & بازگشت مقدار \\
\lr{\textbackslash Comment\{...\}} & کامنت \\
\end{tblr}
\end{matintable}

\subsection{ارجاع به الگوریتم}
\label{sec:code-algorithm-reference}

برای ارجاع به الگوریتم، از \cmd{cref} استفاده
کنید:

\begin{latin}
\begin{minted}{latex}
|\pc{الگوریتم}|\cref{alg:binary}|\pc{فرآیند جستجو را نشان می‌دهد.}|%
\end{minted}
\end{latin}

خروجی: الگوریتم \cref{alg:binary-ch10} فرآیند
جستجو را نشان می‌دهد.

\mnote{برای فهرست الگوریتم‌ها، از \lr{\textbackslash listofalgorithms} در پیش‌متن استفاده کنید.}

\section{تنظیمات پیش‌فرض \lr{minted}}
\label{sec:code-settings}

\lr{MatinBook} تنظیمات زیر را برای \lr{minted}
اعمال می‌کند:

\begin{itemize}
    \item \textbf{فونت:} \lr{JetBrains Mono} با
    اندازه‌ی \lr{footnotesize}.
    \item \textbf{شماره خطوط:} فعال
    (\lr{linenos=true}).
    \item \textbf{شکستن خطوط:} فعال
    (\lr{breaklines=true}).
    \item \textbf{کادر:} خطوط بالا و پایین
    (\lr{frame=lines}).
    \item \textbf{پس‌زمینه:} \lr{matinlightgray!30}.
    \item \textbf{اندازه Tab:} ۴ فاصله
    (\lr{tabsize=4}).
    \item \textbf{حذف تورفتگی مشترک:} فعال
    (\lr{autogobble=true}).
    \item \textbf{استایل رنگ‌آمیزی:} \lr{friendly}.
\end{itemize}

\mnote{برای شخصی‌سازی این تنظیمات، فایل \lr{mb-code.sty} را ویرایش کنید.}

\section{زبان‌های پشتیبانی‌شده}
\label{sec:code-languages}

\lr{minted} از بیش از ۳۰۰ زبان برنامه‌نویسی
پشتیبانی می‌کند. چند زبان پرکاربرد:

\begin{itemize}
    \item \textbf{برنامه‌نویسی:} \lr{python}،
    \lr{cpp}، \lr{c}، \lr{java}،
    \lr{javascript}، \lr{typescript}،
    \lr{rust}، \lr{go}.
    \item \textbf{نشانه‌گذاری:} \lr{html}،
    \lr{xml}، \lr{markdown}، \lr{latex}،
    \lr{yaml}، \lr{json}.
    \item \textbf{پوسته:} \lr{bash}، \lr{zsh}،
    \lr{powershell}.
    \item \textbf{پایگاه داده:} \lr{sql}،
    \lr{postgresql}، \lr{mysql}.
    \item \textbf{سایر:} \lr{text} (ساده)،
    \lr{diff}، \lr{makefile}، \lr{docker}.
\end{itemize}

\mnote{برای دیدن فهرست کامل زبان‌ها، از \lr{pygmentize -L lexers} استفاده کنید.}

\section{خطاهای رایج}
\label{sec:code-mistakes}

\subsection{فراموش کردن \env{latin}}
\label{sec:code-mistake-latin}

اگر کد را داخل \env{latin} قرار ندهید، ممکن
است جهت کد به‌هم بریزد و شماره‌های خطوط
جابه‌جا شوند:

\begin{latin}
\begin{minted}{latex}
% |\pc{اشتباه}|:
\begin{minted}{python}
x = 5  |\pcc{مقدار}|
\end{minted}

% |\pc{درست}|:
\begin{latin}
\begin{minted}{python}
x = 5  |\pcc{مقدار}|
\end{minted}
\end{latin}
\end{minted}
\end{latin}

\subsection{فراموش کردن \cmd{pcc} در پایتون}
\label{sec:code-mistake-pcc}

اگر در پایتون از \lr{\# |\textbackslash pc\{...\}|}
استفاده کنید، \lr{Pygments} علامت \lr{\#} را
به‌عنوان شروع کامنت می‌بیند و \lr{|} را نادیده
می‌گیرد:

\begin{latin}
\begin{minted}{latex}
% |\pc{اشتباه}|:
x = 5  # |\pc{مقدار متغیر}|

% |\pc{درست}|:
x = 5  |\pcc{مقدار متغیر}|
\end{minted}
\end{latin}

\subsection{فراموش کردن \cmd{pc} در \lr{LaTeX}}
\label{sec:code-mistake-pc}

اگر در \lr{LaTeX} متن فارسی را بدون \cmd{pc}
بنویسید، نتیجه به‌صورت مربع‌های خالی نمایش
داده می‌شود:

\begin{latin}
\begin{minted}{latex}
% |\pc{اشتباه}|:
\chapter{فصل نمونه}

% |\pc{درست}|:
\chapter{|\pc{فصل نمونه}|}
\end{minted}
\end{latin}

\section{تمرین‌ها}
\label{sec:code-exercises}

\begin{exercise}
    یک تابع پایتون با کامنت فارسی بنویسید
    و آن را با \cmd{pcc} پیچیده کنید.
    \label{exc:code-python}
\end{exercise}

\begin{exercise}
    یک الگوریتم با عنوان فارسی و متن لاتین
    بنویسید و به آن ارجاع دهید.
    \label{exc:code-algorithm}
\end{exercise}

\begin{exercise}
    یک نمونه‌ی \lr{LaTeX} با متن فارسی بنویسید
    و از \cmd{pc} استفاده کنید.
    \label{exc:code-latex}
\end{exercise}

\begin{exercise}
    یک کد \lr{C++} با کامنت فارسی بنویسید.
    \label{exc:code-cpp}
\end{exercise}

\section{خلاصه}
\label{sec:code-summary}

در این فصل، با کد و الگوریتم در \lr{MatinBook}
آشنا شدیم:

\begin{itemize}
    \item کد با محیط \env{minted} داخل
    \env{latin} نمایش داده می‌شود.

    \item برای کپشن کد، از \cmd{codecaption}
    استفاده کنید.

    \item برای کامنت فارسی در پایتون، از
    \cmd{pcc} داخل \lr{|...|} استفاده کنید.

    \item برای متن فارسی غیرکامنت در
    \lr{LaTeX} یا \lr{bash}، از \cmd{pc} استفاده
    کنید.

    \item برای کد درون‌خطی، از \cmd{inlcode}
    استفاده کنید.

    \item الگوریتم‌ها با محیط \env{algorithm}
    نمایش داده می‌شوند.

    \item کامنت فارسی در الگوریتم با
    \cmd{Comment\{|\cmd{pc}\{...\}|\}} نوشته
    می‌شود.

    \item خطاهای رایج شامل فراموش کردن
    \env{latin}، \cmd{pcc} در پایتون، و \cmd{pc}
    در \lr{LaTeX} است.
\end{itemize}

در فصل بعدی (\cref{chap:graphics})، با نمودار
و شکل در \lr{MatinBook} آشنا می‌شوید و
یاد می‌گیرید که چگونه نمودارهای حرفه‌ای
بسازید.
%%
%% part2-book-components/ch11-graphics.tex
%%
%% Chapter 11: Graphics and Diagrams — Function plots,
%% statistical charts, flowcharts, trees, graphs, 3D
%% surfaces, pie charts, and timelines.
%%

\chapter{نمودار و شکل}
\label{chap:graphics}

\section{چرا این موضوع مهم است؟}
\label{sec:graphics-why}

در کتاب‌های فنی و علمی، نمودارها و شکل‌ها نقش
محوری در انتقال مفاهیم پیچیده دارند. یک نمودار
حرفه‌ای می‌تواند جای چند صفحه متن را بگیرد و
درک مطلب را چند برابر کند.

\lr{MatinBook} از \lr{TikZ}، \lr{PGFPlots} و
\lr{pgf-pie} برای ترسیم نمودارها و شکل‌های
حرفه‌ای استفاده می‌کند. در این فصل، با انواع
نمودارها و روش ساخت آن‌ها آشنا می‌شوید.

\mnote{\lr{TikZ} یک زبان برنامه‌نویسی برای ترسیم شکل‌های برداری است. \lr{PGFPlots} نیز بر پایه‌ی \lr{TikZ} برای رسم نمودارهای داده‌ای ساخته شده است.}

\section{ساختار پایه}
\label{sec:graphics-basic}

هر شکل در \lr{MatinBook} از یک ساختار پایه
پیروی می‌کند:

\begin{itemize}
    \item محیط \env{figure} برای نگهداری شکل.
    \item \cmd{centering} برای وسط‌چین کردن.
    \item محیط \env{latin} برای حفظ جهت لاتین.
    \item محیط \env{tikzpicture} یا
    \env{axis} برای ترسیم.
    \item \cmd{caption} برای کپشن فارسی.
    \item \cmd{label} برای ارجاع.
\end{itemize}

\mnote{محیط \lr{tikzpicture} باید همیشه داخل \lr{\textbackslash begin\{latin\}} قرار گیرد تا جهت ترسیم درست حفظ شود.}

\subsection{نمونه‌ی پایه}
\label{sec:graphics-basic-example}

\begin{latin}
\begin{minted}{latex}
\begin{figure}[h]
\centering
\begin{latin}
\begin{tikzpicture}
    \draw[matinblue, thick] (0,0) -- (3,2);
    \draw[matinred, thick] (0,2) -- (3,0);
\end{tikzpicture}
\end{latin}
\caption{|\pc{یک شکل ساده}|}
\label{fig:simple}
\end{figure}
\end{minted}
\end{latin}

\mnote{دستور \lr{[h]} به \lr{LaTeX} می‌گوید که شکل را «همین‌جا» قرار دهد. اگر فضا کافی نباشد، به صفحه‌ی بعد می‌رود.}

\section{نمودار توابع}
\label{sec:graphics-functions}

پرکاربردترین نوع نمودار در کتاب‌های فنی،
نمودار توابع ریاضی است. \lr{MatinBook} از
\lr{PGFPlots} برای این کار استفاده می‌کند.

\subsection{نمودار تابع ساده}
\label{sec:graphics-simple-function}

\begin{latin}
\begin{minted}{latex}
\begin{figure}[h]
\centering
\begin{latin}
\begin{tikzpicture}
  \begin{axis}[
    width=\textwidth, height=6cm,
    xlabel={$x$},
    ylabel={$f(x)$},
    grid=major,
    domain=-3:3,
    samples=100
  ]
    \addplot[matinblue, thick] {x^2};
    \addlegendentry{$x^2$}
  \end{axis}
\end{tikzpicture}
\end{latin}
\caption{|\pc{نمودار تابع}|\lr{$f(x) = x^{2}$}}
\label{fig:parabola}
\end{figure}
\end{minted}
\end{latin}

خروجی:

\begin{figure}[h]
\centering
\begin{latin}
\begin{tikzpicture}
  \begin{axis}[
    width=0.8\textwidth, height=5cm,
    xlabel={$x$},
    ylabel={$f(x)$},
    grid=major,
    domain=-3:3,
    samples=100
  ]
    \addplot[matinblue, thick] {x^2};
    \addlegendentry{$x^2$}
  \end{axis}
\end{tikzpicture}
\end{latin}
\caption{نمودار تابع \lr{$f(x) = x^{2}$}}
\label{fig:parabola-ch11}
\end{figure}

\mnote{مقدار \lr{width=\textbackslash textwidth} باعث می‌شود نمودار دقیقاً به عرض متن چسبیده و از صفحه سرریز نکند.}

\subsection{نمودار چند تابع}
\label{sec:graphics-multiple-functions}

\begin{latin}
\begin{minted}{latex}
\begin{axis}[
    width=\textwidth, height=6cm,
    xlabel={$x$},
    ylabel={$f(x)$},
    grid=major,
    legend pos=north west,
    domain=-2:2,
    samples=100
]
    \addplot[matinblue, thick] {x^2};
    \addlegendentry{$x^2$}

    \addplot[matinred, thick] {x^3};
    \addlegendentry{$x^3$}

    \addplot[matinorange, thick] {exp(x)};
    \addlegendentry{$e^x$}
\end{axis}
\end{minted}
\end{latin}

خروجی:

\begin{figure}[h]
\centering
\begin{latin}
\begin{tikzpicture}
  \begin{axis}[
    width=0.8\textwidth, height=5cm,
    xlabel={$x$},
    ylabel={$f(x)$},
    grid=major,
    legend pos=north west,
    domain=-2:2,
    samples=100
  ]
    \addplot[matinblue, thick] {x^2};
    \addlegendentry{$x^2$}
    \addplot[matinred, thick] {x^3};
    \addlegendentry{$x^3$}
    \addplot[matinorange, thick] {exp(x)};
    \addlegendentry{$e^x$}
  \end{axis}
\end{tikzpicture}
\end{latin}
\caption{مقایسه‌ی سه تابع}
\label{fig:three-functions}
\end{figure}

\mnote{پالت رنگ \lr{matinbook} به‌طور خودکار برای نمودارهای چندسری اعمال می‌شود. می‌توانید رنگ‌ها را دستی هم مشخص کنید.}

\section{نمودارهای آماری}
\label{sec:graphics-statistical}

\subsection{نمودار میله‌ای}
\label{sec:graphics-bar}

\begin{latin}
\begin{minted}{latex}
\begin{axis}[
    width=\textwidth, height=6cm,
    ybar,
    symbolic x coords={A, B, C, D, E},
    xtick=data,
    ylabel={|\pc{مقدار}|},
    bar width=15pt,
    nodes near coords
]
    \addplot[fill=matinblue!70] coordinates {
        (A, 25) (B, 40) (C, 30) (D, 50) (E, 35)
    };
\end{axis}
\end{minted}
\end{latin}

\mnote{در \lr{PGFPlots}، برچسب‌های محور باید لاتین باشند. برای معادل فارسی، از کپشن یا متن استفاده کنید.}

\subsection{نمودار خطی}
\label{sec:graphics-line}

\begin{latin}
\begin{minted}{latex}
\begin{axis}[
    width=\textwidth, height=6cm,
    xlabel={$n$},
    ylabel={Operations},
    grid=major,
    legend pos=north west,
    domain=0:50,
    samples=100
]
    \addplot[matinred, thick] {x^2};
    \addlegendentry{$O(n^2)$}

    \addplot[matinblue, thick] {x*ln(x)};
    \addlegendentry{$O(n \log n)$}

    \addplot[matingreen, thick] {x};
    \addlegendentry{$O(n)$}
\end{axis}
\end{minted}
\end{latin}

\subsection{نمودار پراکندگی}
\label{sec:graphics-scatter}

\begin{latin}
\begin{minted}{latex}
\begin{axis}[
    width=\textwidth, height=6cm,
    xlabel={X},
    ylabel={Y},
    grid=major
]
    \addplot[only marks, mark=*, color=matinblue] coordinates {
        (1, 45) (2, 52) (3, 58) (4, 65) (5, 70)
    };
\end{axis}
\end{minted}
\end{latin}

\mnote{برای نمودار پراکندگی، از \lr{only marks} استفاده کنید تا نقاط به هم متصل نشوند.}

\section{فلوچارت‌ها}
\label{sec:graphics-flowcharts}

فلوچارت‌ها برای نمایش فرآیندها و الگوریتم‌ها
استفاده می‌شوند. \lr{MatinBook} چند استایل
آماده برای فلوچارت ارائه می‌دهد:

\begin{itemize}
    \item \lr{block}: مستطیل با گوشه‌های گرد.
    \item \lr{decision}: لوزی برای تصمیم.
    \item \lr{arrow}: فلش ضخیم.
\end{itemize}

\begin{latin}
\begin{minted}{latex}
\begin{tikzpicture}[node distance=2cm]
    \node[block] (start) {Start};
    \node[decision, below of=start] (check) {$x > 0$?};
    \node[block, below of=check] (positive) {Positive};
    \node[block, right of=check, node distance=3cm] (negative) {Negative};

    \draw[arrow] (start) -- (check);
    \draw[arrow] (check) -- node[left] {Yes} (positive);
    \draw[arrow] (check) -- node[above] {No} (negative);
\end{tikzpicture}
\end{minted}
\end{latin}

خروجی:

\begin{figure}[h]
\centering
\begin{latin}
\begin{tikzpicture}[node distance=1.5cm, scale=0.9, transform shape]
    \node[block] (start) {Start};
    \node[decision, below of=start] (check) {$x > 0$?};
    \node[block, below of=check] (positive) {Positive};
    \node[block, right of=check, node distance=3cm] (negative) {Negative};

    \draw[arrow] (start) -- (check);
    \draw[arrow] (check) -- node[left] {Yes} (positive);
    \draw[arrow] (check) -- node[above] {No} (negative);
\end{tikzpicture}
\end{latin}
\caption{فلوچارت یک الگوریتم ساده}
\label{fig:flowchart-ch11}
\end{figure}

\mnote{در فلوچارت‌ها، متن داخل گره‌ها باید لاتین باشد. برای متن فارسی، از \lr{\textbackslash rl\{...\}} استفاده کنید.}

\section{درخت‌ها}
\label{sec:graphics-trees}

درخت‌ها برای نمایش ساختارهای سلسله‌مراتبی
استفاده می‌شوند. \lr{MatinBook} یک استایل
\lr{treenode} برای گره‌های درخت ارائه می‌دهد:

\begin{latin}
\begin{minted}{latex}
\begin{tikzpicture}[
    level distance=1.5cm,
    level 1/.style={sibling distance=3cm},
    level 2/.style={sibling distance=1.5cm}
]
    \node[treenode] {8}
        child {node[treenode] {3}
            child {node[treenode] {1}}
            child {node[treenode] {6}}
        }
        child {node[treenode] {10}
            child {node[treenode] {9}}
            child {node[treenode] {14}}
        };
\end{tikzpicture}
\end{minted}
\end{latin}

خروجی:

\begin{figure}[h]
\centering
\begin{latin}
\begin{tikzpicture}[
    level distance=1.5cm,
    scale=0.9,
    transform shape,
    level 1/.style={sibling distance=3cm},
    level 2/.style={sibling distance=1.5cm}
]
    \node[treenode] {8}
        child {node[treenode] {3}
            child {node[treenode] {1}}
            child {node[treenode] {6}}
        }
        child {node[treenode] {10}
            child {node[treenode] {9}}
            child {node[treenode] {14}}
        };
\end{tikzpicture}
\end{latin}
\caption{درخت جستجوی دودویی}
\label{fig:bst-ch11}
\end{figure}

\section{گراف‌ها}
\label{sec:graphics-graphs}

گراف‌ها برای نمایش روابط بین موجودیت‌ها استفاده
می‌شوند. \lr{MatinBook} یک استایل \lr{vertex}
برای رأس‌های گراف ارائه می‌دهد:

\begin{latin}
\begin{minted}{latex}
\begin{tikzpicture}[
    vertex/.style={circle, draw=matindarkblue,
                   fill=matinlightblue!30, minimum size=1cm}
]
    \node[vertex] (A) at (0,0) {A};
    \node[vertex] (B) at (3,0) {B};
    \node[vertex] (C) at (1.5,-2) {C};

    \draw[thick] (A) -- (B);
    \draw[thick] (B) -- (C);
    \draw[thick] (A) -- (C);
\end{tikzpicture}
\end{minted}
\end{latin}

\mnote{برای گراف‌های پیچیده، از کتابخانه‌ی \lr{graphs} در \lr{TikZ} استفاده کنید که امکان تعریف خودکار گراف‌ها را فراهم می‌کند.}

\section{نمودار سه‌بعدی}
\label{sec:graphics-3d}

\lr{PGFPlots} از نمودارهای سه‌بعدی نیز پشتیبانی
می‌کند. برای این کار، باید کتابخانه‌ی
\lr{colormaps} را فعال کنید:

\begin{latin}
\begin{minted}{latex}
\usepgfplotslibrary{colormaps}

\begin{axis}[
    width=\textwidth, height=8cm,
    view={45}{35},
    xlabel={$x$}, ylabel={$y$}, zlabel={$f$}
]
    \addplot3[
        surf, domain=-5:5, domain y=-5:5,
        samples=40, samples y=40,
        colormap/viridis
    ] {sin(deg(sqrt(x^2+y^2)))};
\end{axis}
\end{minted}
\end{latin}

\mnote{نمودارهای سه‌بعدی ممکن است در بعضی از نسخه‌های \lr{PGFPlots} کند باشند. صبور باشید.}

\section{نمودار دایره‌ای}
\label{sec:graphics-pie}

\lr{MatinBook} از \lr{pgf-pie} برای نمودار
دایره‌ای پشتیبانی می‌کند:

\begin{latin}
\begin{minted}{latex}
\begin{tikzpicture}
    \pie[
        radius=3,
        text=legend,
        color={matinblue, matingreen, matinorange, matinred}
    ]{
        35/Python,
        25/JavaScript,
        20/Java,
        12/C++,
        8/Other
    }
\end{tikzpicture}
\end{minted}
\end{latin}

\mnote{در نمودار دایره‌ای، برچسب‌ها باید لاتین باشند. برای متن فارسی، از کپشن استفاده کنید.}

\section{خط زمانی}
\label{sec:graphics-timeline}

خط زمانی برای نمایش رویدادها در طول زمان
استفاده می‌شود:

\begin{latin}
\begin{minted}{latex}
\begin{tikzpicture}[xscale=0.7]
    \draw[thick,->] (0,0) -- (14,0);

    \foreach \x/\year in {0/1990, 3.5/2000, 7/2010, 10.5/2020, 14/2025} {
        \draw[thick] (\x,-0.2) -- (\x,0.2);
        \node[below] at (\x,-0.3) {\small\year};
    }

    \node[event, above] at (0,1.2) {Python\\Created};
    \node[event, above] at (7,1.2) {Rust\\First Stable};
\end{tikzpicture}
\end{minted}
\end{latin}

\section{فاصله‌گذاری و اندازه}
\label{sec:graphics-sizing}

یکی از مشکلات رایج در نمودارها، سرریز به
حاشیه است. برای جلوگیری از این مشکل:

\begin{itemize}
    \item از \lr{width=\textbackslash textwidth} برای
    \env{axis} استفاده کنید.
    \item از \lr{scale=0.9, transform shape}
    برای \env{tikzpicture} استفاده کنید.
    \item عرض ثابت (مثل \lr{width=14cm}) را
    حذف کنید.
\end{itemize}

\mnote{عرض \lr{\textbackslash textwidth} در متن اصلی حدود \lr{۱۰.۳} سانتی‌متر است. اگر عرض نمودار بیشتر باشد، سرریز می‌کند.}

\section{ارجاع به شکل}
\label{sec:graphics-references}

برای ارجاع به شکل، از \cmd{cref} استفاده کنید:

\begin{latin}
\begin{minted}{latex}
|\pc{نمودار}|\cref{fig:parabola}|\pc{تابع درجه دو را نشان می‌دهد.}|%
\end{minted}
\end{latin}

خروجی: نمودار \cref{fig:parabola-ch11} تابع
درجه دو را نشان می‌دهد.

\mnote{نام‌گذاری برچسب شکل‌ها با پیشوند \lr{fig:} انجام می‌شود: \lr{fig:parabola}.}

\section{متن فارسی در \lr{tikzpicture}}
\label{sec:graphics-persian}

برای متن فارسی داخل \env{tikzpicture}، باید از
\cmd{rl} استفاده کنید:

\begin{latin}
\begin{minted}{latex}
\begin{tikzpicture}
    \node[block] (start) {|\rl{شروع}|};
    \node[block, below of=start] (end) {|\rl{پایان}|};
    \draw[arrow] (start) -- (end);
\end{tikzpicture}
\end{minted}
\end{latin}

\mnote{فراموش کردن \lr{\textbackslash rl} باعث می‌شود متن فارسی به‌صورت چپ‌چین و نامفهوم نمایش داده شود.}

\section{خطاهای رایج}
\label{sec:graphics-mistakes}

\subsection{فراموش کردن \env{latin}}
\label{sec:graphics-mistake-latin}

اگر \env{tikzpicture} را داخل \env{latin} قرار
ندهید، ممکن است جهت ترسیم به‌هم بریزد و
شکل در جای اشتباه ظاهر شود.

\textbf{راه‌حل:} همیشه \env{tikzpicture} را داخل
\env{latin} قرار دهید.

\subsection{سرریز نمودار به حاشیه}
\label{sec:graphics-mistake-overflow}

اگر عرض نمودار بیشتر از \cmd{textwidth} باشد،
به حاشیه سرریز می‌کند.

\textbf{راه‌حل:} از \lr{width=\textbackslash textwidth}
استفاده کنید، یا مقیاس را کاهش دهید.

\subsection{فراموش کردن \cmd{rl} در \lr{tikzpicture}}
\label{sec:graphics-mistake-rl}

اگر متن فارسی را داخل \env{tikzpicture} بدون
\cmd{rl} بنویسید، به‌صورت چپ‌چین و نامفهوم
نمایش داده می‌شود.

\textbf{راه‌حل:} متن فارسی را در \lr{\textbackslash rl\{...\}}
قرار دهید.

\section{تمرین‌ها}
\label{sec:graphics-exercises}

\begin{exercise}
    یک نمودار تابع ساده بسازید که عرض آن
    دقیقاً به عرض متن چسبیده باشد.
    \label{exc:graphics-simple}
\end{exercise}

\begin{exercise}
    یک فلوچارت با سه گره بسازید و متن‌های
    فارسی آن را با \cmd{rl} پیچیده کنید.
    \label{exc:graphics-flowchart}
\end{exercise}

\begin{exercise}
    یک درخت دودویی با سه سطح بسازید.
    \label{exc:graphics-tree}
\end{exercise}

\begin{exercise}
    یک نمودار دایره‌ای با چهار بخش بسازید.
    \label{exc:graphics-pie}
\end{exercise}

\section{خلاصه}
\label{sec:graphics-summary}

در این فصل، با نمودار و شکل در \lr{MatinBook}
آشنا شدیم:

\begin{itemize}
    \item \lr{MatinBook} از \lr{TikZ}، \lr{PGFPlots}
    و \lr{pgf-pie} برای ترسیم استفاده می‌کند.

    \item نمودار توابع با \env{axis} و
    \cmd{addplot} ساخته می‌شوند.

    \item نمودارهای آماری شامل میله‌ای، خطی و
    پراکندگی هستند.

    \item فلوچارت‌ها با استایل‌های \lr{block}،
    \lr{decision} و \lr{arrow} ساخته می‌شوند.

    \item درخت‌ها با استایل \lr{treenode} ساخته
    می‌شوند.

    \item گراف‌ها با استایل \lr{vertex} ساخته
    می‌شوند.

    \item \env{tikzpicture} باید همیشه داخل
    \env{latin} قرار گیرد.

    \item متن فارسی داخل \env{tikzpicture} باید
    با \cmd{rl} پیچیده شود.

    \item خطاهای رایج شامل فراموش کردن
    \env{latin}، سرریز نمودار، و فراموش کردن
    \cmd{rl} است.
\end{itemize}

در فصل بعدی (\cref{chap:tables})، با جدول
در \lr{MatinBook} آشنا می‌شوید و یاد می‌گیرید
که چگونه جدول‌های حرفه‌ای بسازید.
%%
%% part2-book-components/ch12-tables.tex
%%
%% Chapter 12: Tables — Modern tables with tabularray,
%% the matintable wrapper, and RTL-aware column types.
%%

\chapter{جدول}
\label{chap:tables}

\section{چرا این موضوع مهم است؟}
\label{sec:tables-why}

جدول‌ها یکی از ابزارهای مهم در کتاب‌های فنی و
علمی هستند. از مقایسه‌ی الگوریتم‌ها تا نمایش
داده‌های تجربی، جدول‌ها به خواننده کمک می‌کنند
تا اطلاعات را سریع‌تر درک کند.

\lr{MatinBook} از بسته‌ی \lr{tabularray} برای
جدول‌ها استفاده می‌کند. این بسته، برخلاف
\lr{tabularx}، با \lr{xepersian} کاملاً سازگار
است و از حروف‌چینی راست‌به‌چپ پشتیبانی می‌کند.

\mnote{بسته‌ی \lr{tabularx} با \lr{xepersian} ناسازگار است و خطای \lr{Illegal pream-token} می‌دهد. هرگز از آن استفاده نکنید.}

\section{محیط \lr{matintable}}
\label{sec:tables-matintable}

ساده‌ترین راه برای ساخت جدول در \lr{MatinBook}،
استفاده از محیط \env{matintable} است. این
محیط، بخش‌های تکراری جدول را خودکار می‌کند:

\begin{latin}
\begin{minted}{latex}
\begin{matintable}{|\pc{عنوان جدول}|}{tab:my-table}
\begin{tblr}{
    width = \textwidth,
    colspec = {l l X[c] X[c]},
    hlines, vlines,
    colsep = 8pt,
    rowsep = 4pt,
    row{1} = {font=\bfseries},
}
|\pc{نوع}| & |\pc{فونت}| & |\pc{وضعیت}| & |\pc{توضیحات}| \\
|\pc{فارسی}| & \lr{XB Niloofar} & |\pc{فعال}| & |\pc{خوانایی بالا}| \\
|\pc{لاتین}| & \lr{Times New Roman} & |\pc{فعال}| & |\pc{استاندارد آکادمیک}| \\
\end{tblr}
\end{matintable}
\end{minted}
\end{latin}

خروجی:

\begin{matintable}{نمونه‌ی جدول با \lr{matintable}}{tab:sample}
\begin{tblr}{
    width = \textwidth,
    colspec = {l l X[c] X[c]},
    hlines, vlines,
    colsep = 8pt,
    rowsep = 4pt,
    row{1} = {font=\bfseries},
}
نوع & فونت & وضعیت & توضیحات \\
فارسی & \lr{XB Niloofar} & فعال & خوانایی بالا \\
لاتین & \lr{Times New Roman} & فعال & استاندارد آکادمیک \\
\end{tblr}
\end{matintable}

\mnote{دستور \lr{\textbackslash matintable} سه بخش تکراری را خودکار می‌کند: وسط‌چین کردن، کپشن، و برچسب.}

\section{ساختار \lr{tblr}}
\label{sec:tables-tblr}

محیط \env{tblr} داخل \env{matintable} قرار
می‌گیرد و تنظیمات جدول را تعیین می‌کند:

\begin{itemize}
    \item \lr{width = \textbackslash textwidth}:
    عرض کل جدول.
    \item \lr{colspec = \{...\}}: نوع ستون‌ها.
    \item \lr{hlines, vlines}: خطوط افقی و
    عمودی.
    \item \lr{colsep = 8pt}: فاصله‌ی داخلی
    ستون‌ها.
    \item \lr{rowsep = 4pt}: فاصله‌ی داخلی
    ردیف‌ها.
    \item \lr{row\{1\} = \{font=\textbackslash bfseries\}}:
    ردیف اول پررنگ.
\end{itemize}

\mnote{تنظیمات \lr{tblr} با کاما از هم جدا می‌شوند. ترتیب آن‌ها مهم نیست.}

\section{انواع ستون}
\label{sec:tables-columns}

\lr{MatinBook} از سه نوع ستون پشتیبانی می‌کند:

\begin{itemize}
    \item \lr{l}: ستون چپ‌چین (برای متن
    لاتین).
    \item \lr{c}: ستون وسط‌چین (برای اعداد
    یا نمادها).
    \item \lr{X[c]}: ستون خودکار با تراز
    وسط (برای متن فارسی).
\end{itemize}

\subsection{ستون \lr{X} در \lr{tabularray}}
\label{sec:tables-x-column}

ستون \lr{X} یک ستون \textbf{خودکار} است که
عرض آن به‌طور خودکار محاسبه می‌شود. می‌توانید
تراز آن را با \lr{[c]}، \lr{[l]} یا \lr{[r]}
تعیین کنید:

\begin{latin}
\begin{minted}{latex}
colspec = {l l X[c] X[c]}
\end{minted}
\end{latin}

\mnote{ستون \lr{X} فقط در \lr{tabularray} وجود دارد و در \lr{tabular} معمولی کار نمی‌کند.}

\subsection{ستون‌های \lr{p} (برای جدول‌های قدیمی)}
\label{sec:tables-p-columns}

\lr{MatinBook} سه نوع ستون \lr{p} هم ارائه می‌دهد
که برای \env{tabular} معمولی کار می‌کنند:

\begin{itemize}
    \item \lr{L\{width\}}: ستون چپ‌چین.
    \item \lr{R\{width\}}: ستون راست‌چین.
    \item \lr{C\{width\}}: ستون وسط‌چین.
\end{itemize}

\mnote{این ستون‌ها فقط برای سازگاری با کدهای قدیمی نگه داشته شده‌اند. برای جدول‌های جدید، از \lr{tblr} و ستون \lr{X} استفاده کنید.}

\section{عرض جدول}
\label{sec:tables-width}

برای اینکه جدول دقیقاً به عرض متن بچسبد، از
\lr{width = \textbackslash textwidth} استفاده
کنید:

\begin{latin}
\begin{minted}{latex}
\begin{tblr}{
    width = \textwidth,
    colspec = {l l X[c] X[c]},
    ...
}
\end{minted}
\end{latin}

\mnote{اگر از \lr{width} استفاده نکنید، عرض جدول به‌طور خودکار محاسبه می‌شود و ممکن است بیشتر از عرض متن شود.}

\subsection{جدول‌های باریک}
\label{sec:tables-narrow}

برای جدول‌های باریک (مثلاً ۲ ستون)، می‌توانید
از عرض کمتر استفاده کنید:

\begin{latin}
\begin{minted}{latex}
\begin{tblr}{
    width = 0.5\textwidth,
    colspec = {c c},
    ...
}
\end{minted}
\end{latin}

\section{خطوط جدول}
\label{sec:tables-rules}

\lr{tabularray} از سه نوع خط پشتیبانی می‌کند:

\begin{itemize}
    \item \lr{hlines}: خطوط افقی بین همه‌ی
    ردیف‌ها.
    \item \lr{vlines}: خطوط عمودی بین همه‌ی
    ستون‌ها.
    \item \lr{hlines, vlines}: هر دو.
\end{itemize}

\subsection{خطوط انتخابی}
\label{sec:tables-selective-rules}

می‌توانید خطوط را برای ردیف یا ستون خاصی تعیین
کنید:

\begin{latin}
\begin{minted}{latex}
\begin{tblr}{
    hline{1,Z} = {1pt},       % |\pc{خط ضخیم در بالا و پایین}|%
    hline{2} = {0.5pt},        % |\pc{خط نازک بعد از ردیف اول}|%
    ...
}
\end{minted}
\end{latin}

\mnote{\lr{Z} به معنای «آخرین ردیف» است. \lr{1,Z} یعنی ردیف اول و آخر.}

\section{رنگ در جدول}
\label{sec:tables-color}

\lr{tabularray} از رنگ‌آمیزی سلول‌ها پشتیبانی
می‌کند:

\begin{latin}
\begin{minted}{latex}
\begin{tblr}{
    row{1} = {bg=matinblue!20},    % |\pc{پس‌زمینه‌ی ردیف اول}|%
    cell{2}{3} = {bg=matingreen!20}, % |\pc{سلول خاص}|%
    ...
}
\end{minted}
\end{latin}

\mnote{برای رنگ‌آمیزی، از پالت \lr{MatinBook} استفاده کنید: \lr{matinblue}، \lr{matingreen}، \lr{matinorange} و...}

\subsection{مثال کامل با رنگ}
\label{sec:tables-color-example}

\begin{matintable}{مقایسه‌ی الگوریتم‌ها با رنگ}{tab:colored}
\begin{tblr}{
    width = \textwidth,
    colspec = {l X[c] X[c]},
    hlines, vlines,
    colsep = 8pt,
    rowsep = 4pt,
    row{1} = {font=\bfseries, bg=matinblue!20},
}
الگوریتم & پیچیدگی & کاربرد \\
مرتب‌سازی سریع & \lr{$O(n \log n)$} & \textcolor{matingreen}{سریع} \\
مرتب‌سازی ادغامی & \lr{$O(n \log n)$} & \textcolor{matingreen}{پایدار} \\
مرتب‌سازی حبابی & \lr{$O(n^2)$} & \textcolor{matinred}{کند} \\
\end{tblr}
\end{matintable}

\section{ادغام سلول‌ها}
\label{sec:tables-merge}

\lr{tabularray} از ادغام سلول‌ها پشتیبانی می‌کند:

\begin{latin}
\begin{minted}{latex}
\begin{tblr}{
    cell{1}{1} = {r=2}{c},   % |\pc{ادغام دو ردیف}|%
    cell{1}{2} = {c=2}{c},   % |\pc{ادغام دو ستون}|%
}
\end{minted}
\end{latin}

\mnote{در \lr{tabularray}، ادغام سلول‌ها با \lr{r} (ردیف) و \lr{c} (ستون) انجام می‌شود.}

\section{جدول‌های چندصفحه‌ای}
\label{sec:tables-multipage}

برای جدول‌های طولانی که بیش از یک صفحه هستند،
از \lr{longtblr} استفاده کنید:

\begin{latin}
\begin{minted}{latex}
\begin{longtblr}{
    width = \textwidth,
    colspec = {l l X[c]},
    hlines, vlines,
    rowhead = 1,
}
|\pc{نوع}| & |\pc{مقدار}| & |\pc{توضیحات}| \\
...
\end{longtblr}
\end{minted}
\end{latin}

\mnote{گزینه‌ی \lr{rowhead = 1} باعث می‌شود ردیف اول (سرصفحه) در هر صفحه‌ی جدید تکرار شود.}

\section{جدول‌های فارسی و لاتین}
\label{sec:tables-persian-latin}

در جدول‌های \lr{MatinBook}، می‌توانید فارسی و
لاتین را با هم ترکیب کنید:

\begin{latin}
\begin{minted}{latex}
\begin{matintable}{|\pc{ترکیب فارسی و لاتین}|}{tab:mixed}
\begin{tblr}{
    width = \textwidth,
    colspec = {l X[c]},
    hlines, vlines,
    colsep = 8pt,
    rowsep = 4pt,
    row{1} = {font=\bfseries},
}
|\pc{مفهوم}| & |\pc{معادل لاتین}| \\
|\pc{الگوریتم}| & \lr{Algorithm} \\
|\pc{برنامه‌نویسی}| & \lr{Programming} \\
\end{tblr}
\end{matintable}
\end{minted}
\end{latin}

خروجی:

\begin{matintable}{ترکیب فارسی و لاتین}{tab:mixed}
\begin{tblr}{
    width = \textwidth,
    colspec = {l X[c]},
    hlines, vlines,
    colsep = 8pt,
    rowsep = 4pt,
    row{1} = {font=\bfseries},
}
مفهوم & معادل لاتین \\
الگوریتم & \lr{Algorithm} \\
برنامه‌نویسی & \lr{Programming} \\
\end{tblr}
\end{matintable}

\mnote{در جدول‌ها، متن لاتین باید داخل \lr{\textbackslash lr\{...\}} قرار گیرد. متن فارسی نیازی به \lr{\textbackslash lr} ندارد.}

\section{جدول در حالت تمام‌عرض}
\label{sec:tables-fullwidth}

اگر جدول شما پهن است و به فضای بیشتری نیاز
دارد، می‌توانید از عرض کامل متن استفاده کنید:

\begin{latin}
\begin{minted}{latex}
\begin{tblr}{
    width = \textwidth,
    colspec = {X[l] X[c] X[c] X[c] X[c]},
    ...
}
\end{minted}
\end{latin}

\mnote{عرض \lr{\textbackslash textwidth} در متن اصلی حدود \lr{۱۰.۳} سانتی‌متر است. در پیش‌متن و پس‌متن، این عرض بیشتر است.}

\section{ارجاع به جدول}
\label{sec:tables-references}

برای ارجاع به جدول، از \cmd{cref} استفاده کنید:

\begin{latin}
\begin{minted}{latex}
|\pc{جدول}|\cref{tab:sample}|\pc{نمونه‌ای از جدول‌های}|\lr{MatinBook}|\pc{است.}|%
\end{minted}
\end{latin}

خروجی: جدول \cref{tab:sample} نمونه‌ای از جدول‌های
\lr{MatinBook} است.

\mnote{نام‌گذاری برچسب جدول‌ها با پیشوند \lr{tab:} انجام می‌شود: \lr{tab:sample}.}

\section{خطاهای رایج}
\label{sec:tables-mistakes}

\subsection{استفاده از \lr{tabularx}}
\label{sec:tables-mistake-tabularx}

\lr{tabularx} با \lr{xepersian} ناسازگار است و
خطای زیر را می‌دهد:

\begin{latin}
\begin{minted}{text}
Package array: Illegal pream-token
  (\TX@col@width): `c' used.
\end{minted}
\end{latin}

\textbf{راه‌حل:} به‌جای \lr{tabularx} از
\lr{tabularray} استفاده کنید.

\subsection{عرض ثابت در ستون}
\label{sec:tables-mistake-fixed-width}

اگر از عرض ثابت (مثل \lr{C\{3cm\}}) استفاده
کنید، ممکن است جدول سرریز کند.

\textbf{راه‌حل:} از ستون \lr{X[c]} استفاده کنید
که عرض آن خودکار است.

\subsection{فراموش کردن \cmd{lr} در سلول‌های لاتین}
\label{sec:tables-mistake-lr}

اگر متن لاتین در سلول جدول را بدون \cmd{lr}
بنویسید، به‌صورت نامفهوم نمایش داده می‌شود:

\begin{latin}
\begin{minted}{latex}
% |\pc{اشتباه}|:
|\pc{الگوریتم}| & Algorithm \\

% |\pc{درست}|:
|\pc{الگوریتم}| & \lr{Algorithm} \\
\end{minted}
\end{latin}

\section{تمرین‌ها}
\label{sec:tables-exercises}

\begin{exercise}
    یک جدول دو ستونی بسازید که عرض آن دقیقاً
    به عرض متن بچسبد.
    \label{exc:tables-basic}
\end{exercise}

\begin{exercise}
    یک جدول با چهار ستون بسازید که دو ستون
    آن لاتین و دو ستون فارسی باشد.
    \label{exc:tables-mixed}
\end{exercise}

\begin{exercise}
    یک جدول با رنگ‌آمیزی ردیف اول و یک سلول
    خاص بسازید.
    \label{exc:tables-color}
\end{exercise}

\begin{exercise}
    یک جدول طولانی بسازید که به دو صفحه
    تقسیم شود و از \lr{longtblr} استفاده کنید.
    \label{exc:tables-long}
\end{exercise}

\section{خلاصه}
\label{sec:tables-summary}

در این فصل، با جدول در \lr{MatinBook} آشنا
شدیم:

\begin{itemize}
    \item \lr{MatinBook} از \lr{tabularray}
    استفاده می‌کند که با \lr{xepersian} سازگار
    است.

    \item محیط \env{matintable} بخش‌های تکراری
    جدول را خودکار می‌کند.

    \item محیط \env{tblr} تنظیمات جدول را تعیین
    می‌کند: عرض، ستون‌ها، خطوط، فاصله‌ها.

    \item سه نوع ستون: \lr{l}، \lr{c}، \lr{X[c]}.

    \item برای عرض دقیق، از
    \lr{width = \textbackslash textwidth} استفاده
    کنید.

    \item خطوط جدول با \lr{hlines} و \lr{vlines}
    تعیین می‌شوند.

    \item رنگ‌آمیزی با \lr{bg=...} انجام می‌شود.

    \item جدول‌های طولانی با \lr{longtblr} ساخته
    می‌شوند.

    \item خطاهای رایج شامل استفاده از
    \lr{tabularx}، عرض ثابت، و فراموش کردن
    \cmd{lr} است.
\end{itemize}

در فصل بعدی (\cref{chap:margin-notes})، با
یادداشت حاشیه‌ای در \lr{MatinBook} آشنا
می‌شوید و یاد می‌گیرید که چگونه از فضای
حاشیه‌ی پهن استفاده کنید.
%%
%% part2-book-components/ch13-margin-notes.tex
%%
%% Chapter 13: Margin Notes — RTL margin notes and
%% margin figures in the wide outer margin.
%%

\chapter{یادداشت حاشیه}
\label{chap:margin-notes}

\section{چرا این موضوع مهم است؟}
\label{sec:margin-notes-why}

یکی از ویژگی‌های متمایزکننده‌ی کتاب‌های درسی
حرفه‌ای، استفاده از \textbf{حاشیه‌ی بیرونی پهن}
برای یادداشت‌های توضیحی، شکل‌های کوچک و
فرمول‌های مکمل است. این طراحی، از کتاب‌های
\lr{Boyer} و \lr{Stewart} الهام گرفته شده
است.

\lr{MatinBook} یک حاشیه‌ی بیرونی پهن (\pnum{۵.۹}
سانتی‌متر) ارائه می‌دهد که مخصوص این کار طراحی
شده است. در این فصل، با دو دستور \cmd{mnote}
و \cmd{marginfig} آشنا می‌شوید.

\mnote{حاشیه‌ی پهن در \lr{frontmatter} و \lr{backmatter} حذف می‌شود و فقط در \lr{mainmatter} فعال است.}

\section{دستور \cmd{mnote}}
\label{sec:margin-notes-mnote}

دستور \cmd{mnote} برای افزودن یک یادداشت متنی
ساده در حاشیه استفاده می‌شود:

\begin{latin}
\begin{minted}{latex}
|\pc{این یک متن اصلی است.}|\mnote{|\pc{این یک یادداشت حاشیه‌ای است.}|}%
|\pc{و این ادامه‌ی متن اصلی است.}|%
\end{minted}
\end{latin}

خروجی:

این یک متن اصلی است.\mnote{این یک یادداشت حاشیه‌ای است.}
و این ادامه‌ی متن اصلی است.

\mnote{یادداشت‌های حاشیه‌ای باید کوتاه و مکمل متن اصلی باشند. طول ایده‌آل: ۱ تا ۳ خط.}

\subsection{ویژگی‌های \cmd{mnote}}
\label{sec:margin-notes-mnote-features}

دستور \cmd{mnote} ویژگی‌های زیر را دارد:

\begin{itemize}
    \item \textbf{بدون شماره:} یادداشت‌ها
    شماره‌گذاری نمی‌شوند (برخلاف پانویس).
    \item \textbf{راست‌به‌چپ:} متن یادداشت
    به‌طور خودکار راست‌به‌چپ است.
    \item \textbf{حاشیه‌ی بیرونی:} در صفحات فرد،
    در سمت راست؛ در صفحات زوج، در سمت چپ.
    \item \textbf{اندازه‌ی کوچک:} متن یادداشت با
    \cmd{footnotesize} نمایش داده می‌شود.
\end{itemize}

\mnote{یادداشت‌های حاشیه‌ای با \lr{\textbackslash marginpar} پیاده‌سازی شده‌اند که بخشی از هسته‌ی \lr{LaTeX} است.}

\section{دستور \cmd{marginfig}}
\label{sec:margin-notes-marginfig}

دستور \cmd{marginfig} برای افزودن یک شکل کوچک
با کپشن در حاشیه استفاده می‌شود:

\begin{latin}
\begin{minted}{latex}
|\pc{این یک متن اصلی است.}|%
\marginfig{figures/plot.png}{|\pc{نمودار نمونه}|}%
\end{minted}
\end{latin}

این دستور دو آرگومان می‌گیرد:

\begin{enumerate}
    \item \textbf{نام فایل:} مسیر فایل تصویر.
    \item \textbf{کپشن:} متن توضیح زیر شکل.
\end{enumerate}

\mnote{عرض شکل در حاشیه، به‌طور خودکار با \lr{\textbackslash marginparwidth} تنظیم می‌شود که \lr{۳.۵} سانتی‌متر است.}

\section{محدودیت‌ها}
\label{sec:margin-notes-limitations}

یادداشت‌های حاشیه‌ای چند محدودیت دارند:

\subsection{عدم استفاده در \env{tcolorbox}}
\label{sec:margin-notes-limitation-tcolorbox}

دستور \cmd{mnote} \textbf{نمی‌تواند} داخل
\env{tcolorbox} (مثل \env{theorem}، \env{definition}
و...) استفاده شود، چون \cmd{marginpar} در
محیط‌های شناور کار نمی‌کند.

\textbf{راه‌حل:} یادداشت را بیرون از جعبه قرار دهید.

\subsection{عدم استفاده در \env{figure}}
\label{sec:margin-notes-limitation-figure}

دستور \cmd{mnote} \textbf{نمی‌تواند} داخل
\env{figure} استفاده شود.

\textbf{راه‌حل:} یادداشت را در متن اصلی، بیرون
از \env{figure}، قرار دهید.

\subsection{نیاز به هندسه‌ی \lr{mainmatter}}
\label{sec:margin-notes-limitation-geometry}

یادداشت‌های حاشیه‌ای فقط در \lr{mainmatter} کار
می‌کنند، چون در \lr{frontmatter} و \lr{backmatter}
حاشیه‌ی پهن حذف می‌شود.

\textbf{راه‌حل:} از \cmd{mainmattergeometry} برای
فعال‌سازی حاشیه‌ی پهن استفاده کنید.

\subsection{کامپایل چندباره}
\label{sec:margin-notes-limitation-compile}

یادداشت‌های حاشیه‌ای برای موقعیت‌یابی پایدار،
حداقل \textbf{۲ بار کامپایل} نیاز دارند.

\textbf{راه‌حل:} چرخه‌ی کامل کامپایل را اجرا
کنید.

\mnote{اگر یادداشت‌ها در جای اشتباه ظاهر شدند، یک بار دیگر کامپایل کنید.}

\section{جهت متن در صفحات زوج و فرد}
\label{sec:margin-notes-direction}

یادداشت‌های حاشیه‌ای در \lr{MatinBook} در
هر دو نوع صفحه (فرد و زوج) راست‌به‌چپ هستند.
این کار با استفاده از \cmd{RL} از \lr{xepersian}
انجام می‌شود.

جدول \cref{tab:margin-direction} رفتار یادداشت‌ها
را در صفحات مختلف نشان می‌دهد.

\begin{matintable}{رفتار یادداشت‌های حاشیه‌ای}{tab:margin-direction}
\begin{tblr}{
    width = \textwidth,
    colspec = {l X[c] X[c]},
    hlines, vlines,
    colsep = 8pt,
    rowsep = 4pt,
    row{1} = {font=\bfseries},
}
نوع صفحه & سمت حاشیه & جهت متن \\
فرد (راست) & راست & راست‌به‌چپ \\
زوج (چپ) & چپ & راست‌به‌چپ \\
\end{tblr}
\end{matintable}

\mnote{در \lr{MatinBook}، از \lr{\textbackslash reversemarginpar} استفاده نمی‌شود، چون جهت را دو بار معکوس می‌کند.}

\section{مثال‌های کاربردی}
\label{sec:margin-notes-examples}

\subsection{یادداشت توضیحی}
\label{sec:margin-notes-example-explanation}

یادداشت‌ها برای توضیحات مکمل مناسب هستند:

\begin{latin}
\begin{minted}{latex}
|\pc{الگوریتم دایکسترا برای یافتن کوتاه‌ترین مسیر در گراف‌های وزن‌دار استفاده می‌شود.}|%
\mnote{|\pc{پیچیدگی این الگوریتم}|\lr{$O((V+E) \log V)$}|\pc{است.}|}%
\end{minted}
\end{latin}

خروجی:

الگوریتم دایکسترا برای یافتن کوتاه‌ترین مسیر در
گراف‌های وزن‌دار استفاده می‌شود.\mnote{پیچیدگی این الگوریتم \lr{$O((V+E) \log V)$} است.}

\subsection{یادداشت هشدار}
\label{sec:margin-notes-example-warning}

یادداشت‌ها برای هشدارها نیز مناسب هستند:

\begin{latin}
\begin{minted}{latex}
|\pc{تقسیم بر صفر تعریف نشده است.}|%
\mnote{|\pc{هنگام ساده‌سازی عبارات کسری، همواره فرض کنید مخرج غیرصفر است.}|}%
\end{minted}
\end{latin}

خروجی:

تقسیم بر صفر تعریف نشده است.\mnote{هنگام ساده‌سازی عبارات کسری، همواره فرض کنید مخرج غیرصفر است.}

\subsection{یادداشت ارجاع}
\label{sec:margin-notes-example-reference}

یادداشت‌ها برای ارجاع به منابع نیز مناسب هستند:

\begin{latin}
\begin{minted}{latex}
|\pc{این قضیه اولین بار توسط اویلر اثبات شد.}|%
\mnote{|\pc{برای مطالعه‌ی بیشتر، به}|\cite{knuth1984}|\pc{مراجعه کنید.}|}%
\end{minted}
\end{latin}

\section{شخصی‌سازی}
\label{sec:margin-notes-customization}

برای شخصی‌سازی یادداشت‌های حاشیه‌ای، می‌توانید
پارامترهای زیر را در فایل \file{mb-layout.sty}
تغییر دهید:

\begin{matintable}{پارامترهای قابل تنظیم یادداشت حاشیه}{tab:margin-params}
\begin{tblr}{
    width = \textwidth,
    colspec = {l X[l] X[c]},
    hlines, vlines,
    colsep = 6pt,
    rowsep = 4pt,
    row{1} = {font=\bfseries},
}
پارامتر & توضیح & مقدار پیش‌فرض \\
\lr{marginparwidth} & عرض ناحیه‌ی یادداشت & \pnum{۳.۵} سانتی‌متر \\
\lr{marginparsep} & فاصله از متن اصلی & \pnum{۰.۶} سانتی‌متر \\
\lr{outer} & حاشیه‌ی بیرونی کل & \pnum{۵.۹} سانتی‌متر \\
\end{tblr}
\end{matintable}

\mnote{تغییر این پارامترها، بر چیدمان کل کتاب تأثیر می‌گذارد. با احتیاط عمل کنید.}

\section{تفاوت با پانویس}
\label{sec:margin-notes-vs-footnote}

یادداشت‌های حاشیه‌ای با پانویس تفاوت‌هایی
دارند. جدول \cref{tab:margin-vs-footnote} این
تفاوت‌ها را نشان می‌دهد.

\begin{matintable}{تفاوت یادداشت حاشیه و پانویس}{tab:margin-vs-footnote}
\begin{tblr}{
    width = \textwidth,
    colspec = {l X[c] X[c]},
    hlines, vlines,
    colsep = 6pt,
    rowsep = 4pt,
    row{1} = {font=\bfseries},
}
ویژگی & یادداشت حاشیه & پانویس \\
موقعیت & حاشیه‌ی صفحه & پایین صفحه \\
شماره‌گذاری & ندارد & دارد \\
طول & کوتاه (۱-۳ خط) & متوسط (۱-۵ خط) \\
کاربرد & توضیح مکمل & ارجاع یا توضیح \\
دستور & \cmd{mnote} & \cmd{footnote} \\
\end{tblr}
\end{matintable}

\mnote{از یادداشت حاشیه برای توضیحات کوتاه و مکمل استفاده کنید، و از پانویس برای ارجاعات و توضیحات بلندتر.}

\section{خطاهای رایج}
\label{sec:margin-notes-mistakes}

\subsection{استفاده در \env{tcolorbox}}
\label{sec:margin-notes-mistake-tcolorbox}

\textbf{مشکل:} \cmd{mnote} داخل \env{tcolorbox}
(مثل \env{theorem}، \env{definition}) کار نمی‌کند
و خطا می‌دهد.

\textbf{راه‌حل:} یادداشت را بیرون از جعبه قرار
دهید.

\subsection{استفاده در \env{figure}}
\label{sec:margin-notes-mistake-figure}

\textbf{مشکل:} \cmd{mnote} داخل \env{figure} کار
نمی‌کند.

\textbf{راه‌حل:} یادداشت را در متن اصلی قرار
دهید.

\subsection{یادداشت بسیار طولانی}
\label{sec:margin-notes-mistake-long}

\textbf{مشکل:} اگر یادداشت بیش از ۳ خط باشد،
ممکن است از حاشیه سرریز کند یا به صفحه‌ی بعد
برود.

\textbf{راه‌حل:} یادداشت‌ها را کوتاه نگه دارید
(حداکثر ۳ خط). برای توضیحات بلندتر، از متن اصلی
یا پانویس استفاده کنید.

\subsection{فراموش کردن هندسه‌ی \lr{mainmatter}}
\label{sec:margin-notes-mistake-geometry}

\textbf{مشکل:} اگر \cmd{mainmattergeometry} را
فراخوانی نکنید، حاشیه‌ی پهن وجود ندارد و
یادداشت‌ها در جای اشتباه ظاهر می‌شوند.

\textbf{راه‌حل:} همیشه پس از \cmd{mainmatter}،
دستور \cmd{mainmattergeometry} را فراخوانی کنید.

\section{تمرین‌ها}
\label{sec:margin-notes-exercises}

\begin{exercise}
    یک متن ساده با دو یادداشت حاشیه‌ای بنویسید.
    آیا یادداشت‌ها در حاشیه‌ی درست ظاهر
    می‌شوند؟
    \label{exc:margin-notes-basic}
\end{exercise}

\begin{exercise}
    یک یادداشت هشدار در حاشیه بنویسید که
    به یک بخش دیگر کتاب ارجاع دهد.
    \label{exc:margin-notes-warning}
\end{exercise}

\begin{exercise}
    یک \env{theorem} بنویسید و یک یادداشت
    حاشیه‌ای در بیرون آن (نه داخلش) قرار
    دهید.
    \label{exc:margin-notes-outside}
\end{exercise}

\begin{exercise}
    یک یادداشت بسیار طولانی بنویسید و ببینید
    آیا سرریز می‌کند. سپس آن را کوتاه کنید.
    \label{exc:margin-notes-long}
\end{exercise}

\section{خلاصه}
\label{sec:margin-notes-summary}

در این فصل، با یادداشت حاشیه‌ای در \lr{MatinBook}
آشنا شدیم:

\begin{itemize}
    \item \lr{MatinBook} یک حاشیه‌ی بیرونی پهن
    (\pnum{۵.۹} سانتی‌متر) برای یادداشت‌ها ارائه
    می‌دهد.

    \item دو دستور برای یادداشت:
    \cmd{mnote} (متنی) و \cmd{marginfig} (شکل).

    \item یادداشت‌ها بدون شماره هستند و جهت
    راست‌به‌چپ دارند.

    \item یادداشت‌ها در حاشیه‌ی بیرونی قرار
    می‌گیرند (راست در صفحات فرد، چپ در
    صفحات زوج).

    \item محدودیت‌ها: عدم استفاده در
    \env{tcolorbox} و \env{figure}، نیاز به
    \cmd{mainmattergeometry}، و کامپایل چندباره.

    \item یادداشت‌ها با پانویس تفاوت دارند:
    بدون شماره، کوتاه‌تر، و در حاشیه.

    \item خطاهای رایج شامل استفاده در
    \env{tcolorbox}، استفاده در \env{figure}،
    یادداشت بسیار طولانی، و فراموش کردن
    \cmd{mainmattergeometry} است.
\end{itemize}

در فصل بعدی (\cref{chap:bibliography})، با
منابع و کتاب‌نامه در \lr{MatinBook} آشنا
می‌شوید و یاد می‌گیرید که چگونه به منابع
ارجاع دهید.
%%
%% part2-book-components/ch14-bibliography.tex
%%
%% Chapter 14: Bibliography — biblatex + biber, citations,
%% and the bibliography display.
%%

\chapter{منابع و کتاب‌نامه}
\label{chap:bibliography}

\section{چرا این موضوع مهم است؟}
\label{sec:bibliography-why}

استناددهی یکی از ارکان اساسی نگارش علمی است.
هر ادعای علمی باید به منبع معتبر ارجاع داده شود
و خواننده باید بتواند به‌راحتی منابع را پیدا
کند.

\lr{MatinBook} از \lr{biblatex} با موتور
\lr{biber} برای مدیریت منابع استفاده می‌کند.
این ترکیب، استاندارد مدرن \lr{LaTeX} است و از
تمام قابلیت‌های پیشرفته‌ی کتاب‌نامه پشتیبانی
می‌کند.

\mnote{موتور \lr{biber} جایگزین \lr{bibtex} قدیمی است و از \lr{Unicode} و فارسی پشتیبانی می‌کند.}

\section{فایل \lr{.bib}}
\label{sec:bibliography-bibfile}

منابع در یک فایل جداگانه با پسوند \file{.bib}
ذخیره می‌شوند. هر منبع، یک ورودی \lr{BibTeX}
دارد:

\begin{latin}
\begin{minted}{bibtex}
@book{knuth1984,
    author    = {Donald E. Knuth},
    title     = {The {\TeX}book},
    year      = {1984},
    publisher = {Addison-Wesley},
    address   = {Reading, MA}
}

@article{einstein1905,
    author  = {Albert Einstein},
    title   = {On the Electrodynamics of Moving Bodies},
    journal = {Annalen der Physik},
    volume  = {17},
    pages   = {891--921},
    year    = {1905}
}
\end{minted}
\end{latin}

\mnote{کلید هر منبع (مثل \lr{knuth1984}) باید یکتا باشد. برای استناد، از همین کلید استفاده می‌کنید.}

\section{بارگذاری فایل \lr{.bib}}
\label{sec:bibliography-loading}

برای بارگذاری فایل \lr{.bib}، از دستور
\cmd{addbibresource} در پیش از
\cmd{begin\{document\}} استفاده کنید:

\begin{latin}
\begin{minted}{latex}
\documentclass{matinbook}

\addbibresource{references.bib}

\begin{document}
...
\end{document}
\end{minted}
\end{latin}

\mnote{فایل \lr{.bib} باید در همان پوشه‌ی \lr{main.tex} یا در مسیر مشخص‌شده باشد.}

\section{استناد}
\label{sec:bibliography-citation}

\lr{biblatex} چهار دستور اصلی برای استناد ارائه
می‌دهد:

\begin{itemize}
    \item \cmd{cite\{key\}}: استناد ساده.
    \item \cmd{parencite\{key\}}: استناد پرانتزی.
    \item \cmd{textcite\{key\}}: استناد متنی.
    \item \cmd{autocite\{key\}}: استناد خودکار.
\end{itemize}

\subsection{استناد ساده}
\label{sec:bibliography-cite}

\begin{latin}
\begin{minted}{latex}
|\pc{کتاب}|\cite{knuth1984}|\pc{مرجع استاندارد}|\lr{\TeX}|\pc{است.}|%
\end{minted}
\end{latin}

خروجی: کتاب \cite{knuth1984} مرجع استاندارد
\lr{\TeX} است.

\mnote{دستور \lr{\textbackslash cite} شماره‌ی منبع را در براکت نمایش می‌دهد: \lr{[1]}.}

\subsection{استناد پرانتزی}
\label{sec:bibliography-parencite}

\begin{latin}
\begin{minted}{latex}
|\pc{این الگوریتم در}|\parencite{cormen2009}|\pc{معرفی شده است.}|%
\end{minted}
\end{latin}

خروجی: این الگوریتم در \parencite{cormen2009}
معرفی شده است.

\mnote{در سبک عددی، \lr{\textbackslash parencite} و \lr{\textbackslash cite} یکسان نمایش داده می‌شوند.}

\subsection{استناد متنی}
\label{sec:bibliography-textcite}

\begin{latin}
\begin{minted}{latex}
\textcite{cormen2009}|\pc{این الگوریتم را معرفی کرد.}|%
\end{minted}
\end{latin}

خروجی: \textcite{cormen2009} این الگوریتم را
معرفی کرد.

\mnote{دستور \lr{\textbackslash textcite} نام نویسنده را در متن و شماره‌ی منبع را در براکت نمایش می‌دهد.}

\subsection{استناد به چند منبع}
\label{sec:bibliography-multiple}

\begin{latin}
\begin{minted}{latex}
|\pc{منابع}|\cite{knuth1984,lamport1994}|\pc{برای مطالعه توصیه می‌شوند.}|%
\end{minted}
\end{latin}

خروجی: منابع \cite{knuth1984,lamport1994} برای
مطالعه توصیه می‌شوند.

\mnote{برای استناد به چند منبع، کلیدها را با کاما جدا کنید.}

\section{نمایش کتاب‌نامه}
\label{sec:bibliography-print}

برای نمایش کتاب‌نامه، از دستور
\cmd{printbibliography} در \lr{backmatter}
استفاده کنید:

\begin{latin}
\begin{minted}{latex}
\backmatter
\frontmattergeometry

\printbibliography[title={|\pc{منابع و مراجع}|}]
\end{minted}
\end{latin}

\mnote{کتاب‌نامه را \textbf{هرگز} داخل \lr{\textbackslash begin\{latin\}} قرار ندهید. \lr{xepersian} جهت متن را به‌درستی مدیریت می‌کند.}

\subsection{عنوان فارسی}
\label{sec:bibliography-persian-title}

عنوان کتاب‌نامه را با آرگومان \lr{title} تعیین
کنید:

\begin{latin}
\begin{minted}{latex}
\printbibliography[title={|\pc{منابع و مراجع}|}]
\end{minted}
\end{latin}

\mnote{عنوان پیش‌فرض کتاب‌نامه «منابع» است. با آرگومان \lr{title} می‌توانید آن را تغییر دهید.}

\section{سبک کتاب‌نامه}
\label{sec:bibliography-style}

\lr{MatinBook} از سبک \textbf{عددی} (\lr{numeric})
استفاده می‌کند. یعنی استنادها به‌صورت شماره
نمایش داده می‌شوند: \lr{[1]}، \lr{[2]} و...

\begin{latin}
\begin{minted}{latex}
\RequirePackage[
    backend=biber,
    style=numeric,
    sorting=none
]{biblatex}
\end{minted}
\end{latin}

\mnote{ترتیب \lr{sorting=none} یعنی منابع به ترتیب استناد در کتاب‌نامه ظاهر می‌شوند، نه به ترتیب الفبایی.}

\subsection{سبک‌های جایگزین}
\label{sec:bibliography-alternative-styles}

اگر سبک دیگری می‌خواهید، می‌توانید آن را در
فایل \file{mb-core.sty} تغییر دهید:

\begin{itemize}
    \item \lr{style=numeric}: عددی (پیش‌فرض).
    \item \lr{style=alphabetic}: الفبایی-عددی
    (مثل \lr{[Knu84]}).
    \item \lr{style=authoryear}: نویسنده-سال
    (مثل \lr{(Knuth, 1984)}).
\end{itemize}

\mnote{برای کتاب‌های فنی، سبک عددی توصیه می‌شود چون فشرده‌تر است.}

\section{انواع ورودی \lr{BibTeX}}
\label{sec:bibliography-entry-types}

جدول \cref{tab:bib-types} انواع اصلی ورودی
\lr{BibTeX} را نشان می‌دهد.

\begin{matintable}{انواع ورودی \lr{BibTeX}}{tab:bib-types}
\begin{tblr}{
    width = \textwidth,
    colspec = {l X[l]},
    hlines, vlines,
    colsep = 8pt,
    rowsep = 4pt,
    row{1} = {font=\bfseries},
}
نوع & کاربرد \\
\lr{@book} & کتاب \\
\lr{@article} & مقاله‌ی ژورنالی \\
\lr{@inproceedings} & مقاله‌ی کنفرانسی \\
\lr{@online} & وب‌سایت \\
\lr{@phdthesis} & پایان‌نامه‌ی دکتری \\
\lr{@mastersthesis} & پایان‌نامه‌ی کارشناسی ارشد \\
\lr{@techreport} & گزارش فنی \\
\lr{@misc} & سایر \\
\end{tblr}
\end{matintable}

\mnote{هر نوع ورودی، فیلدهای مخصوص خود را دارد. برای اطلاعات کامل، به مستندات \lr{biblatex} مراجعه کنید.}

\section{منابع فارسی}
\label{sec:bibliography-persian}

\lr{biblatex} با \lr{biber} از منابع فارسی
پشتیبانی می‌کند. برای یک منبع فارسی، از فیلد
\lr{langid} استفاده کنید:

\begin{latin}
\begin{minted}{bibtex}
@book{persian-book,
    author    = {|\pc{نام نویسنده}|},
    title     = {|\pc{عنوان کتاب}|},
    year      = {|\pnum{۱۴۰۴}|},
    publisher = {|\pc{نام ناشر}|},
    langid    = {persian}
}
\end{minted}
\end{latin}

\mnote{فیلد \lr{langid = \{persian\}} باعث می‌شود \lr{biblatex} جهت متن را راست‌به‌چپ کند.}

\section{چرخه‌ی کامپایل}
\label{sec:bibliography-compile}

برای ساخت کتاب‌نامه، باید چرخه‌ی چهار مرحله‌ای
را اجرا کنید:

\begin{latin}
\begin{minted}{bash}
# |\pc{مرحله‌ی اول: کامپایل اول}|%
xelatex -shell-escape main.tex

# |\pc{مرحله‌ی دوم: پردازش کتاب‌نامه}|%
biber main

# |\pc{مرحله‌ی سوم: کامپایل دوم}|%
xelatex -shell-escape main.tex

# |\pc{مرحله‌ی چهارم: کامپایل سوم}|%
xelatex -shell-escape main.tex
\end{minted}
\end{latin}

\mnote{اگر یکی از این مراحل را فراموش کنید، کتاب‌نامه خالی خواهد بود یا استنادها به‌صورت «؟؟» نمایش داده می‌شوند.}

\section{استناد به منابع بدون استناد}
\label{sec:bibliography-nocite}

اگر می‌خواهید منبعی در کتاب‌نامه ظاهر شود
بدون اینکه در متن به آن استناد کنید، از دستور
\cmd{nocite} استفاده کنید:

\begin{latin}
\begin{minted}{latex}
\nocite{knuth1984}
\nocite{*}
\end{minted}
\end{latin}

\begin{itemize}
    \item \cmd{nocite\{key\}}: یک منبع خاص.
    \item \cmd{nocite\{*\}}: همه‌ی منابع فایل
    \lr{.bib}.
\end{itemize}

\mnote{استفاده از \lr{\textbackslash nocite\{*\}} برای کتاب‌های مرجع مفید است، چون همه‌ی منابع را در کتاب‌نامه می‌آورد.}

\section{خطاهای رایج}
\label{sec:bibliography-mistakes}

\subsection{فراموش کردن \cmd{addbibresource}}
\label{sec:bibliography-mistake-addbibresource}

\textbf{مشکل:} اگر \cmd{addbibresource} را فراموش
کنید، کتاب‌نامه خالی خواهد بود.

\textbf{راه‌حل:} همیشه \cmd{addbibresource} را
در پیش از \cmd{begin\{document\}} قرار دهید.

\subsection{فراموش کردن \lr{biber}}
\label{sec:bibliography-mistake-biber}

\textbf{مشکل:} اگر \lr{biber} را اجرا نکنید،
استنادها به‌صورت «؟؟» نمایش داده می‌شوند.

\textbf{راه‌حل:} چرخه‌ی کامل چهار مرحله‌ای را
اجرا کنید.

\subsection{قرار دادن کتاب‌نامه در \env{latin}}
\label{sec:bibliography-mistake-latin}

\textbf{مشکل:} اگر \cmd{printbibliography} را
داخل \cmd{begin\{latin\}} قرار دهید، عنوان
فارسی «منابع و مراجع» به‌هم می‌ریزد.

\textbf{راه‌حل:} کتاب‌نامه را \textbf{هرگز} داخل
\env{latin} قرار ندهید.

\subsection{کلید تکراری}
\label{sec:bibliography-mistake-duplicate}

\textbf{مشکل:} اگر دو منبع کلید یکسان داشته
باشند، \lr{biber} خطا می‌دهد.

\textbf{راه‌حل:} کلیدهای یکتا برای هر منبع
تعریف کنید.

\section{تمرین‌ها}
\label{sec:bibliography-exercises}

\begin{exercise}
    سه منبع (یک کتاب، یک مقاله، یک وب‌سایت)
    در فایل \lr{.bib} تعریف کنید.
    \label{exc:bibliography-bibfile}
\end{exercise}

\begin{exercise}
    به هر سه منبع استناد کنید و کتاب‌نامه را
    نمایش دهید.
    \label{exc:bibliography-cite}
\end{exercise}

\begin{exercise}
    یک منبع فارسی با فیلد \lr{langid} تعریف
    کنید و به آن استناد دهید.
    \label{exc:bibliography-persian}
\end{exercise}

\begin{exercise}
    از \cmd{nocite\{*\}} استفاده کنید و همه‌ی
    منابع را در کتاب‌نامه نمایش دهید.
    \label{exc:bibliography-nocite}
\end{exercise}

\section{خلاصه}
\label{sec:bibliography-summary}

در این فصل، با منابع و کتاب‌نامه در \lr{MatinBook}
آشنا شدیم:

\begin{itemize}
    \item \lr{MatinBook} از \lr{biblatex} با
    موتور \lr{biber} استفاده می‌کند.

    \item منابع در فایل \lr{.bib} ذخیره می‌شوند.

    \item فایل \lr{.bib} با \cmd{addbibresource}
    بارگذاری می‌شود.

    \item چهار دستور استناد: \cmd{cite}،
    \cmd{parencite}، \cmd{textcite}، \cmd{autocite}.

    \item کتاب‌نامه با \cmd{printbibliography}
    نمایش داده می‌شود.

    \item سبک پیش‌فرض، عددی (\lr{numeric}) است.

    \item منابع فارسی با \lr{langid = \{persian\}}
    تعریف می‌شوند.

    \item چرخه‌ی کامپایل چهار مرحله‌ای: \lr{xelatex}
    + \lr{biber} + \lr{xelatex} + \lr{xelatex}.

    \item خطاهای رایج شامل فراموش کردن
    \cmd{addbibresource}، فراموش کردن \lr{biber}،
    قرار دادن کتاب‌نامه در \env{latin}، و کلید
    تکراری است.
\end{itemize}

در فصل بعدی (\cref{chap:index})، با نمایه در
\lr{MatinBook} آشنا می‌شوید و یاد می‌گیرید
که چگونه یک نمایه‌ی حرفه‌ای بسازید.
%%
%% part2-book-components/ch15-index.tex
%%
%% Chapter 15: Index — xindy with Persian sorting, main
%% entries, sub-entries, and cross-references.
%%

\chapter{نمایه}
\label{chap:index}

\section{چرا این موضوع مهم است؟}
\label{sec:index-why}

نمایه یکی از ابزارهای مهم ناوبری در کتاب‌های
فنی و علمی است. یک نمایه‌ی حرفه‌ای به خواننده
کمک می‌کند که مفاهیم مورد نظر را سریع پیدا
کند و مسیر خود را در کتاب گم نکند.

\lr{MatinBook} از \lr{xindy} با ماژول
\lr{persian-variant2} برای ساخت نمایه استفاده
می‌کند. این ترکیب، مرتب‌سازی صحیح فارسی را
تضمین می‌کند.

\mnote{بسته‌ی \lr{makeindex} نمی‌تواند فارسی را درست مرتب کند و حروف پ، چ، ژ، گ، ک را اشتباه می‌چیند. به همین دلیل \lr{MatinBook} از \lr{xindy} استفاده می‌کند.}

\section{ساختار نمایه}
\label{sec:index-structure}

نمایه از چند نوع مدخل تشکیل شده است:

\begin{itemize}
    \item \textbf{مدخل اصلی:} یک واژه یا عبارت
    که در نمایه ظاهر می‌شود.
    \item \textbf{زیرمدخل:} یک مدخل فرعی که
    زیر مدخل اصلی قرار می‌گیرد.
    \item \textbf{زیرزیرمدخل:} عمیق‌ترین سطح
    مدخل.
    \item \textbf{ارجاع «نگاه کنید به»:} ارجاع
    از یک واژه به واژه‌ی دیگر.
    \item \textbf{ارجاع «همچنین نگاه کنید به»:}
    ارجاع تکمیلی.
\end{itemize}

\mnote{مرتب‌سازی نمایه در \lr{MatinBook} بر اساس الفبای فارسی انجام می‌شود و با \lr{xindy} به‌درستی کار می‌کند.}

\section{مدخل اصلی}
\label{sec:index-main}

برای ایجاد یک مدخل اصلی، از دستور \cmd{index}
استفاده کنید:

\begin{latin}
\begin{minted}{latex}
|\pc{الگوریتم}|\index{|\pc{الگوریتم}|}|\pc{یک مفهوم اساسی است.}|%
\end{minted}
\end{latin}

خروجی: الگوریتم\index{الگوریتم} یک مفهوم اساسی است.

\mnote{دستور \lr{\textbackslash index} واژه را در نمایه ثبت می‌کند ولی در متن چیزی نمایش نمی‌دهد.}

\section{زیرمدخل}
\label{sec:index-sub}

برای ایجاد یک زیرمدخل، از علامت \lr{!} استفاده
کنید:

\begin{latin}
\begin{minted}{latex}
|\pc{برنامه‌نویسی}|\index{|\pc{برنامه‌نویسی!پایتون}|}|\pc{یک زبان محبوب است.}|%
\end{minted}
\end{latin}

خروجی: برنامه‌نویسی\index{برنامه‌نویسی!پایتون} یک
زبان محبوب است.

\mnote{در نمایه، این مدخل به‌صورت «برنامه‌نویسی ← پایتون» نمایش داده می‌شود.}

\subsection{زیرزیرمدخل}
\label{sec:index-subsub}

برای زیرزیرمدخل، دو علامت \lr{!} استفاده
کنید:

\begin{latin}
\begin{minted}{latex}
|\pc{ریاضیات}|\index{|\pc{ریاضیات!جبر!گروه}|}|\pc{شاخه‌ای از ریاضیات است.}|%
\end{minted}
\end{latin}

\mnote{توصیه می‌شود بیش از دو سطح زیرمدخل نداشته باشید، چون نمایه پیچیده می‌شود.}

\section{دستورات کمکی}
\label{sec:index-helpers}

\lr{MatinBook} چند دستور کمکی برای نمایه ارائه
می‌دهد. جدول \cref{tab:index-helpers} این
دستورات را خلاصه می‌کند.

\begin{matintable}{دستورات کمکی نمایه}{tab:index-helpers}
\begin{tblr}{
    width = \textwidth,
    colspec = {l X[l]},
    hlines, vlines,
    colsep = 8pt,
    rowsep = 4pt,
    row{1} = {font=\bfseries},
}
دستور & کاربرد \\
\lr{\textbackslash idxbold\{...\}} & مدخل با فونت سیاه \\
\lr{\textbackslash idxitalic\{...\}} & مدخل با فونت مورب \\
\lr{\textbackslash idxsee\{...\}\{...\}} & ارجاع «نگاه کنید به» \\
\lr{\textbackslash idxseealso\{...\}\{...\}} & ارجاع «همچنین نگاه کنید به» \\
\lr{\textbackslash idxsub\{...\}\{...\}} & زیرمدخل سریع \\
\lr{\textbackslash idxsubsub\{...\}\{...\}\{...\}} & زیرزیرمدخل سریع \\
\end{tblr}
\end{matintable}

\subsection{مدخل با فونت سیاه}
\label{sec:index-bold}

برای نمایش یک مدخل با فونت سیاه، از
\cmd{idxbold} استفاده کنید:

\begin{latin}
\begin{minted}{latex}
|\pc{مفهوم مهم}|\idxbold{|\pc{مفهوم مهم}|}%
\end{minted}
\end{latin}

\mnote{مدخل‌های سیاه برای مفاهیم کلیدی استفاده می‌شوند و در نمایه برجسته‌تر دیده می‌شوند.}

\subsection{ارجاع «نگاه کنید به»}
\label{sec:index-see}

برای ارجاع از یک واژه به واژه‌ی دیگر، از
\cmd{idxsee} استفاده کنید:

\begin{latin}
\begin{minted}{latex}
\idxsee{|\pc{پایتون}|}{|\pc{زبان برنامه‌نویسی}|}%
\end{minted}
\end{latin}

خروجی: در نمایه، این مدخل به‌صورت «پایتون ← نگاه
کنید به زبان برنامه‌نویسی» نمایش داده می‌شود.

\mnote{ارجاع «نگاه کنید به» برای واژه‌های مترادف یا مفاهیم مرتبط استفاده می‌شود.}

\subsection{ارجاع «همچنین نگاه کنید به»}
\label{sec:index-seealso}

برای ارجاع تکمیلی، از \cmd{idxseealso} استفاده
کنید:

\begin{latin}
\begin{minted}{latex}
\idxseealso{|\pc{برنامه‌نویسی}|}{|\pc{الگوریتم}|}%
\end{minted}
\end{latin}

\mnote{ارجاع «همچنین نگاه کنید به» برای مفاهیم مرتبط ولی نه مترادف استفاده می‌شود.}

\subsection{زیرمدخل سریع}
\label{sec:index-subfast}

برای ایجاد زیرمدخل سریع، از \cmd{idxsub}
استفاده کنید:

\begin{latin}
\begin{minted}{latex}
\idxsub{|\pc{ریاضیات}|}{|\pc{آنالیز}|}%
\end{minted}
\end{latin}

خروجی: این دستور معادل
\cmd{index\{ریاضیات!آنالیز\}} است.

\mnote{دستور \lr{\textbackslash idxsub} معادل \lr{\textbackslash index\{اصلی!زیر\}} است و کوتاه‌تر نوشته می‌شود.}

\section{مرتب‌سازی فارسی}
\label{sec:index-sorting}

مرتب‌سازی نمایه در \lr{MatinBook} با \lr{xindy}
و ماژول \lr{persian-variant2} انجام می‌شود.
این ماژول، ترتیب صحیح حروف فارسی را رعایت
می‌کند.

جدول \cref{tab:index-sorting} ترتیب حروف
مشکل‌دار در فارسی را نشان می‌دهد.

\begin{matintable}{ترتیب حروف مشکل‌دار در فارسی}{tab:index-sorting}
\begin{tblr}{
    width = \textwidth,
    colspec = {c X[l] X[l]},
    hlines, vlines,
    colsep = 8pt,
    rowsep = 4pt,
    row{1} = {font=\bfseries},
}
حرف & ترتیب صحیح & ترتیب \lr{makeindex} \\
پ & بعد از ب & اشتباه \\
چ & بعد از ج & اشتباه \\
ژ & بعد از ز & اشتباه \\
گ & بعد از ک & اشتباه \\
ک & بعد از گ & اشتباه \\
\end{tblr}
\end{matintable}

\mnote{اگر از \lr{makeindex} استفاده کنید، این حروف در جای اشتباه قرار می‌گیرند و نمایه نامرتب می‌شود.}

\section{چاپ نمایه}
\label{sec:index-print}

برای چاپ نمایه، از دستور \cmd{printindex} در
\lr{backmatter} استفاده کنید:

\begin{latin}
\begin{minted}{latex}
\backmatter
\frontmattergeometry

\printbibliography[title={|\pc{منابع و مراجع}|}]

\printindex
\end{minted}
\end{latin}

\mnote{گزینه‌ی \lr{intoc} در \lr{mb-index.sty} باعث می‌شود نمایه در فهرست مطالب ظاهر شود.}

\section{چرخه‌ی کامپایل}
\label{sec:index-compile}

برای ساخت نمایه، باید چرخه‌ی کامل کامپایل را
اجرا کنید:

\begin{latin}
\begin{minted}{bash}
# |\pc{مرحله‌ی اول: کامپایل اول}|%
xelatex -shell-escape main.tex

# |\pc{مرحله‌ی دوم: پردازش نمایه}|%
# |\pc{\lr{xindy} به‌طور خودکار توسط \lr{imakeidx} اجرا می‌شود}|%

# |\pc{مرحله‌ی سوم: کامپایل دوم}|%
xelatex -shell-escape main.tex

# |\pc{مرحله‌ی چهارم: کامپایل سوم}|%
xelatex -shell-escape main.tex
\end{minted}
\end{latin}

\mnote{اگر از \lr{-shell-escape} استفاده نکنید، \lr{xindy} اجرا نمی‌شود و نمایه خالی می‌ماند.}

\section{تنظیمات نمایه}
\label{sec:index-settings}

تنظیمات نمایه در فایل \file{mb-index.sty} تعریف
شده است:

\begin{latin}
\begin{minted}{latex}
\makeindex[
    intoc,
    program=xindy,
    options={-L persian-variant2 -C utf8
             -M texindy -M page-ranges}
]
\end{minted}
\end{latin}

\mnote{گزینه‌ی \lr{intoc} نمایه را به فهرست مطالب اضافه می‌کند. \lr{persian-variant2} نیز مرتب‌سازی صحیح فارسی را تضمین می‌کند.}

\section{مثال کامل}
\label{sec:index-example}

در ادامه، یک مثال کامل از استفاده‌ی نمایه در
یک فصل آورده شده است:

\begin{latin}
\begin{minted}{latex}
\chapter{|\pc{مفاهیم برنامه‌نویسی}|}

\section{|\pc{زبان‌های برنامه‌نویسی}|}

|\pc{پایتون}|\index{|\pc{پایتون}|}|\pc{یک زبان سطح بالا است.}|%

|\pc{برنامه‌نویسی شیءگرا}|\index{|\pc{برنامه‌نویسی!شیءگرا}|}%
|\pc{یک پارادایم است.}|%

\idxsee{|\pc{جاوااسکریپت}|}{|\pc{زبان برنامه‌نویسی وب}|}%

\idxseealso{|\pc{الگوریتم}|}{|\pc{ساختمان داده}|}%
\end{minted}
\end{latin}

\mnote{هر چه مدخل‌های بیشتری در کتاب داشته باشید، نمایه مفصل‌تر و کاربردی‌تر خواهد بود.}

\section{خطاهای رایج}
\label{sec:index-mistakes}

\subsection{استفاده از \lr{makeindex} به‌جای \lr{xindy}}
\label{sec:index-mistake-makeindex}

\textbf{مشکل:} اگر از \lr{makeindex} استفاده
کنید، حروف پ، چ، ژ، گ، ک اشتباه مرتب می‌شوند.

\textbf{راه‌حل:} همیشه از \lr{xindy} استفاده
کنید. این تنظیم در \file{mb-index.sty} از پیش
انجام شده است.

\subsection{فراموش کردن \lr{-shell-escape}}
\label{sec:index-mistake-shellescape}

\textbf{مشکل:} اگر \lr{-shell-escape} را
فراموش کنید، \lr{xindy} اجرا نمی‌شود و نمایه
خالی می‌ماند.

\textbf{راه‌حل:} همیشه از
\lr{xelatex -shell-escape} استفاده کنید.

\subsection{مدخل‌های تکراری}
\label{sec:index-mistake-duplicate}

\textbf{مشکل:} اگر یک واژه را چند بار با
\cmd{index} علامت‌گذاری کنید، در نمایه چند
بار ظاهر می‌شود.

\textbf{راه‌حل:} واژه را فقط در محل اصلی
تعریف کنید و در جاهای دیگر به آن ارجاع دهید.

\subsection{مدخل‌های بسیار زیاد}
\label{sec:index-mistake-toomany}

\textbf{مشکل:} اگر مدخل‌های نمایه بسیار زیاد
باشند، نمایه از حد معمول بزرگ‌تر می‌شود و
کاربردی‌بودن خود را از دست می‌دهد.

\textbf{راه‌حل:} فقط مفاهیم کلیدی را در نمایه
بگنجانید. قاعده‌ی تجربی: ۵ تا ۱۰ مدخل در هر
صفحه.

\section{تمرین‌ها}
\label{sec:index-exercises}

\begin{exercise}
    یک فصل ساده با سه مدخل اصلی و دو زیرمدخل
    بنویسید و نمایه را چاپ کنید.
    \label{exc:index-basic}
\end{exercise}

\begin{exercise}
    یک مدخل با فونت سیاه و یک مدخل با فونت
    مورب بسازید.
    \label{exc:index-bold}
\end{exercise}

\begin{exercise}
    از \cmd{idxsee} و \cmd{idxseealso} برای
    ارجاع بین مدخل‌ها استفاده کنید.
    \label{exc:index-see}
\end{exercise}

\begin{exercise}
    نمایه‌ی خود را با نمایه‌ی یک کتاب فارسی
    دیگر مقایسه کنید. آیا ترتیب حروف
    یکسان است؟
    \label{exc:index-compare}
\end{exercise}

\section{خلاصه}
\label{sec:index-summary}

در این فصل، با نمایه در \lr{MatinBook} آشنا
شدیم:

\begin{itemize}
    \item \lr{MatinBook} از \lr{xindy} با
    ماژول \lr{persian-variant2} برای مرتب‌سازی
    صحیح فارسی استفاده می‌کند.

    \item \lr{makeindex} نمی‌تواند فارسی را
    درست مرتب کند و باید از \lr{xindy} استفاده
    کرد.

    \item مدخل اصلی با \cmd{index\{...\}}.

    \item زیرمدخل با \cmd{index\{اصلی!زیر\}}.

    \item زیرزیرمدخل با
    \cmd{index\{اصلی!زیر!زیر\}}.

    \item دستورات کمکی: \cmd{idxbold}،
    \cmd{idxitalic}، \cmd{idxsee}،
    \cmd{idxseealso}، \cmd{idxsub}،
    \cmd{idxsubsub}.

    \item نمایه با \cmd{printindex} در
    \lr{backmatter} چاپ می‌شود.

    \item خطاهای رایج شامل استفاده از
    \lr{makeindex}، فراموش کردن
    \lr{-shell-escape}، مدخل‌های تکراری، و
    مدخل‌های بسیار زیاد است.
\end{itemize}

در فصل بعدی (\cref{chap:cross-refs})، با
ارجاع هوشمند در \lr{MatinBook} آشنا می‌شوید
و یاد می‌گیرید که چگونه به عناصر مختلف کتاب
ارجاع دهید.
%%
%% part2-book-components/ch16-cross-refs.tex
%%
%% Chapter 16: Cross-References — cleveref, hyperref,
%% label conventions, and all reference commands.
%%

\chapter{ارجاع هوشمند}
\label{chap:cross-refs}

\section{چرا این موضوع مهم است؟}
\label{sec:cross-refs-why}

در کتاب‌های فنی و علمی، ارجاع دقیق به عناصر
مختلف (معادله، قضیه، شکل، جدول، بخش) اهمیت
ویژه‌ای دارد. خواننده باید بتواند به‌سرعت به
مرجع مورد نظر برسد و از پیوستگی مطالب مطمئن
شود.

\lr{MatinBook} از \lr{cleveref} برای ارجاع
هوشمند استفاده می‌کند. این بسته، نوع مرجع را
به‌طور خودکار تشخیص می‌دهد و پیشوند مناسب
(قضیه، شکل، جدول) را اضافه می‌کند.

\mnote{پیش از \lr{cleveref}، برای هر نوع مرجع باید دستور جداگانه‌ای می‌نوشتیم: \lr{\textbackslash ref}، \lr{\textbackslash figref}، \lr{\textbackslash tabref} و... حالا فقط \lr{\textbackslash cref} کافی است.}

\section{دستور \cmd{cref}}
\label{sec:cross-refs-cref}

دستور \cmd{cref} پرکاربردترین دستور ارجاع در
\lr{MatinBook} است. این دستور، نوع مرجع را
خودکار تشخیص می‌دهد:

\begin{latin}
\begin{minted}{latex}
|\pc{طبق}|\cref{thm:pythagoras-ch8}|\pc{، در هر مثلث قائم‌الزاویه...}|%
\end{minted}
\end{latin}

خروجی: طبق \cref{thm:pythagoras-ch8}، در هر
مثلث قائم‌الزاویه...

\mnote{دستور \lr{\textbackslash cref} پیشوند مناسب را از فایل \lr{fa-IR.sty} می‌گیرد. مثلاً برای قضیه، «قضیه»؛ برای شکل، «شکل».}

\subsection{مراجع مختلف}
\label{sec:cross-refs-types}

جدول \cref{tab:cref-types} پیشوندهای
\cmd{cref} برای انواع مختلف مرجع را نشان
می‌دهد.

\begin{matintable}{پیشوندهای \lr{cref} برای انواع مرجع}{tab:cref-types}
\begin{tblr}{
    width = \textwidth,
    colspec = {l X[l] X[l]},
    hlines, vlines,
    colsep = 8pt,
    rowsep = 4pt,
    row{1} = {font=\bfseries},
}
نوع & برچسب & خروجی \\
معادله & \lr{eq:label} & معادله ۱.۱ \\
قضیه & \lr{thm:label} & قضیه ۱.۱ \\
لم & \lr{lem:label} & لم ۱.۱ \\
نتیجه & \lr{cor:label} & نتیجه ۱.۱ \\
گزاره & \lr{prop:label} & گزاره ۱.۱ \\
تعریف & \lr{def:label} & تعریف ۱.۱ \\
مثال & \lr{ex:label} & مثال ۱.۱ \\
نکته & \lr{rem:label} & نکته ۱.۱ \\
تمرین & \lr{exc:label} & تمرین ۱.۱ \\
شکل & \lr{fig:label} & شکل ۱.۱ \\
جدول & \lr{tab:label} & جدول ۱.۱ \\
کد & \lr{code:label} & کد ۱.۱ \\
الگوریتم & \lr{alg:label} & الگوریتم ۱.۱ \\
فصل & \lr{chap:label} & فصل ۱ \\
بخش & \lr{sec:label} & بخش ۱.۱ \\
\end{tblr}
\end{matintable}

\section{دستورات کمکی}
\label{sec:cross-refs-helpers}

\lr{MatinBook} چند دستور کمکی برای ارجاع ارائه
می‌دهد که همه به \cmd{cref} واگذار می‌شوند:

\begin{itemize}
    \item \cmd{thmref\{label\}}: معادل
    \cmd{cref\{label\}} برای قضیه.
    \item \cmd{lemref\{label\}}: برای لم.
    \item \cmd{corref\{label\}}: برای نتیجه.
    \item \cmd{propref\{label\}}: برای گزاره.
    \item \cmd{defref\{label\}}: برای تعریف.
    \item \cmd{exref\{label\}}: برای مثال.
    \item \cmd{remref\{label\}}: برای نکته.
    \item \cmd{excref\{label\}}: برای تمرین.
\end{itemize}

\mnote{این دستورات از نظر عملکردی معادل \lr{\textbackslash cref} هستند و فقط برای خوانایی بیشتر سورس کد نگه داشته شده‌اند.}

\section{ارجاع به چند مرجع}
\label{sec:cross-refs-multiple}

\subsection{ارجاع گروهی}
\label{sec:cross-refs-group}

برای ارجاع به چند مرجع به‌طور هم‌زمان، کلیدها
را با کاما جدا کنید:

\begin{latin}
\begin{minted}{latex}
|\pc{طبق}|\cref{thm:pythagoras-ch8,def:graph-ch8}|\pc{، ...}|%
\end{minted}
\end{latin}

خروجی: طبق \cref{thm:pythagoras-ch8,def:graph-ch8}، ...

\mnote{در فارسی، \lr{\textbackslash cref\{a,b\}} به‌صورت «قضیه ۱.۱ و تعریف ۱.۲» نمایش داده می‌شود.}

\subsection{ارجاع بازه‌ای}
\label{sec:cross-refs-range}

برای ارجاع به یک بازه‌ی پیوسته، از
\cmd{crefrange} استفاده کنید:

\begin{latin}
\begin{minted}{latex}
|\pc{معادلات}|\crefrange{eq:abs}{eq:inner}|\pc{...}|%
\end{minted}
\end{latin}

خروجی: معادلات \crefrange{eq:abs}{eq:inner}...

\mnote{دستور \lr{\textbackslash crefrange} فقط برای مراجعی کار می‌کند که در یک دسته هستند (مثلاً سه معادله‌ی پیوسته).}

\subsection{ارجاع با \cmd{Cref}}
\label{sec:cross-refs-Cref}

دستور \cmd{Cref} نسخه‌ی «حرف اول بزرگ» است.
در فارسی تفاوتی با \cmd{cref} ندارد، چون
فارسی حالت بزرگ و کوچک ندارد.

\begin{latin}
\begin{minted}{latex}
\Cref{thm:pythagoras-ch8}
\end{minted}
\end{latin}

\mnote{در انگلیسی، \lr{\textbackslash cref} می‌شود «theorem» و \lr{\textbackslash Cref} می‌شود «Theorem». در فارسی هر دو یکسان هستند.}

\section{نام‌گذاری برچسب‌ها}
\label{sec:cross-refs-labels}

برای اینکه برچسب‌ها معنادار و یکتا باشند،
از پیشوندهای استاندارد استفاده کنید:

\begin{matintable}{پیشوندهای استاندارد برچسب}{tab:label-prefixes}
\begin{tblr}{
    width = \textwidth,
    colspec = {l X[l] X[l]},
    hlines, vlines,
    colsep = 8pt,
    rowsep = 4pt,
    row{1} = {font=\bfseries},
}
نوع & پیشوند & مثال \\
معادله & \lr{eq:} & \lr{eq:einstein} \\
قضیه & \lr{thm:} & \lr{thm:pythagoras} \\
لم & \lr{lem:} & \lr{lem:euclid} \\
نتیجه & \lr{cor:} & \lr{cor:monotone} \\
گزاره & \lr{prop:} & \lr{prop:q-dense} \\
تعریف & \lr{def:} & \lr{def:graph} \\
مثال & \lr{ex:} & \lr{ex:integral} \\
نکته & \lr{rem:} & \lr{rem:division-zero} \\
تمرین & \lr{exc:} & \lr{exc:sum-even} \\
شکل & \lr{fig:} & \lr{fig:parabola} \\
جدول & \lr{tab:} & \lr{tab:comparison} \\
کد & \lr{code:} & \lr{code:factorial} \\
الگوریتم & \lr{alg:} & \lr{alg:binary} \\
فصل & \lr{chap:} & \lr{chap:introduction} \\
بخش & \lr{sec:} & \lr{sec:motivation} \\
زیربخش & \lr{subsec:} & \lr{subsec:history} \\
\end{tblr}
\end{matintable}

\mnote{برچسب‌های معنادار، خواندن فایل‌های \lr{tex} را آسان‌تر می‌کنند و از خطاهای ارجاع جلوگیری می‌کنند.}

\section{لینک‌های \lr{hyperref}}
\label{sec:cross-refs-hyperref}

تمام ارجاعات در \lr{MatinBook} به‌صورت
\textbf{لینک فعال} هستند. یعنی خواننده
می‌تواند روی شماره‌ی مرجع کلیک کند و به آن
برود.

\begin{latin}
\begin{minted}{latex}
\RequirePackage[bookmarks]{hyperref}

\hypersetup{
    colorlinks=true,
    linkcolor=blue,
    citecolor=blue,
    urlcolor=blue
}
\end{minted}
\end{latin}

\mnote{رنگ لینک‌ها در \lr{MatinBook} آبی است. برای تغییر، فایل \lr{mb-references.sty} را ویرایش کنید.}

\subsection{بوکمارک‌های PDF}
\label{sec:cross-refs-bookmarks}

\lr{hyperref} به‌طور خودکار بوکمارک‌های
\lr{PDF} را از فصل‌ها و بخش‌ها می‌سازد. این
بوکمارک‌ها در پنل کنار \lr{PDF} نمایش داده
می‌شوند.

\mnote{گزینه‌ی \lr{bookmarksnumbered=true} در \lr{mb-references.sty} باعث می‌شود شماره‌ها در بوکمارک‌ها هم نمایش داده شوند.}

\section{مثال‌های کاربردی}
\label{sec:cross-refs-examples}

\subsection{ارجاع در متن علمی}
\label{sec:cross-refs-example-scientific}

\begin{latin}
\begin{minted}{latex}
|\pc{طبق}|\cref{thm:fundamental}|\pc{، هر عدد طبیعی بزرگ‌تر از ۱ به عوامل اول تجزیه می‌شود.}|%
|\pc{این قضیه در}|\cref{chap:numbertheory}|\pc{به‌طور کامل بررسی شده است.}|%
|\pc{برای مثال،}|\cref{ex:primes-under-50}|\pc{اعداد اول کوچک‌تر از ۵۰ را نشان می‌دهد.}|%
\end{minted}
\end{latin}

\subsection{ارجاع در متن فنی}
\label{sec:cross-refs-example-technical}

\begin{latin}
\begin{minted}{latex}
|\pc{الگوریتم}|\cref{alg:binary}|\pc{روش جستجوی دودویی را نشان می‌دهد.}|%
|\pc{پیچیدگی این الگوریتم}|\lr{$O(\log n)$}|\pc{است، همان‌طور که در}|\cref{tab:complexity}|\pc{نشان داده شده.}|%
|\pc{پیاده‌سازی آن در}|\cref{code:binary-search}|\pc{آمده است.}|%
\end{minted}
\end{latin}

\mnote{استفاده از \lr{\textbackslash cref} در متن‌های علمی و فنی، خوانایی و دقت ارجاعات را افزایش می‌دهد.}

\section{خطاهای رایج}
\label{sec:cross-refs-mistakes}

\subsection{استفاده از \cmd{ref} به‌جای \cmd{cref}}
\label{sec:cross-refs-mistake-ref}

\textbf{مشکل:} \cmd{ref} فقط شماره‌ی مرجع را
نمایش می‌دهد، بدون پیشوند:

\begin{latin}
\begin{minted}{latex}
% |\pc{اشتباه}|:
|\pc{طبق معادله}|\ref{eq:einstein}|\pc{...}|%

% |\pc{درست}|:
|\pc{طبق}|\cref{eq:einstein}|\pc{...}|%
\end{minted}
\end{latin}

\textbf{راه‌حل:} همیشه از \cmd{cref} استفاده
کنید.

\subsection{قرار دادن \cmd{label} بیرون از محیط}
\label{sec:cross-refs-mistake-label}

\textbf{مشکل:} اگر \cmd{label} را بیرون از
محیط (مثل \env{theorem}) قرار دهید، ارجاع
کار نمی‌کند.

\textbf{راه‌حل:} \cmd{label} را همیشه داخل محیط
قرار دهید.

\subsection{برچسب تکراری}
\label{sec:cross-refs-mistake-duplicate}

\textbf{مشکل:} اگر دو عنصر برچسب یکسان داشته
باشند، \lr{LaTeX} خطا می‌دهد و ارجاع به یکی
از آن‌ها ممکن است اشتباه باشد.

\textbf{راه‌حل:} از برچسب‌های یکتا استفاده
کنید.

\subsection{ارجاع به عنصر قبل از تعریف}
\label{sec:cross-refs-mistake-forward}

\textbf{مشکل:} اگر به عنصری ارجاع دهید که
هنوز تعریف نشده، خروجی «؟؟» نمایش داده
می‌شود.

\textbf{راه‌حل:} حداقل دو بار کامپایل کنید.
اگر پس از دو بار همچنان «؟؟» بود، برچسب
را بررسی کنید.

\mnote{اگر خروجی «؟؟» دیدید، اول مطمئن شوید برچسب وجود دارد، سپس دوباره کامپایل کنید.}

\section{تمرین‌ها}
\label{sec:cross-refs-exercises}

\begin{exercise}
    به یک معادله، یک قضیه و یک شکل ارجاع
    دهید و پیشوندهای خروجی را بررسی کنید.
    \label{exc:cross-refs-basic}
\end{exercise}

\begin{exercise}
    به سه معادله به‌طور هم‌زمان ارجاع دهید
    و سپس به یک بازه‌ی پیوسته ارجاع دهید.
    \label{exc:cross-refs-multiple}
\end{exercise}

\begin{exercise}
    از دستورات کمکی (\cmd{thmref}،
    \cmd{defref} و...) استفاده کنید و خروجی
    را با \cmd{cref} مقایسه کنید.
    \label{exc:cross-refs-helpers}
\end{exercise}

\begin{exercise}
    یک برچسب تکراری بسازید و خطای آن را
    ببینید. سپس آن را اصلاح کنید.
    \label{exc:cross-refs-duplicate}
\end{exercise}

\section{خلاصه}
\label{sec:cross-refs-summary}

در این فصل، با ارجاع هوشمند در \lr{MatinBook}
آشنا شدیم:

\begin{itemize}
    \item \lr{MatinBook} از \lr{cleveref} برای
    ارجاع هوشمند استفاده می‌کند.

    \item دستور اصلی، \cmd{cref} است که نوع
    مرجع را خودکار تشخیص می‌دهد.

    \item پیشوندها از فایل \lr{fa-IR.sty}
    می‌آیند (قضیه، شکل، جدول، ...).

    \item \cmd{Cref} نسخه‌ی «حرف اول بزرگ» است
    ولی در فارسی تفاوتی با \cmd{cref} ندارد.

    \item \cmd{crefrange} برای ارجاع بازه‌ای
    و \cmd{cref\{a,b\}} برای ارجاع گروهی.

    \item دستورات کمکی: \cmd{thmref}،
    \cmd{lemref}، \cmd{corref}، \cmd{propref}،
    \cmd{defref}، \cmd{exref}، \cmd{remref}،
    \cmd{excref}.

    \item برچسب‌ها با پیشوندهای استاندارد:
    \lr{eq:}، \lr{thm:}، \lr{fig:}، \lr{tab:}،
    ...

    \item لینک‌های \lr{hyperref} آبی هستند
    و بوکمارک‌های \lr{PDF} خودکار ساخته
    می‌شوند.

    \item خطاهای رایج شامل استفاده از
    \cmd{ref} به‌جای \cmd{cref}، قرار دادن
    \cmd{label} بیرون از محیط، برچسب
    تکراری، و ارجاع به عنصر قبل از تعریف
    است.
\end{itemize}

این فصل، آخرین فصل از بخش دوم (اجزای کتاب)
بود. در بخش سوم (شخصی‌سازی)، یاد می‌گیرید
که چگونه \lr{MatinBook} را برای نیازهای
خود تنظیم کنید.

%========================================================================
% PART 3 — CUSTOMIZATION
%========================================================================
\part{شخصی‌سازی}
\label{part:customization}

%%
%% part3-customization/ch17-options.tex
%%
%% Chapter 17: Class Options — draft/final modes, font
%% selection, and one/two-sided layout.
%%

\chapter{آپشن‌های کلاس}
\label{chap:options}

\section{چرا این موضوع مهم است؟}
\label{sec:options-why}

آپشن‌های کلاس، به شما امکان می‌دهند که
\lr{MatinBook} را برای نیازهای خاص خود تنظیم
کنید. از انتخاب فونت تا حالت پیش‌نویس، همه
با یک خط \cmd{documentclass} قابل کنترل
هستند.

\lr{MatinBook} از سیستم \textbf{کلید-مقدار}
\lr{LaTeX 2023} استفاده می‌کند که بسیار
مدرن‌تر و پایدارتر از روش‌های قدیمی است.
در این فصل، با تمام آپشن‌های کلاس آشنا
می‌شوید.

\mnote{سیستم \lr{\textbackslash DeclareKeys} در \lr{LaTeX 2023} جایگزین \lr{\textbackslash DeclareOption} قدیمی شده و از سینتکس \lr{key=value} پشتیبانی می‌کند.}

\section{آپشن‌های حالت}
\label{sec:options-mode}

\subsection{حالت \lr{final}}
\label{sec:options-final}

حالت پیش‌فرض \lr{MatinBook}، \textbf{حالت
نهایی} است. در این حالت، کتاب به‌صورت
حرفه‌ای و بدون هیچ نشانه‌ی اضافی چاپ
می‌شود.

\begin{latin}
\begin{minted}{latex}
\documentclass[final]{matinbook}
\end{minted}
\end{latin}

\mnote{حالت \lr{final} نیازی به ذکر ندارد، چون پیش‌فرض است. یعنی \lr{\textbackslash documentclass\{matinbook\}} معادل \lr{\textbackslash documentclass[final]\{matinbook\}} است.}

\subsection{حالت \lr{draft}}
\label{sec:options-draft}

در حالت \textbf{پیش‌نویس}، \lr{MatinBook}
رفتارهای زیر را تغییر می‌دهد:

\begin{itemize}
    \item \textbf{واترمارک:} عبارت \lr{DRAFT}
    با شفافیت \pnum{۹۰٪} در پس‌زمینه‌ی هر
    صفحه نمایش داده می‌شود.
    \item \textbf{جعبه‌های سرریز:} خطوطی که از
    حاشیه خارج می‌شوند، با خط مشکی \pnum{۵pt}
    مشخص می‌شوند.
    \item \textbf{تصاویر:} به‌جای نمایش تصاویر،
    فقط کادر آن‌ها رسم می‌شود (سرعت کامپایل
    بالاتر).
\end{itemize}

\begin{latin}
\begin{minted}{latex}
\documentclass[draft]{matinbook}
\end{minted}
\end{latin}

\mnote{حالت \lr{draft} برای مراحل اولیه‌ی نوشتن بسیار مفید است، ولی قبل از چاپ نهایی حتماً به \lr{final} برگردید.}

\section{آپشن‌های فونت}
\label{sec:options-fonts}

\lr{MatinBook} از پنج فونت فارسی پشتیبانی
می‌کند. برای انتخاب فونت، نام آن را به‌عنوان
آپشن به کلاس بدهید:

\begin{latin}
\begin{minted}{latex}
% |\pc{فونت پیش‌فرض}|%
\documentclass[niloofar]{matinbook}

% |\pc{فونت وزیرمتن}|%
\documentclass[vazirmatn]{matinbook}

% |\pc{فونت ساحل}|%
\documentclass[sahel]{matinbook}

% |\pc{فونت ایران‌لاتوس}|%
\documentclass[irlotus]{matinbook}

% |\pc{فونت بی‌نازنین}|%
\documentclass[bnazanin]{matinbook}
\end{minted}
\end{latin}

جدول \cref{tab:font-options} این پنج فونت
را مقایسه می‌کند.

\begin{matintable}{فونت‌های پشتیبانی‌شده‌ی \lr{MatinBook}}{tab:font-options}
\begin{tblr}{
    width = \textwidth,
    colspec = {l X[l] X[c]},
    hlines, vlines,
    colsep = 8pt,
    rowsep = 4pt,
    row{1} = {font=\bfseries},
}
آپشن & فونت & پیشنهاد \\
\lr{niloofar} & \lr{XB Niloofar} & پیش‌فرض \\
\lr{vazirmatn} & \lr{Vazirmatn} & مدرن \\
\lr{sahel} & \lr{Sahel} & ساده \\
\lr{irlotus} & \lr{IR Lotus} & سنتی \\
\lr{bnazanin} & \lr{B Nazanin} & تجاری \\
\end{tblr}
\end{matintable}

\mnote{فونت \lr{B Nazanin} تجاری است و ممکن است روی همه‌ی سیستم‌ها نصب نباشد. برای کتاب‌های فنی، \lr{XB Niloofar} یا \lr{Vazirmatn} توصیه می‌شود.}

\subsection{ترکیب چند آپشن}
\label{sec:options-combine}

می‌توانید چند آپشن را با کاما از هم جدا کنید:

\begin{latin}
\begin{minted}{latex}
\documentclass[draft, vazirmatn]{matinbook}
\end{minted}
\end{latin}

\mnote{ترتیب آپشن‌ها مهم نیست. \lr{[draft, vazirmatn]} معادل \lr{[vazirmatn, draft]} است.}

\section{آپشن‌های چیدمان}
\label{sec:options-layout}

\subsection{حالت دوطرفه و یک‌طرفه}
\label{sec:options-twoside}

\lr{MatinBook} به‌طور پیش‌فرض از حالت
\textbf{دوطرفه} (\lr{twoside}) استفاده
می‌کند. یعنی صفحات فرد و زوج، چیدمان
آینه‌ای دارند (حاشیه‌ی پهن در بیرون).

\begin{latin}
\begin{minted}{latex}
% |\pc{حالت دوطرفه (پیش‌فرض)}|%
\documentclass[twoside]{matinbook}

% |\pc{حالت یک‌طرفه}|%
\documentclass[oneside]{matinbook}
\end{minted}
\end{latin}

\mnote{حالت \lr{twoside} برای کتاب‌های چاپی مناسب است. حالت \lr{oneside} برای اسناد دیجیتال یا فایل‌های خواندنی روی صفحه‌نمایش.}

\subsection{آپشن‌های \lr{geometry}}
\label{sec:options-geometry}

برای تغییر هندسه‌ی صفحه، می‌توانید از
دستور \cmd{geometry} در پیش از
\cmd{begin\{document\}} استفاده کنید:

\begin{latin}
\begin{minted}{latex}
\usepackage{geometry}
\geometry{
    paperwidth=17.8cm,
    paperheight=25.4cm,
    outer=5.9cm,
    inner=1.5cm
}
\end{minted}
\end{latin}

\mnote{برای جزئیات بیشتر درباره‌ی چیدمان، به \cref{chap:layout} مراجعه کنید.}

\section{آپشن‌های سفارشی}
\label{sec:options-custom}

\lr{MatinBook} از آپشن‌های سفارشی نیز پشتیبانی
می‌کند. برای افزودن آپشن جدید، به فایل
\file{matinbook.cls} بروید و از
\cmd{DeclareKeys} استفاده کنید.

\subsection{مثال: افزودن آپشن \lr{print}}
\label{sec:options-custom-example}

فرض کنید می‌خواهید یک آپشن \lr{print} اضافه
کنید که رنگ‌ها را برای چاپ سیاه‌وسفید تنظیم
کند:

\begin{latin}
\begin{minted}{latex}
% |\pc{در فایل}|\lr{matinbook.cls}:
\newif\if@matin@print
\@matin@printfalse

\DeclareKeys{
    print .code = \@matin@printtrue,
}

\ProcessKeyOptions
\end{minted}
\end{latin}

سپس در ماژول‌های دیگر، از این فلگ استفاده
کنید:

\begin{latin}
\begin{minted}{latex}
\if@matin@print
    % |\pc{تنظیمات چاپ سیاه‌وسفید}|%
\fi
\end{minted}
\end{latin}

\mnote{این یک مثال پیشرفته است. برای کارهای معمول، نیازی به افزودن آپشن جدید نیست.}

\section{آپشن‌ها در فایل‌های مختلف}
\label{sec:options-files}

اگر پروژه‌ی شما چند فایل دارد، آپشن‌ها
فقط در \file{main.tex} تعریف می‌شوند:

\begin{latin}
\begin{minted}{latex}
% |\pc{در فایل}|\lr{main.tex}:
\documentclass[draft, vazirmatn]{matinbook}

\begin{document}
\input{chapters/chapter-01}
\end{document}
\end{minted}
\end{latin}

\mnote{فایل‌های فصل (\lr{chapter-01.tex} و...) نباید \lr{\textbackslash documentclass} داشته باشند. آن‌ها فقط محتوا هستند.}

\section{جدول خلاصه‌ی آپشن‌ها}
\label{sec:options-summary-table}

جدول \cref{tab:all-options} تمام آپشن‌های
\lr{MatinBook} را خلاصه می‌کند.

\begin{matintable}{خلاصه‌ی آپشن‌های \lr{MatinBook}}{tab:all-options}
\begin{tblr}{
    width = \textwidth,
    colspec = {l X[l] X[c]},
    hlines, vlines,
    colsep = 8pt,
    rowsep = 4pt,
    row{1} = {font=\bfseries},
}
آپشن & توضیح & پیش‌فرض \\
\lr{draft} & حالت پیش‌نویس & خیر \\
\lr{final} & حالت نهایی & بله \\
\lr{niloofar} & فونت \lr{XB Niloofar} & بله \\
\lr{vazirmatn} & فونت \lr{Vazirmatn} & خیر \\
\lr{sahel} & فونت \lr{Sahel} & خیر \\
\lr{irlotus} & فونت \lr{IR Lotus} & خیر \\
\lr{bnazanin} & فونت \lr{B Nazanin} & خیر \\
\lr{twoside} & چیدمان دوطرفه & بله \\
\lr{oneside} & چیدمان یک‌طرفه & خیر \\
\end{tblr}
\end{matintable}

\section{خطاهای رایج}
\label{sec:options-mistakes}

\subsection{آپشن ناشناخته}
\label{sec:options-mistake-unknown}

\textbf{مشکل:} اگر آپشن ناشناخته‌ای بدهید،
\lr{LaTeX} آن را به کلاس \lr{book} می‌فرستد
و ممکن است خطا بدهد.

\textbf{راه‌حل:} فقط از آپشن‌های مستند استفاده
کنید.

\subsection{ترکیب آپشن‌های متناقض}
\label{sec:options-mistake-conflict}

\textbf{مشکل:} اگر \lr{draft} و \lr{final} را
با هم بدهید، آخرین آپشن اعمال می‌شود.

\textbf{راه‌حل:} فقط یکی از آپشن‌های متناقض
را بدهید.

\subsection{تعریف آپشن در فایل فصل}
\label{sec:options-mistake-chapter}

\textbf{مشکل:} اگر \cmd{documentclass} را در
فایل فصل بنویسید، خطا می‌دهد.

\textbf{راه‌حل:} \cmd{documentclass} فقط در
\file{main.tex}.

\section{تمرین‌ها}
\label{sec:options-exercises}

\begin{exercise}
    کتاب خود را در حالت \lr{draft} کامپایل
    کنید و واترمارک را ببینید.
    \label{exc:options-draft}
\end{exercise}

\begin{exercise}
    فونت را به \lr{vazirmatn} تغییر دهید و
    نتیجه را مقایسه کنید.
    \label{exc:options-font}
\end{exercise}

\begin{exercise}
    یک کتاب با حالت \lr{oneside} بسازید و
    تفاوت آن را با \lr{twoside} ببینید.
    \label{exc:options-side}
\end{exercise}

\begin{exercise}
    یک آپشن سفارشی (مثلاً \lr{print}) به
    \file{matinbook.cls} اضافه کنید.
    \label{exc:options-custom}
\end{exercise}

\section{خلاصه}
\label{sec:options-summary}

در این فصل، با آپشن‌های کلاس \lr{MatinBook}
آشنا شدیم:

\begin{itemize}
    \item \lr{MatinBook} از سیستم کلید-مقدار
    \lr{LaTeX 2023} استفاده می‌کند.

    \item حالت \lr{final} پیش‌فرض است و
    کتاب را بدون نشانه‌ی اضافی چاپ می‌کند.

    \item حالت \lr{draft} واترمارک، جعبه‌های
    سرریز و تصاویر ساده‌شده را فعال می‌کند.

    \item پنج فونت فارسی پشتیبانی می‌شود:
    \lr{niloofar}، \lr{vazirmatn}، \lr{sahel}،
    \lr{irlotus}، \lr{bnazanin}.

    \item حالت \lr{twoside} پیش‌فرض است و
    برای کتاب‌های چاپی مناسب است.

    \item آپشن‌های سفارشی با
    \cmd{DeclareKeys} قابل افزودن هستند.

    \item آپشن‌ها فقط در \file{main.tex}
    تعریف می‌شوند.

    \item خطاهای رایج شامل آپشن ناشناخته،
    ترکیب آپشن‌های متناقض، و تعریف آپشن در
    فایل فصل است.
\end{itemize}

در فصل بعدی (\cref{chap:colors})، با سیستم
رنگ \lr{MatinBook} آشنا می‌شوید و یاد
می‌گیرید که چگونه رنگ‌ها را شخصی‌سازی
کنید.
%%
%% part3-customization/ch18-colors.tex
%%
%% Chapter 18: Color System — CMYK palette, five color
%% families, and customization.
%%

\chapter{سیستم رنگ}
\label{chap:colors}

\section{چرا این موضوع مهم است؟}
\label{sec:colors-why}

رنگ‌ها در \lr{MatinBook} فقط برای زیبایی نیستند؛
آن‌ها نقش معنایی دارند. هر رنگ، نوع خاصی از
محتوا را نشان می‌دهد و به خواننده کمک می‌کند
که سریع‌تر نوع مطلب را تشخیص دهد.

\lr{MatinBook} از پالت \textbf{CMYK} استفاده
می‌کند که برای چاپ حرفه‌ای طراحی شده است. در
این فصل، با این پالت آشنا می‌شوید و یاد
می‌گیرید که چگونه آن را شخصی‌سازی کنید.

\mnote{مدل \lr{CMYK} برای چاپ، و مدل \lr{RGB} برای نمایش دیجیتال مناسب است. \lr{MatinBook} از \lr{CMYK} استفاده می‌کند چون کتاب‌ها معمولاً چاپ می‌شوند.}

\section{پالت رنگی \lr{MatinBook}}
\label{sec:colors-palette}

\lr{MatinBook} ۲۰ رنگ را در ۵ خانواده سازمان‌دهی
می‌کند. هر خانواده، نقش معنایی مشخصی دارد:

\begin{itemize}
    \item \textbf{آبی:} قضیه، لم، نتیجه، گزاره.
    \item \textbf{سبز:} تعریف، راه‌حل.
    \item \textbf{نارنجی:} مثال، نکته.
    \item \textbf{قرمز:} هشدار، نکته.
    \item \textbf{بنفش:} تمرین.
\end{itemize}

\mnote{رنگ \lr{matinblue} رنگ غالب سیستم است. سایر رنگ‌ها یا سایه‌های آبی هستند یا رنگ‌های عملکردی.}

\section{خانواده‌ی آبی}
\label{sec:colors-blue}

خانواده‌ی آبی، رنگ غالب سیستم است و برای
محیط‌های نظری استفاده می‌شود. جدول
\cref{tab:colors-blue} این خانواده را نشان
می‌دهد.

\begin{matintable}{خانواده‌ی آبی}{tab:colors-blue}
\begin{tblr}{
    width = \textwidth,
    colspec = {l X[l] X[c]},
    hlines, vlines,
    colsep = 8pt,
    rowsep = 4pt,
    row{1} = {font=\bfseries},
}
نام & کاربرد & \lr{CMYK} \\
\lr{matinblue} & قضیه & \lr{(80, 30, 0, 30)} \\
\lr{matindarkblue} & آبی تیره & \lr{(80, 30, 0, 60)} \\
\lr{matinlightblue} & آبی روشن & \lr{(15, 5, 0, 5)} \\
\end{tblr}
\end{matintable}

\begin{quote}
    \textcolor{matinblue}{\textbf{متن آبی}} —
    \textcolor{matindarkblue}{\textbf{متن آبی تیره}} —
    \colorbox{matinlightblue}{\textbf{پس‌زمینه‌ی آبی روشن}}
\end{quote}

\section{خانواده‌ی سبز}
\label{sec:colors-green}

خانواده‌ی سبز به تعاریف و راه‌حل‌ها اختصاص
دارد. جدول \cref{tab:colors-green} این
خانواده را نشان می‌دهد.

\begin{matintable}{خانواده‌ی سبز}{tab:colors-green}
\begin{tblr}{
    width = \textwidth,
    colspec = {l X[l] X[c]},
    hlines, vlines,
    colsep = 8pt,
    rowsep = 4pt,
    row{1} = {font=\bfseries},
}
نام & کاربرد & \lr{CMYK} \\
\lr{matingreen} & تعریف & \lr{(80, 0, 45, 30)} \\
\lr{matindarkgreen} & سبز تیره & \lr{(80, 0, 45, 65)} \\
\lr{matinlightgreen} & سبز روشن & \lr{(10, 0, 5, 5)} \\
\end{tblr}
\end{matintable}

\begin{quote}
    \textcolor{matingreen}{\textbf{متن سبز}} —
    \textcolor{matindarkgreen}{\textbf{متن سبز تیره}} —
    \colorbox{matinlightgreen}{\textbf{پس‌زمینه‌ی سبز روشن}}
\end{quote}

\section{خانواده‌ی نارنجی}
\label{sec:colors-orange}

خانواده‌ی نارنجی برای مثال‌ها استفاده می‌شود.
جدول \cref{tab:colors-orange} این خانواده
را نشان می‌دهد.

\begin{matintable}{خانواده‌ی نارنجی}{tab:colors-orange}
\begin{tblr}{
    width = \textwidth,
    colspec = {l X[l] X[c]},
    hlines, vlines,
    colsep = 8pt,
    rowsep = 4pt,
    row{1} = {font=\bfseries},
}
نام & کاربرد & \lr{CMYK} \\
\lr{matinorange} & مثال & \lr{(0, 45, 85, 10)} \\
\lr{matindarkorange} & نارنجی تیره & \lr{(0, 45, 85, 55)} \\
\lr{matinlightorange} & نارنجی روشن & \lr{(0, 5, 15, 0)} \\
\end{tblr}
\end{matintable}

\begin{quote}
    \textcolor{matinorange}{\textbf{متن نارنجی}} —
    \textcolor{matindarkorange}{\textbf{متن نارنجی تیره}} —
    \colorbox{matinlightorange}{\textbf{پس‌زمینه‌ی نارنجی روشن}}
\end{quote}

\section{خانواده‌ی قرمز}
\label{sec:colors-red}

خانواده‌ی قرمز به نکته‌ها و هشدارها اختصاص
دارد. جدول \cref{tab:colors-red} این
خانواده را نشان می‌دهد.

\begin{matintable}{خانواده‌ی قرمز}{tab:colors-red}
\begin{tblr}{
    width = \textwidth,
    colspec = {l X[l] X[c]},
    hlines, vlines,
    colsep = 8pt,
    rowsep = 4pt,
    row{1} = {font=\bfseries},
}
نام & کاربرد & \lr{CMYK} \\
\lr{matinred} & هشدار & \lr{(0, 65, 75, 10)} \\
\lr{matindarkred} & قرمز تیره & \lr{(0, 65, 75, 55)} \\
\lr{matinlightred} & قرمز روشن & \lr{(0, 15, 15, 0)} \\
\end{tblr}
\end{matintable}

\begin{quote}
    \textcolor{matinred}{\textbf{متن قرمز}} —
    \textcolor{matindarkred}{\textbf{متن قرمز تیره}} —
    \colorbox{matinlightred}{\textbf{پس‌زمینه‌ی قرمز روشن}}
\end{quote}

\section{خانواده‌ی بنفش}
\label{sec:colors-purple}

خانواده‌ی بنفش برای تمرین‌ها استفاده می‌شود.
جدول \cref{tab:colors-purple} این خانواده
را نشان می‌دهد.

\begin{matintable}{خانواده‌ی بنفش}{tab:colors-purple}
\begin{tblr}{
    width = \textwidth,
    colspec = {l X[l] X[c]},
    hlines, vlines,
    colsep = 8pt,
    rowsep = 4pt,
    row{1} = {font=\bfseries},
}
نام & کاربرد & \lr{CMYK} \\
\lr{matinpurple} & تمرین & \lr{(20, 60, 0, 30)} \\
\lr{matindarkpurple} & بنفش تیره & \lr{(20, 60, 0, 65)} \\
\lr{matinlightpurple} & بنفش روشن & \lr{(5, 10, 0, 5)} \\
\end{tblr}
\end{matintable}

\begin{quote}
    \textcolor{matinpurple}{\textbf{متن بنفش}} —
    \textcolor{matindarkpurple}{\textbf{متن بنفش تیره}} —
    \colorbox{matinlightpurple}{\textbf{پس‌زمینه‌ی بنفش روشن}}
\end{quote}

\section{رنگ‌های خنثی}
\label{sec:colors-neutral}

رنگ‌های خنثی برای متن، حاشیه و پس‌زمینه
استفاده می‌شوند. جدول \cref{tab:colors-neutral}
این رنگ‌ها را نشان می‌دهد.

\begin{matintable}{رنگ‌های خنثی}{tab:colors-neutral}
\begin{tblr}{
    width = \textwidth,
    colspec = {l X[l] X[c]},
    hlines, vlines,
    colsep = 8pt,
    rowsep = 4pt,
    row{1} = {font=\bfseries},
}
نام & کاربرد & \lr{CMYK} \\
\lr{matindarkgray} & متن اصلی & \lr{(45, 25, 0, 70)} \\
\lr{matinmediumgray} & متن فرعی & \lr{(10, 0, 0, 35)} \\
\lr{matinlightgray} & پس‌زمینه & \lr{(0, 0, 0, 5)} \\
\lr{matinbordergray} & حاشیه & \lr{(5, 0, 0, 20)} \\
\lr{matincream} & پس‌زمینه‌ی گرم & \lr{(0, 0, 0, 0)} \\
\end{tblr}
\end{matintable}

\mnote{متن اصلی نه با سیاه مطلق، بلکه با خاکستری تیره‌ی نزدیک به سیاه حروف‌چینی می‌شود تا کنتراست چشم‌نوازتری ایجاد شود.}

\section{استفاده از رنگ‌ها}
\label{sec:colors-usage}

\lr{MatinBook} سه دستور اصلی برای استفاده
از رنگ‌ها ارائه می‌دهد:

\begin{itemize}
    \item \cmd{textcolor\{رنگ\}\{متن\}}:
    متن رنگی.
    \item \cmd{colorbox\{رنگ\}\{متن\}}:
    پس‌زمینه‌ی رنگی.
    \item \cmd{color\{رنگ\}}: تغییر رنگ جاری.
\end{itemize}

\begin{latin}
\begin{minted}{latex}
\textcolor{matinblue}{|\pc{متن آبی}|}%
\colorbox{matinlightblue}{|\pc{پس‌زمینه‌ی آبی}|}%
{\color{matinred}|\pc{متن قرمز}|}%
\end{minted}
\end{latin}

\mnote{برای رنگ‌آمیزی متن، از \lr{\textbackslash textcolor} استفاده کنید. برای پس‌زمینه، از \lr{\textbackslash colorbox}.}

\subsection{ترکیب رنگ‌ها}
\label{sec:colors-mixing}

می‌توانید سایه‌های مختلفی از هر رنگ بسازید:

\begin{latin}
\begin{minted}{latex}
\textcolor{matinblue!20}{|\pc{آبی خیلی روشن}|}%
\textcolor{matinblue!50}{|\pc{آبی متوسط}|}%
\textcolor{matinblue!80}{|\pc{آبی تیره}|}%
\end{minted}
\end{latin}

خروجی:

\begin{quote}
    \textcolor{matinblue!20}{آبی خیلی روشن} —
    \textcolor{matinblue!50}{آبی متوسط} —
    \textcolor{matinblue!80}{آبی تیره}
\end{quote}

\mnote{عدد بعد از \lr{!} درصد ترکیب با رنگ پایه را نشان می‌دهد. \lr{matinblue!20} یعنی ۲۰٪ آبی و ۸۰٪ رنگ پایه.}

\section{قواعد استفاده از رنگ}
\label{sec:colors-rules}

\lr{MatinBook} چند قاعده برای استفاده‌ی
حرفه‌ای از رنگ‌ها توصیه می‌کند:

\begin{enumerate}
    \item \textbf{حداکثر ۳ رنگ در هر صفحه}
    (بدون احتساب رنگ جعبه‌ها).
    \item \textbf{استفاده‌ی معنایی، نه تزئینی.}
    آبی برای قضیه، سبز برای تعریف، نارنجی
    برای مثال.
    \item \textbf{کنتراست ۴.۵:۱} برای
    دسترس‌پذیری.
    \item \textbf{پرهیز از سیاه مطلق روی
    سفید مطلق.}
    \item \textbf{رنگ غالب، آبی است.} سایر
    رنگ‌ها عملکردی هستند.
\end{enumerate}

\mnote{استفاده‌ی بیش از حد از رنگ، کتاب را پرزرق‌وبرق می‌کند و تمرکز خواننده را از محتوا می‌گیرد.}

\section{شخصی‌سازی پالت}
\label{sec:colors-customization}

برای شخصی‌سازی پالت رنگ، فایل
\file{mb-theme-colors.sty} را ویرایش کنید.
این فایل، همه‌ی رنگ‌ها را با
\cmd{providecolor} تعریف می‌کند:

\begin{latin}
\begin{minted}{latex}
\providecolor{matinblue}{cmyk}{0.80, 0.30, 0.00, 0.30}
\providecolor{matingreen}{cmyk}{0.80, 0.00, 0.45, 0.30}
\providecolor{matinorange}{cmyk}{0.00, 0.45, 0.85, 0.10}
\end{minted}
\end{latin}

\mnote{چون از \lr{\textbackslash providecolor} استفاده شده، می‌توانید رنگ‌ها را با \lr{\textbackslash definecolor} قبل از بارگذاری کلاس بازنویسی کنید.}

\subsection{بازنویسی رنگ در سند}
\label{sec:colors-override}

برای بازنویسی یک رنگ خاص در سند خود، آن
را قبل از \cmd{documentclass} تعریف کنید:

\begin{latin}
\begin{minted}{latex}
\documentclass{matinbook}

\usepackage{xcolor}
\definecolor{matinblue}{cmyk}{0.70, 0.20, 0.00, 0.20}

\begin{document}
...
\end{document}
\end{minted}
\end{latin}

\mnote{چون \lr{MatinBook} از \lr{\textbackslash providecolor} استفاده می‌کند، تعریف شما اولویت دارد.}

\section{خطاهای رایج}
\label{sec:colors-mistakes}

\subsection{استفاده از رنگ تعریف‌نشده}
\label{sec:colors-mistake-undefined}

\textbf{مشکل:} اگر رنگی تعریف نشده باشد،
خطای \lr{Undefined color} می‌دهد.

\textbf{راه‌حل:} از پالت \lr{MatinBook} استفاده
کنید یا رنگ جدید را با \cmd{definecolor}
تعریف کنید.

\subsection{استفاده‌ی بیش از حد از رنگ}
\label{sec:colors-mistake-toomuch}

\textbf{مشکل:} اگر در هر صفحه بیش از ۳ رنگ
استفاده کنید، صفحه شلوغ به نظر می‌رسد.

\textbf{راه‌حل:} از رنگ‌ها به‌صورت معنایی
استفاده کنید، نه تزئینی.

\subsection{کنتراست پایین}
\label{sec:colors-mistake-contrast}

\textbf{مشکل:} اگر متن رنگی روی پس‌زمینه‌ی
روشن نوشته شود، خواندنش سخت می‌شود.

\textbf{راه‌حل:} از کنتراست ۴.۵:۱ استفاده
کنید. پالت \lr{MatinBook} این کنتراست را
تضمین می‌کند.

\section{تمرین‌ها}
\label{sec:colors-exercises}

\begin{exercise}
    یک متن ساده با هر پنج خانواده‌ی رنگی
    بنویسید.
    \label{exc:colors-basic}
\end{exercise}

\begin{exercise}
    از ترکیب \lr{matinblue!40} استفاده کنید
    و با \lr{matinlightblue} مقایسه کنید.
    \label{exc:colors-mixing}
\end{exercise}

\begin{exercise}
    رنگ \lr{matinblue} را در سند خود بازنویسی
    کنید و کتاب را کامپایل کنید.
    \label{exc:colors-override}
\end{exercise}

\begin{exercise}
    یک صفحه بسازید که حداکثر ۳ رنگ داشته
    باشد. سپس رنگی اضافه کنید و تفاوت را
    ببینید.
    \label{exc:colors-limit}
\end{exercise}

\section{خلاصه}
\label{sec:colors-summary}

در این فصل، با سیستم رنگ \lr{MatinBook}
آشنا شدیم:

\begin{itemize}
    \item \lr{MatinBook} از پالت \lr{CMYK}
    استفاده می‌کند که برای چاپ مناسب است.

    \item ۲۰ رنگ در ۵ خانواده سازمان‌دهی
    شده‌اند: آبی، سبز، نارنجی، قرمز، بنفش.

    \item رنگ \lr{matinblue} رنگ غالب
    سیستم است.

    \item رنگ‌های خنثی شامل خاکستری تیره،
    متوسط، روشن، مرزی و کِرِم هستند.

    \item سه دستور اصلی: \cmd{textcolor}،
    \cmd{colorbox}، \cmd{color}.

    \item ترکیب رنگ‌ها با \lr{!} امکان‌پذیر
    است: \lr{matinblue!40}.

    \item قواعد استفاده: حداکثر ۳ رنگ در
    صفحه، استفاده‌ی معنایی، کنتراست ۴.۵:۱.

    \item شخصی‌سازی با ویرایش
    \file{mb-theme-colors.sty}.

    \item خطاهای رایج شامل رنگ تعریف‌نشده،
    استفاده‌ی بیش از حد، و کنتراست پایین
    است.
\end{itemize}

در فصل بعدی (\cref{chap:fonts})، با فونت‌ها
در \lr{MatinBook} آشنا می‌شوید و یاد می‌گیرید
که چگونه فونت‌ها را شخصی‌سازی کنید.
%%
%% part3-customization/ch19-fonts.tex
%%
%% Chapter 19: Fonts — The five Persian fonts, the Latin
%% font, the math font, and the monospace font.
%%

\chapter{فونت‌ها}
\label{chap:fonts}

\section{چرا این موضوع مهم است؟}
\label{sec:fonts-why}

فونت، یکی از مهم‌ترین عوامل در خوانایی و
زیبایی کتاب است. یک فونت خوب، چشم خواننده را
خسته نمی‌کند و درک مطلب را آسان‌تر می‌سازد.

\lr{MatinBook} از پنج فونت فارسی پشتیبانی
می‌کند و همچنین فونت‌های لاتین، ریاضی و
مونواسپیس را به‌طور دقیق تنظیم می‌کند. در
این فصل، با این فونت‌ها آشنا می‌شوید و
یاد می‌گیرید که چگونه آن‌ها را شخصی‌سازی
کنید.

\mnote{\lr{MatinBook} از \lr{fontspec} و \lr{unicode-math} برای مدیریت فونت‌ها استفاده می‌کند که از فونت‌های \lr{OpenType} پشتیبانی می‌کنند.}

\section{فونت‌های فارسی}
\label{sec:fonts-persian}

\lr{MatinBook} پنج فونت فارسی را پشتیبانی
می‌کند که با آپشن کلاس قابل انتخاب هستند.
جدول \cref{tab:persian-fonts} این فونت‌ها
را خلاصه می‌کند.

\begin{matintable}{فونت‌های فارسی \lr{MatinBook}}{tab:persian-fonts}
\begin{tblr}{
    width = \textwidth,
    colspec = {l X[l] X[c] X[c]},
    hlines, vlines,
    colsep = 6pt,
    rowsep = 4pt,
    row{1} = {font=\bfseries},
}
آپشن & فونت & مقیاس & پیش‌فرض \\
\lr{niloofar} & \lr{XB Niloofar} & \pnum{۱.۲} & بله \\
\lr{vazirmatn} & \lr{Vazirmatn} & \pnum{۱.۲} & خیر \\
\lr{sahel} & \lr{Sahel} & \pnum{۱.۲} & خیر \\
\lr{irlotus} & \lr{IR Lotus} & \pnum{۱.۲} & خیر \\
\lr{bnazanin} & \lr{B Nazanin} & \pnum{۱.۰} & خیر \\
\end{tblr}
\end{matintable}

\subsection{فونت \lr{XB Niloofar}}
\label{sec:fonts-niloofar}

\lr{XB Niloofar} فونت پیش‌فرض \lr{MatinBook}
است. این فونت به‌دلیل خوانایی بالا در متن‌های
فنی و پوشش کامل حروف فارسی انتخاب شده است.

\begin{latin}
\begin{minted}{latex}
\documentclass[niloofar]{matinbook}
\end{minted}
\end{latin}

مقیاس این فونت \pnum{۱.۲} است، یعنی ۲۰٪
بزرگ‌تر از فونت پایه‌ی سند (\pnum{۱۰pt}).

\mnote{مقیاس \pnum{۱.۲} برای فونت فارسی، استاندارد نشر فنی ایران است. این تنظیم باعث می‌شود فونت فارسی با فونت لاتین تعادل بصری داشته باشد.}

\subsection{فونت \lr{Vazirmatn}}
\label{sec:fonts-vazirmatn}

\lr{Vazirmatn} یک فونت فارسی مدرن با طراحی
هندسی است که خوانایی بالایی در نمایش دیجیتال
دارد:

\begin{latin}
\begin{minted}{latex}
\documentclass[vazirmatn]{matinbook}
\end{minted}
\end{latin}

\mnote{فونت \lr{Vazirmatn} برای کتاب‌های دیجیتال و وب مناسب‌تر است، ولی در چاپ هم کیفیت خوبی دارد.}

\subsection{فونت \lr{Sahel}}
\label{sec:fonts-sahel}

\lr{Sahel} یک فونت فارسی ساده و خوانا است
که برای متن‌های طولانی مناسب است:

\begin{latin}
\begin{minted}{latex}
\documentclass[sahel]{matinbook}
\end{minted}
\end{latin}

\subsection{فونت \lr{IR Lotus}}
\label{sec:fonts-irlotus}

\lr{IR Lotus} یک فونت فارسی سنتی است که
حس کلاسیک به کتاب می‌دهد:

\begin{latin}
\begin{minted}{latex}
\documentclass[irlotus]{matinbook}
\end{minted}
\end{latin}

\subsection{فونت \lr{B Nazanin}}
\label{sec:fonts-bnazanin}

\lr{B Nazanin} یک فونت تجاری است که پوشش
حروف آن کامل نیست و برای کتاب‌های فنی
توصیه نمی‌شود:

\begin{latin}
\begin{minted}{latex}
\documentclass[bnazanin]{matinbook}
\end{minted}
\end{latin}

\mnote{فونت \lr{B Nazanin} تجاری است و ممکن است روی همه‌ی سیستم‌ها نصب نباشد. برای کتاب‌های فنی، \lr{XB Niloofar} یا \lr{Vazirmatn} توصیه می‌شود.}

\section{فونت لاتین}
\label{sec:fonts-latin}

فونت لاتین پیش‌فرض \lr{MatinBook}، \lr{Times
New Roman} است. اگر این فونت روی سیستم
نصب نباشد، \lr{TeX Gyre Termes} به‌عنوان
جایگزین استفاده می‌شود.

\begin{latin}
\begin{minted}{latex}
% |\pc{در فایل}|\lr{matinbook.cls}:
\IfFontExistsTF{Times New Roman}{
    \setlatintextfont{Times New Roman}
}{
    \setlatintextfont{TeX Gyre Termes}
    \PackageWarning{matinbook}{Times New Roman not found}
}
\end{minted}
\end{latin}

\mnote{فونت \lr{TeX Gyre Termes} یک فونت رایگان و متن‌باز است که ظاهری مشابه \lr{Times New Roman} دارد.}

\section{فونت ریاضی}
\label{sec:fonts-math}

فونت ریاضی پیش‌فرض \lr{MatinBook}، \lr{XITS
Math} است. این فونت، پوشش گسترده‌ای از
نمادهای ریاضی دارد و با \lr{unicode-math}
سازگار است.

\begin{latin}
\begin{minted}{latex}
% |\pc{در فایل}|\lr{matinbook.cls}:
\IfFontExistsTF{XITS Math}{
    \setmathfont{XITS Math}
}{
    \IfFontExistsTF{XITSMath-Regular.otf}{
        \setmathfont{XITSMath-Regular.otf}
    }{
        \IfFontExistsTF{Latin Modern Math}{
            \setmathfont{Latin Modern Math}
        }{
            \PackageWarning{matinbook}{No math font found}
        }
    }
}
\end{minted}
\end{latin}

\mnote{فونت \lr{XITS Math} از \lr{OpenType} پشتیبانی می‌کند و نمادهای ریاضی را با کیفیت بالا نمایش می‌دهد.}

\section{فونت مونواسپیس}
\label{sec:fonts-monospace}

فونت مونواسپیس پیش‌فرض \lr{MatinBook}،
\lr{DejaVu Sans Mono} است. این فونت برای
نمایش کد، اعداد و متن‌های فنی مناسب است.

\begin{latin}
\begin{minted}{latex}
% |\pc{در فایل}|\lr{matinbook.cls}:
\IfFontExistsTF{DejaVu Sans Mono}{
    \setmonofont{DejaVu Sans Mono}[Scale=0.9]
}{
    \IfFontExistsTF{Courier New}{
        \setmonofont{Courier New}[Scale=0.9]
    }{
        \PackageWarning{matinbook}{No monospace font found}
    }
}
\end{minted}
\end{latin}

\mnote{فونت \lr{DejaVu Sans Mono} روی اکثر توزیع‌های \lr{Linux} به‌طور پیش‌فرض نصب است.}

\section{فونت کد}
\label{sec:fonts-code}

فونت کد پیش‌فرض \lr{MatinBook}، \lr{JetBrains
Mono} است. این فونت، لیگچرهای مخصوص
برنامه‌نویسی دارد که خوانایی کد را افزایش
می‌دهند.

\begin{latin}
\begin{minted}{latex}
% |\pc{در فایل}|\lr{matinbook.cls}:
\IfFontExistsTF{JetBrains Mono}{
    \newfontfamily\codefont{JetBrains Mono}[Scale=0.9,Ligatures=TeX]
}{
    \IfFontExistsTF{DejaVu Sans Mono}{
        \newfontfamily\codefont{DejaVu Sans Mono}[Scale=0.9,Ligatures=TeX]
    }{
        \newfontfamily\codefont{Courier New}[Scale=0.9]
    }
}
\end{minted}
\end{latin}

\mnote{اگر \lr{JetBrains Mono} نصب نباشد، \lr{DejaVu Sans Mono} به‌عنوان جایگزین استفاده می‌شود.}

\section{فونت فارسی در کد}
\label{sec:fonts-persian-code}

برای نمایش متن فارسی داخل کد، \lr{MatinBook}
از فونت \lr{DejaVu Sans Mono} استفاده می‌کند
که از حروف فارسی پشتیبانی می‌کند.

\begin{latin}
\begin{minted}{latex}
% |\pc{در فایل}|\lr{mb-code.sty}:
\defpersianfont\persiancodingfont{DejaVu Sans Mono}[Script=Arabic,Scale=0.9]
\newcommand{\pc}[1]{\rl{\persiancodingfont\rl{#1}}}
\end{minted}
\end{latin}

\mnote{دستور \lr{\textbackslash pc} برای متن فارسی غیرکامنت و \lr{\textbackslash pcc} برای کامنت فارسی در پایتون استفاده می‌شود.}

\section{نصب فونت‌ها}
\label{sec:fonts-installation}

\subsection{نصب \lr{XB Niloofar}}
\label{sec:fonts-install-niloofar}

\begin{latin}
\begin{minted}{bash}
mkdir -p ~/.local/share/fonts
cp fonts/niloofar/*.ttf ~/.local/share/fonts/
fc-cache -fv
\end{minted}
\end{latin}

\subsection{نصب \lr{XITS Math}}
\label{sec:fonts-install-xits}

\begin{latin}
\begin{minted}{bash}
# Debian/Ubuntu
sudo apt install fonts-xits

# macOS
brew install --cask font-xits
\end{minted}
\end{latin}

\subsection{نصب \lr{DejaVu Sans Mono}}
\label{sec:fonts-install-dejavu}

\begin{latin}
\begin{minted}{bash}
# Debian/Ubuntu
sudo apt install fonts-dejavu
\end{minted}
\end{latin}

\mnote{برای نصب فونت‌ها روی ویندوز، روی فایل \lr{.ttf} راست‌کلیک کرده و \lr{Install} بزنید.}

\section{تغییر فونت پیش‌فرض}
\label{sec:fonts-change-default}

اگر می‌خواهید فونت پیش‌فرض را برای همیشه
تغییر دهید، فایل \file{matinbook.cls} را
ویرایش کنید:

\begin{latin}
\begin{minted}{latex}
% |\pc{مقدار پیش‌فرض}|\lr{XBNiloofar}|%
\def\matin@fontname{XBNiloofar}

% |\pc{تغییر به}|\lr{Vazirmatn}|%
\def\matin@fontname{Vazirmatn}
\end{minted}
\end{latin}

\mnote{با این تغییر، فونت پیش‌فرض کلاس \lr{Vazirmatn} می‌شود و نیازی به آپشن در هر سند نیست.}

\section{خطاهای رایج}
\label{sec:fonts-mistakes}

\subsection{فونت پیدا نشد}
\label{sec:fonts-mistake-notfound}

\textbf{مشکل:} اگر فونت روی سیستم نصب نباشد،
خطای \lr{Font X not found} می‌دهد.

\textbf{راه‌حل:} فونت را نصب کنید و
\lr{fc-cache -fv} را اجرا کنید.

\subsection{فونت تجاری}
\label{sec:fonts-mistake-commercial}

\textbf{مشکل:} فونت \lr{B Nazanin} تجاری
است و ممکن است روی همه‌ی سیستم‌ها نصب
نباشد.

\textbf{راه‌حل:} از فونت‌های رایگان مانند
\lr{XB Niloofar} یا \lr{Vazirmatn}
استفاده کنید.

\subsection{فونت فارسی در کد}
\label{sec:fonts-mistake-persian-code}

\textbf{مشکل:} اگر برای متن فارسی در کد از
فونت لاتین استفاده کنید، نتیجه مربع‌های
خالی می‌شود.

\textbf{راه‌حل:} از دستورات \cmd{pc} و
\cmd{pcc} استفاده کنید که فونت فارسی
مناسب را فعال می‌کنند.

\section{تمرین‌ها}
\label{sec:fonts-exercises}

\begin{exercise}
    یک کتاب با فونت \lr{Vazirmatn} بسازید و
    آن را با \lr{XB Niloofar} مقایسه کنید.
    \label{exc:fonts-vazirmatn}
\end{exercise}

\begin{exercise}
    فونت پیش‌فرض کلاس را به \lr{Sahel}
    تغییر دهید.
    \label{exc:fonts-default}
\end{exercise}

\begin{exercise}
    یک نمونه کد با کامنت فارسی بنویسید و
    فونت \lr{DejaVu Sans Mono} را ببینید.
    \label{exc:fonts-code}
\end{exercise}

\begin{exercise}
    نصب بودن فونت‌های اصلی را با \lr{fc-list}
    بررسی کنید.
    \label{exc:fonts-check}
\end{exercise}

\section{خلاصه}
\label{sec:fonts-summary}

در این فصل، با فونت‌های \lr{MatinBook}
آشنا شدیم:

\begin{itemize}
    \item \lr{MatinBook} پنج فونت فارسی
    را پشتیبانی می‌کند: \lr{niloofar}،
    \lr{vazirmatn}، \lr{sahel}، \lr{irlotus}،
    \lr{bnazanin}.

    \item فونت پیش‌فرض، \lr{XB Niloofar} با
    مقیاس \pnum{۱.۲} است.

    \item فونت لاتین، \lr{Times New Roman}
    (با جایگزین \lr{TeX Gyre Termes}).

    \item فونت ریاضی، \lr{XITS Math}.

    \item فونت مونواسپیس، \lr{DejaVu Sans Mono}.

    \item فونت کد، \lr{JetBrains Mono} با
    لیگچرهای برنامه‌نویسی.

    \item متن فارسی در کد با \lr{DejaVu Sans Mono}
    و دستورات \cmd{pc} و \cmd{pcc} نمایش داده
    می‌شود.

    \item خطاهای رایج شامل فونت پیدا نشد،
    فونت تجاری، و فونت فارسی در کد است.
\end{itemize}

در فصل بعدی (\cref{chap:layout})، با
صفحه‌آرایی در \lr{MatinBook} آشنا می‌شوید
و یاد می‌گیرید که چگونه چیدمان کتاب را
شخصی‌سازی کنید.
%%
%% part3-customization/ch20-layout.tex
%%
%% Chapter 20: Page Layout — Geometry, margins, line
%% spacing, headers, footers, and paragraph settings.
%%

\chapter{صفحه‌آرایی}
\label{chap:layout}

\section{چرا این موضوع مهم است؟}
\label{sec:layout-why}

صفحه‌آرایی، پایه‌ی کیفیت بصری یک کتاب است.
حاشیه‌های درست، فاصله‌ی خطوط مناسب، سرصفحه‌ی
خوانا و پاصفحه‌ی تمیز، همگی به خوانایی و
زیبایی کتاب کمک می‌کنند.

\lr{MatinBook} یک چیدمان حرفه‌ای بر اساس
استاندارد \lr{Boyer/Stewart} ارائه می‌دهد.
در این فصل، با اجزای این چیدمان آشنا
می‌شوید و یاد می‌گیرید که چگونه آن را
شخصی‌سازی کنید.

\mnote{چیدمان \lr{MatinBook} از کتاب‌های درسی دانشگاهی مانند \lr{Boyer} و \lr{Stewart} الهام گرفته شده است.}

\section{هندسه‌ی صفحه}
\label{sec:layout-geometry}

\subsection{ابعاد کاغذ}
\label{sec:layout-paper}

\lr{MatinBook} از ابعاد \pnum{۱۷.۸} ×
\pnum{۲۵.۴} سانتی‌متر (نزدیک به \lr{7 × 10}
اینچ) استفاده می‌کند. این ابعاد، استاندارد
نشر آکادمیک بین‌المللی است.

\begin{latin}
\begin{minted}{latex}
% |\pc{در فایل}|\lr{mb-layout.sty}:
\geometry{
    paperwidth=17.8cm,
    paperheight=25.4cm,
    ...
}
\end{minted}
\end{latin}

\mnote{این ابعاد، نزدیک به نسبت طلایی است و برای کتاب‌های فنی، خوانایی بالایی دارد.}

\subsection{حاشیه‌ها در \lr{mainmatter}}
\label{sec:layout-margins-main}

در متن اصلی، \lr{MatinBook} از حاشیه‌ی
بیرونی پهن (\pnum{۵.۹} سانتی‌متر) استفاده
می‌کند:

\begin{latin}
\begin{minted}{latex}
\geometry{
    paperwidth=17.8cm,
    paperheight=25.4cm,
    top=1.7cm,
    bottom=1.7cm,
    inner=1.5cm,
    outer=5.9cm,              % |\pc{حاشیه‌ی پهن}|%
    marginparwidth=3.5cm,     % |\pc{عرض یادداشت}|%
    marginparsep=0.6cm,       % |\pc{فاصله از متن}|%
    headheight=15pt,
    headsep=5pt,
    footskip=18pt,
    bindingoffset=0.4cm
}
\end{minted}
\end{latin}

جدول \cref{tab:layout-margins} این مقادیر
را خلاصه می‌کند.

\begin{matintable}{حاشیه‌های متن اصلی}{tab:layout-margins}
\begin{tblr}{
    width = \textwidth,
    colspec = {l X[c] X[l]},
    hlines, vlines,
    colsep = 8pt,
    rowsep = 4pt,
    row{1} = {font=\bfseries},
}
حاشیه & مقدار & کاربرد \\
بالا & \pnum{۱.۷} سانتی‌متر & فاصله از لبه‌ی بالا \\
پایین & \pnum{۱.۷} سانتی‌متر & فاصله از لبه‌ی پایین \\
داخلی & \pnum{۱.۵} سانتی‌متر & نزدیک عطف \\
بیرونی & \pnum{۵.۹} سانتی‌متر & برای یادداشت حاشیه \\
عرض یادداشت & \pnum{۳.۵} سانتی‌متر & ناحیه‌ی یادداشت \\
فاصله از متن & \pnum{۰.۶} سانتی‌متر & بین متن و یادداشت \\
فضای صحافی & \pnum{۰.۴} سانتی‌متر & برای صحافی \\
\end{tblr}
\end{matintable}

\mnote{عرض متن در این حالت، حدود \pnum{۱۰.۳} سانتی‌متر است که برای متن فارسی با فونت \lr{XB Niloofar} بهینه است.}

\subsection{حاشیه‌ها در \lr{frontmatter} و \lr{backmatter}}
\label{sec:layout-margins-front}

در پیش‌متن و پس‌متن، هندسه‌ی صفحه متقارن
می‌شود (بدون حاشیه‌ی پهن):

\begin{latin}
\begin{minted}{latex}
% |\pc{در فایل}|\lr{mb-layout.sty}:
\newcommand{\frontmattergeometry}{%
    \newgeometry{
        paperwidth=17.8cm,
        paperheight=25.4cm,
        top=1.7cm,
        bottom=1.7cm,
        inner=1.5cm,
        outer=1.5cm,          % |\pc{حاشیه‌ی متقارن}|%
        marginparwidth=0pt,
        ...
    }%
}
\end{minted}
\end{latin}

\mnote{دلیل حذف حاشیه‌ی پهن در پیش‌متن و پس‌متن این است که این بخش‌ها به یادداشت حاشیه‌ای نیاز ندارند.}

\section{فاصله‌ی خطوط}
\label{sec:layout-spacing}

\lr{MatinBook} از فاصله‌ی خطوط \pnum{۱.۱۵}
استفاده می‌کند که مطابق سبک کتاب‌های درسی
\lr{Boyer/Stewart} است:

\begin{latin}
\begin{minted}{latex}
\setstretch{1.15}
\end{minted}
\end{latin}

\mnote{فاصله‌ی \pnum{۱.۱۵} برای متن فارسی با فونت \lr{XB Niloofar} بهینه است و خوانایی را افزایش می‌دهد.}

\subsection{تنظیمات پاراگراف}
\label{sec:layout-paragraph}

\begin{latin}
\begin{minted}{latex}
\setlength{\parindent}{1em}
\setlength{\parskip}{0pt}
\end{minted}
\end{latin}

\begin{itemize}
    \item \cmd{parindent}: تورفتگی پاراگراف.
    \item \cmd{parskip}: فاصله‌ی بین پاراگراف‌ها.
\end{itemize}

\mnote{در فارسی، تورفتگی \lr{۱em} استاندارد است. پاراگراف اول هر بخش، تورفتگی ندارد.}

\section{سرصفحه و پاصفحه}
\label{sec:layout-headers}

\lr{MatinBook} از \lr{fancyhdr} برای مدیریت
سرصفحه و پاصفحه استفاده می‌کند.

\subsection{صفحات فرد و زوج}
\label{sec:layout-headers-oddeven}

\begin{itemize}
    \item \textbf{صفحات فرد:} نام فصل در بالا
    راست، شماره صفحه در پایین راست.
    \item \textbf{صفحات زوج:} نام فصل در بالا
    چپ، شماره صفحه در پایین چپ.
\end{itemize}

\begin{latin}
\begin{minted}{latex}
\pagestyle{fancy}
\fancyhf{}

\fancyhead[LE]{%
    \ifnum\c@chapter>0
        \footnotesize\textbf{\leftmark}%
    \fi
}

\fancyhead[RO]{%
    \ifnum\c@chapter>0
        \footnotesize\textbf{\leftmark}%
    \fi
}

\fancyfoot[LE,RO]{\footnotesize\textbf{\thepage}}
\end{minted}
\end{latin}

\mnote{در \lr{MatinBook}، از \lr{\textbackslash leftmark} برای هر دو سرصفحه استفاده می‌شود. یعنی نام فصل در هر دو صفحه نمایش داده می‌شود.}

\subsection{صفحات شروع فصل}
\label{sec:layout-headers-plain}

صفحات شروع فصل، از استایل \lr{plain} استفاده
می‌کنند و فقط شماره صفحه را در پاصفحه
نمایش می‌دهند:

\begin{latin}
\begin{minted}{latex}
\fancypagestyle{plain}{%
    \fancyhf{}
    \fancyfoot[LE,RO]{\footnotesize\textbf{\thepage}}
    \renewcommand{\headrulewidth}{0pt}
    \renewcommand{\footrulewidth}{0pt}
}
\end{minted}
\end{latin}

\mnote{این طراحی تمیز، توجه خواننده را به ابتدای فصل جلب می‌کند و فضای بیشتری برای عنوان فصل فراهم می‌آورد.}

\subsection{خط جداکننده‌ی سرصفحه}
\label{sec:layout-headers-rule}

\begin{latin}
\begin{minted}{latex}
\renewcommand{\headrulewidth}{0.4pt}
\renewcommand{\footrulewidth}{0pt}
\end{minted}
\end{latin}

\mnote{خط جداکننده‌ی سرصفحه با ضخامت \pnum{۰.۴pt} طراحی شده است. پاصفحه خط ندارد.}

\section{تنظیمات کپشن}
\label{sec:layout-captions}

\begin{latin}
\begin{minted}{latex}
\captionsetup{
    font=footnotesize,
    labelfont=bf,
    skip=8pt,
    justification=centering
}
\end{minted}
\end{latin}

\mnote{کپشن‌ها با فونت \lr{footnotesize} و برچسب پررنگ نمایش داده می‌شوند.}

\section{تنظیمات پانویس}
\label{sec:layout-footnotes}

\begin{latin}
\begin{minted}{latex}
\setlength{\footnotesep}{8pt}
\renewcommand{\footnoterule}{%
    \kern-3pt
    \hrule width 0.3\textwidth height 0.4pt
    \kern2.6pt
}
\end{minted}
\end{latin}

\mnote{خط جداکننده‌ی پانویس، ۳۰٪ عرض متن است. این طراحی، استاندارد کتاب‌های آکادمیک است.}

\section{حالت \lr{draft}}
\label{sec:layout-draft}

در حالت \lr{draft}، \lr{MatinBook} واترمارک
\lr{DRAFT} را فعال می‌کند:

\begin{latin}
\begin{minted}{latex}
\if@matin@draft
    \AtBeginDocument{%
        \RequirePackage{draftwatermark}%
        \SetWatermarkText{\lr{DRAFT}}%
        \SetWatermarkScale{1.5}%
        \SetWatermarkColor[gray]{0.9}%
    }
\fi
\end{minted}
\end{latin}

\mnote{واترمارک در حالت \lr{draft} به‌عنوان یادآوری برای بازگشت به حالت \lr{final} عمل می‌کند.}

\section{شخصی‌سازی چیدمان}
\label{sec:layout-customization}

برای شخصی‌سازی چیدمان، فایل
\file{mb-layout.sty} را ویرایش کنید. سه
راه اصلی وجود دارد:

\begin{enumerate}
    \item \textbf{تغییر هندسه:} مقادیر
    \cmd{geometry} را تغییر دهید.
    \item \textbf{تغییر سرصفحه:} تنظیمات
    \cmd{fancyhdr} را تغییر دهید.
    \item \textbf{تغییر فاصله‌ها:} مقادیر
    \cmd{setstretch} و \cmd{parindent} را
    تغییر دهید.
\end{enumerate}

\mnote{تغییر هندسه‌ی صفحه، بر کل کتاب تأثیر می‌گذارد. قبل از تغییر، از کتاب پشتیبان بگیرید.}

\section{خطاهای رایج}
\label{sec:layout-mistakes}

\subsection{فراموش کردن \cmd{frontmattergeometry}}
\label{sec:layout-mistake-frontgeo}

\textbf{مشکل:} اگر \cmd{frontmattergeometry}
را فراخوانی نکنید، پیش‌متن با حاشیه‌ی پهن
نمایش داده می‌شود.

\textbf{راه‌حل:} همیشه پس از \cmd{frontmatter}،
دستور \cmd{frontmattergeometry} را فراخوانی
کنید.

\subsection{فراموش کردن \cmd{mainmattergeometry}}
\label{sec:layout-mistake-maingeo}

\textbf{مشکل:} اگر \cmd{mainmattergeometry}
را فراخوانی نکنید، متن اصلی با حاشیه‌ی
متقارن نمایش داده می‌شود.

\textbf{راه‌حل:} همیشه پس از \cmd{mainmatter}،
دستور \cmd{mainmattergeometry} را فراخوانی
کنید.

\subsection{سرریز متن به حاشیه}
\label{sec:layout-mistake-overflow}

\textbf{مشکل:} اگر متن یا شکل از حاشیه سرریز
کند، خطای \lr{Overfull hbox} می‌دهد.

\textbf{راه‌حل:} از \cmd{emergencystretch} و
\cmd{tolerance} استفاده کنید (که در
\lr{MatinBook} از پیش تنظیم شده است).

\section{تمرین‌ها}
\label{sec:layout-exercises}

\begin{exercise}
    حاشیه‌ی بیرونی را از \pnum{۵.۹} به
    \pnum{۶.۵} سانتی‌متر تغییر دهید و کتاب
    را دوباره کامپایل کنید.
    \label{exc:layout-margins}
\end{exercise}

\begin{exercise}
    فاصله‌ی خطوط را از \pnum{۱.۱۵} به
    \pnum{۱.۳} تغییر دهید و نتیجه را ببینید.
    \label{exc:layout-spacing}
\end{exercise}

\begin{exercise}
    رنگ سرصفحه را تغییر دهید.
    \label{exc:layout-headers}
\end{exercise}

\begin{exercise}
    کتاب را در حالت \lr{draft} کامپایل کنید
    و واترمارک را ببینید.
    \label{exc:layout-draft}
\end{exercise}

\section{خلاصه}
\label{sec:layout-summary}

در این فصل، با صفحه‌آرایی \lr{MatinBook}
آشنا شدیم:

\begin{itemize}
    \item ابعاد کاغذ \pnum{۱۷.۸} ×
    \pnum{۲۵.۴} سانتی‌متر.

    \item حاشیه‌ی بیرونی پهن \pnum{۵.۹}
    سانتی‌متر در متن اصلی.

    \item حاشیه‌ی متقارن \pnum{۱.۵} سانتی‌متر
    در پیش‌متن و پس‌متن.

    \item فاصله‌ی خطوط \pnum{۱.۱۵}.

    \item تورفتگی پاراگراف \pnum{۱em}.

    \item سرصفحه با نام فصل، پاصفحه با شماره
    صفحه.

    \item خط جداکننده‌ی سرصفحه با ضخامت
    \pnum{۰.۴pt}.

    \item تنظیمات کپشن و پانویس.

    \item حالت \lr{draft} با واترمارک.

    \item خطاهای رایج شامل فراموش کردن
    \cmd{frontmattergeometry}، فراموش کردن
    \cmd{mainmattergeometry}، و سرریز متن
    است.
\end{itemize}

در فصل بعدی (\cref{chap:themes})، با تم‌ها
در \lr{MatinBook} آشنا می‌شوید و یاد
می‌گیرید که چگونه تم سفارشی بسازید.
%%
%% part3-customization/ch21-themes.tex
%%
%% Chapter 21: Themes — The default theme, its components,
%% and how to switch themes.
%%

\chapter{تم‌ها}
\label{chap:themes}

\section{چرا این موضوع مهم است؟}
\label{sec:themes-why}

یک تم، مجموعه‌ای از تصمیمات طراحی است که
هویت بصری کتاب را تعیین می‌کند: رنگ‌ها،
فونت‌ها، جلد، و چیدمان. با تغییر تم،
می‌توانید کتاب خود را کاملاً متفاوت جلوه
دهید.

\lr{MatinBook} یک سیستم تم ماژولار ارائه
می‌دهد که به شما امکان می‌دهد تم پیش‌فرض را
استفاده کنید یا تم سفارشی بسازید. در این
فصل، با این سیستم آشنا می‌شوید.

\mnote{تم پیش‌فرض \lr{MatinBook}، \lr{default} نام دارد و تمام قابلیت‌های کلاس را فعال می‌کند.}

\section{تم پیش‌فرض}
\label{sec:themes-default}

تم پیش‌فرض \lr{MatinBook}، \lr{default} است.
این تم از سه فایل تشکیل شده است:

\begin{itemize}
    \item \file{mb-theme-colors.sty}: پالت
    رنگی \lr{CMYK}.
    \item \file{mb-theme-cover.sty}: طراحی
    جلد.
    \item \file{mb-theme-default.sty}:
    بارگذارنده‌ی تم.
\end{itemize}

\mnote{تم \lr{default} از ۲۰ رنگ \lr{CMYK} در ۵ خانواده استفاده می‌کند که مطابق استاندارد نشر آموزشی ایران است.}

\subsection{ساختار فایل‌های تم}
\label{sec:themes-structure}

تم‌ها در پوشه‌ی \file{tex/themes/} قرار
دارند:

\begin{latin}
\begin{minted}{text}
tex/themes/
└── default/
    ├── mb-theme-colors.sty
    ├── mb-theme-cover.sty
    └── mb-theme-default.sty
\end{minted}
\end{latin}

\mnote{هر تم، یک پوشه‌ی جداگانه با سه فایل \lr{.sty} دارد.}

\section{بارگذاری تم}
\label{sec:themes-loading}

\subsection{تم پیش‌فرض}
\label{sec:themes-loading-default}

تم پیش‌فرض به‌طور خودکار در \file{matinbook.cls}
بارگذاری می‌شود:

\begin{latin}
\begin{minted}{latex}
% |\pc{در فایل}|\lr{matinbook.cls}:
\RequirePackage{mb-theme-colors}
...
\RequirePackage{mb-theme-cover}
\end{minted}
\end{latin}

\mnote{ترتیب بارگذاری مهم است: \lr{mb-theme-colors} باید \lr{قبل از} مصرف‌کنندگان (مثل \lr{mb-boxes}) بارگذاری شود.}

\subsection{تم سفارشی}
\label{sec:themes-loading-custom}

برای بارگذاری یک تم سفارشی، فایل
\file{matinbook.cls} را ویرایش کنید:

\begin{latin}
\begin{minted}{latex}
% |\pc{به‌جای}|%
\RequirePackage{mb-theme-colors}
\RequirePackage{mb-theme-cover}

% |\pc{بنویسید}|%
\RequirePackage{tex/themes/mytheme/mb-theme-default}
\end{minted}
\end{latin}

\mnote{مسیر تم سفارشی باید نسبت به پوشه‌ی پروژه باشد.}

\section{ساختار یک تم}
\label{sec:themes-components}

هر تم از سه جزء تشکیل شده است:

\subsection{پالت رنگ}
\label{sec:themes-colors}

فایل \file{mb-theme-colors.sty} رنگ‌های تم
را تعریف می‌کند:

\begin{latin}
\begin{minted}{latex}
\providecolor{matinblue}{cmyk}{0.80, 0.30, 0.00, 0.30}
\providecolor{matingreen}{cmyk}{0.80, 0.00, 0.45, 0.30}
\providecolor{matinorange}{cmyk}{0.00, 0.45, 0.85, 0.10}
\providecolor{matinred}{cmyk}{0.00, 0.65, 0.75, 0.10}
\providecolor{matinpurple}{cmyk}{0.20, 0.60, 0.00, 0.30}
\end{minted}
\end{latin}

\mnote{استفاده از \lr{\textbackslash providecolor} به‌جای \lr{\textbackslash definecolor} امکان بازنویسی رنگ‌ها را در سند فراهم می‌کند.}

\subsection{جلد}
\label{sec:themes-cover}

فایل \file{mb-theme-cover.sty} جلد جلو و
عقب را طراحی می‌کند:

\begin{latin}
\begin{minted}{latex}
\newcommand{\makecover}[4]{%
    \begin{titlepage}\thispagestyle{empty}
        \mb@frontart
        ...
    \end{titlepage}%
}

\newcommand{\makebackcover}[1]{%
    \clearpage
    \thispagestyle{empty}%
    \mb@backart
    ...
}
\end{minted}
\end{latin}

\mnote{جلد با \lr{TikZ} ساخته می‌شود و به حداقل دو بار کامپایل نیاز دارد.}

\subsection{بارگذارنده‌ی تم}
\label{sec:themes-loader}

فایل \file{mb-theme-default.sty} اجزای تم
را بارگذاری می‌کند:

\begin{latin}
\begin{minted}{latex}
\RequirePackage{mb-theme-cover}
\end{minted}
\end{latin}

\mnote{فایل \lr{mb-theme-colors} اینجا بارگذاری نمی‌شود، چون در \lr{matinbook.cls} بارگذاری می‌شود.}

\section{ایجاد تم سفارشی}
\label{sec:themes-creating}

برای ایجاد یک تم سفارشی، این مراحل را
دنبال کنید:

\begin{enumerate}
    \item یک پوشه‌ی جدید در
    \file{tex/themes/} بسازید (مثلاً
    \file{mytheme/}).

    \item سه فایل تم را کپی کنید:

    \begin{latin}
\begin{minted}{bash}
cp tex/themes/default/*.sty tex/themes/mytheme/
\end{minted}
    \end{latin}

    \item فایل‌های تم را ویرایش کنید.

    \item در \file{matinbook.cls}، تم جدید را
    بارگذاری کنید.
\end{enumerate}

\mnote{برای شروع، از تم \lr{default} کپی بگیرید و آن را تغییر دهید. این ساده‌تر از نوشتن تم از صفر است.}

\subsection{تغییر رنگ‌ها}
\label{sec:themes-custom-colors}

ساده‌ترین راه برای تغییر تم، تغییر رنگ‌ها
است. مثلاً برای یک تم با رنگ اصلی سبز:

\begin{latin}
\begin{minted}{latex}
% |\pc{در فایل}|\lr{mb-theme-colors.sty}:
\providecolor{matinblue}{cmyk}{0.80, 0.00, 0.45, 0.30}  % |\pc{سبز}|%
\providecolor{matindarkblue}{cmyk}{0.80, 0.00, 0.45, 0.60}
\providecolor{matinlightblue}{cmyk}{0.10, 0.00, 0.05, 0.05}
\end{minted}
\end{latin}

\mnote{به‌جای تغییر نام رنگ‌ها، فقط مقادیر \lr{CMYK} آن‌ها را تغییر دهید. این کار، سازگاری با کلاس را حفظ می‌کند.}

\subsection{تغییر جلد}
\label{sec:themes-custom-cover}

برای تغییر جلد، فایل \file{mb-theme-cover.sty}
را ویرایش کنید. مثلاً برای تغییر رنگ
پس‌زمینه‌ی جلد:

\begin{latin}
\begin{minted}{latex}
\providecolor{matinnavy}{RGB}{30, 42, 58}
% |\pc{تغییر به قرمز تیره}|%
\providecolor{matinnavy}{RGB}{120, 30, 30}
\end{minted}
\end{latin}

\mnote{برای تغییرات ساده، فقط رنگ‌ها را تغییر دهید. برای طراحی کاملاً جدید، عناصر \lr{TikZ} را جابه‌جا یا حذف کنید.}

\section{چند تم نمونه}
\label{sec:themes-examples}

\subsection{تم \lr{minimal}}
\label{sec:themes-minimal}

یک تم مینیمال که فقط از دو رنگ استفاده
می‌کند:

\begin{latin}
\begin{minted}{latex}
\providecolor{matinblue}{cmyk}{0.00, 0.00, 0.00, 0.80}
\providecolor{matingreen}{cmyk}{0.00, 0.00, 0.00, 0.60}
\providecolor{matinorange}{cmyk}{0.00, 0.00, 0.00, 0.50}
\end{minted}
\end{latin}

\mnote{تم‌های مینیمال برای کتاب‌های آکادمیک و رسمی مناسب‌ترند.}

\subsection{تم \lr{vibrant}}
\label{sec:themes-vibrant}

یک تم پرانرژی با رنگ‌های زنده:

\begin{latin}
\begin{minted}{latex}
\providecolor{matinblue}{cmyk}{1.00, 0.00, 0.00, 0.00}
\providecolor{matingreen}{cmyk}{1.00, 0.00, 1.00, 0.00}
\providecolor{matinorange}{cmyk}{0.00, 0.50, 1.00, 0.00}
\end{minted}
\end{latin}

\mnote{تم‌های پرانرژی برای کتاب‌های آموزشی و مخاطبان جوان مناسب‌ترند.}

\section{خطاهای رایج}
\label{sec:themes-mistakes}

\subsection{عدم بارگذاری \file{mb-theme-colors}}
\label{sec:themes-mistake-colors}

\textbf{مشکل:} اگر \file{mb-theme-colors} را
بارگذاری نکنید، رنگ‌ها تعریف نمی‌شوند و
خطای \lr{Undefined color} می‌دهید.

\textbf{راه‌حل:} مطمئن شوید که
\file{mb-theme-colors} در فاز ۱ \lr{matinbook.cls}
بارگذاری می‌شود.

\subsection{ترتیب اشتباه بارگذاری}
\label{sec:themes-mistake-order}

\textbf{مشکل:} اگر \file{mb-theme-cover} را
قبل از \file{mb-theme-colors} بارگذاری کنید،
خطا می‌دهید.

\textbf{راه‌حل:} ترتیب صحیح:
\file{mb-theme-colors} → \file{mb-theme-cover}.

\subsection{تغییر نام رنگ‌ها}
\label{sec:themes-mistake-rename}

\textbf{مشکل:} اگر نام رنگ‌ها را تغییر دهید،
ماژول‌های مصرف‌کننده (مثل \file{mb-boxes})
خطا می‌دهند.

\textbf{راه‌حل:} نام‌ها را تغییر ندهید؛ فقط
مقادیر \lr{CMYK} را تغییر دهید.

\section{تمرین‌ها}
\label{sec:themes-exercises}

\begin{exercise}
    یک تم سفارشی بسازید که از رنگ‌های
    متفاوتی استفاده کند.
    \label{exc:themes-custom}
\end{exercise}

\begin{exercise}
    رنگ اصلی تم را از آبی به سبز تغییر
    دهید.
    \label{exc:themes-green}
\end{exercise}

\begin{exercise}
    یک تم مینیمال بسازید که فقط از سه
    رنگ خاکستری استفاده کند.
    \label{exc:themes-minimal}
\end{exercise}

\begin{exercise}
    تم خود را در \file{matinbook.cls} بارگذاری
    کنید و کتاب را کامپایل کنید.
    \label{exc:themes-load}
\end{exercise}

\section{خلاصه}
\label{sec:themes-summary}

در این فصل، با تم‌ها در \lr{MatinBook}
آشنا شدیم:

\begin{itemize}
    \item تم پیش‌فرض \lr{default} است و از
    سه فایل تشکیل شده.

    \item هر تم در پوشه‌ی
    \file{tex/themes/} قرار دارد.

    \item اجزای تم: \file{mb-theme-colors}،
    \file{mb-theme-cover}،
    \file{mb-theme-default}.

    \item برای ایجاد تم سفارشی، پوشه‌ی جدید
    بسازید و سه فایل را کپی کنید.

    \item ساده‌ترین راه تغییر تم، تغییر
    رنگ‌ها است.

    \item نام رنگ‌ها را تغییر ندهید؛ فقط
    مقادیر \lr{CMYK} را تغییر دهید.

    \item خطاهای رایج شامل عدم بارگذاری
    \file{mb-theme-colors}، ترتیب اشتباه
    بارگذاری، و تغییر نام رنگ‌ها است.
\end{itemize}

در فصل بعدی (\cref{chap:localization})،
با لوکال و زبان در \lr{MatinBook} آشنا
می‌شوید و یاد می‌گیرید که چگونه بین فارسی
و انگلیسی سوئیچ کنید.
%%
%% part3-customization/ch22-localization.tex
%%
%% Chapter 22: Localization — Persian (fa-IR) and English
%% (en-US) locales, language switching, and number/date
%% formatting.
%%

\chapter{لوکال و زبان}
\label{chap:localization}

\section{چرا این موضوع مهم است؟}
\label{sec:localization-why}

\lr{MatinBook} برای کتاب‌های فنی فارسی طراحی
شده است، ولی از کتاب‌های انگلیسی نیز پشتیبانی
می‌کند. لوکال‌ها (فایل‌های ترجمه) نام عناصر
مختلف کتاب را به زبان مورد نظر ترجمه می‌کنند.

در این فصل، با دو لوکال \lr{fa-IR} و
\lr{en-US} آشنا می‌شوید و یاد می‌گیرید که
چگونه بین آن‌ها سوئیچ کنید.

\mnote{لوکال \lr{fa-IR} به‌طور پیش‌فرض در \lr{MatinBook} بارگذاری می‌شود. لوکال \lr{en-US} برای کتاب‌های انگلیسی است.}

\section{لوکال فارسی (\lr{fa-IR})}
\label{sec:localization-fa}

لوکال \lr{fa-IR} نام عناصر \lr{LaTeX} را
به فارسی ترجمه می‌کند. جدول
\cref{tab:fa-names} این ترجمه‌ها را نشان
می‌دهد.

\begin{matintable}{نام‌های فارسی عناصر \lr{LaTeX}}{tab:fa-names}
\begin{tblr}{
    width = \textwidth,
    colspec = {l X[l] X[l]},
    hlines, vlines,
    colsep = 6pt,
    rowsep = 4pt,
    row{1} = {font=\bfseries},
}
نام \lr{LaTeX} & ترجمه‌ی فارسی & کاربرد \\
\lr{contentsname} & فهرست مطالب & فهرست \\
\lr{listfigurename} & فهرست تصاویر & فهرست شکل‌ها \\
\lr{listtablename} & فهرست جداول & فهرست جدول‌ها \\
\lr{bibname} & منابع و مراجع & کتاب‌نامه \\
\lr{indexname} & نمایه & نمایه \\
\lr{chaptername} & فصل & سرصفحه \\
\lr{appendixname} & پیوست & پیوست‌ها \\
\lr{proofname} & اثبات & محیط اثبات \\
\lr{figurename} & شکل & کپشن شکل \\
\lr{tablename} & جدول & کپشن جدول \\
\end{tblr}
\end{matintable}

\mnote{این نام‌ها در فایل \lr{fa-IR.sty} تعریف شده‌اند و با \lr{\textbackslash renewcommand} بازنویسی می‌شوند.}

\subsection{نام‌های \lr{cleveref}}
\label{sec:localization-fa-cref}

لوکال فارسی، نام‌های \lr{cleveref} را نیز
ترجمه می‌کند. جدول \cref{tab:fa-cref-names}
این ترجمه‌ها را نشان می‌دهد.

\begin{matintable}{نام‌های فارسی \lr{cleveref}}{tab:fa-cref-names}
\begin{tblr}{
    width = \textwidth,
    colspec = {l X[l] X[l]},
    hlines, vlines,
    colsep = 6pt,
    rowsep = 4pt,
    row{1} = {font=\bfseries},
}
نوع & مفرد & جمع \\
معادله & معادله & معادلات \\
فصل & فصل & فصول \\
بخش & بخش & بخش‌ها \\
شکل & شکل & اشکال \\
جدول & جدول & جداول \\
الگوریتم & الگوریتم & الگوریتم‌ها \\
قضیه & قضیه & قضایا \\
لم & لم & لم‌ها \\
نتیجه & نتیجه & نتایج \\
تعریف & تعریف & تعاریف \\
مثال & مثال & مثال‌ها \\
نکته & نکته & نکات \\
تمرین & تمرین & تمرین‌ها \\
گزاره & گزاره & گزاره‌ها \\
\end{tblr}
\end{matintable}

\mnote{این نام‌ها با \lr{\textbackslash crefname} تعریف می‌شوند و در \lr{mb-theorem.sty} و \lr{fa-IR.sty} تنظیم شده‌اند.}

\subsection{نام‌های محیط‌های علمی}
\label{sec:localization-fa-theorems}

نام‌های محیط‌های علمی نیز در لوکال فارسی
تعریف شده‌اند:

\begin{itemize}
    \item \cmd{theoremname}: قضیه.
    \item \cmd{lemmaname}: لم.
    \item \cmd{corollaryname}: نتیجه.
    \item \cmd{propositionname}: گزاره.
    \item \cmd{definitionname}: تعریف.
    \item \cmd{examplename}: مثال.
    \item \cmd{remarkname}: نکته.
    \item \cmd{exercisename}: تمرین.
    \item \cmd{solutionname}: راه‌حل.
\end{itemize}

\mnote{این نام‌ها در عنوان جعبه‌های علمی استفاده می‌شوند و در فایل \lr{fa-IR.sty} تعریف شده‌اند.}

\section{لوکال انگلیسی (\lr{en-US})}
\label{sec:localization-en}

لوکال \lr{en-US} نسخه‌ی انگلیسی
\lr{fa-IR} است و برای کتاب‌های انگلیسی
استفاده می‌شود.

\begin{latin}
\begin{minted}{latex}
% |\pc{در فایل}|\lr{matinbook.cls}:
% |\pc{به‌جای}|%
\RequirePackage{fa-IR}

% |\pc{بنویسید}|%
\RequirePackage{en-US}
\end{minted}
\end{latin}

\mnote{لوکال \lr{en-US} از ساختار مشابه \lr{fa-IR} پیروی می‌کند ولی ترجمه‌ها انگلیسی هستند.}

\section{سوئیچ بین زبان‌ها}
\label{sec:localization-switching}

\lr{MatinBook} در حال حاضر از سوئیچ زبان
در زمان اجرا پشتیبانی نمی‌کند. یعنی
نمی‌توانید در وسط سند، زبان را از فارسی به
انگلیسی تغییر دهید.

برای سوئیچ، باید کلاس را با لوکال انگلیسی
کامپایل کنید:

\begin{latin}
\begin{minted}{latex}
% |\pc{در فایل}|\lr{matinbook.cls}:
\RequirePackage{en-US}
\end{minted}
\end{latin}

\mnote{پشتیبانی از سوئیچ زبان در زمان اجرا، یکی از برنامه‌های آینده‌ی \lr{MatinBook} است.}

\section{متن دوزبانه}
\label{sec:localization-bilingual}

در کتاب‌های دوزبانه، می‌توانید متن فارسی و
انگلیسی را با هم ترکیب کنید:

\begin{latin}
\begin{minted}{latex}
|\pc{این کتاب درباره}|\lr{machine learning}|\pc{و}|\lr{deep learning}|\pc{است.}|%
\end{minted}
\end{latin}

خروجی: این کتاب درباره \lr{machine learning}
و \lr{deep learning} است.

\mnote{برای متن لاتین در متن فارسی، از \lr{\textbackslash lr} استفاده کنید. برای پاراگراف کامل لاتین، از محیط \lr{latin}.}

\section{قالب اعداد}
\label{sec:localization-numbers}

\subsection{اعداد در متن فارسی}
\label{sec:localization-numbers-fa}

در متن فارسی، اعداد به‌طور خودکار به ارقام
فارسی (۰-۹) تبدیل می‌شوند:

\begin{latin}
\begin{minted}{latex}
|\pc{این کتاب شامل ۱۲ فصل است.}|%
\end{minted}
\end{latin}

خروجی: این کتاب شامل ۱۲ فصل است.

\subsection{اعداد در متن لاتین}
\label{sec:localization-numbers-latin}

در متن لاتین، اعداد به‌صورت لاتین (0-9)
نمایش داده می‌شوند:

\begin{latin}
\begin{minted}{latex}
\begin{latin}
The book has 12 chapters.
\end{latin}
\end{minted}
\end{latin}

خروجی:

\begin{latin}
The book has 12 chapters.
\end{latin}

\subsection{اعداد اعشاری فارسی}
\label{sec:localization-numbers-decimal}

برای اعداد اعشاری فارسی، از دستور \cmd{pnum}
استفاده کنید:

\begin{latin}
\begin{minted}{latex}
|\pc{مقیاس}|\pnum{۱.۲}|\pc{به این معناست...}|%
\end{minted}
\end{latin}

خروجی: مقیاس \pnum{۱.۲} به این معناست...

\mnote{دستور \lr{\textbackslash pnum} از باگ معکوس‌شدن اعداد اعشاری در \lr{bidi} جلوگیری می‌کند.}

\subsection{اعداد در ریاضیات}
\label{sec:localization-numbers-math}

در ریاضیات، اعداد از بافت پیروی می‌کنند:

\begin{itemize}
    \item در متن فارسی: اعداد ریاضی به‌صورت
    فارسی نمایش داده می‌شوند.
    \item در \cmd{lr}: اعداد ریاضی به‌صورت
    لاتین نمایش داده می‌شوند.
\end{itemize}

\begin{latin}
\begin{minted}{latex}
|\pc{معادله}|\lr{$E = mc^{2}$}|\pc{...}|%
\end{minted}
\end{latin}

\mnote{برای نمایش اعداد لاتین در ریاضیات، از \lr{\textbackslash lr\{...\}} استفاده کنید.}

\section{قالب تاریخ}
\label{sec:localization-dates}

\subsection{تاریخ فارسی}
\label{sec:localization-dates-fa}

برای تاریخ فارسی، از تاریخ تقویم هجری
شمسی استفاده کنید:

\begin{latin}
\begin{minted}{latex}
\date{|\pc{مهر ۱۴۰۵}|}
\end{minted}
\end{latin}

خروجی: مهر ۱۴۰۵.

\mnote{دستور \lr{\textbackslash today} در حالت فارسی، تاریخ هجری شمسی را نمایش می‌دهد.}

\subsection{تاریخ انگلیسی}
\label{sec:localization-dates-en}

در حالت انگلیسی، از تاریخ میلادی استفاده
کنید:

\begin{latin}
\begin{minted}{latex}
\date{October 2026}
\end{minted}
\end{latin}

\section{شخصی‌سازی لوکال}
\label{sec:localization-custom}

برای شخصی‌سازی نام‌ها، فایل \file{fa-IR.sty}
یا \file{en-US.sty} را ویرایش کنید. مثلاً
برای تغییر «فهرست مطالب» به «محتویات»:

\begin{latin}
\begin{minted}{latex}
% |\pc{در فایل}|\lr{fa-IR.sty}:
\renewcommand{\contentsname}{|\pc{محتویات}|}
\end{minted}
\end{latin}

\mnote{برای تغییرات موقت، می‌توانید در سند خود نیز این دستورات را بازنویسی کنید.}

\section{خطاهای رایج}
\label{sec:localization-mistakes}

\subsection{فراموش کردن \cmd{lr}}
\label{sec:localization-mistake-lr}

\textbf{مشکل:} اگر متن لاتین را در متن فارسی
بدون \cmd{lr} بنویسید، جهت متن به‌هم
می‌ریزد.

\textbf{راه‌حل:} همیشه از \cmd{lr} برای متن
لاتین در متن فارسی استفاده کنید.

\subsection{فراموش کردن \cmd{pnum}}
\label{sec:localization-mistake-pnum}

\textbf{مشکل:} اگر اعداد اعشاری فارسی را
بدون \cmd{pnum} بنویسید، به‌صورت معکوس
نمایش داده می‌شوند (مثلاً \lr{۱.۲} می‌شود
\lr{۲.۱}).

\textbf{راه‌حل:} از \cmd{pnum} برای همه‌ی
اعداد اعشاری فارسی استفاده کنید.

\subsection{ترکیب لوکال‌ها}
\label{sec:localization-mistake-mixed}

\textbf{مشکل:} اگر هر دو لوکال \lr{fa-IR} و
\lr{en-US} را بارگذاری کنید، لوکال دوم
اولی را بازنویسی می‌کند.

\textbf{راه‌حل:} فقط یکی از لوکال‌ها را
بارگذاری کنید.

\section{تمرین‌ها}
\label{sec:localization-exercises}

\begin{exercise}
    یک کتاب با لوکال \lr{en-US} بسازید و
    ببینید که نام‌ها به انگلیسی نمایش
    داده می‌شوند.
    \label{exc:localization-en}
\end{exercise}

\begin{exercise}
    نام «فهرست مطالب» را به «محتویات» تغییر
    دهید.
    \label{exc:localization-custom}
\end{exercise}

\begin{exercise}
    یک پاراگراف فارسی با سه واژه‌ی لاتین
    بنویسید و از \cmd{lr} استفاده کنید.
    \label{exc:localization-lr}
\end{exercise}

\begin{exercise}
    یک عدد اعشاری فارسی را با و بدون
    \cmd{pnum} بنویسید و تفاوت را ببینید.
    \label{exc:localization-pnum}
\end{exercise}

\section{خلاصه}
\label{sec:localization-summary}

در این فصل، با لوکال و زبان در \lr{MatinBook}
آشنا شدیم:

\begin{itemize}
    \item \lr{MatinBook} از دو لوکال پشتیبانی
    می‌کند: \lr{fa-IR} (فارسی) و \lr{en-US}
    (انگلیسی).

    \item لوکال فارسی، نام‌های \lr{LaTeX} و
    \lr{cleveref} را به فارسی ترجمه می‌کند.

    \item لوکال انگلیسی، نسخه‌ی انگلیسی
    \lr{fa-IR} است.

    \item سوئیچ بین زبان‌ها با تغییر
    \cmd{RequirePackage\{fa-IR\}} به
    \cmd{RequirePackage\{en-US\}} در
    \file{matinbook.cls} انجام می‌شود.

    \item در متن دوزبانه، از \cmd{lr} برای
    متن لاتین در متن فارسی استفاده کنید.

    \item اعداد در متن فارسی به‌طور خودکار
    فارسی می‌شوند.

    \item برای اعداد اعشاری فارسی، از
    \cmd{pnum} استفاده کنید.

    \item تاریخ فارسی با تقویم هجری شمسی.

    \item خطاهای رایج شامل فراموش کردن
    \cmd{lr}، فراموش کردن \cmd{pnum}، و
    ترکیب لوکال‌ها است.
\end{itemize}

این فصل، آخرین فصل از بخش سوم (شخصی‌سازی)
بود. در بخش چهارم (مباحث پیشرفته)، با
ساختار ماژولار، تم سفارشی، دستورات سفارشی
و پروژه‌های بزرگ آشنا می‌شوید.

%========================================================================
% PART 4 — ADVANCED
%========================================================================
\part{مباحث پیشرفته}
\label{part:advanced}

%%
%% part4-advanced/ch23-modules.tex
%%
%% Chapter 23: Modular Structure — The 19 modules of
%% MatinBook, their roles, and the loading order.
%%

\chapter{ساختار ماژولار}
\label{chap:modules}

\section{چرا این موضوع مهم است؟}
\label{sec:modules-why}

\lr{MatinBook} یک کلاس ماژولار است. یعنی
به‌جای یک فایل بزرگ، از ۱۹ ماژول مستقل
تشکیل شده که هر یک مسئولیت مشخصی دارند.
این معماری، نگهداری، توسعه و شخصی‌سازی
کلاس را بسیار آسان‌تر می‌کند.

درک ساختار ماژولار، برای شخصی‌سازی پیشرفته
و افزودن قابلیت‌های جدید ضروری است. در این
فصل، با این ساختار آشنا می‌شوید.

\mnote{معماری ماژولار یعنی اگر فقط می‌خواهید یک ماژول را تغییر دهید، نیازی به دست زدن به بقیه‌ی کلاس ندارید.}

\section{فهرست ماژول‌ها}
\label{sec:modules-list}

\lr{MatinBook} ۱۹ ماژول دارد که در سه دسته
تقسیم می‌شوند:

\subsection{هسته (۲ ماژول)}
\label{sec:modules-core}

\begin{itemize}
    \item \file{mb-core.sty}: بارگذارنده‌ی
    بسته‌های پایه.
    \item \file{mb-utils.sty}: دستورات کمکی.
\end{itemize}

\subsection{ماژول‌های کاربردی (۱۰ ماژول)}
\label{sec:modules-feature}

\begin{itemize}
    \item \file{mb-math.sty}: ریاضیات.
    \item \file{mb-boxes.sty}: استایل‌های
    جعبه.
    \item \file{mb-theorem.sty}: محیط‌های
    قضیه.
    \item \file{mb-code.sty}: کد.
    \item \file{mb-layout.sty}: چیدمان.
    \item \file{mb-headings.sty}: عناوین.
    \item \file{mb-references.sty}: ارجاعات.
    \item \file{mb-graphics.sty}: گرافیک.
    \item \file{mb-index.sty}: نمایه.
    \item \file{mb-table.sty}: جداول.
\end{itemize}

\subsection{تم و لوکال (۷ ماژول)}
\label{sec:modules-theme}

\begin{itemize}
    \item \file{mb-theme-colors.sty}: پالت
    رنگی.
    \item \file{mb-theme-cover.sty}: جلد.
    \item \file{mb-theme-default.sty}:
    بارگذارنده‌ی تم.
    \item \file{mb-typography.sty}:
    تایپوگرافی.
    \item \file{fa-IR.sty}: لوکال فارسی.
    \item \file{en-US.sty}: لوکال انگلیسی.
\end{itemize}

\mnote{جمع این ماژول‌ها ۱۹ عدد است: ۲ هسته + ۱۰ کاربردی + ۳ تم + ۲ لوکال + ۱ تایپوگرافی + ۱ کلاس اصلی.}

\section{نقش هر ماژول}
\label{sec:modules-roles}

جدول \cref{tab:module-roles} نقش هر ماژول
را خلاصه می‌کند.

\begin{matintable}{نقش ماژول‌های \lr{MatinBook}}{tab:module-roles}
\begin{tblr}{
    width = \textwidth,
    colspec = {l X[l]},
    hlines, vlines,
    colsep = 8pt,
    rowsep = 4pt,
    row{1} = {font=\bfseries},
}
ماژول & نقش \\
\lr{mb-core} & بارگذاری بسته‌های پایه \\
\lr{mb-utils} & دستورات کمکی \\
\lr{mb-math} & عملگرها و جداکننده‌های ریاضی \\
\lr{mb-boxes} & استایل‌های \lr{tcolorbox} \\
\lr{mb-theorem} & محیط‌های قضیه \\
\lr{mb-code} & \lr{minted} و پشتیبانی فارسی \\
\lr{mb-layout} & هندسه، سرصفحه، پاصفحه \\
\lr{mb-headings} & سبک عناوین \\
\lr{mb-references} & \lr{hyperref} و \lr{cleveref} \\
\lr{mb-graphics} & \lr{TikZ} و \lr{PGFPlots} \\
\lr{mb-index} & \lr{xindy} \\
\lr{mb-table} & \lr{tabularray} \\
\lr{mb-theme-colors} & پالت رنگی \lr{CMYK} \\
\lr{mb-theme-cover} & جلد جلو و عقب \\
\lr{mb-theme-default} & بارگذارنده‌ی تم \\
\lr{mb-typography} & \lr{microtype} و کشیده \\
\lr{fa-IR} & لوکال فارسی \\
\lr{en-US} & لوکال انگلیسی \\
\lr{matinbook.cls} & کلاس اصلی \\
\end{tblr}
\end{matintable}

\section{ترتیب بارگذاری}
\label{sec:modules-order}

ترتیب بارگذاری ماژول‌ها \textbf{حیاتی}
است. \lr{MatinBook} از سه فاز استفاده می‌کند:

\subsection{فاز ۱ — قبل از \lr{xepersian}}
\label{sec:modules-phase1}

\begin{latin}
\begin{minted}{text}
1.  mb-core
2.  mb-theme-colors
3.  mb-table
4.  mb-math
5.  mb-boxes
6.  mb-theorem
7.  mb-code
8.  mb-layout
9.  mb-headings
10. mb-references
11. mb-graphics
12. mb-index
13. mb-utils
14. mb-theme-cover
\end{minted}
\end{latin}

\mnote{در فاز ۱، ترتیب دقیق مهم است: \lr{mb-theme-colors} باید قبل از مصرف‌کنندگان (مثل \lr{mb-boxes}) بارگذاری شود.}

\subsection{فاز ۲ — \lr{xepersian}}
\label{sec:modules-phase2}

\begin{latin}
\begin{minted}{text}
15. xepersian
\end{minted}
\end{latin}

\mnote{\lr{xepersian} باید \textbf{آخرین} بسته‌ی فاز ۲ باشد. اگر بسته‌ای بعد از آن بارگذاری شود، تعاریف \lr{bidi} را بازنویسی می‌کند.}

\subsection{فاز ۳ — بعد از \lr{xepersian}}
\label{sec:modules-phase3}

\begin{latin}
\begin{minted}{text}
16. mb-typography
17. fa-IR
18. mb-theme-default
\end{minted}
\end{latin}

\mnote{در فاز ۳، فونت‌ها، تایپوگرافی و لوکال بارگذاری می‌شوند. این‌ها نیاز به \lr{xepersian} دارند.}

\section{چرا این ترتیب مهم است؟}
\label{sec:modules-why-order}

\begin{enumerate}
    \item \textbf{\lr{xepersian} آخرین بسته}:
    چون \lr{bidi} و سایر بسته‌ها را بازنویسی
    می‌کند. اگر بسته‌ای بعد از آن بارگذاری
    شود، تعاریف آن از بین می‌رود.

    \item \textbf{رنگ‌ها قبل از مصرف‌کنندگان}:
    \file{mb-theme-colors} قبل از
    \file{mb-boxes}، \file{mb-code}،
    \file{mb-layout}، \file{mb-headings} و
    \file{mb-graphics} بارگذاری می‌شود.

    \item \textbf{فونت‌ها بعد از \lr{xepersian}}:
    \cmd{settextfont} توسط \lr{xepersian}
    تعریف می‌شود، پس فونت‌ها باید در فاز ۳
    تنظیم شوند.
\end{enumerate}

\mnote{تغییر این ترتیب، باعث خطاهای عجیب و غریب می‌شود. اگر ماژول جدیدی اضافه می‌کنید، حتماً آن را در فاز مناسب قرار دهید.}

\section{ساختار پوشه‌ها}
\label{sec:modules-structure}

ماژول‌ها در پوشه‌ی \file{tex/} قرار دارند:

\begin{latin}
\begin{minted}{text}
tex/
├── core/
│   ├── mb-core.sty
│   └── mb-utils.sty
├── locales/
│   ├── fa-IR.sty
│   └── en-US.sty
├── modules/
│   ├── boxes/mb-boxes.sty
│   ├── code/mb-code.sty
│   ├── graphics/mb-graphics.sty
│   ├── index/mb-index.sty
│   ├── layout/
│   │   ├── mb-layout.sty
│   │   └── mb-headings.sty
│   ├── math/mb-math.sty
│   ├── references/mb-references.sty
│   ├── table/mb-table.sty
│   ├── theorem/mb-theorem.sty
│   └── typography/mb-typography.sty
└── themes/
    └── default/
        ├── mb-theme-colors.sty
        ├── mb-theme-cover.sty
        └── mb-theme-default.sty
\end{minted}
\end{latin}

\mnote{هر ماژول در پوشه‌ی مخصوص خود قرار دارد. این ساختار، پیدا کردن ماژول‌ها را آسان می‌کند.}

\section{افزودن ماژول جدید}
\label{sec:modules-adding}

برای افزودن یک ماژول جدید، این مراحل را
دنبال کنید:

\begin{enumerate}
    \item یک پوشه‌ی جدید در
    \file{tex/modules/} بسازید (مثلاً
    \file{mymodule/}).

    \item یک فایل \file{.sty} در آن بسازید
    (مثلاً \file{mb-mymodule.sty}).

    \item فایل را با ساختار استاندارد
    بنویسید.

    \item در \file{matinbook.cls}، ماژول را
    در فاز مناسب بارگذاری کنید.
\end{enumerate}

\subsection{ساختار یک ماژول}
\label{sec:modules-template}

\begin{latin}
\begin{minted}{latex}
\NeedsTeXFormat{LaTeX2e}[2023/06/01]
\ProvidesPackage{mb-mymodule}[2026/10/08 v1.2 MatinBook My Module]

%========================================================================
% MB-MYMODULE — DESCRIPTION
%========================================================================

% |\pc{کد ماژول...}|%

\PackageInfo{mb-mymodule}{My module loaded}

\endinput
\end{minted}
\end{latin}

\mnote{هر ماژول باید \lr{\textbackslash NeedsTeXFormat}، \lr{\textbackslash ProvidesPackage}، \lr{\textbackslash PackageInfo} و \lr{\textbackslash endinput} داشته باشد.}

\section{ساختار کلاس اصلی}
\label{sec:modules-cls}

فایل \file{matinbook.cls} سه فاز دارد:

\begin{latin}
\begin{minted}{latex}
% |\pc{فاز ۱: بسته‌های پایه}|%
\RequirePackage{mb-core}
\RequirePackage{mb-theme-colors}
\RequirePackage{mb-table}
\RequirePackage{mb-math}
\RequirePackage{mb-boxes}
\RequirePackage{mb-theorem}
\RequirePackage{mb-code}
\RequirePackage{mb-layout}
\RequirePackage{mb-headings}
\RequirePackage{mb-references}
\RequirePackage{mb-graphics}
\RequirePackage{mb-index}
\RequirePackage{mb-utils}
\RequirePackage{mb-theme-cover}

% |\pc{فاز ۲:}|\lr{xepersian}|%
\RequirePackage{xepersian}

% |\pc{فاز ۳: تنظیمات فارسی}|%
% |\pc{فونت‌ها...}|%
\RequirePackage{mb-typography}
\RequirePackage{fa-IR}
\end{minted}
\end{latin}

\mnote{در \lr{matinbook.cls}، هیچ ماژولی نباید خارج از این سه فاز بارگذاری شود.}

\section{اشکال‌زدایی}
\label{sec:modules-debugging}

اگر خطایی رخ داد، این مراحل را انجام
دهید:

\begin{enumerate}
    \item \textbf{بررسی فاز:} آیا ماژول در
    فاز درست بارگذاری شده؟

    \item \textbf{بررسی ترتیب:} آیا ترتیب
    ماژول‌ها درست است؟

    \item \textbf{بررسی وابستگی:} آیا ماژول
    به بسته‌ای نیاز دارد که بارگذاری
    نشده؟

    \item \textbf{بررسی نسخه:} آیا نسخه‌ی
    \lr{LaTeX} و بسته‌ها سازگار هستند؟
\end{enumerate}

\mnote{برای اشکال‌زدایی، می‌توانید از \lr{\textbackslash listfiles} در ابتدای سند استفاده کنید تا فهرست همه‌ی فایل‌های بارگذاری‌شده را ببینید.}

\section{خطاهای رایج}
\label{sec:modules-mistakes}

\subsection{بارگذاری ماژول در فاز اشتباه}
\label{sec:modules-mistake-phase}

\textbf{مشکل:} اگر ماژولی که به
\lr{xepersian} نیاز دارد را در فاز ۱
بارگذاری کنید، خطا می‌دهد.

\textbf{راه‌حل:} ماژول‌های وابسته به
\lr{xepersian} را در فاز ۳ بارگذاری کنید.

\subsection{بارگذاری بسته بعد از \lr{xepersian}}
\label{sec:modules-mistake-after}

\textbf{مشکل:} اگر بسته‌ای را بعد از
\lr{xepersian} بارگذاری کنید، تعاریف
\lr{bidi} بازنویسی می‌شوند و خطا می‌دهید.

\textbf{راه‌حل:} \lr{xepersian} را آخرین
بسته‌ی فاز ۲ نگه دارید.

\subsection{فراموش کردن \cmd{RequirePackage}}
\label{sec:modules-mistake-require}

\textbf{مشکل:} اگر ماژول را در
\file{matinbook.cls} بارگذاری نکنید،
استفاده نمی‌شود.

\textbf{راه‌حل:} همیشه ماژول‌های جدید را
در \file{matinbook.cls} اضافه کنید.

\section{تمرین‌ها}
\label{sec:modules-exercises}

\begin{exercise}
    فهرست ۱۹ ماژول \lr{MatinBook} را بنویسید
    و نقش هر یک را توضیح دهید.
    \label{exc:modules-list}
\end{exercise}

\begin{exercise}
    یک ماژول ساده بسازید که یک دستور جدید
    تعریف کند و آن را در فاز ۱ بارگذاری
    کنید.
    \label{exc:modules-add}
\end{exercise}

\begin{exercise}
    یک ماژول بسازید که به \lr{xepersian}
    نیاز دارد و آن را در فاز ۳ بارگذاری
    کنید.
    \label{exc:modules-phase3}
\end{exercise}

\begin{exercise}
    ترتیب بارگذاری ماژول‌ها را عوض کنید و
    خطای آن را ببینید. سپس آن را اصلاح
    کنید.
    \label{exc:modules-order}
\end{exercise}

\section{خلاصه}
\label{sec:modules-summary}

در این فصل، با ساختار ماژولار \lr{MatinBook}
آشنا شدیم:

\begin{itemize}
    \item \lr{MatinBook} از ۱۹ ماژول مستقل
    تشکیل شده که در سه دسته تقسیم می‌شوند:
    هسته، کاربردی، تم و لوکال.

    \item ماژول‌ها در سه فاز بارگذاری
    می‌شوند: قبل از \lr{xepersian}،
    \lr{xepersian}، بعد از \lr{xepersian}.

    \item \lr{xepersian} باید آخرین بسته‌ی
    فاز ۲ باشد.

    \item رنگ‌ها باید قبل از مصرف‌کنندگان
    بارگذاری شوند.

    \item فونت‌ها باید بعد از \lr{xepersian}
    تنظیم شوند.

    \item ماژول جدید با ساختار استاندارد
    (\cmd{NeedsTeXFormat}،
    \cmd{ProvidesPackage}، \cmd{PackageInfo}،
    \cmd{endinput}) ساخته می‌شود.

    \item خطاهای رایج شامل بارگذاری ماژول
    در فاز اشتباه، بارگذاری بسته بعد از
    \lr{xepersian}، و فراموش کردن
    \cmd{RequirePackage} است.
\end{itemize}

در فصل بعدی (\cref{chap:custom-theme})،
با ساخت تم سفارشی آشنا می‌شوید و یاد
می‌گیرید که چگونه یک تم کاملاً جدید طراحی
کنید.
%%
%% part4-advanced/ch24-custom-theme.tex
%%
%% Chapter 24: Custom Theme — Design a complete theme from
%% scratch: colors, cover, and loader.
%%

\chapter{تم سفارشی}
\label{chap:custom-theme}

\section{چرا این موضوع مهم است؟}
\label{sec:custom-theme-why}

تم پیش‌فرض \lr{MatinBook} برای بیشتر کتاب‌ها
مناسب است، ولی گاهی نیاز دارید که تم را
کاملاً تغییر دهید. مثلاً:

\begin{itemize}
    \item کتاب شما باید هویت بصری ناشر را
    داشته باشد.
    \item می‌خواهید از رنگ‌های سازمانی خود
    استفاده کنید.
    \item نیاز به طراحی جلد کاملاً جدید دارید.
    \item می‌خواهید ظاهر کتاب با کتاب‌های دیگر
    مجموعه هماهنگ باشد.
\end{itemize}

در این فصل، یاد می‌گیرید که چگونه یک تم
سفارشی از صفر بسازید.

\mnote{برای شروع، از تم \lr{default} کپی بگیرید و آن را تغییر دهید. این ساده‌تر از نوشتن تم از صفر است.}

\section{مراحل ساخت تم}
\label{sec:custom-theme-steps}

ساخت یک تم سفارشی در پنج مرحله انجام
می‌شود:

\begin{enumerate}
    \item \textbf{ایجاد پوشه:} یک پوشه‌ی
    جدید در \file{tex/themes/}.
    \item \textbf{کپی فایل‌ها:} سه فایل تم
    پیش‌فرض را کپی کنید.
    \item \textbf{ویرایش رنگ‌ها:} پالت رنگی
    را تغییر دهید.
    \item \textbf{ویرایش جلد:} جلد را طراحی
    کنید.
    \item \textbf{بارگذاری:} تم جدید را در
    \file{matinbook.cls} بارگذاری کنید.
\end{enumerate}

\subsection{مرحله ۱ — ایجاد پوشه}
\label{sec:custom-theme-step1}

\begin{latin}
\begin{minted}{bash}
cd ~/Desktop/matinbook/tex/themes
mkdir mytheme
\end{minted}
\end{latin}

\mnote{نام پوشه را با حروف کوچک و بدون فاصله بنویسید. مثلاً \lr{mytheme}، \lr{corporate}، \lr{university}.}

\subsection{مرحله ۲ — کپی فایل‌ها}
\label{sec:custom-theme-step2}

\begin{latin}
\begin{minted}{bash}
cp default/mb-theme-colors.sty mytheme/
cp default/mb-theme-cover.sty mytheme/
cp default/mb-theme-default.sty mytheme/
\end{minted}
\end{latin}

\mnote{سه فایل تم پیش‌فرض را کپی می‌کنیم تا ساختار درست را داشته باشیم.}

\subsection{مرحله ۳ — ویرایش رنگ‌ها}
\label{sec:custom-theme-step3}

فایل \file{mb-theme-colors.sty} را باز کنید
و رنگ‌ها را تغییر دهید. مهم‌ترین نکته:
\textbf{نام رنگ‌ها را تغییر ندهید}، فقط
مقادیر \lr{CMYK} را عوض کنید.

\begin{latin}
\begin{minted}{latex}
% |\pc{قبل: تم پیش‌فرض}|%
\providecolor{matinblue}{cmyk}{0.80, 0.30, 0.00, 0.30}

% |\pc{بعد: تم سفارشی}|%
\providecolor{matinblue}{cmyk}{0.70, 0.20, 0.00, 0.20}
\end{minted}
\end{latin}

\mnote{چرا نام‌ها را تغییر ندهیم؟ چون ماژول‌های دیگر (مثل \lr{mb-boxes} و \lr{mb-code}) به این نام‌ها وابسته‌اند.}

\subsection{مرحله ۴ — ویرایش جلد}
\label{sec:custom-theme-step4}

فایل \file{mb-theme-cover.sty} را باز کنید
و جلد را طراحی کنید. سه سطح تغییر وجود
دارد:

\begin{enumerate}
    \item \textbf{تغییر رنگ‌ها:} ساده‌ترین
    سطح.
    \item \textbf{تغییر متن:} تغییر عنوان‌ها
    و اطلاعات.
    \item \textbf{تغییر عناصر:} حذف یا
    افزودن عناصر \lr{TikZ}.
\end{enumerate}

\subsection{مرحله ۵ — بارگذاری}
\label{sec:custom-theme-step5}

در \file{matinbook.cls}، خط بارگذاری تم را
تغییر دهید:

\begin{latin}
\begin{minted}{latex}
% |\pc{قبل}|%
\RequirePackage{mb-theme-colors}
...
\RequirePackage{mb-theme-cover}

% |\pc{بعد}|%
\RequirePackage{tex/themes/mytheme/mb-theme-colors}
...
\RequirePackage{tex/themes/mytheme/mb-theme-cover}
\end{minted}
\end{latin}

\mnote{مسیر تم سفارشی نسبت به پوشه‌ی پروژه است.}

\section{طراحی پالت رنگی}
\label{sec:custom-theme-palette}

\subsection{انتخاب رنگ اصلی}
\label{sec:custom-theme-primary}

رنگ اصلی تم، هویت بصری کتاب را تعیین
می‌کند. پیشنهاد می‌شود رنگ اصلی را از
رنگ‌های زیر انتخاب کنید:

\begin{matintable}{پیشنهاد رنگ اصلی برای تم سفارشی}{tab:theme-primary}
\begin{tblr}{
    width = \textwidth,
    colspec = {l X[c] X[l]},
    hlines, vlines,
    colsep = 8pt,
    rowsep = 4pt,
    row{1} = {font=\bfseries},
}
رنگ & \lr{CMYK} & حس بصری \\
سرمه‌ای & \lr{(100, 80, 30, 30)} & رسمی، آکادمیک \\
سبز تیره & \lr{(90, 30, 80, 40)} & آرام، طبیعی \\
بنفش & \lr{(60, 80, 0, 20)} & خلاق، نوآور \\
قرمز تیره & \lr{(20, 100, 80, 20)} & قدرتمند، توجه‌برانگیز \\
نارنجی & \lr{(0, 60, 100, 0)} & گرم، دوستانه \\
\end{tblr}
\end{matintable}

\mnote{برای کتاب‌های دانشگاهی، رنگ‌های سرد (سرمه‌ای، سبز تیره) مناسب‌ترند. برای کتاب‌های آموزشی، رنگ‌های گرم (نارنجی، قرمز) مناسب‌ترند.}

\subsection{ساخت سایه‌ها}
\label{sec:custom-theme-shades}

برای هر رنگ اصلی، سه سایه بسازید:

\begin{itemize}
    \item \textbf{رنگ اصلی:} برای عنوان.
    \item \textbf{رنگ تیره:} برای متن روی
    پس‌زمینه‌ی روشن.
    \item \textbf{رنگ روشن:} برای پس‌زمینه.
\end{itemize}

\begin{latin}
\begin{minted}{latex}
% |\pc{رنگ اصلی}|%
\providecolor{matinblue}{cmyk}{0.80, 0.30, 0.00, 0.30}

% |\pc{رنگ تیره: افزایش}|\lr{K}|%
\providecolor{matindarkblue}{cmyk}{0.80, 0.30, 0.00, 0.60}

% |\pc{رنگ روشن: کاهش همه}|%
\providecolor{matinlightblue}{cmyk}{0.15, 0.05, 0.00, 0.05}
\end{minted}
\end{latin}

\mnote{برای ساخت سایه‌ی روشن، مقادیر \lr{CMYK} را به حدود \lr{۱۰-۲۰٪} کاهش دهید. برای سایه‌ی تیره، مقدار \lr{K} را افزایش دهید.}

\section{طراحی جلد}
\label{sec:custom-theme-cover}

\subsection{جلد جلو}
\label{sec:custom-theme-front}

جلد جلو در \file{mb-theme-cover.sty} با
\cmd{mb@frontart} تعریف شده است. برای
طراحی جدید:

\begin{latin}
\begin{minted}{latex}
\newcommand{\mb@frontart}{%
    \begin{tikzpicture}[remember picture,overlay]
        % |\pc{پس‌زمینه}|%
        \fill[coverprimary]
            (current page.north west) rectangle
            (current page.south east);

        % |\pc{عناصر تزئینی...}|%

        % |\pc{متن}|%
        \node[white, font=\fontsize{34}{43}\selectfont\bfseries]
            at (current page.center) {|\rl{عنوان کتاب}|};
    \end{tikzpicture}%
}
\end{minted}
\end{latin}

\mnote{در \lr{tikzpicture}، متن فارسی باید داخل \lr{\textbackslash rl\{...\}} باشد.}

\subsection{جلد عقب}
\label{sec:custom-theme-back}

جلد عقب در \file{mb-theme-cover.sty} با
\cmd{mb@backart} و \cmd{makebackcover} تعریف
شده است. این جلد از دو بلوک ایستاتیک تشکیل
شده:

\begin{latin}
\begin{minted}{latex}
\newcommand{\makebackcover}[1]{%
    \clearpage
    \thispagestyle{empty}%
    \mb@backart
    \vspace*{6.5cm}%
    \noindent
    \begin{minipage}[t][\dimexpr\paperheight-8.9cm\relax][s]{\textwidth}
        % |\pc{بلوک اول: درباره این کتاب}|%
        ...

        % |\pc{فاصله‌ی ۰.۷ سانتی‌متر}|%
        \vspace{0.7cm}

        % |\pc{بلوک دوم: درباره این نویسنده}|%
        ...
    \end{minipage}%
    \clearpage
}
\end{minted}
\end{latin}

\mnote{ارتفاع \lr{minipage} به‌صورت \lr{\textbackslash dimexpr\textbackslash paperheight-8.9cm\textbackslash relax} تعیین شده تا از سرریز به صفحه‌ی دوم جلوگیری شود.}

\section{مثال کامل: تم \lr{university}}
\label{sec:custom-theme-example}

در ادامه، یک تم کامل با رنگ سرمه‌ای و
جلد ساده ارائه می‌شود.

\subsection{فایل \file{mb-theme-colors.sty}}
\label{sec:custom-theme-example-colors}

\begin{latin}
\begin{minted}{latex}
% |\pc{خانواده‌ی سرمه‌ای}|%
\providecolor{matinblue}{cmyk}{1.00, 0.80, 0.30, 0.30}
\providecolor{matindarkblue}{cmyk}{1.00, 0.80, 0.30, 0.60}
\providecolor{matinlightblue}{cmyk}{0.15, 0.10, 0.00, 0.05}

% |\pc{خانواده‌ی طلایی}|%
\providecolor{matingreen}{cmyk}{0.00, 0.20, 0.60, 0.10}
\providecolor{matindarkgreen}{cmyk}{0.00, 0.20, 0.60, 0.50}
\providecolor{matinlightgreen}{cmyk}{0.00, 0.05, 0.15, 0.00}
\end{minted}
\end{latin}

\subsection{فایل \file{mb-theme-default.sty}}
\label{sec:custom-theme-example-loader}

\begin{latin}
\begin{minted}{latex}
\ProvidesPackage{mb-theme-default}[2026/10/08 v1.2 University Theme]
\RequirePackage{mb-theme-cover}
\PackageInfo{mb-theme-default}{University theme loaded}
\endinput
\end{minted}
\end{latin}

\mnote{فایل \lr{mb-theme-colors} اینجا بارگذاری نمی‌شود، چون در \lr{matinbook.cls} بارگذاری می‌شود.}

\section{خطاهای رایج}
\label{sec:custom-theme-mistakes}

\subsection{تغییر نام رنگ‌ها}
\label{sec:custom-theme-mistake-rename}

\textbf{مشکل:} اگر نام رنگ‌ها را تغییر
دهید (مثلاً \lr{matinblue} را به
\lr{mycolor} تغییر دهید)، ماژول‌های
مصرف‌کننده خطا می‌دهند.

\textbf{راه‌حل:} فقط مقادیر \lr{CMYK} را
تغییر دهید، نام‌ها را نگه دارید.

\subsection{فراموش کردن بارگذاری تم}
\label{sec:custom-theme-mistake-loading}

\textbf{مشکل:} اگر تم جدید را در
\file{matinbook.cls} بارگذاری نکنید، تم
پیش‌فرض استفاده می‌شود.

\textbf{راه‌حل:} مسیر تم جدید را در
\file{matinbook.cls} بنویسید.

\subsection{فراموش کردن \file{mb-theme-colors}}
\label{sec:custom-theme-mistake-colors}

\textbf{مشکل:} اگر \file{mb-theme-colors}
را در فاز ۱ بارگذاری نکنید، رنگ‌ها
تعریف نمی‌شوند.

\textbf{راه‌حل:} \file{mb-theme-colors}
باید در فاز ۱ \lr{matinbook.cls} باشد.

\subsection{استفاده‌ی مستقیم از رنگ‌های سفارشی}
\label{sec:custom-theme-mistake-direct}

\textbf{مشکل:} اگر در سند خود مستقیماً از
رنگ‌های \lr{RGB} استفاده کنید، با پالت
\lr{CMYK} تم هماهنگ نمی‌شود.

\textbf{راه‌حل:} از نام‌های استاندارد
\lr{matin*} استفاده کنید.

\section{تمرین‌ها}
\label{sec:custom-theme-exercises}

\begin{exercise}
    یک تم سفارشی با رنگ اصلی بنفش بسازید.
    \label{exc:custom-theme-purple}
\end{exercise}

\begin{exercise}
    جلد جلو را با یک عنصر جدید (مثلاً یک
    دایره‌ی طلایی) طراحی کنید.
    \label{exc:custom-theme-cover}
\end{exercise}

\begin{exercise}
    تم سفارشی خود را در \file{matinbook.cls}
    بارگذاری کنید و کتاب را کامپایل کنید.
    \label{exc:custom-theme-load}
\end{exercise}

\begin{exercise}
    یک تم مینیمال با سه رنگ خاکستری طراحی
    کنید.
    \label{exc:custom-theme-minimal}
\end{exercise}

\section{خلاصه}
\label{sec:custom-theme-summary}

در این فصل، با ساخت تم سفارشی در \lr{MatinBook}
آشنا شدیم:

\begin{itemize}
    \item ساخت تم سفارشی در پنج مرحله:
    ایجاد پوشه، کپی فایل‌ها، ویرایش رنگ‌ها،
    ویرایش جلد، بارگذاری.

    \item نام رنگ‌ها را تغییر ندهید؛ فقط
    مقادیر \lr{CMYK} را تغییر دهید.

    \item برای هر رنگ، سه سایه بسازید: اصلی،
    تیره، روشن.

    \item جلد جلو در \cmd{mb@frontart} و جلد
    عقب در \cmd{mb@backart} تعریف شده‌اند.

    \item در \lr{tikzpicture}، متن فارسی باید
    داخل \cmd{rl} باشد.

    \item جلد عقب از دو بلوک ایستاتیک با
    فاصله‌ی \pnum{۰.۷} سانتی‌متر تشکیل شده.

    \item خطاهای رایج شامل تغییر نام رنگ‌ها،
    فراموش کردن بارگذاری تم، فراموش کردن
    \file{mb-theme-colors}، و استفاده‌ی
    مستقیم از رنگ‌های سفارشی است.
\end{itemize}

در فصل بعدی (\cref{chap:custom-commands})،
با دستورات سفارشی آشنا می‌شوید و یاد
می‌گیرید که چگونه دستورات و محیط‌های جدید
تعریف کنید.
%%
%% part4-advanced/ch25-custom-commands.tex
%%
%% Chapter 25: Custom Commands — How to define new commands
%% and environments with \newcommand, \NewDocumentCommand,
%% and \NewDocumentEnvironment.
%%

\chapter{دستورات سفارشی}
\label{chap:custom-commands}

\section{چرا این موضوع مهم است؟}
\label{sec:custom-commands-why}

در نوشتن کتاب‌های فنی، بارها پیش می‌آید که
یک دستور یا محیط را چندین بار تکرار می‌کنید.
تعریف دستورات سفارشی، این تکرار را حذف
می‌کند و خوانایی سورس را افزایش می‌دهد.

\lr{MatinBook} خودش از دستورات سفارشی زیادی
استفاده می‌کند (مثل \cmd{cref}، \cmd{pnum}،
\cmd{mnote}). در این فصل، یاد می‌گیرید که
چگونه دستورات و محیط‌های سفارشی خود را
تعریف کنید.

\mnote{تعریف دستور سفارشی، یکی از مهارت‌های کلیدی در \lr{LaTeX} است که به شما امکان می‌دهد سورس کتاب را کوتاه‌تر و خواناتر کنید.}

\section{روش‌های تعریف دستور}
\label{sec:custom-commands-methods}

\lr{LaTeX} سه روش اصلی برای تعریف دستور
ارائه می‌دهد:

\begin{enumerate}
    \item \cmd{newcommand}: روش سنتی.
    \item \cmd{renewcommand}: بازنویسی دستور
    موجود.
    \item \cmd{NewDocumentCommand}: روش مدرن
    \lr{LaTeX3}.
\end{enumerate}

\mnote{\lr{LaTeX3} روش‌های مدرن‌تری برای تعریف دستور ارائه می‌دهد که از تعداد متغیر آرگومان‌ها و آرگومان‌های اختیاری پشتیبانی می‌کنند.}

\section{دستور ساده}
\label{sec:custom-commands-simple}

\subsection{با \cmd{newcommand}}
\label{sec:custom-commands-newcommand}

ساده‌ترین روش تعریف دستور، \cmd{newcommand}
است:

\begin{latin}
\begin{minted}{latex}
\newcommand{\mycommand}{|\pc{متن ثابت}|}
\end{minted}
\end{latin}

این دستور، یک دستور بدون آرگومان تعریف
می‌کند. برای استفاده:

\begin{latin}
\begin{minted}{latex}
\mycommand
\end{minted}
\end{latin}

\mnote{اگر دستور از قبل تعریف شده باشد، \lr{\textbackslash newcommand} خطا می‌دهد. برای اجتناب، از \lr{\textbackslash providecommand} استفاده کنید.}

\subsection{با آرگومان}
\label{sec:custom-commands-args}

برای تعریف دستور با آرگومان:

\begin{latin}
\begin{minted}{latex}
\newcommand{\myemph}[1]{\textbf{#1}}
\end{minted}
\end{latin}

استفاده:

\begin{latin}
\begin{minted}{latex}
\myemph{|\pc{متن پررنگ}|}
\end{minted}
\end{latin}

خروجی: \textbf{متن پررنگ}.

\mnote{تعداد آرگومان‌ها با عدد داخل \lr{[]} تعیین می‌شود. حداکثر ۹ آرگومان امکان‌پذیر است.}

\subsection{با آرگومان اختیاری}
\label{sec:custom-commands-optional}

برای تعریف دستور با آرگومان اختیاری:

\begin{latin}
\begin{minted}{latex}
\newcommand{\mybox}[2][matinblue]{%
    \colorbox{#1}{#2}%
}
\end{minted}
\end{latin}

استفاده:

\begin{latin}
\begin{minted}{latex}
\mybox{|\pc{متن پیش‌فرض}|}%
\mybox[matinred]{|\pc{متن قرمز}|}%
\end{minted}
\end{latin}

\mnote{آرگومان اختیاری با \lr{[...]} و مقدار پیش‌فرض داخل \lr{[]} تعریف می‌شود.}

\section{دستور مدرن (\lr{LaTeX3})}
\label{sec:custom-commands-latex3}

روش مدرن \lr{LaTeX3} از \cmd{NewDocumentCommand}
استفاده می‌کند:

\begin{latin}
\begin{minted}{latex}
\NewDocumentCommand{\myemph}{m}{%
    \textbf{#1}%
}
\end{minted}
\end{latin}

\mnote{روش \lr{LaTeX3} ایمن‌تر است و از انواع آرگومان (اجباری، اختیاری، ستاره‌دار) پشتیبانی می‌کند.}

\subsection{انواع آرگومان}
\label{sec:custom-commands-argtypes}

جدول \cref{tab:arg-types} انواع آرگومان در
\cmd{NewDocumentCommand} را نشان می‌دهد.

\begin{matintable}{انواع آرگومان در \lr{NewDocumentCommand}}{tab:arg-types}
\begin{tblr}{
    width = \textwidth,
    colspec = {c X[l] X[l]},
    hlines, vlines,
    colsep = 8pt,
    rowsep = 4pt,
    row{1} = {font=\bfseries},
}
نوع & توضیح & مثال \\
\lr{m} & آرگومان اجباری & \lr{\{متن\}} \\
\lr{o} & آرگومان اختیاری & \lr{[متن]} \\
\lr{O\{default\}} & اختیاری با پیش‌فرض & \lr{[متن]} \\
\lr{s} & ستاره‌ی اختیاری & \lr{*} \\
\lr{v} & آرگومان verbatim & \lr{|متن|} \\
\end{tblr}
\end{matintable}

\mnote{برای دستورات ساده، \lr{\textbackslash newcommand} کافی است. برای دستورات پیچیده، از \lr{\textbackslash NewDocumentCommand} استفاده کنید.}

\subsection{مثال: دستور با ستاره}
\label{sec:custom-commands-star}

\begin{latin}
\begin{minted}{latex}
\NewDocumentCommand{\myemph}{s m}{%
    \IfBooleanTF{#1}%
        {\textbf{#2}}%
        {\textit{#2}}%
}
\end{minted}
\end{latin}

استفاده:

\begin{latin}
\begin{minted}{latex}
\myemph{|\pc{متن مورب}|}%
\myemph*{|\pc{متن پررنگ}|}%
\end{minted}
\end{latin}

\mnote{این دستور با \lr{\textbackslash myemph} متن را مورب و با \lr{\textbackslash myemph*} متن را پررنگ می‌کند.}

\section{تعریف محیط سفارشی}
\label{sec:custom-commands-env}

\subsection{با \cmd{newenvironment}}
\label{sec:custom-commands-newenv}

برای تعریف محیط سفارشی:

\begin{latin}
\begin{minted}{latex}
\newenvironment{myenv}{%
    \begin{quote}\itshape
}{%
    \end{quote}
}
\end{minted}
\end{latin}

استفاده:

\begin{latin}
\begin{minted}{latex}
\begin{myenv}
|\pc{متن داخل محیط}|%
\end{myenv}
\end{minted}
\end{latin}

\mnote{محیط سفارشی با \lr{\textbackslash newenvironment} تعریف می‌شود و دو بخش دارد: کد شروع و کد پایان.}

\subsection{با \cmd{NewDocumentEnvironment}}
\label{sec:custom-commands-newdocenv}

روش مدرن:

\begin{latin}
\begin{minted}{latex}
\NewDocumentEnvironment{myenv}{m}{%
    \begin{quote}\itshape
    \textbf{#1}\par
}{%
    \end{quote}
}
\end{minted}
\end{latin}

استفاده:

\begin{latin}
\begin{minted}{latex}
\begin{myenv}{|\pc{عنوان}|}
|\pc{متن داخل محیط}|%
\end{myenv}
\end{minted}
\end{latin}

\mnote{روش \lr{LaTeX3} از آرگومان‌های محیط پشتیبانی می‌کند.}

\section{دستورات مفید \lr{MatinBook}}
\label{sec:custom-commands-matinbook}

\lr{MatinBook} چند دستور کمکی برای ساده‌سازی
نوشتن متن ارائه می‌دهد:

\begin{matintable}{دستورات کمکی \lr{MatinBook}}{tab:helper-commands}
\begin{tblr}{
    width = \textwidth,
    colspec = {l X[l] X[l]},
    hlines, vlines,
    colsep = 6pt,
    rowsep = 4pt,
    row{1} = {font=\bfseries},
}
دستور & کاربرد & معادل \\
\lr{\textbackslash cmd\{...\}} & نام دستور & \lr{\textbackslash texttt\{\textbackslash ...\}} \\
\lr{\textbackslash env\{...\}} & نام محیط & \lr{\textbackslash texttt\{...\}} \\
\lr{\textbackslash pkg\{...\}} & نام بسته & \lr{\textbackslash texttt\{...\}} \\
\lr{\textbackslash file\{...\}} & نام فایل & \lr{\textbackslash texttt\{...\}} \\
\lr{\textbackslash opt\{...\}} & آپشن & \lr{\textbackslash texttt\{...\}} \\
\lr{\textbackslash pnum\{...\}} & عدد فارسی & \lr{\textbackslash lr\{...\}} \\
\lr{\textbackslash pc\{...\}} & فارسی در کد & \lr{\textbackslash rl\{...\}} \\
\lr{\textbackslash pcc\{...\}} & کامنت فارسی & — \\
\end{tblr}
\end{matintable}

\mnote{دستورات \lr{\textbackslash cmd}، \lr{\textbackslash env}، \lr{\textbackslash pkg}، \lr{\textbackslash file} و \lr{\textbackslash opt} در \lr{main.tex} مستندات تعریف شده‌اند (نه در کلاس).}

\section{مثال‌های کاربردی}
\label{sec:custom-commands-examples}

\subsection{دستور برای نمایش واژه‌های کلیدی}
\label{sec:custom-commands-keyword}

\begin{latin}
\begin{minted}{latex}
\newcommand{\keyword}[1]{\textbf{\textcolor{matinblue}{#1}}}
\end{minted}
\end{latin}

استفاده:

\begin{latin}
\begin{minted}{latex}
\keyword{|\pc{الگوریتم}|}|\pc{یک مفهوم اساسی است.}|%
\end{minted}
\end{latin}

خروجی: \textbf{\textcolor{matinblue}{الگوریتم}}
یک مفهوم اساسی است.

\subsection{دستور برای نمایش معادل لاتین}
\label{sec:custom-commands-latin-equiv}

\begin{latin}
\begin{minted}{latex}
\newcommand{\latinterm}[2]{#1\lr{(#2)}}
\end{minted}
\end{latin}

استفاده:

\begin{latin}
\begin{minted}{latex}
\latinterm{|\pc{الگوریتم}|}{algorithm}
\end{minted}
\end{latin}

خروجی: الگوریتم\lr{(algorithm)}.

\mnote{این دستور برای نمایش واژه‌ی فارسی به‌همراه معادل لاتین آن در پرانتز مفید است.}

\subsection{محیط برای نکات مهم}
\label{sec:custom-commands-important}

\begin{latin}
\begin{minted}{latex}
\NewDocumentEnvironment{important}{m}{%
    \begin{tcolorbox}[
        colback=matinred!10,
        colframe=matinred,
        title=|\rl{#1}|%
    ]
}{%
    \end{tcolorbox}
}
\end{minted}
\end{latin}

\mnote{برای محیط‌های سفارشی پیچیده، از \lr{tcolorbox} استفاده کنید که امکانات بیشتری دارد.}

\section{خطاهای رایج}
\label{sec:custom-commands-mistakes}

\subsection{تعریف دستور موجود}
\label{sec:custom-commands-mistake-exists}

\textbf{مشکل:} اگر دستوری را که از قبل
وجود دارد با \cmd{newcommand} تعریف کنید،
خطا می‌دهد.

\textbf{راه‌حل:} از \cmd{renewcommand} برای
بازنویسی یا \cmd{providecommand} برای
تعریف امن استفاده کنید.

\subsection{تعداد اشتباه آرگومان‌ها}
\label{sec:custom-commands-mistake-args}

\textbf{مشکل:} اگر تعداد آرگومان‌های تعریف‌شده
با تعداد آرگومان‌های استفاده‌شده یکسان
نباشد، خطا می‌دهد.

\textbf{راه‌حل:} تعداد آرگومان‌ها را در
تعریف و استفاده بررسی کنید.

\subsection{فراموش کردن \cmd{makeatletter}}
\label{sec:custom-commands-mistake-makeatletter}

\textbf{مشکل:} اگر دستوری با \lr{@} در نام
تعریف کنید (مثل \cmd{my@command})،
\lr{LaTeX} خطا می‌دهد.

\textbf{راه‌حل:} از \cmd{makeatletter} و
\cmd{makeatother} استفاده کنید، یا از
\lr{@} در نام دستور اجتناب کنید.

\mnote{در فایل‌های \lr{.sty}، استفاده از \lr{@} در نام دستورات داخلی مرسوم است، ولی در سند اصلی باید از \lr{\textbackslash makeatletter} استفاده کنید.}

\section{تمرین‌ها}
\label{sec:custom-commands-exercises}

\begin{exercise}
    یک دستور ساده تعریف کنید که متن را
    پررنگ کند.
    \label{exc:custom-commands-bold}
\end{exercise}

\begin{exercise}
    یک دستور با دو آرگومان تعریف کنید
    که یک واژه و معادل لاتین آن را نمایش
    دهد.
    \label{exc:custom-commands-two-args}
\end{exercise}

\begin{exercise}
    یک دستور با آرگومان اختیاری تعریف کنید
    که رنگ متن را تغییر دهد.
    \label{exc:custom-commands-optional}
\end{exercise}

\begin{exercise}
    یک محیط سفارشی تعریف کنید که متن را
    در یک جعبه‌ی رنگی قرار دهد.
    \label{exc:custom-commands-env}
\end{exercise}

\section{خلاصه}
\label{sec:custom-commands-summary}

در این فصل، با دستورات سفارشی در \lr{MatinBook}
آشنا شدیم:

\begin{itemize}
    \item سه روش تعریف دستور:
    \cmd{newcommand}، \cmd{renewcommand}،
    \cmd{NewDocumentCommand}.

    \item دستور ساده بدون آرگومان:
    \cmd{newcommand\{\textbackslash mycmd\}\{...\}}.

    \item دستور با آرگومان:
    \cmd{newcommand\{\textbackslash mycmd\}[1]\{...\}}.

    \item دستور با آرگومان اختیاری:
    \cmd{newcommand\{\textbackslash mycmd\}[2][default]\{...\}}.

    \item روش مدرن \lr{LaTeX3} از انواع آرگومان
    (اجباری، اختیاری، ستاره‌دار) پشتیبانی
    می‌کند.

    \item محیط سفارشی با \cmd{newenvironment}
    یا \cmd{NewDocumentEnvironment}.

    \item دستورات کمکی \lr{MatinBook}:
    \cmd{cmd}، \cmd{env}، \cmd{pkg}،
    \cmd{file}، \cmd{opt}، \cmd{pnum}،
    \cmd{pc}، \cmd{pcc}.

    \item خطاهای رایج شامل تعریف دستور موجود،
    تعداد اشتباه آرگومان‌ها، و فراموش کردن
    \cmd{makeatletter} است.
\end{itemize}

در فصل بعدی (\cref{chap:large-projects})،
با مدیریت پروژه‌های بزرگ آشنا می‌شوید و
یاد می‌گیرید که چگونه کتاب‌های چندفایلی
را سازمان‌دهی کنید.
%%
%% part4-advanced/ch26-large-projects.tex
%%
%% Chapter 26: Large Projects — How to organize books with
%% multiple files, parts, and contributors.
%%

\chapter{پروژه‌های بزرگ}
\label{chap:large-projects}

\section{چرا این موضوع مهم است؟}
\label{sec:large-projects-why}

وقتی کتاب شما به بیش از ۱۰۰ صفحه می‌رسد،
نوشتن همه‌ی محتوا در یک فایل \file{.tex}
غیرعملی می‌شود. کامپایل کند می‌شود، پیدا
کردن خطاها سخت می‌شود، و همکاری با دیگران
دشوار می‌شود.

\lr{MatinBook} به‌طور خاص برای پروژه‌های بزرگ
طراحی شده است. در این فصل، یاد می‌گیرید که
چگونه یک کتاب چندفایلی را سازمان‌دهی کنید.

\mnote{کتاب‌های حرفه‌ای معمولاً بیش از ۲۰۰ صفحه هستند و به‌طور طبیعی به چندین فایل تقسیم می‌شوند.}

\section{ساختار پیشنهادی}
\label{sec:large-projects-structure}

\subsection{ساختار ساده}
\label{sec:large-projects-simple}

برای کتاب‌های کوچک تا متوسط (کمتر از ۱۰
فصل):

\begin{latin}
\begin{minted}{text}
my-book/
├── main.tex
├── references.bib
├── chapters/
│   ├── chapter-01.tex
│   ├── chapter-02.tex
│   ├── chapter-03.tex
│   └── ...
├── figures/
│   ├── figure-01.pdf
│   └── ...
└── code/
    ├── example-01.py
    └── ...
\end{minted}
\end{latin}

\mnote{هر فصل در یک فایل جداگانه قرار می‌گیرد. این کار، کامپایل را سریع‌تر و مدیریت را آسان‌تر می‌کند.}

\subsection{ساختار پیشرفته}
\label{sec:large-projects-advanced}

برای کتاب‌های بزرگ (بیش از ۱۰ فصل یا
چند بخش):

\begin{latin}
\begin{minted}{text}
my-book/
├── main.tex
├── references.bib
├── frontmatter/
│   ├── cover.tex
│   ├── colophon.tex
│   ├── preface.tex
│   └── ...
├── part1-basics/
│   ├── chapter-01.tex
│   ├── chapter-02.tex
│   └── ...
├── part2-advanced/
│   ├── chapter-05.tex
│   ├── chapter-06.tex
│   └── ...
├── appendices/
│   ├── app-a.tex
│   └── ...
├── backmatter/
│   ├── bibliography.tex
│   └── index.tex
├── figures/
│   ├── chapter-01/
│   ├── chapter-02/
│   └── ...
└── code/
    ├── chapter-01/
    └── ...
\end{minted}
\end{latin}

\mnote{در ساختار پیشرفته، هر بخش از کتاب یک پوشه‌ی جداگانه دارد. این کار، مدیریت پروژه‌های بزرگ را آسان‌تر می‌کند.}

\section{فایل اصلی}
\label{sec:large-projects-main}

فایل \file{main.tex}، نقطه‌ی ورود پروژه
است. این فایل شامل:

\begin{itemize}
    \item \cmd{documentclass}.
    \item فراداده (\cmd{title}، \cmd{author}،
    \cmd{date}).
    \item \cmd{addbibresource}.
    \item \cmd{begin\{document\}}.
    \item جلد، پیش‌متن، متن اصلی، پس‌متن.
    \item \cmd{input} برای هر فصل.
    \item \cmd{end\{document\}}.
\end{itemize}

\begin{latin}
\begin{minted}{latex}
\documentclass{matinbook}

\addbibresource{references.bib}

\title{|\pc{عنوان کتاب}|}
\author{|\pc{نام نویسنده}|}
\date{|\pc{مهر ۱۴۰۵}|}

\booktitle{|\pc{عنوان کتاب}|}

\begin{document}

% |\pc{جلد جلو}|%
\makecover{...}{...}{...}{...}

% |\pc{پیش‌متن}|%
\frontmatter
\frontmattergeometry

%%
%% frontmatter/cover.tex
%%
%% The front cover is generated directly by \makecover in main.tex.
%% This file is intentionally empty and kept as a placeholder so
%% that the directory structure stays uniform.
%%
%% If you want to customize the cover, edit the \makecover call
%% in main.tex, or override the cover art in
%% tex/themes/default/mb-theme-cover.sty.
%%
%%
%% frontmatter/preface.tex
%%
%% Preface: Why MatinBook, who is this book for, what is
%% its philosophy.
%%

\chapter*{پیشگفتار}
\addcontentsline{toc}{chapter}{پیشگفتار}
\label{ch:preface}

\section*{چرا \lr{MatinBook}؟}

حروف‌چینی کتاب‌های فنی فارسی، سال‌هاست با چالش‌های
متعددی روبه‌روست. ابزارهای موجود، یا برای زبان فارسی
بهینه نشده‌اند، یا از نظر بصری کیفیت لازم را ندارند،
یا برای پروژه‌های بزرگ مقیاس‌پذیر نیستند. این وضعیت،
باعث شده که بسیاری از کتاب‌های فنی فارسی، با وجود
محتوای ارزشمند، از نظر ظاهری فاصله‌ی زیادی با
استانداردهای نشر جهانی داشته باشند.

\lr{MatinBook} با هدف پر کردن این شکاف طراحی شده است.
این کلاس \lr{LaTeX}، تمام ابزارهای لازم برای تولید
یک کتاب فنی حرفه‌ای فارسی را در یک بسته‌ی منسجم فراهم
می‌کند: از حروف‌چینی راست‌به‌چپ و پشتیبانی کامل از
فونت‌های فارسی، تا محیط‌های علمی رنگ‌آمیزی‌شده،
نمودارهای پیچیده، و مدیریت حرفه‌ای منابع و نمایه.

\section*{این کتاب برای چه کسانی است؟}

این راهنما برای چهار گروه اصلی نوشته شده است:

\begin{itemize}
    \item \textbf{نویسندگان:} کسانی که می‌خواهند کتاب
    فنی یا علمی فارسی بنویسند و به دنبال ابزاری حرفه‌ای
    برای حروف‌چینی هستند.

    \item \textbf{پژوهشگران:} کسانی که می‌خواهند مقالات،
    پایان‌نامه‌ها یا گزارش‌های فنی خود را با کیفیت
    بالا منتشر کنند.

    \item \textbf{دانشجویان:} کسانی که می‌خواهند با
    ابزارهای حرفه‌ای حروف‌چینی آشنا شوند و مهارتی
    ارزشمند برای آینده‌ی حرفه‌ای خود کسب کنند.

    \item \textbf{مدرسان:} کسانی که می‌خواهند جزوه‌ها،
    کتاب‌های درسی یا مطالب آموزشی خود را با استاندارد
    بالایی آماده کنند.
\end{itemize}

\section*{فلسفه‌ی طراحی}

\lr{MatinBook} بر سه اصل بنیادین استوار است:

\subsection*{۱. وضوح بر تزئین}

هر عنصر بصری در \lr{MatinBook}، نقشی معنایی دارد.
رنگ‌ها، جعبه‌ها، و نمودارها، نه برای زیبایی صرف،
بلکه برای کمک به درک بهتر محتوا طراحی شده‌اند.
این اصل، از فلسفه‌ی نشر \lr{MIT Press} و
\lr{Springer} الهام گرفته است.

\subsection*{۲. سازگاری بر نوآوری}

هر عنصر در \lr{MatinBook}، یک شکل استاندارد دارد
و در تمام کتاب، یکسان ظاهر می‌شود. این سازگاری،
باعث می‌شود خواننده بدون حواس‌پرتی، روی محتوا
تمرکز کند.

\subsection*{۳. خویشتن‌داری بر افزونگی}

فضای خالی، یک عنصر طراحی است. \lr{MatinBook}
از پر کردن هر اینچ صفحه پرهیز می‌کند و به متن
اجازه می‌دهد نفس بکشد. این اصل، خوانایی را
افزایش می‌دهد و خستگی چشم را کاهش می‌دهد.

\section*{ساختار این راهنما}

این راهنما در شش بخش اصلی و پنج پیوست سازماندهی
شده است:

\begin{enumerate}
    \item \textbf{شروع به کار:} معرفی، نصب، اولین
    کتاب، ساختار کلی.

    \item \textbf{اجزای کتاب:} جلد، فصل، متن، محیط‌های
    علمی، ریاضیات، کد، نمودار، جدول، یادداشت حاشیه،
    منابع، نمایه، ارجاع.

    \item \textbf{شخصی‌سازی:} آپشن‌ها، رنگ‌ها، فونت‌ها،
    صفحه‌آرایی، تم‌ها، زبان.

    \item \textbf{مباحث پیشرفته:} ماژول‌ها، تم سفارشی،
    دستورات سفارشی، پروژه‌های بزرگ.

    \item \textbf{مثال‌های کامل:} سه فصل نمونه از
    کتاب‌های ریاضی، برنامه‌نویسی و فیزیک.

    \item \textbf{مرجع دستورات:} مرجع کامل کلاس،
    محیط‌ها و دستورات.
\end{enumerate}

هر فصل با یک مثال ساده شروع می‌شود، به تدریج پیچیده‌تر
می‌شود، و با تمرین‌های عملی به پایان می‌رسد. خواننده
می‌تواند فصل‌ها را به ترتیب بخواند، یا هر فصل را
به‌عنوان یک مرجع مستقل استفاده کند.

\section*{سپاسگزاری}

\lr{MatinBook} بدون پشتیبانی جامعه‌ی \lr{LaTeX} فارسی
و تلاش‌های بی‌دریغ توسعه‌دهندگان بسته‌های متن‌باز،
امکان‌پذیر نبود. از \lr{Vafa Khalighi} برای
\lr{XePersian}، از \lr{Hassan Mesgarha} برای
\lr{XePersian-HM}، از \lr{Thomas F. Sturm} برای
\lr{tcolorbox}، از \lr{Geoffrey Poore} برای
\lr{minted}، و از تمام کسانی که در این مسیر
کمک کرده‌اند، صمیمانه سپاسگزارم.

\vspace{1cm}

\begin{flushleft}
    \textbf{متین محمودی} \\
    مهر ۱۴۰۵
\end{flushleft}

\clearpage

% |\pc{متن اصلی}|%
\mainmatter
\mainmattergeometry

\part{|\pc{مبانی}|}
\input{part1-basics/chapter-01}
\input{part1-basics/chapter-02}

\part{|\pc{مباحث پیشرفته}|}
\input{part2-advanced/chapter-05}
\input{part2-advanced/chapter-06}

% |\pc{پس‌متن}|%
\backmatter
\frontmattergeometry

\printbibliography[title={|\pc{منابع و مراجع}|}]
\printindex

% |\pc{جلد عقب}|%
\makebackcover{...}

\end{document}
\end{minted}
\end{latin}

\mnote{فایل \lr{main.tex} باید کوتاه بماند. اگر بیش از ۱۰۰ خط شد، احتمالاً باید بخشی از محتوا را به فایل‌های جداگانه منتقل کنید.}

\section{فایل‌های فصل}
\label{sec:large-projects-chapters}

هر فایل فصل، فقط شامل محتوای همان فصل
است. \textbf{بدون} \cmd{documentclass} و
\textbf{بدون} \cmd{begin\{document\}}:

\begin{latin}
\begin{minted}{latex}
% |\pc{در فایل}|\lr{chapters/chapter-01.tex}:

\chapter{|\pc{مقدمه}|}
\label{chap:introduction}

\section{|\pc{انگیزه}|}
\label{sec:motivation}

|\pc{متن فصل...}|%

\begin{theorem}[|\pc{قضیه‌ی نمونه}|]
    |\pc{متن قضیه...}|%
    \label{thm:sample}
\end{theorem}

\begin{proof}
    |\pc{اثبات...}|%
\end{proof}
\end{minted}
\end{latin}

\mnote{فایل‌های فصل نباید \lr{\textbackslash documentclass} داشته باشند. آن‌ها فقط محتوا هستند و با \lr{\textbackslash input} وارد می‌شوند.}

\section{\cmd{input} در برابر \cmd{include}}
\label{sec:large-projects-input-include}

\lr{LaTeX} دو دستور برای وارد کردن فایل
ارائه می‌دهد:

\begin{itemize}
    \item \cmd{input\{file\}}: فایل را در
    همان‌جا وارد می‌کند.
    \item \cmd{include\{file\}}: فایل را
    در یک صفحه‌ی جدید وارد می‌کند.
\end{itemize}

\subsection{تفاوت‌ها}
\label{sec:large-projects-differences}

جدول \cref{tab:input-vs-include} تفاوت‌های
\cmd{input} و \cmd{include} را نشان می‌دهد.

\begin{matintable}{تفاوت \lr{input} و \lr{include}}{tab:input-vs-include}
\begin{tblr}{
    width = \textwidth,
    colspec = {l X[c] X[c]},
    hlines, vlines,
    colsep = 8pt,
    rowsep = 4pt,
    row{1} = {font=\bfseries},
}
ویژگی & \lr{\textbackslash input} & \lr{\textbackslash include} \\
صفحه‌ی جدید & خیر & بله \\
فایل \lr{.aux} جدا & خیر & بله \\
سرعت & سریع & کندتر \\
مناسب برای & فصل‌های کوتاه & فصل‌های بلند \\
\end{tblr}
\end{matintable}

\mnote{برای کتاب‌ها، معمولاً از \lr{\textbackslash input} استفاده می‌شود. \lr{\textbackslash include} فقط برای پروژه‌های بسیار بزرگ مفید است.}

\subsection{مثال}
\label{sec:large-projects-example}

\begin{latin}
\begin{minted}{latex}
% |\pc{با}|\lr{\textbackslash input}|%
\input{chapters/chapter-01}
\input{chapters/chapter-02}

% |\pc{با}|\lr{\textbackslash include}|%
\include{chapters/chapter-01}
\include{chapters/chapter-02}
\end{minted}
\end{latin}

\section{مدیریت شکل‌ها}
\label{sec:large-projects-figures}

برای پروژه‌های بزرگ، بهتر است شکل‌ها را در
پوشه‌های جداگانه سازمان‌دهی کنید:

\begin{latin}
\begin{minted}{text}
figures/
├── chapter-01/
│   ├── figure-01.pdf
│   ├── figure-02.pdf
│   └── ...
├── chapter-02/
│   └── ...
└── cover/
    └── cover-image.pdf
\end{minted}
\end{latin}

سپس در سند:

\begin{latin}
\begin{minted}{latex}
\includegraphics{figures/chapter-01/figure-01.pdf}
\end{minted}
\end{latin}

\mnote{برای کوتاه‌تر کردن مسیرها، می‌توانید از \lr{\textbackslash graphicspath} استفاده کنید که در \lr{MatinBook} از پیش تنظیم شده است.}

\section{مدیریت کد}
\label{sec:large-projects-code}

برای کدهای طولانی، بهتر است آن‌ها را در
فایل‌های جداگانه ذخیره کنید و با
\cmd{inputminted} وارد کنید:

\begin{latin}
\begin{minted}{text}
code/
├── chapter-01/
│   ├── example-01.py
│   └── example-02.py
└── chapter-02/
    └── ...
\end{minted}
\end{latin}

سپس در سند:

\begin{latin}
\begin{minted}{latex}
\begin{latin}
\inputminted{python}{code/chapter-01/example-01.py}
\end{latin}
\codecaption{|\pc{مثال اول}|}
\label{code:example-01}
\end{minted}
\end{latin}

\mnote{دستور \lr{\textbackslash inputminted} کد را از یک فایل خارجی وارد می‌کند. این کار، سورس کتاب را کوتاه‌تر می‌کند.}

\section{مدیریت کتاب‌نامه}
\label{sec:large-projects-bibliography}

برای کتاب‌های بزرگ، بهتر است کتاب‌نامه را
در چند فایل تقسیم کنید:

\begin{latin}
\begin{minted}{text}
bibliography/
├── books.bib
├── articles.bib
├── online.bib
└── persian.bib
\end{minted}
\end{latin}

سپس در \file{main.tex}:

\begin{latin}
\begin{minted}{latex}
\addbibresource{bibliography/books.bib}
\addbibresource{bibliography/articles.bib}
\addbibresource{bibliography/online.bib}
\addbibresource{bibliography/persian.bib}
\end{minted}
\end{latin}

\mnote{می‌توانید چند فایل \lr{.bib} را با \lr{\textbackslash addbibresource} بارگذاری کنید. \lr{biber} همه را با هم پردازش می‌کند.}

\section{همکاری با دیگران}
\label{sec:large-projects-collaboration}

اگر با دیگران روی یک کتاب کار می‌کنید،
این نکات را رعایت کنید:

\begin{enumerate}
    \item \textbf{استفاده از \lr{Git}:} نسخه‌بندی
    و همکاری.
    \item \textbf{تقسیم فایل‌ها:} هر نفر یک
    فصل.
    \item \textbf{قواعد نگارشی:} یکدست‌سازی
    سبک نوشتار.
    \item \textbf{برچسب‌های یکتا:} جلوگیری از
    تداخل برچسب‌ها.
    \item \textbf{کامپایل مکرر:} بررسی
    خطاها پس از هر تغییر.
\end{enumerate}

\mnote{در همکاری گروهی، \lr{Git} ابزار ضروری است. هر نفر روی برنچ خود کار می‌کند و سپس تغییرات را ادغام می‌کند.}

\section{کامپایل پروژه‌های بزرگ}
\label{sec:large-projects-compile}

\subsection{کامپایل کامل}
\label{sec:large-projects-compile-full}

\begin{latin}
\begin{minted}{bash}
# |\pc{مرحله‌ی اول: کامپایل اول}|%
xelatex -shell-escape main.tex

# |\pc{مرحله‌ی دوم: کتاب‌نامه}|%
biber main

# |\pc{مرحله‌ی سوم: کامپایل دوم}|%
xelatex -shell-escape main.tex

# |\pc{مرحله‌ی چهارم: کامپایل سوم}|%
xelatex -shell-escape main.tex
\end{minted}
\end{latin}

\mnote{برای پروژه‌های بزرگ، کامپایل کامل ممکن است چند دقیقه طول بکشد. از \lr{latexmk} برای خودکارسازی استفاده کنید.}

\subsection{استفاده از \lr{latexmk}}
\label{sec:large-projects-compile-latexmk}

\begin{latin}
\begin{minted}{bash}
latexmk -xelatex -shell-escape main.tex
\end{minted}
\end{latin}

\mnote{\lr{latexmk} به‌طور خودکار تعداد کامپایل‌های لازم را تشخیص می‌دهد و کل چرخه را اجرا می‌کند.}

\subsection{کامپایل فصل به فصل}
\label{sec:large-projects-compile-chapter}

برای صرفه‌جویی در زمان، می‌توانید فقط
یک فصل را کامپایل کنید:

\begin{latin}
\begin{minted}{bash}
# |\pc{در}|\lr{main.tex}|%
% |\pc{فقط فصل اول را کامپایل کن}|%
\includeonly{chapters/chapter-01}
\end{minted}
\end{latin}

\mnote{دستور \lr{\textbackslash includeonly} فقط با \lr{\textbackslash include} کار می‌کند، نه \lr{\textbackslash input}.}

\section{خطاهای رایج}
\label{sec:large-projects-mistakes}

\subsection{فراموش کردن \cmd{input}}
\label{sec:large-projects-mistake-input}

\textbf{مشکل:} اگر \cmd{input} را برای یک
فصل فراموش کنید، آن فصل در کتاب ظاهر
نمی‌شود.

\textbf{راه‌حل:} فهرست \cmd{input}ها را
با فهرست فصل‌ها مقایسه کنید.

\subsection{برچسب تکراری}
\label{sec:large-projects-mistake-duplicate}

\textbf{مشکل:} اگر در دو فصل برچسب یکسان
استفاده کنید، خطای \lr{Multiply defined labels}
می‌دهید.

\textbf{راه‌حل:} از پیشوندهای یکتا برای
هر فصل استفاده کنید (مثلاً \lr{ch01:}،
\lr{ch02:}).

\subsection{مسیر اشتباه}
\label{sec:large-projects-mistake-path}

\textbf{مشکل:} اگر مسیر فایل‌ها اشتباه
باشد، خطای \lr{File not found} می‌دهید.

\textbf{راه‌حل:} مسیرها را نسبت به
\file{main.tex} بررسی کنید.

\subsection{کامپایل ناقص}
\label{sec:large-projects-mistake-incomplete}

\textbf{مشکل:} اگر فقط یک بار کامپایل
کنید، ارجاعات و کتاب‌نامه کامل نمی‌شوند.

\textbf{راه‌حل:} چرخه‌ی کامل کامپایل را
اجرا کنید (حداقل ۲ بار \lr{xelatex} +
\lr{biber}).

\section{تمرین‌ها}
\label{sec:large-projects-exercises}

\begin{exercise}
    یک کتاب با ساختار پیشرفته (پوشه‌بندی
    شده) بسازید و آن را کامپایل کنید.
    \label{exc:large-projects-structure}
\end{exercise}

\begin{exercise}
    یک کتاب با دو بخش و پنج فصل بسازید.
    از \cmd{input} برای وارد کردن فصل‌ها
    استفاده کنید.
    \label{exc:large-projects-parts}
\end{exercise}

\begin{exercise}
    شکل‌ها را در پوشه‌های جداگانه سازمان‌دهی
    کنید و از \cmd{includegraphics} با
    مسیر کامل استفاده کنید.
    \label{exc:large-projects-figures}
\end{exercise}

\begin{exercise}
    یک کتاب با \cmd{includeonly} کامپایل
    کنید تا فقط یک فصل را ببینید.
    \label{exc:large-projects-includeonly}
\end{exercise}

\section{خلاصه}
\label{sec:large-projects-summary}

در این فصل، با مدیریت پروژه‌های بزرگ در
\lr{MatinBook} آشنا شدیم:

\begin{itemize}
    \item برای کتاب‌های بزرگ، محتوا را به
    فایل‌های جداگانه تقسیم کنید.

    \item ساختار ساده: \file{main.tex} +
    \file{chapters/} + \file{figures/} +
    \file{code/}.

    \item ساختار پیشرفته: تقسیم به بخش‌ها
    (\file{part1-*}، \file{part2-*} و...).

    \item \file{main.tex} باید کوتاه باشد
    (کمتر از ۱۰۰ خط).

    \item فایل‌های فصل نباید
    \cmd{documentclass} داشته باشند.

    \item \cmd{input} فایل را در همان‌جا
    وارد می‌کند؛ \cmd{include} صفحه‌ی
    جدید ایجاد می‌کند.

    \item برای کدهای طولانی، از
    \cmd{inputminted} استفاده کنید.

    \item کتاب‌نامه را می‌توان در چند فایل
    \lr{.bib} تقسیم کرد.

    \item در همکاری گروهی، از \lr{Git}
    و قواعد یکدست استفاده کنید.

    \item خطاهای رایج شامل فراموش کردن
    \cmd{input}، برچسب تکراری، مسیر
    اشتباه، و کامپایل ناقص است.
\end{itemize}

این فصل، آخرین فصل از بخش چهارم (مباحث
پیشرفته) بود. در بخش پنجم (مثال‌های کامل)،
سه فصل کامل از کتاب‌های ریاضی، برنامه‌نویسی
و فیزیک ارائه می‌شود.

%========================================================================
% PART 5 — EXAMPLES
%========================================================================
\part{مثال‌های کامل}
\label{part:examples}

%%
%% part5-examples/ch27-math-book.tex
%%
%% Chapter 27: A Math Book — A complete sample chapter from
%% a mathematics textbook, demonstrating all MatinBook
%% features in a real-world context.
%%

\chapter{کتاب ریاضی}
\label{chap:math-book}

\section{چرا این موضوع مهم است؟}
\label{sec:math-book-why}

در فصل‌های قبل، با تمام قابلیت‌های \lr{MatinBook}
آشنا شدید. ولی دانستن هر قابلیت به‌تنهایی کافی
نیست؛ مهم این است که بتوانید همه‌ی آن‌ها را
در یک کتاب واقعی ترکیب کنید.

در این فصل، یک فصل کامل از یک کتاب ریاضی
خیالی را ارائه می‌دهیم که همه‌ی اجزای
\lr{MatinBook} را در بافت واقعی نشان می‌دهد.
هدف، الگوبرداری برای نوشتن کتاب ریاضی خودتان
است.

\mnote{این فصل، یک «مطالعه‌ی موردی» است. به‌جای توضیح قابلیت‌ها، آن‌ها را در عمل می‌بینید.}

\section{ساختار فصل نمونه}
\label{sec:math-book-structure}

فصل نمونه، از یک کتاب ریاضی خیالی با عنوان
«مقدمه‌ای بر آنالیز حقیقی» است. این فصل
شامل بخش‌های زیر است:

\begin{itemize}
    \item مقدمه.
    \item تعاریف.
    \item قضایا و اثبات‌ها.
    \item مثال‌ها.
    \item تمرین‌ها.
    \item راه‌حل‌ها.
    \item نمودارها.
\end{itemize}

\mnote{این فصل را می‌توانید به‌عنوان الگو برای نوشتن فصل‌های کتاب ریاضی خودتان استفاده کنید.}

\section{مقدمه}
\label{sec:math-book-intro}

آنالیز حقیقی\index{آنالیز حقیقی}، شاخه‌ای از
ریاضیات است که به مطالعه‌ی اعداد حقیقی،
دنباله‌ها، سری‌ها، حد، پیوستگی و مشتق
می‌پردازد. این شاخه، پایه‌ی بسیاری از
شاخه‌های دیگر ریاضیات است.

در این فصل، با مفاهیم اساسی آنالیز حقیقی
آشنا می‌شویم و قضایای مهم آن را بررسی
می‌کنیم.

\mnote{آنالیز حقیقی، پیش‌نیاز آنالیز مختلط، آنالیز تابعی و نظریه‌ی اندازه است.}

\section{دنباله‌ها و همگرایی}
\label{sec:math-book-sequences}

\subsection{تعریف دنباله}
\label{sec:math-book-sequence-def}

\begin{definition}[دنباله]
    یک \textbf{دنباله}\index{دنباله} در
    مجموعه‌ی \lr{$\R$}، تابعی است از
    \lr{$\N$} به \lr{$\R$}:
    \[
        a: \N \to \R, \quad n \mapsto a_n
    \]
    که در آن \lr{$a_n$} را جمله‌ی \lr{$n$}-ام
    دنباله می‌نامیم.
    \label{def:sequence}
\end{definition}

طبق \cref{def:sequence}، یک دنباله را
می‌توان به‌صورت \lr{$(a_n)_{n=1}^{\infty}$}
نمایش داد.

\begin{example}
    دنباله‌ی \lr{$a_n = \frac{1}{n}$}:
    \[
        1, \frac{1}{2}, \frac{1}{3}, \frac{1}{4}, \ldots
    \]
    این دنباله به صفر نزدیک می‌شود، ولی
    هرگز به آن نمی‌رسد.
    \label{ex:harmonic}
\end{example}

\subsection{همگرایی}
\label{sec:math-book-convergence}

\begin{definition}[همگرایی دنباله]
    دنباله‌ی \lr{$(a_n)$} به عدد \lr{$L$}
    \textbf{همگرا}\index{همگرایی} است اگر:
    \[
        \forall \varepsilon > 0, \;
        \exists N \in \N, \;
        \forall n \geq N: \abs{a_n - L} < \varepsilon
    \]
    در این صورت می‌نویسیم
    \lr{$\lim_{n \to \infty} a_n = L$}.
    \label{def:convergence}
\end{definition}

\mnote{تعریف همگرایی، یکی از مهم‌ترین تعاریف در آنالیز حقیقی است. این تعریف، مفهوم «نزدیک شدن» را دقیق می‌کند.}

\section{قضایای مهم}
\label{sec:math-book-theorems}

\subsection{یکتایی حد}
\label{sec:math-book-unique-limit}

\begin{theorem}[یکتایی حد]
    اگر دنباله‌ای همگرا باشد، حد آن یکتاست.
    \label{thm:unique-limit}
\end{theorem}

\begin{proof}
    فرض کنید \lr{$\lim a_n = L_1$} و
    \lr{$\lim a_n = L_2$} با
    \lr{$L_1 \neq L_2$}. قرار دهید
    \lr{$\varepsilon = \frac{\abs{L_1 - L_2}}{2} > 0$}.

    طبق \cref{def:convergence}، عدد \lr{$N_1$}
    وجود دارد که برای \lr{$n \geq N_1$}،
    \lr{$\abs{a_n - L_1} < \varepsilon$}.
    همچنین عدد \lr{$N_2$} وجود دارد که برای
    \lr{$n \geq N_2$}،
    \lr{$\abs{a_n - L_2} < \varepsilon$}.

    برای \lr{$n \geq \max(N_1, N_2)$}:
    \begin{align*}
        \abs{L_1 - L_2}
        &= \abs{(L_1 - a_n) + (a_n - L_2)} \\
        &\leq \abs{L_1 - a_n} + \abs{a_n - L_2} \\
        &< \varepsilon + \varepsilon = \abs{L_1 - L_2}
    \end{align*}
    که تناقض است. بنابراین \lr{$L_1 = L_2$}.
\end{proof}

\subsection{قضیه‌ی فشردگی}
\label{sec:math-book-squeeze}

\begin{theorem}[قضیه‌ی فشردگی]
    اگر \lr{$a_n \leq b_n \leq c_n$} برای
    همه‌ی \lr{$n \geq N_0$} و
    \lr{$\lim a_n = \lim c_n = L$}، آنگاه
    \lr{$\lim b_n = L$}.
    \label{thm:squeeze}
\end{theorem}

\begin{proof}
    طبق \cref{def:convergence}، برای هر
    \lr{$\varepsilon > 0$}، اعداد \lr{$N_1$}
    و \lr{$N_2$} وجود دارند که:
    \begin{align*}
        n \geq N_1 &\implies \abs{a_n - L} < \varepsilon \\
        n \geq N_2 &\implies \abs{c_n - L} < \varepsilon
    \end{align*}

    برای \lr{$n \geq \max(N_0, N_1, N_2)$}:
    \[
        L - \varepsilon < a_n \leq b_n \leq c_n < L + \varepsilon
    \]
    بنابراین \lr{$\abs{b_n - L} < \varepsilon$}،
    که نتیجه می‌دهد \lr{$\lim b_n = L$}.
\end{proof}

\begin{remark}
    قضیه‌ی فشردگی (\cref{thm:squeeze}) یکی از
    ابزارهای قدرتمند برای محاسبه‌ی حد
    دنباله‌های پیچیده است. برای مثال، برای
    محاسبه‌ی \lr{$\lim \frac{\sin n}{n}$}،
    از نامساوی \lr{$-\frac{1}{n} \leq \frac{\sin n}{n} \leq \frac{1}{n}$}
    استفاده می‌کنیم.
    \label{rem:squeeze-application}
\end{remark}

\subsection{قضیه‌ی بولتزانو-وایرشتراس}
\label{sec:math-book-bolzano}

\begin{theorem}[بولتزانو-وایرشتراس]
    هر دنباله‌ی کراندار در \lr{$\R$} دارای
    زیردنباله‌ی همگراست.
    \label{thm:bolzano-weierstrass}
\end{theorem}

\mnote{قضیه‌ی بولتزانو-وایرشتراس، یکی از نتایج بنیادین در آنالیز حقیقی است. اثبات آن از اصل کامل بودن \lr{$\R$} استفاده می‌کند.}

\section{نمودارها}
\label{sec:math-book-figures}

نمودار \cref{fig:convergence-plot} همگرایی
دنباله‌ی \lr{$a_n = \frac{1}{n}$} را نشان
می‌دهد.

\begin{figure}[h]
\centering
\begin{latin}
\begin{tikzpicture}
  \begin{axis}[
    width=0.8\textwidth, height=6cm,
    xlabel={$n$},
    ylabel={$a_n$},
    grid=major,
    domain=1:20,
    samples=20,
    ymin=0, ymax=1.1
  ]
    \addplot[only marks, mark=*, color=matinblue, mark size=2pt]
        {1/x};
    \addlegendentry{$a_n = 1/n$}

    \addplot[matinred, thick, dashed, domain=1:20]
        {0};
    \addlegendentry{$\lim a_n = 0$}
  \end{axis}
\end{tikzpicture}
\end{latin}
\caption{همگرایی دنباله‌ی \lr{$a_n = \frac{1}{n}$} به صفر}
\label{fig:convergence-plot}
\end{figure}

\mnote{نمودار \lr{\cref{fig:convergence-plot}} به‌وضوح نشان می‌دهد که دنباله به صفر نزدیک می‌شود ولی هرگز به آن نمی‌رسد.}

\section{تمرین‌ها}
\label{sec:math-book-exercises}

\begin{exercise}
    ثابت کنید دنباله‌ی \lr{$a_n = \frac{n}{n+1}$}
    به \lr{$1$} همگراست.
    \label{exc:convergence-n-over-n1}
\end{exercise}

\begin{solution}
    باید نشان دهیم که برای هر
    \lr{$\varepsilon > 0$}، عدد \lr{$N$}
    وجود دارد که برای \lr{$n \geq N$}:
    \[
        \abs{\frac{n}{n+1} - 1} = \frac{1}{n+1} < \varepsilon
    \]
    با انتخاب \lr{$N > \frac{1}{\varepsilon} - 1$}،
    این نامساوی برقرار است.
\end{solution}

\begin{exercise}
    با استفاده از \cref{thm:squeeze}، حد
    دنباله‌ی \lr{$a_n = \frac{\cos n}{n}$} را
    محاسبه کنید.
    \label{exc:squeeze-cos}
\end{exercise}

\begin{solution}
    از نامساوی
    \lr{$-\frac{1}{n} \leq \frac{\cos n}{n} \leq \frac{1}{n}$}
    و \cref{thm:squeeze} استفاده می‌کنیم.
    چون \lr{$\lim \frac{1}{n} = \lim \left(-\frac{1}{n}\right) = 0$}،
    داریم \lr{$\lim \frac{\cos n}{n} = 0$}.
\end{solution}

\begin{exercise}
    آیا دنباله‌ی \lr{$a_n = (-1)^n$} همگراست؟
    چرا؟
    \label{exc:alternating}
\end{exercise}

\section{خلاصه}
\label{sec:math-book-summary}

در این فصل، یک فصل کامل از یک کتاب ریاضی
را دیدیم که شامل:

\begin{itemize}
    \item \textbf{تعاریف}: دنباله، همگرایی.
    \item \textbf{قضایا}: یکتایی حد، فشردگی،
    بولتزانو-وایرشتراس.
    \item \textbf{اثبات‌ها}: با استدلال ریاضی.
    \item \textbf{مثال‌ها}: دنباله‌ی هارمونیک.
    \item \textbf{نکته‌ها}: کاربرد قضیه‌ی
    فشردگی.
    \item \textbf{نمودارها}: همگرایی دنباله.
    \item \textbf{تمرین‌ها}: با راه‌حل.
    \item \textbf{نمایه}: با \cmd{index}.
    \item \textbf{ارجاعات}: با \cmd{cref}.
\end{itemize}

\mnote{این فصل، الگویی برای نوشتن فصل‌های کتاب ریاضی شماست. می‌توانید ساختار آن را برای موضوعات دیگر تطبیق دهید.}

در فصل بعدی (\cref{chap:programming-book})،
یک فصل نمونه از یک کتاب برنامه‌نویسی را
می‌بینیم که بر کد و الگوریتم تمرکز دارد.
%%
%% part5-examples/ch28-programming-book.tex
%%
%% Chapter 28: A Programming Book — A complete sample
%% chapter from a programming textbook, focusing on code,
%% algorithms, and data structures.
%%

\chapter{کتاب برنامه‌نویسی}
\label{chap:programming-book}

\section{چرا این موضوع مهم است؟}
\label{sec:programming-book-why}

در فصل قبل، یک فصل نمونه از کتاب ریاضی
دیدیم که بر قضیه، اثبات و معادلات تمرکز
داشت. در این فصل، یک فصل نمونه از کتاب
برنامه‌نویسی را می‌بینیم که بر کد، الگوریتم
و ساختار داده تمرکز دارد.

هدف این فصل، نشان دادن این است که چگونه
می‌توان از قابلیت‌های \lr{MatinBook} برای
نوشتن کتاب‌های برنامه‌نویسی استفاده کرد.

\mnote{کتاب‌های برنامه‌نویسی چالش‌های خاص خود را دارند: کد، الگوریتم، و توضیحات فارسی باید در کنار هم هماهنگ باشند.}

\section{ساختار فصل نمونه}
\label{sec:programming-book-structure}

فصل نمونه، از یک کتاب برنامه‌نویسی خیالی
با عنوان «مقدمه‌ای بر ساختمان داده‌ها»
است. این فصل شامل بخش‌های زیر است:

\begin{itemize}
    \item مقدمه.
    \item کد نمونه (پایتون و \lr{C++}).
    \item الگوریتم‌ها.
    \item جدول مقایسه.
    \item تمرین‌ها.
    \item راه‌حل‌ها.
\end{itemize}

\section{مقدمه}
\label{sec:programming-book-intro}

ساختمان داده\index{ساختمان داده}، روشی برای
سازمان‌دهی و ذخیره‌سازی داده‌ها در حافظه‌ی
کامپیوتر است. انتخاب ساختمان داده‌ی مناسب،
تأثیر مستقیمی بر کارایی الگوریتم‌ها دارد.

در این فصل، با ساختمان داده‌ی \textbf{پشته}
و کاربردهای آن آشنا می‌شویم.

\mnote{پشته یکی از ساده‌ترین و در عین حال پرکاربردترین ساختمان داده‌هاست.}

\section{پشته}
\label{sec:programming-book-stack}

\subsection{تعریف پشته}
\label{sec:programming-book-stack-def}

\begin{definition}[پشته]
    \textbf{پشته}\index{پشته} یک ساختمان
    داده‌ی خطی است که در آن درج و حذف
    عناصر از یک سر انجام می‌شود. این
    سر را \textbf{بالای پشته} می‌نامیم.
    \label{def:stack}
\end{definition}

طبق \cref{def:stack}، پشته از اصل
\textbf{\lr{LIFO}}\index{LIFO} (آخرین
ورودی، اولین خروجی) پیروی می‌کند.

\begin{example}
    عملیات اصلی پشته:
    \begin{itemize}
        \item \lr{push(x)}: درج \lr{x} در
        بالای پشته.
        \item \lr{pop()}: حذف و بازگشت عنصر
        بالای پشته.
        \item \lr{top()}: بازگشت عنصر بالای
        پشته بدون حذف آن.
        \item \lr{empty()}: بررسی خالی بودن
        پشته.
    \end{itemize}
    \label{ex:stack-operations}
\end{example}

\subsection{پیاده‌سازی در پایتون}
\label{sec:programming-book-python}

پیاده‌سازی پشته در پایتون با لیست:

\begin{latin}
\begin{minted}{python}
class Stack:
    """A simple stack implementation using a Python list."""

    def __init__(self) -> None:
        self._data: list = []

    def push(self, item) -> None:
        """Add an item to the top of the stack."""
        self._data.append(item)

    def pop(self):
        """Remove and return the top item."""
        if self.is_empty():
            raise IndexError("pop from empty stack")
        return self._data.pop()

    def peek(self):
        """Return the top item without removing it."""
        if self.is_empty():
            raise IndexError("peek from empty stack")
        return self._data[-1]

    def is_empty(self) -> bool:
        """Return True if the stack is empty."""
        return len(self._data) == 0

    def size(self) -> int:
        """Return the number of items in the stack."""
        return len(self._data)


# |\pc{استفاده از پشته}|
stack = Stack()
stack.push(10)
stack.push(20)
stack.push(30)
print(stack.pop())   # 30
print(stack.peek())  # 20
\end{minted}
\end{latin}
\codecaption{پیاده‌سازی پشته در پایتون}
\label{code:stack-python}

\mnote{در پایتون، لیست به‌عنوان ساختار پایه برای پشته استفاده می‌شود. عملیات \lr{append} و \lr{pop} هر دو \lr{$O(1)$} هستند.}

\subsection{پیاده‌سازی در \lr{C++}}
\label{sec:programming-book-cpp}

پیاده‌سازی پشته در \lr{C++} با قالب:

\begin{latin}
\begin{minted}{cpp}
#include <vector>
#include <stdexcept>

template <typename T>
class Stack {
private:
    std::vector<T> data;

public:
    void push(const T& item) {
        data.push_back(item);
    }

    T pop() {
        if (data.empty()) {
            throw std::runtime_error("pop from empty stack");
        }
        T item = data.back();
        data.pop_back();
        return item;
    }

    const T& peek() const {
        if (data.empty()) {
            throw std::runtime_error("peek from empty stack");
        }
        return data.back();
    }

    bool empty() const {
        return data.empty();
    }

    size_t size() const {
        return data.size();
    }
};
\end{minted}
\end{latin}
\codecaption{پیاده‌سازی پشته در \lr{C++}}
\label{code:stack-cpp}

\mnote{در \lr{C++}، از قالب (\lr{template}) استفاده می‌کنیم تا پشته با هر نوع داده‌ای کار کند.}

\section{الگوریتم‌ها}
\label{sec:programming-book-algorithms}

\subsection{بررسی توازن پرانتزها}
\label{sec:programming-book-parentheses}

یکی از کاربردهای کلاسیک پشته، بررسی
توازن پرانتزها در یک عبارت است.

\begin{latin}
\begin{algorithm}
\caption{بررسی توازن پرانتزها}
\label{alg:balanced-parentheses}
\begin{algorithmic}[1]
    \Require A string $s$ containing parentheses
    \Ensure \textbf{true} if the parentheses are balanced,
            \textbf{false} otherwise
    \State Create an empty stack $S$
    \For{each character $c$ in $s$}
        \If{$c$ is an opening bracket}
            \State $S.\text{push}(c)$
        \ElsIf{$c$ is a closing bracket}
            \If{$S$ is empty}
                \State \Return \textbf{false}
            \EndIf
            \State $top \gets S.\text{pop}()$
            \If{$top$ does not match $c$}
                \State \Return \textbf{false}
            \EndIf
        \EndIf
    \EndFor
    \State \Return $S$ is empty
\end{algorithmic}
\end{algorithm}
\end{latin}

\mnote{این الگوریتم در کامپایلرها، ویرایشگرهای متن و تحلیل‌گرهای نحوی کاربرد دارد.}

\subsection{پیچیدگی}
\label{sec:programming-book-complexity}

\begin{theorem}[پیچیدگی بررسی توازن]
    الگوریتم \cref{alg:balanced-parentheses}
    دارای پیچیدگی زمانی \lr{$O(n)$} و
    پیچیدگی فضایی \lr{$O(n)$} است، که
    \lr{$n$} طول رشته‌ی ورودی است.
    \label{thm:balanced-complexity}
\end{theorem}

\begin{proof}
    هر کاراکتر رشته دقیقاً یک بار پردازش
    می‌شود (حلقه‌ی \lr{for} در خط ۲). عملیات
    \lr{push} و \lr{pop} هر دو \lr{$O(1)$}
    هستند. بنابراین پیچیدگی زمانی
    \lr{$O(n)$} است.

    در بدترین حالت (رشته‌ای با \lr{$n$}
    پرانتز باز)، پشته دارای \lr{$n$} عنصر
    خواهد بود. بنابراین پیچیدگی فضایی
    \lr{$O(n)$} است.
\end{proof}

\section{جدول مقایسه}
\label{sec:programming-book-comparison}

جدول \cref{tab:stack-comparison} پیاده‌سازی
پشته در دو زبان را مقایسه می‌کند.

\begin{matintable}{مقایسه‌ی پیاده‌سازی پشته در پایتون و \lr{C++}}{tab:stack-comparison}
\begin{tblr}{
    width = \textwidth,
    colspec = {l X[c] X[c]},
    hlines, vlines,
    colsep = 8pt,
    rowsep = 4pt,
    row{1} = {font=\bfseries},
}
ویژگی & پایتون & \lr{C++} \\
نوع‌دهی & پویا & ایستا \\
ساختار پایه & \lr{list} & \lr{vector} \\
قالب & ندارد & دارد \\
مدیریت خطا & \lr{IndexError} & \lr{std::runtime\_error} \\
کارایی & متوسط & بالا \\
\end{tblr}
\end{matintable}

\mnote{در پایتون، نوع‌دهی پویا و در \lr{C++}، نوع‌دهی ایستا است. هر دو رویکرد مزایا و معایب خود را دارند.}

\section{تمرین‌ها}
\label{sec:programming-book-exercises}

\begin{exercise}
    یک تابع بنویسید که با استفاده از پشته،
    یک رشته را معکوس کند.
    \label{exc:reverse-string}
\end{exercise}

\begin{solution}
    \begin{latin}
\begin{minted}{python}
def reverse_string(s: str) -> str:
    """Reverse a string using a stack."""
    stack = Stack()
    for char in s:
        stack.push(char)

    result = []
    while not stack.is_empty():
        result.append(stack.pop())

    return "".join(result)


# |\pc{تست}|
print(reverse_string("hello"))  # "olleh"
\end{minted}
    \end{latin}
\end{solution}

\begin{exercise}
    با استفاده از پشته، یک ماشین حساب ساده
    برای عبارات پسوندی (\lr{postfix}) بنویسید.
    \label{exc:postfix-calculator}
\end{exercise}

\begin{exercise}
    تفاوت پشته و صف را توضیح دهید و یک
    کاربرد واقعی برای هر یک بنویسید.
    \label{exc:stack-vs-queue}
\end{exercise}

\section{خلاصه}
\label{sec:programming-book-summary}

در این فصل، یک فصل کامل از یک کتاب
برنامه‌نویسی را دیدیم که شامل:

\begin{itemize}
    \item \textbf{تعاریف}: پشته، \lr{LIFO}.
    \item \textbf{کد}: پیاده‌سازی در پایتون
    و \lr{C++}.
    \item \textbf{الگوریتم}: بررسی توازن
    پرانتزها.
    \item \textbf{قضیه}: پیچیدگی الگوریتم.
    \item \textbf{اثبات}: تحلیل پیچیدگی.
    \item \textbf{جدول}: مقایسه‌ی دو زبان.
    \item \textbf{تمرین‌ها}: با راه‌حل.
\end{itemize}

\mnote{این فصل، الگویی برای نوشتن فصل‌های کتاب برنامه‌نویسی شماست. می‌توانید ساختار آن را برای موضوعات دیگر (مثل درخت، گراف، مرتب‌سازی) تطبیق دهید.}

در فصل بعدی (\cref{chap:physics-book})،
یک فصل نمونه از یک کتاب فیزیک را می‌بینیم
که بر معادلات و نمودارهای فیزیکی تمرکز
دارد.
%%
%% part5-examples/ch29-physics-book.tex
%%
%% Chapter 29: A Physics Book — A complete sample chapter
%% from a physics textbook, focusing on equations, units,
%% and experimental data.
%%

\chapter{کتاب فیزیک}
\label{chap:physics-book}

\section{چرا این موضوع مهم است؟}
\label{sec:physics-book-why}

در دو فصل قبل، نمونه‌هایی از کتاب ریاضی
(با تمرکز بر قضیه و اثبات) و کتاب
برنامه‌نویسی (با تمرکز بر کد و الگوریتم)
دیدیم. در این فصل، نمونه‌ای از کتاب
فیزیک را می‌بینیم که چالش‌های خاص خود
را دارد: معادلات، یکاها، داده‌های تجربی
و نمودارها.

هدف این فصل، نشان دادن این است که چگونه
می‌توان از \lr{MatinBook} برای نوشتن
کتاب‌های فیزیک استفاده کرد.

\mnote{کتاب‌های فیزیک معمولاً ترکیبی از معادلات پیچیده، جداول داده‌های تجربی و نمودارهای دقیق هستند.}

\section{ساختار فصل نمونه}
\label{sec:physics-book-structure}

فصل نمونه، از یک کتاب فیزیک خیالی با
عنوان «مبانی مکانیک کلاسیک» است. این
فصل شامل بخش‌های زیر است:

\begin{itemize}
    \item مقدمه.
    \item معادلات حرکت.
    \item یکاها و ابعاد.
    \item داده‌های تجربی.
    \item نمودارها.
    \item تمرین‌ها.
\end{itemize}

\section{مقدمه}
\label{sec:physics-book-intro}

مکانیک کلاسیک\index{مکانیک کلاسیک}،
شاخه‌ای از فیزیک است که به مطالعه‌ی حرکت
اجسام می‌پردازد. این شاخه، پایه‌ی بسیاری
از شاخه‌های دیگر فیزیک است.

در این فصل، با قوانین نیوتن و معادلات
حرکت آشنا می‌شویم.

\mnote{مکانیک کلاسیک توسط \lr{Isaac Newton} در قرن هفدهم پایه‌گذاری شد.}

\section{قوانین نیوتن}
\label{sec:physics-book-newton}

\subsection{قانون اول}
\label{sec:physics-book-newton-first}

\begin{theorem}[قانون اول نیوتن]
    هر جسم در حال سکون یا حرکت یکنواخت
    در خط راست، به همان حالت باقی می‌ماند
    مگر اینکه نیروی خارجی به آن وارد شود.
    \label{thm:newton-first}
\end{theorem}

\mnote{قانون اول نیوتن را «قانون اینرسی» یا «قانون لختی» نیز می‌نامند.}

\subsection{قانون دوم}
\label{sec:physics-book-newton-second}

\begin{theorem}[قانون دوم نیوتن]
    شتاب جسم با نیروی خالص وارد بر آن
    متناسب و با جرم آن نسبت عکس دارد:
    \[
        \vec{F}_{\text{net}} = m \vec{a}
    \]
    \label{thm:newton-second}
\end{theorem}

طبق \cref{thm:newton-second}، اگر نیروی
خالص وارد بر جسمی صفر باشد، شتاب آن نیز
صفر است.

\begin{example}
    جسمی به جرم \lr{$2\,\text{kg}$} تحت
    نیروی \lr{$10\,\text{N}$} قرار دارد.
    شتاب آن برابر است با:
    \[
        a = \frac{F}{m} = \frac{10}{2} = 5\,\text{m/s}^2
    \]
    \label{ex:newton-second}
\end{example}

\subsection{قانون سوم}
\label{sec:physics-book-newton-third}

\begin{theorem}[قانون سوم نیوتن]
    برای هر کنش، واکنشی مساوی و در جهت
    مخالف وجود دارد:
    \[
        \vec{F}_{12} = -\vec{F}_{21}
    \]
    \label{thm:newton-third}
\end{theorem}

\mnote{قوانین نیوتن \cref{thm:newton-first,thm:newton-second,thm:newton-third} پایه‌ی مکانیک کلاسیک هستند.}

\section{معادلات حرکت}
\label{sec:physics-book-equations}

\subsection{حرکت با شتاب ثابت}
\label{sec:physics-book-constant-accel}

برای حرکت با شتاب ثابت \lr{$a$}، معادلات
حرکت به‌صورت زیر هستند:

\begin{align}
    v &= v_0 + at \label{eq:velocity} \\
    x &= x_0 + v_0 t + \frac{1}{2} a t^2 \label{eq:position} \\
    v^2 &= v_0^2 + 2a(x - x_0) \label{eq:velocity-position}
\end{align}

معادلات \cref{eq:velocity,eq:position,eq:velocity-position}
به‌عنوان معادلات سینماتیک شناخته می‌شوند.

\mnote{این معادلات فقط برای شتاب ثابت معتبرند. برای شتاب متغیر، باید از حساب دیفرانسیل استفاده کرد.}

\begin{example}
    خودرویی با سرعت اولیه‌ی
    \lr{$20\,\text{m/s}$} با شتاب
    \lr{$-5\,\text{m/s}^2$} ترمز می‌کند.
    مدت زمان تا توقف:
    \[
        t = \frac{v - v_0}{a} = \frac{0 - 20}{-5} = 4\,\text{s}
    \]
    \label{ex:braking}
\end{example}

\subsection{حرکت پرتابی}
\label{sec:physics-book-projectile}

برای حرکت پرتابی در دو بُعد، معادلات
به‌صورت زیر هستند:

\begin{align}
    x(t) &= v_0 \cos\theta \cdot t \\
    y(t) &= v_0 \sin\theta \cdot t - \frac{1}{2} g t^2
\end{align}

\mnote{در حرکت پرتابی، مؤلفه‌ی افقی سرعت ثابت است ولی مؤلفه‌ی عمودی تحت تأثیر گرانش تغییر می‌کند.}

\section{یکاها و ابعاد}
\label{sec:physics-book-units}

در فیزیک، استفاده‌ی درست از یکاها
ضروری است. جدول \cref{tab:si-units}
یکاهای اصلی \lr{SI} را نشان می‌دهد.

\begin{matintable}{یکاهای اصلی \lr{SI}}{tab:si-units}
\begin{tblr}{
    width = \textwidth,
    colspec = {l X[l] X[l] X[c]},
    hlines, vlines,
    colsep = 8pt,
    rowsep = 4pt,
    row{1} = {font=\bfseries},
}
کمیت & نام یکا & نماد & بُعد \\
طول & متر & \lr{m} & \lr{L} \\
جرم & کیلوگرم & \lr{kg} & \lr{M} \\
زمان & ثانیه & \lr{s} & \lr{T} \\
جریان & آمپر & \lr{A} & \lr{I} \\
دما & کلوین & \lr{K} & \lr{Θ} \\
مقدار ماده & مول & \lr{mol} & \lr{N} \\
شدت روشنایی & کندلا & \lr{cd} & \lr{J} \\
\end{tblr}
\end{matintable}

\mnote{کمیت‌های فرعی (مثل نیرو، انرژی، توان) از ترکیب یکاهای اصلی ساخته می‌شوند.}

\subsection{تحلیل ابعادی}
\label{sec:physics-book-dimensional}

\begin{example}
    بررسی ابعادی معادله‌ی \lr{$v = v_0 + at$}:
    \begin{itemize}
        \item \lr{$v$} و \lr{$v_0$}: \lr{$LT^{-1}$}.
        \item \lr{$at$}: \lr{$LT^{-2} \cdot T = LT^{-1}$}.
    \end{itemize}
    بنابراین دو طرف معادله، ابعاد یکسان
    دارند.
    \label{ex:dimensional-analysis}
\end{example}

\mnote{تحلیل ابعادی، ابزاری قدرتمند برای بررسی درستی معادلات فیزیکی است.}

\section{داده‌های تجربی}
\label{sec:physics-book-data}

جدول \cref{tab:free-fall-data} داده‌های
تجربی مربوط به سقوط آزاد را نشان می‌دهد.

\begin{matintable}{داده‌های تجربی سقوط آزاد}{tab:free-fall-data}
\begin{tblr}{
    width = \textwidth,
    colspec = {c c c},
    hlines, vlines,
    colsep = 8pt,
    rowsep = 4pt,
    row{1} = {font=\bfseries},
}
زمان \lr{(s)} & سرعت \lr{(m/s)} & ارتفاع \lr{(m)} \\
۰ & ۰٫۰ & ۰٫۰ \\
۰٫۵ & ۴٫۹ & ۱٫۲ \\
۱٫۰ & ۹٫۸ & ۴٫۹ \\
۱٫۵ & ۱۴٫۷ & ۱۱٫۰ \\
۲٫۰ & ۱۹٫۶ & ۱۹٫۶ \\
\end{tblr}
\end{matintable}

\mnote{داده‌های تجربی معمولاً با عدم قطعیت همراه هستند. در این جدول، مقادیر با یک رقم اعشار گزارش شده‌اند.}

\section{نمودارها}
\label{sec:physics-book-graphs}

نمودار \cref{fig:free-fall} داده‌های
سقوط آزاد را نمایش می‌دهد.

\begin{figure}[h]
\centering
\begin{latin}
\begin{tikzpicture}
  \begin{axis}[
    width=0.8\textwidth, height=6cm,
    xlabel={$t$ (s)},
    ylabel={$v$ (m/s)},
    grid=major,
    legend pos=north west
  ]
    \addplot[only marks, mark=*, color=matinblue, mark size=2.5pt]
        coordinates {(0,0) (0.5,4.9) (1,9.8) (1.5,14.7) (2,19.6)};
    \addlegendentry{|\pc{داده‌های تجربی}|}

    \addplot[matinred, thick, domain=0:2] {9.8*x};
    \addlegendentry{$v = gt$}
  \end{axis}
\end{tikzpicture}
\end{latin}
\caption{نمودار سرعت-زمان در سقوط آزاد}
\label{fig:free-fall}
\end{figure}

\mnote{نمودار \lr{\cref{fig:free-fall}} نشان می‌دهد که سرعت با زمان خطی افزایش می‌یابد و شیب خط، شتاب گرانش \lr{$g$} است.}

\section{تمرین‌ها}
\label{sec:physics-book-exercises}

\begin{exercise}
    جسمی از ارتفاع \lr{$45\,\text{m}$}
    رها می‌شود. با فرض \lr{$g = 10\,\text{m/s}^2$}،
    سرعت آن هنگام برخورد با زمین چقدر
    است؟
    \label{exc:free-fall}
\end{exercise}

\begin{solution}
    از معادله‌ی \cref{eq:velocity-position}:
    \[
        v^2 = v_0^2 + 2g h = 0 + 2 \cdot 10 \cdot 45 = 900
    \]
    بنابراین \lr{$v = 30\,\text{m/s}$}.
\end{solution}

\begin{exercise}
    خودرویی با سرعت \lr{$72\,\text{km/h}$}
    در حال حرکت است. راننده ترمز می‌کند و
    خودرو پس از \lr{$40\,\text{m}$} متوقف
    می‌شود. شتاب ترمز چقدر است؟
    \label{exc:braking-distance}
\end{exercise}

\begin{exercise}
    ابعاد کمیت \lr{$F = ma$} را با استفاده
    از یکاهای اصلی \lr{SI} بیان کنید.
    \label{exc:force-dimension}
\end{exercise}

\section{خلاصه}
\label{sec:physics-book-summary}

در این فصل، یک فصل کامل از یک کتاب
فیزیک را دیدیم که شامل:

\begin{itemize}
    \item \textbf{قضایا}: سه قانون نیوتن.
    \item \textbf{معادلات}: سینماتیک،
    حرکت پرتابی.
    \item \textbf{اثبات‌ها}: برای قضایا.
    \item \textbf{مثال‌ها}: با محاسبات عددی.
    \item \textbf{جداول}: یکاهای \lr{SI}،
    داده‌های تجربی.
    \item \textbf{نمودارها}: سرعت-زمان.
    \item \textbf{تمرین‌ها}: با راه‌حل.
\end{itemize}

\mnote{این فصل، الگویی برای نوشتن فصل‌های کتاب فیزیک شماست. می‌توانید ساختار آن را برای موضوعات دیگر (مثل الکترومغناطیس، ترمودینامیک) تطبیق دهید.}

این فصل، آخرین فصل از بخش پنجم (مثال‌های
کامل) بود. در بخش ششم (مرجع دستورات)،
مرجع کامل کلاس، محیط‌ها و دستورات ارائه
می‌شود.

%========================================================================
% PART 6 — REFERENCE
%========================================================================
\part{مرجع دستورات}
\label{part:reference}

\input{part6-reference/ch30-class-options}
\input{part6-reference/ch31-environments}
\input{part6-reference/ch32-commands}

%========================================================================
% APPENDICES
%========================================================================
\appendix

\input{appendices/app-a-troubleshooting}
\input{appendices/app-b-tools}
\input{appendices/app-c-resources}
\input{appendices/app-d-changelog}
\input{appendices/app-e-license}

%------------------------------------------------------------------------
% BACK MATTER (symmetric margins)
%------------------------------------------------------------------------
\backmatter
\frontmattergeometry

% Bibliography
\printbibliography[title={منابع و مراجع}]

% Index
\printindex

%------------------------------------------------------------------------
% BACK COVER
%------------------------------------------------------------------------
\authorbio{%
    مatin محمودی، پژوهشگر و توسعه‌دهنده‌ی ابزارهای حروف‌چینی
    فارسی است. او سال‌ها در زمینه‌ی نشر دیجیتال و طراحی
    کلاس‌های \lr{LaTeX} برای زبان فارسی فعالیت کرده است.
    \lr{MatinBook} حاصل تجربه‌ی او در این حوزه است.
}

\makebackcover{%
    \textbf{راهنمای جامع MatinBook} مرجعی کامل برای
    نوشتن کتاب‌های فنی و علمی فارسی با استفاده از کلاس
    \lr{MatinBook} است.

    \vspace{0.4cm}

    \textbf{در این راهنما می‌آموزید:}
    \begin{itemize}
        \item نصب و راه‌اندازی MatinBook
        \item ساختار یک کتاب حرفه‌ای
        \item محیط‌های علمی، ریاضی و کد
        \item نمودار، جدول و شکل
        \item شخصی‌سازی رنگ، فونت و چیدمان
        \item ایجاد تم سفارشی
        \item انتشار کتاب نهایی
    \end{itemize}

    \vspace{0.4cm}

    مناسب برای نویسندگان، محققان، دانشجویان و
    مدرسانی که می‌خواهند کتاب‌های فنی فارسی
    با کیفیت حرفه‌ای تولید کنند.
}

\end{document}